% Generated by the blueprint tool (blueprint-pdf/1) from a blueprint snapshot; edit the
% Markdown sources, not this file.  Compile with `latexmk -pdf` or by running
% `pdflatex` twice (the second run settles the contents and the page references).
% Images, if any, are read from assets/ next to this file.
\documentclass[11pt]{article}
\usepackage[T1]{fontenc}
\usepackage[utf8]{inputenc}
\usepackage{lmodern}
\usepackage{textcomp}
\usepackage[a4paper,margin=2.5cm,headheight=14pt]{geometry}
\usepackage{amsmath,amssymb,amsthm}
\usepackage{amscd}  % KaTeX's CD diagrams; part of amsmath, so always there
\usepackage{mathtools}
\usepackage{bm}
\usepackage{graphicx}
\usepackage{xcolor}
\usepackage{array,booktabs,tabularx}
\usepackage{fancyvrb}
\usepackage[normalem]{ulem}
\usepackage{microtype}
% Optional, and only with their fonts: a small TeX installation has the
% packages without the fonts, which fails at the first formula.
\IfFileExists{mathrsfs.sty}{\IfFileExists{rsfs10.tfm}{\usepackage{mathrsfs}}{}}{}
\IfFileExists{stmaryrd.sty}{\IfFileExists{stmary10.tfm}{\usepackage{stmaryrd}}{}}{}
\IfFileExists{cancel.sty}{\usepackage{cancel}}{}
\usepackage[hyphens]{url}
\usepackage[pdfusetitle,bookmarksnumbered=false]{hyperref}
\usepackage{bookmark}
\definecolor{bplink}{HTML}{1F4E9C}
\definecolor{bpmuted}{HTML}{5F676E}
\hypersetup{colorlinks=true,linkcolor=bplink,urlcolor=bplink,citecolor=bplink}
\setlength{\parindent}{0pt}
\setlength{\parskip}{0.55em plus 0.2em minus 0.1em}
% Lists as tight as the page's.
\makeatletter
\def\@listi{\leftmargin\leftmargini\parsep 0pt\topsep 0.3em plus 0.1em\itemsep 0.1em plus 0.1em}
\let\@listI\@listi
\def\@listii{\leftmargin\leftmarginii\labelwidth\leftmarginii\advance\labelwidth-\labelsep\parsep 0pt\topsep 0.1em\itemsep 0.1em}
% An ordered list that does not start at 1.
\newcommand{\bpstartat}[1]{\setcounter{\@enumctr}{\numexpr#1-1\relax}}
\makeatother
\setcounter{tocdepth}{3}
\emergencystretch=3em
\allowdisplaybreaks

% A numbered statement, set as a paper sets one: "Theorem 2.1 (Title)."
% #1 is the lead ("Theorem~2.1"), #2 what follows it in roman.
\newenvironment{bpstatement}[2]{\par\addvspace{0.6\baselineskip}\noindent\textbf{#1}#2\textbf{.}\enspace\ignorespaces}
  {\par\addvspace{0.4\baselineskip}}
\newenvironment{bpstatementit}[2]{\begin{bpstatement}{#1}{#2}\itshape\ignorespaces}{\end{bpstatement}}
% A step: a relation with prose of its own, read under the object it starts from.
\newenvironment{bpstep}[1]{\par\addvspace{0.3\baselineskip}\begin{list}{}{\leftmargin=1.6em\rightmargin=0pt\topsep=0pt}\item\relax
  {\small\color{bpmuted}#1}\par\nopagebreak}{\end{list}\par}
% Notes about where an object is written out.
\newenvironment{bpnote}{\par\small\color{bpmuted}\itshape}{\par}
\newcommand{\bpomitted}[1]{\par{\small\color{bpmuted}\itshape #1\par}}
\newcommand{\bpbodyheading}[1]{\par\addvspace{0.4\baselineskip}\noindent\textbf{#1}\par\nopagebreak}
\newcommand{\bprule}{\par\medskip\noindent\hfil\rule{0.3\linewidth}{0.4pt}\hfil\par\medskip}
\newcommand{\bpboundary}[1]{\par\nopagebreak{\small\color{bpmuted}#1\par}\nopagebreak}
% A reference to something this export leaves out, and one that resolves nowhere.
\newcommand{\bpout}[1]{#1\textsuperscript{\hyperlink{bp-outside}{\ensuremath{\ast}}}}
\newcommand{\bpbroken}[1]{\textcolor{red!70!black}{[#1?]}}
\newcommand{\bpslug}[1]{\texttt{#1}}
% A bibliography entry: label, anchor, entry; the label hangs in the margin.
\newcommand{\bpbibitem}[3]{\par\setbox0\hbox{[#1]\enspace}\ifdim\wd0<5em\setbox0\hbox to 5em{[#1]\hfil}\fi
  \hangindent=\wd0 \hangafter=1 \noindent\box0\phantomsection\label{#2}#3\par\smallskip}
% A character the export knows no way to set.
\newcommand{\bpglyph}[1]{\fbox{\scriptsize #1}}
% Images, at their own size unless they are wider than the line.
\newsavebox{\bpimagebox}
\newcommand{\bpimage}[1]{\sbox{\bpimagebox}{\includegraphics{#1}}%
  \ifdim\wd\bpimagebox>\linewidth\resizebox{\linewidth}{!}{\usebox{\bpimagebox}}\else\usebox{\bpimagebox}\fi}
\newcommand{\bpmissingimage}[1]{\fbox{\parbox{0.8\linewidth}{\centering\small\color{bpmuted}Image not included: #1}}}
\DefineVerbatimEnvironment{bpverbatim}{Verbatim}{fontsize=\small,xleftmargin=1em}
\renewcommand{\qedsymbol}{\ensuremath{\blacksquare}}
\hypersetup{pdftitle={Weil Conjectures — Part Ⅲ Weights and the Riemann hypothesis},pdfsubject={Part Ⅲ Weights and the Riemann hypothesis; all levels},pdfcreator={blueprint}}
% KaTeX commands the website knows
\providecommand{\R}{\mathbb{R}}\providecommand{\N}{\mathbb{N}}\providecommand{\Z}{\mathbb{Z}}
\providecommand{\Q}{\mathbb{Q}}\providecommand{\C}{\mathbb{C}}
\providecommand{\Reals}{\mathbb{R}}\providecommand{\reals}{\mathbb{R}}\providecommand{\Complex}{\mathbb{C}}
\providecommand{\cnums}{\mathbb{C}}\providecommand{\natnums}{\mathbb{N}}\providecommand{\Bbb}[1]{\mathbb{#1}}
\providecommand{\bold}[1]{\mathbf{#1}}\providecommand{\frak}[1]{\mathfrak{#1}}\providecommand{\mathbbm}[1]{\mathbb{#1}}
\providecommand{\lang}{\langle}\providecommand{\rang}{\rangle}\providecommand{\bra}[1]{\langle #1\rvert}
\providecommand{\ket}[1]{\lvert #1\rangle}\providecommand{\braket}[1]{\langle #1\rangle}
\providecommand{\set}[1]{\{#1\}}\providecommand{\Set}[1]{\left\{#1\right\}}
\providecommand{\infin}{\infty}\providecommand{\isin}{\in}\providecommand{\plusmn}{\pm}\providecommand{\sub}{\subset}
\providecommand{\sube}{\subseteq}\providecommand{\supe}{\supseteq}\providecommand{\exist}{\exists}
\providecommand{\bull}{\bullet}\providecommand{\larr}{\leftarrow}\providecommand{\rarr}{\rightarrow}
\providecommand{\lrarr}{\leftrightarrow}\providecommand{\Larr}{\Leftarrow}\providecommand{\Rarr}{\Rightarrow}
\providecommand{\Lrarr}{\Leftrightarrow}\providecommand{\uarr}{\uparrow}\providecommand{\darr}{\downarrow}
\providecommand{\harr}{\leftrightarrow}\providecommand{\lArr}{\Leftarrow}\providecommand{\rArr}{\Rightarrow}
\providecommand{\hArr}{\Leftrightarrow}\providecommand{\alef}{\aleph}\providecommand{\alefsym}{\aleph}
\providecommand{\real}{\Re}\providecommand{\image}{\Im}\providecommand{\weierp}{\wp}
\providecommand{\Alpha}{\mathrm{A}}\providecommand{\Beta}{\mathrm{B}}\providecommand{\Epsilon}{\mathrm{E}}
\providecommand{\Zeta}{\mathrm{Z}}\providecommand{\Eta}{\mathrm{H}}\providecommand{\Iota}{\mathrm{I}}
\providecommand{\Kappa}{\mathrm{K}}\providecommand{\Mu}{\mathrm{M}}\providecommand{\Nu}{\mathrm{N}}
\providecommand{\Omicron}{\mathrm{O}}\providecommand{\Rho}{\mathrm{P}}\providecommand{\Tau}{\mathrm{T}}
\providecommand{\Chi}{\mathrm{X}}\providecommand{\omicron}{o}\providecommand{\gt}{>}\providecommand{\lt}{<}
\providecommand{\And}{\&}\providecommand{\hdashline}{\hline}\providecommand{\operatornamewithlimits}[1]{\operatorname*{#1}}
\providecommand{\llbracket}{[\![}\providecommand{\rrbracket}{]\!]}\providecommand{\cancel}[1]{#1}
\providecommand{\mathscr}{\mathcal}
% the characters beyond ASCII this document uses
\DeclareUnicodeCharacter{2160}{\textup{I}}
\DeclareUnicodeCharacter{2161}{\textup{II}}
\DeclareUnicodeCharacter{2162}{\textup{III}}
% the project's macros ([katex.macros] in blueprint.toml)
\providecommand{\A}{}\renewcommand{\A}{\mathbf A}
\providecommand{\AEn}{}\renewcommand{\AEn}{A_{E,n}}
\providecommand{\Aaf}{}\renewcommand{\Aaf}{\mathbf A}
\providecommand{\Ab}{}\renewcommand{\Ab}{\mathsf{Ab}}
\providecommand{\AffSch}{}\renewcommand{\AffSch}{\operatorname{AffSch}}
\providecommand{\Aone}{}\renewcommand{\Aone}{\mathbf A^1}
\providecommand{\Bl}{}\renewcommand{\Bl}{\operatorname{Bl}}
\providecommand{\Br}{}\renewcommand{\Br}{\mathrm{Br}}
\providecommand{\C}{}\renewcommand{\C}{\mathbb C}
\providecommand{\Cech}{}\renewcommand{\Cech}{\check{H}}
\providecommand{\Coext}{}\renewcommand{\Coext}{\operatorname{Coext}}
\providecommand{\Cone}{}\renewcommand{\Cone}{\operatorname{Cone}}
\providecommand{\Cons}{}\renewcommand{\Cons}{\operatorname{Cons}}
\providecommand{\Cref}{}\renewcommand{\Cref}[1]{\text{#1}}
\providecommand{\D}{}\renewcommand{\D}{D}
\providecommand{\DLAFT}{}\renewcommand{\DLAFT}{\operatorname{FT}}
\providecommand{\DLAFr}{}\renewcommand{\DLAFr}{\operatorname{Fr}}
\providecommand{\DLAQ}{}\renewcommand{\DLAQ}{{\overline{\mathbf Q}_\ell}}
\providecommand{\DLRQ}{}\renewcommand{\DLRQ}{{\overline{\mathbf Q}_\ell}}
\providecommand{\Dbc}{}\renewcommand{\Dbc}{D^b_c}
\providecommand{\Dc}{}\renewcommand{\Dc}{D_{\mathrm{cons}}}
\providecommand{\Dctf}{}\renewcommand{\Dctf}{D^b_{ctf}}
\providecommand{\Dic}{}\renewcommand{\Dic}{D_{\IndCons}}
\providecommand{\Dplus}{}\renewcommand{\Dplus}{D^{+}}
\providecommand{\End}{}\renewcommand{\End}{\operatorname{End}}
\providecommand{\Et}{}\renewcommand{\Et}{\mathrm{\acute{e}t}}
\providecommand{\Ex}{}\renewcommand{\Ex}{\mathrm{Ex}}
\providecommand{\Ext}{}\renewcommand{\Ext}{\operatorname{Ext}}
\providecommand{\F}{}\renewcommand{\F}{\mathbb F}
\providecommand{\FT}{}\renewcommand{\FT}{\operatorname{FT}}
\providecommand{\Fbar}{}\renewcommand{\Fbar}{{\overline{\mathbf F}_q}}
\providecommand{\Fell}{}\renewcommand{\Fell}{\mathbf F_\ell}
\providecommand{\Fib}{}\renewcommand{\Fib}{\operatorname{Fib}}
\providecommand{\Fix}{}\renewcommand{\Fix}{\operatorname{Fix}}
\providecommand{\Fp}{}\renewcommand{\Fp}{{\mathbf F_{p}}}
\providecommand{\Fq}{}\renewcommand{\Fq}{{\mathbf F_q}}
\providecommand{\Fqn}{}\renewcommand{\Fqn}{{\mathbf F_{q^{n}}}}
\providecommand{\Frob}{}\renewcommand{\Frob}{\operatorname{Frob}}
\providecommand{\GL}{}\renewcommand{\GL}{\operatorname{GL}}
\providecommand{\Gal}{}\renewcommand{\Gal}{\operatorname{Gal}}
\providecommand{\Hens}{}\renewcommand{\Hens}{\operatorname{Hens}}
\providecommand{\Hilb}{}\renewcommand{\Hilb}{\operatorname{Hilb}}
\providecommand{\Hom}{}\renewcommand{\Hom}{\operatorname{Hom}}
\providecommand{\IndCons}{}\renewcommand{\IndCons}{\operatorname{IndCons}}
\providecommand{\Kar}{}\renewcommand{\Kar}{\operatorname{Kar}}
\providecommand{\LT}{}\renewcommand{\LT}{\operatorname{LT}}
\providecommand{\Lie}{}\renewcommand{\Lie}{\operatorname{Lie}}
\providecommand{\Loc}{}\renewcommand{\Loc}{\operatorname{Loc}}
\providecommand{\Mod}{}\renewcommand{\Mod}{\mathsf{Mod}}
\providecommand{\N}{}\renewcommand{\N}{\mathbb{N}}
\providecommand{\Needspace}{}\renewcommand{\Needspace}{\par}
\providecommand{\O}{}\renewcommand{\O}{{\mathcal{O}}}
\providecommand{\Ob}{}\renewcommand{\Ob}{\mathrm{Ob}}
\providecommand{\Oe}{}\renewcommand{\Oe}{\mathcal O_E}
\providecommand{\PP}{}\renewcommand{\PP}{\mathbf{P}}
\providecommand{\Perf}{}\renewcommand{\Perf}{\operatorname{Perf}}
\providecommand{\Pic}{}\renewcommand{\Pic}{\mathrm{Pic}}
\providecommand{\Pm}{}\renewcommand{\Pm}{\mathbf P^{m}}
\providecommand{\Pone}{}\renewcommand{\Pone}{\mathbf P^{1}}
\providecommand{\Pqn}{}\renewcommand{\Pqn}{\mathbf P^{n}}
\providecommand{\ProFin}{}\renewcommand{\ProFin}{\operatorname{ProFin}}
\providecommand{\Proet}{}\renewcommand{\Proet}{\mathrm{pro\acute{e}t}}
\providecommand{\Q}{}\renewcommand{\Q}{\mathbb Q}
\providecommand{\Qb}{}\renewcommand{\Qb}{\overline{\mathbf Q}_{\ell}}
\providecommand{\Qbar}{}\renewcommand{\Qbar}{\overline{\mathbf Q}}
\providecommand{\QbarX}{}\renewcommand{\QbarX}{\overline{\mathbf Q}_{\ell,X}}
\providecommand{\Qbarl}{}\renewcommand{\Qbarl}{{\overline{\mathbf Q}_{\ell}}}
\providecommand{\Qell}{}\renewcommand{\Qell}{{\mathbf{Q}_{\ell}}}
\providecommand{\Ql}{}\renewcommand{\Ql}{\mathbf Q_{\ell}}
\providecommand{\Qlb}{}\renewcommand{\Qlb}{{\overline{\mathbf Q}_\ell}}
\providecommand{\Qlbar}{}\renewcommand{\Qlbar}{\overline{\mathbf Q}_\ell}
\providecommand{\R}{}\renewcommand{\R}{\mathbb{R}}
\providecommand{\RG}{}\renewcommand{\RG}{\operatorname{\mathrm{R}\Gamma}}
\providecommand{\RGamma}{}\renewcommand{\RGamma}{R\Gamma}
\providecommand{\RGammac}{}\renewcommand{\RGammac}{R\Gamma_c}
\providecommand{\RHom}{}\renewcommand{\RHom}{R\mathcal{H}om}
\providecommand{\RHomX}{}\renewcommand{\RHomX}{R\mathcal H om_X}
\providecommand{\Rat}{}\renewcommand{\Rat}{\operatorname{IsRat}}
\providecommand{\Res}{}\renewcommand{\Res}{\operatorname{Res}}
\providecommand{\Rlim}{}\renewcommand{\Rlim}{R\!\varprojlim}
\providecommand{\SIm}{}\renewcommand{\SIm}{\operatorname{SIm}}
\providecommand{\Sch}{}\renewcommand{\Sch}{\mathsf{Sch}}
\providecommand{\Sh}{}\renewcommand{\Sh}{\operatorname{Sh}}
\providecommand{\Span}{}\renewcommand{\Span}{\operatorname{span}}
\providecommand{\Spec}{}\renewcommand{\Spec}{\operatorname{Spec}}
\providecommand{\Tr}{}\renewcommand{\Tr}{\operatorname{Tr}}
\providecommand{\Trsm}{}\renewcommand{\Trsm}{\operatorname{Tr}^{\mathrm{sm}}}
\providecommand{\Wt}{}\renewcommand{\Wt}{\operatorname{Wt}}
\providecommand{\Z}{}\renewcommand{\Z}{\mathbb{Z}}
\providecommand{\Zar}{}\renewcommand{\Zar}{\mathrm{zar}}
\providecommand{\Zell}{}\renewcommand{\Zell}{{\mathbf{Z}_{\ell}}}
\providecommand{\Zl}{}\renewcommand{\Zl}{{\mathbf Z_\ell}}
\providecommand{\affet}{}\renewcommand{\affet}[1]{{#1}^{\mathrm{aff}}_{\mathrm{ét}}}
\providecommand{\affproet}{}\renewcommand{\affproet}[1]{{#1}^{\mathrm{aff}}_{\mathrm{proét}}}
\providecommand{\boxt}{}\renewcommand{\boxt}{\boxtimes_A^{\mathrm L}}
\providecommand{\cF}{}\renewcommand{\cF}{\mathcal{F}}
\providecommand{\cHom}{}\renewcommand{\cHom}{R\mathcal{H}\!om}
\providecommand{\cI}{}\renewcommand{\cI}{\mathcal{I}}
\providecommand{\cL}{}\renewcommand{\cL}{\mathcal L}
\providecommand{\cM}{}\renewcommand{\cM}{\mathcal M}
\providecommand{\cN}{}\renewcommand{\cN}{\mathcal N}
\providecommand{\cO}{}\renewcommand{\cO}{\mathcal O}
\providecommand{\cT}{}\renewcommand{\cT}{\mathcal T}
\providecommand{\checked}{}\renewcommand{\checked}{\color{red}\aftergroup\restorecolor}
\providecommand{\coeff}{}\renewcommand{\coeff}{\operatorname{coeff}}
\providecommand{\coker}{}\renewcommand{\coker}{\operatorname{coker}}
\providecommand{\colim}{}\renewcommand{\colim}{\operatorname{colim}}
\providecommand{\constructedby}{}\renewcommand{\constructedby}[1]{\par\smallskip\noindent\textcolor{purple!60!black}{\textsf{Constructed by:} }
  \footnotesize\sloppy\nolinkurl{#1}\par}
\providecommand{\deepresult}{}\renewcommand{\deepresult}[1]{\par\smallskip\noindent\textcolor{orange!45!black}{
  \textsf{Blueprint role: named deep theorem node; see the cluster proof architecture.}}
  \par\noindent\textcolor{orange!45!black}{\textsf{Mathematical source:} }#1\par}
\providecommand{\dt}{}\renewcommand{\dt}{\mathbin{\otimes^{\mathbf L}}}
\providecommand{\dw}{}\renewcommand{\dw}{\operatorname{dw}}
\providecommand{\et}{}\renewcommand{\et}[1]{{#1}_{\mathrm{ét}}}
\providecommand{\ev}{}\renewcommand{\ev}{\operatorname{ev}}
\providecommand{\fib}{}\renewcommand{\fib}{\operatorname{fib}}
\providecommand{\formalizationnote}{}\renewcommand{\formalizationnote}[1]{\begin{quote}\small\raggedright\textbf{Formalization note.} #1\end{quote}}
\providecommand{\fpqc}{}\renewcommand{\fpqc}{X_{\mathrm{fpqc}}}
\providecommand{\gammacont}{}\renewcommand{\gammacont}{\Gamma_{\mathrm{cont}}}
\providecommand{\gm}{}\renewcommand{\gm}{\mathbf G_m}
\providecommand{\hocolim}{}\renewcommand{\hocolim}{\operatorname{hocolim}}
\providecommand{\holim}{}\renewcommand{\holim}{\operatorname{holim}}
\providecommand{\hrB}{}\renewcommand{\hrB}{\@ifnextchar]{\hrEnd}{\hrStep}}
\providecommand{\hrEnd}{}\renewcommand{\hrEnd}{\@firstoftwo{}}
\providecommand{\hrStep}{}\renewcommand{\hrStep}{\@firstoftwo{\hrB}}
\providecommand{\id}{}\renewcommand{\id}{\operatorname{id}}
\providecommand{\kunFl}{}\renewcommand{\kunFl}{\mathbf F_\ell}
\providecommand{\kunZn}{}\renewcommand{\kunZn}{\mathbf Z/n\mathbf Z}
\providecommand{\kuncd}{}\renewcommand{\kuncd}{\operatorname{cd}}
\providecommand{\kundim}{}\renewcommand{\kundim}{\dim}
\providecommand{\kunotimes}{}\renewcommand{\kunotimes}[1]{\otimes^{\mathbf L}_{#1}}
\providecommand{\kunsupp}{}\renewcommand{\kunsupp}{\operatorname{Supp}}
\providecommand{\lotimes}{}\renewcommand{\lotimes}{\otimes^{\mathrm L}}
\providecommand{\m}{}\renewcommand{\m}{{\mathfrak{m}}}
\providecommand{\n}{}\renewcommand{\n}{{\mathfrak{n}}}
\providecommand{\nagataGap}{}\renewcommand{\nagataGap}[1]{\par\noindent\textbf{Formalization gap.} #1\par}
\providecommand{\openobligation}{}\renewcommand{\openobligation}[1]{\par\smallskip\noindent\textcolor{red!65!black}{\textsf{Open obligation:} }
  \footnotesize\sloppy\nolinkurl{#1}\par}
\providecommand{\p}{}\renewcommand{\p}{{\mathfrak{p}}}
\providecommand{\pr}{}\renewcommand{\pr}{\operatorname{pr}}
\providecommand{\preshv}{}\renewcommand{\preshv}[1]{\mathrm{PreShv}(#1)}
\providecommand{\pro}{}\renewcommand{\pro}[1]{\mathrm{Pro}(#1)}
\providecommand{\proet}{}\renewcommand{\proet}{X_{\mathrm{proét}}}
\providecommand{\q}{}\renewcommand{\q}{{\mathfrak{q}}}
\providecommand{\rad}{}\renewcommand{\rad}{\operatorname{rad}}
\providecommand{\rank}{}\renewcommand{\rank}{\operatorname{rank}}
\providecommand{\red}{}\renewcommand{\red}{\color{red}}
\providecommand{\restorecolor}{}\renewcommand{\restorecolor}{\color{black}}
\providecommand{\retiredobligation}{}\renewcommand{\retiredobligation}[1]{\par\smallskip\noindent\textcolor{red!75!black}{\textsf{Retired invalid route:} }
  \footnotesize\sloppy\nolinkurl{#1}\par}
\providecommand{\rk}{}\renewcommand{\rk}{\operatorname{rk}}
\providecommand{\sep}{}\renewcommand{\sep}{\textnormal{sep}}
\providecommand{\shv}{}\renewcommand{\shv}[1]{\mathrm{Shv}(#1)}
\providecommand{\spec}{}\renewcommand{\spec}[1]{\mathrm{Spec}(#1)}
\providecommand{\stacks}{}\renewcommand{\stacks}[1]{\href{https://stacks.math.columbia.edu/tag/#1}{\textnormal{Stacks, Tag~#1}}}
\providecommand{\stacksproject}{}\renewcommand{\stacksproject}[1]{{\cite[\href{https://stacks.math.columbia.edu/tag/#1}{Tag #1}]{stacks-project}}}
\providecommand{\statusnote}{}\renewcommand{\statusnote}[1]{\begin{quote}\small\raggedright\textbf{Status.} #1\end{quote}}
\providecommand{\tildeideal}{}\renewcommand{\tildeideal}[2]{\widetilde{(#1, #2)}}
\providecommand{\tildering}{}\renewcommand{\tildering}[3]{{#1}_{\widetilde{(#2, #3)}}}
\providecommand{\tr}{}\renewcommand{\tr}{\operatorname{tr}}
\providecommand{\wt}{}\renewcommand{\wt}{\operatorname{wt}}

\begin{document}
\pagestyle{plain}
\begin{center}
{\LARGE\bfseries Weil Conjectures\par}
\vspace{0.7em}
{\large Part Ⅲ: Weights and the Riemann hypothesis\par}
\vspace{0.5em}
{\small\color{bpmuted}Snapshot 2026-09-26, revision 8e4cb88e2a\par}
\end{center}
\vspace{0.5em}
{\small\tableofcontents}
\bigskip\hrule\bigskip

% object: part-Weights-and-the-Riemann-hypothesis (blueprint/part-Weights-and-the-Riemann-hypothesis/_section.md)
\part*{\texorpdfstring{Part Ⅲ\quad Weights and the Riemann hypothesis}{Part Ⅲ Weights and the Riemann hypothesis}}
\phantomsection\label{bp1}%
\addcontentsline{toc}{part}{\texorpdfstring{Part Ⅲ\quad Weights and the Riemann hypothesis}{Part Ⅲ Weights and the Riemann hypothesis}}
% text: prose-P3-01
\phantomsection\label{bp2}%
After Parts I and II we can write the number of points of a smooth proper variety as an alternating sum of powers of Frobenius eigenvalues. But we know nothing about the size of these eigenvalues, so the formula alone estimates nothing. For a curve of genus \(g\) the formula reads \(N_m=1+q^m-\sum_j\alpha_j^m\). It gives Weil's bound \(|N_m-1-q^m|\le2gq^{m/2}\) only once we know that \(|\alpha_j|=q^{1/2}\). This Part proves the analogous statement in every dimension, Weil's axiom (d) (Corollary \hyperref[bp326]{19.2}): for smooth proper \(X_0\) of pure dimension \(d\), every eigenvalue \(\alpha\) of Frobenius on \(H^i(X,\mathbf Q_\ell)\) satisfies \(|\iota(\alpha)|=q^{i/2}\) for every complex embedding \(\iota\). Theorem \hyperref[bp342]{19.7} deduces from it the Riemann hypothesis, the integrality of the factors \(P_i\) and their independence of \(\ell\).

For curves Weil proved this with the Jacobian. There the eigenvalues on \(H^1\) are those of Frobenius on a Tate module, and the positivity of the Rosati involution forces \(|\alpha|=q^{1/2}\). In higher dimension, however, no such structure is available. Instead, Deligne's first proof [\hyperref[bp350]{Del74}] fibred \(X_0\) by a Lefschetz pencil and used the monodromy of the pencil. In his second paper [\hyperref[bp351]{Del80}], by contrast, he proved a theorem about all sheaves:

\begin{quote}
if \(K\) has weights at most \(w\), then so does \(Rf_!K\)
\end{quote}

(Theorem \hyperref[bp38]{13.11}; weights are defined in Definition \hyperref[bp9]{13.1}). In particular, for the constant sheaf on a smooth proper \(X_0\), the theorem gives the upper bound \(|\iota(\alpha)|\le q^{i/2}\) on \(H^i\). Since Poincaré duality pairs \(\alpha\) with \(q^d/\alpha\) on \(H^{2d-i}\), the upper bound on \(H^{2d-i}\) gives the lower bound on \(H^i\). Here we follow Laumon's proof [\hyperref[bp352]{Lau87}] of the Deligne–Laumon theorem, which replaces the Lefschetz pencil by the \(\ell\)-adic Fourier transform.

The theorem has to be about sheaves because the proof fibres the variety, locally, in curves over a variety of smaller dimension, and along such a fibration the direct image of the constant sheaf is no longer constant. Everything therefore comes down to a pure lisse sheaf \(F\) of weight \(w\) on a curve \(C\). There we must show that \(H^r_c(C,F)\) has weights at most \(w+r\) (Theorem \hyperref[bp269]{17.15}), and only \(H^1_c\) is hard.

On a curve the one tool is the trace formula. It writes the \(L\)-function, a product of local factors over the closed points, as the ratio of the characteristic polynomials on \(H^1_c\) and \(H^2_c\). The difficulty is cancellation: a large local eigenvalue gives a local factor with a pole near the origin, but the other factors might cancel it. Rankin's trick from the theory of modular forms rules this out when the traces are real, because then the even tensor powers have Euler factors with nonnegative coefficients. By Lemma \hyperref[bp113]{14.24}, no factor of such a product can have a pole closer to the origin than the product itself. On an affine curve, moreover, only \(H^2_c\) contributes poles, and their position is controlled by the determinants. These determinants are pure because a lisse sheaf of rank one has finite geometric monodromy (Section \hyperref[bp41]{14.1}). The outcome is the bound \(|\iota(\alpha)|^2\le q_x^\beta\), where \(\beta\) is the largest determinant weight (Proposition \hyperref[bp116]{14.25}). Applied to exterior powers, this bound shows that every irreducible constituent of a real sheaf is pure (Theorem \hyperref[bp120]{14.26}, Chapter \hyperref[bp39]{14}).

This bounds eigenvalues at points of a curve. What we must bound, however, is a cohomology group \(H^1_c(\mathbf A^1,j_*M)\), for a pure lisse sheaf \(M\) on a dense open \(j:U_0\hookrightarrow\mathbf A^1\). The Fourier transform (Chapter \hyperref[bp167]{16}) turns this group into the stalk of a sheaf on the dual line, and a twist of that sheaf is a quotient of a real sheaf. The Plancherel formula, together with two facts about pure sheaves on curves (Chapter \hyperref[bp122]{15}), shows that the sheaf has weight \(w+1\).

In higher dimension the same positivity shows that real sheaves on any base are mixed (Theorem \hyperref[bp295]{18.8}). This lets us apply the curve case along a fibration (Chapter \hyperref[bp271]{18}). The Weil conjectures are then a short deduction (Chapter \hyperref[bp319]{19}).

% object: chap-Weights (blueprint/chap-Weights/_section.md), content hidden
\section*{\texorpdfstring{13\quad Weights}{13 Weights}}
\phantomsection\label{bp3}%
\addcontentsline{toc}{section}{\texorpdfstring{13\quad Weights}{13 Weights}}
% text: prose-13-01
\phantomsection\label{bp4}%
For the elliptic curve \(y^2=x^3-x\) over \(\mathbf F_3\) of the introduction, Frobenius acts on \(H^1\) with eigenvalues \(\pm i\sqrt3\). The Riemann hypothesis for this curve is then the statement \(|\pm i\sqrt3|=3^{1/2}\). But the eigenvalues live in \(\overline{\mathbf Q}_\ell\), and an \(\ell\)-adic number has no complex absolute value until we choose a field embedding \(\iota:\overline{\mathbf Q}_\ell\hookrightarrow\mathbf C\). Since no choice is preferred, we fix one \(\iota\) throughout, and the final theorems hold for every choice.

An eigenvalue \(\alpha\) of Frobenius at a closed point of degree \(n\) has weight \(w\) if \(|\iota(\alpha)|=q^{nw/2}\). The factor \(n\) makes the weight independent of the field over which we look at the point: over \(\mathbf F_{q^{mn}}\) the eigenvalue becomes \(\alpha^m\) and \(q^n\) becomes \(q^{mn}\). For example, on \(\mathbf P^1\) Frobenius acts on \(H^0\) by \(1\), of weight \(0\), and on \(H^2=\overline{\mathbf Q}_\ell(-1)\) by \(q\), of weight \(2\). Similarly, the eigenvalues on \(H^1\) of our elliptic curve have weight \(1\). Weil's purity axiom says that this pattern persists: \(H^i\) of a smooth proper variety is pure of weight \(i\). The name comes from Hodge theory, where \(H^i\) of a compact Kähler manifold carries a Hodge structure of weight \(i\).

The proof cannot stay with smooth proper varieties. It fibres them in curves, and the cohomology of the fibres forms sheaves on the base. So we need weights for sheaves, and for cohomology that is not pure. Without smoothness and properness only an upper bound survives. For example, let \(\mathbf G_m=\mathbf P^1\setminus\{0,\infty\}\). The localization sequence gives \(H^1_c(\mathbf G_m,\overline{\mathbf Q}_\ell)=\overline{\mathbf Q}_\ell\) with trivial Frobenius and \(H^2_c(\mathbf G_m,\overline{\mathbf Q}_\ell)=\overline{\mathbf Q}_\ell(-1)\). This agrees with the trace formula \(q-1=\#\mathbf G_m(\mathbf F_q)\). Thus \(H^1_c\) has weight \(0\), strictly less than \(1\).

A sheaf is pure of weight \(w\) if all its eigenvalues at all closed points have weight \(w\). It is mixed if it has a finite filtration whose quotients are pure. A bounded complex \(K\) has weights at most \(w\) if each \(\mathcal H^i(K)\) is mixed with all weights at most \(w+i\) (Definition \hyperref[bp9]{13.1}). Here the shift by \(i\) puts the constant sheaf and its cohomology into one statement: \(R\Gamma_c(X,\overline{\mathbf Q}_\ell)\) has weights at most \(0\) exactly when \(H^i_c\) has weights at most \(i\), as on \(\mathbf G_m\). In these terms the Deligne–Laumon theorem (Theorem \hyperref[bp38]{13.11}) says that \(Rf_!\) does not increase weights.

Two warnings will matter later. The first is that weights depend on \(\iota\). For example, consider the rank-one Weil sheaf on \(\operatorname{Spec}\mathbf F_q\) on which Frobenius acts by \(1+\sqrt2\in\overline{\mathbf Q}_\ell\). It has weight \(2\log_q(1+\sqrt2)>0\) for an embedding with \(\iota(\sqrt2)=\sqrt2\), and weight \(2\log_q(\sqrt2-1)=-2\log_q(1+\sqrt2)\) for the other one. For this reason Theorem \hyperref[bp342]{19.7} needs the Deligne–Laumon theorem for every \(\iota\): only then do the reciprocal roots of \(P_i\) have absolute value \(q^{i/2}\) under every complex embedding.

The second warning is that being mixed is not a pointwise condition: it asks for a filtration with pure quotients, and such a filtration has to be constructed. Most of Chapter \hyperref[bp271]{18} is spent constructing one for direct images.

We define weights on Weil sheaves, sheaves on \(X\) with a Frobenius structure, rather than on sheaves on \(X_0\). The reason is that normalizing weights requires twisting by the geometrically constant rank-one Weil sheaves \(\mathcal B_b\), on which a degree-one Frobenius acts by a scalar \(b\in\overline{\mathbf Q}_\ell^\times\). These do not come from \(X_0\) when \(b\) is not an \(\ell\)-adic unit.

% object: sec-Weil-sheaves-and-weights (blueprint/chap-Weights/sec-Weil-sheaves-and-weights/_section.md), content hidden
\subsection*{\texorpdfstring{13.1\quad Weil sheaves and weights}{13.1 Weil sheaves and weights}}
\phantomsection\label{bp5}%
\addcontentsline{toc}{subsection}{\texorpdfstring{13.1\quad Weil sheaves and weights}{13.1 Weil sheaves and weights}}
% text: prose-13-02
\phantomsection\label{bp6}%
Weights are defined by the absolute values of the Frobenius eigenvalues on stalks at closed points, so we first fix how Frobenius acts on sheaves.

On a smooth curve \(U_0\) we want to treat lisse Weil sheaves as representations, just as a lisse sheaf on \(U_0\) is a continuous representation of \(\pi_1(U_0)\). However, the profinite group \(\pi_1(U_0)\) is an extension of \(\widehat{\mathbf Z}\) by \(\pi_1(U)\), and it has too few representations for this. Indeed, a continuous representation of a profinite group preserves a lattice, so a lift \(s\) of Frobenius acts with eigenvalues that are \(\ell\)-adic units. Hence the rank-one Weil sheaf \(\mathcal B_\ell\), on which \(s\) acts by \(\ell\) and \(\pi_1(U)\) acts trivially, is not a representation of \(\pi_1(U_0)\). Concretely, \(n\mapsto\ell^n\) does not extend continuously to \(\widehat{\mathbf Z}\): the integers \(m!\) tend to \(0\) in \(\widehat{\mathbf Z}\), whereas \(\ell^{m!}\) tends to \(0\) in \(\mathbf Z_\ell\) and not to \(1\).

The Weil group keeps \(\pi_1(U)\) with its topology and replaces \(\widehat{\mathbf Z}\) by the discrete group \(\mathbf Z\) of integer powers of Frobenius. Continuity then constrains only the geometric part. These representations of the Weil group are the language of Theorem \hyperref[bp58]{14.5} and of the monodromy groups of Section \hyperref[bp69]{14.2}.

% text: situation-weil
\phantomsection\label{bp7}%
\textbf{Convention.} Put \(k_0=\mathbf F_q\), \(k=\overline{\mathbf F}_q\), and fix a prime \(\ell\nmid q\) and a field embedding \(\iota:\overline{\mathbf Q}_\ell\hookrightarrow\mathbf C\). For \(X_0/k_0\) separated of finite type, we write \(X=X_0\times_{k_0}k\), and for a closed point \(x\) of \(X_0\) write \(\deg(x)=[\kappa(x):k_0]\) for its degree and \(q_x=q^{\deg(x)}=\#\kappa(x)\). Moreover, geometric Frobenius, Weil complexes and their Frobenius operators are those of Section \bpout{10.5}. As there, a Weil sheaf is a Weil complex concentrated in degree zero, and the six operations carry Weil structures to Weil structures. In this Part Weil sheaves and complexes have coefficients in \(\overline{\mathbf Q}_\ell\); a Weil complex \(K\) with coefficients in a finite extension \(E\subset\overline{\mathbf Q}_\ell\) of \(\mathbf Q_\ell\) is called pure, mixed, or of weights at most \(w\) if \(K\otimes_E\overline{\mathbf Q}_\ell\) is (see Lemma \hyperref[bp11]{13.2} and the remark after it for why scalars must be extended).

On a connected normal scheme of finite type over a field, we identify a lisse \(E\)-sheaf with the corresponding continuous representation of the profinite fundamental group, by Proposition \hyperref[bp34]{13.10}, proved in Section \hyperref[bp26]{13.2} below. This covers all arguments with representations below, since they take place on smooth schemes. Moreover, a lisse Weil \(\overline{\mathbf Q}_\ell\)-sheaf has a lisse Weil model over some finite extension \(E/\mathbf Q_\ell\), and two models become identified after a finite extension of scalars. This model includes the Weil isomorphism, which need not preserve a geometric \(\mathcal O_E\)-lattice.

For a geometrically connected smooth curve \(U_0\), put \(U=U_0\times_{k_0}k\) and choose a geometric point \(\bar\eta\) of \(U\). We then write \(\pi_1(U)\) and \(\pi_1(U_0)\) for \(\pi_1(U,\bar\eta)\) and \(\pi_1(U_0,\bar\eta)\).  The inverse image of the integer powers of geometric Frobenius in \(\pi_1(U_0)\) is its Weil group \(W(U_0)\).  It contains the geometric fundamental group, with an exact sequence

\[
1\longrightarrow \pi_1(U)\longrightarrow W(U_0)
\xrightarrow{\deg}\mathbf Z\longrightarrow1.
\]

A degree-one lift \(s\) splits this sequence as a group extension. In these terms, lisse Weil sheaves are the same as representations \(\rho\) of \(W(U_0)\), defined over a finite extension of \(\mathbf Q_\ell\), whose restriction to \(\pi_1(U)\) is continuous. For \(x\in|U_0|\), the local Frobenius \(\operatorname{Frob}_x\) at \(x\) (Theorem \bpout{12.30}) is a conjugacy class of degree \(\deg(x)\).

% text: prose-13-03
\phantomsection\label{bp8}%
Every sheaf defined over \(k_0\) is a Weil sheaf. By means of the Weil isomorphism, Frobenius acts on the stalk at a geometric point over each closed point \(x\), so we can speak of the eigenvalues of \(\operatorname{Frob}_x\). In the description by representations the Weil isomorphism is transport by \(s\), and conversely it determines the action of \(s\).

The definition below contains two kinds of bounds. The first, being punctually of weights \(\le w\), bounds the eigenvalues at each closed point; it is what the curve case yields on each stalk of a direct image. But it cannot serve as input for the curve case, which needs pure lisse sheaves. The second, mixedness, has the opposite virtue: a mixed sheaf can be treated one pure quotient at a time, and for a mixed sheaf the bound on weights is again a pointwise condition (Lemma \hyperref[bp14]{13.3}). In Chapter \hyperref[bp271]{18} we use this to turn the stalkwise bounds into the bound of Theorem \hyperref[bp38]{13.11}.

% object: definition-weights (blueprint/chap-Weights/sec-Weil-sheaves-and-weights/definition-weights/_section.md)
\phantomsection\label{bp9}%
\begin{bpstatement}{Definition~13.1}{ (Weights, purity and mixedness)}
Let \(K\) be a constructible Weil sheaf on \(X\), let \(x\) be a closed point of \(X_0\), and let \(\alpha\) be an eigenvalue of \(\operatorname{Frob}_x\) on \(K_{\bar x}\). We say \(\alpha\) has \emph{\(\iota\)-weight \(w\)} if \(|\iota(\alpha)|=q_x^{w/2}\).

\begin{enumerate}
\item \(K\) is \emph{punctually \(\iota\)-pure of weight \(w\)}, or simply \emph{pure of weight \(w\)}, if every such eigenvalue, at every closed point, has \(\iota\)-weight \(w\).
\item \(K\) is \emph{punctually of weights \(\le w\)} if \(|\iota(\alpha)|\le q_x^{w/2}\) for every such eigenvalue at every closed point.
\item \(K\) is \emph{\(\iota\)-mixed} if it has a finite filtration by Weil subsheaves whose nonzero quotients are pure, of possibly different real weights. It is \emph{mixed of weights at most \(w\)}, or has \emph{weights at most \(w\)}, if it has such a filtration whose nonzero quotients have weights at most \(w\).
\end{enumerate}

A bounded constructible Weil complex \(K\) has \emph{weights at most \(w\)} if \(\mathcal H^i(K)\) has weights at most \(w+i\) for all \(i\), and it is \emph{pure of weight \(w\)} if each nonzero \(\mathcal H^i(K)\) is pure of weight \(w+i\).
\end{bpstatement}

% text: prose-13-04
\phantomsection\label{bp10}%
Deligne calls a complex \(K\) pure of weight \(w\) if \(K\) has weights at most \(w\) and its Verdier dual \(\mathbb D K\) has weights at most \(-w\). Over a point the two notions agree. In general neither implies the other, and we use ours only for sheaves and over a point.

The cohomology groups of Theorem \hyperref[bp38]{13.11} are Weil sheaves on \(\operatorname{Spec}k_0\). For them, having weights at most \(w\) is just a condition on the eigenvalues of Frobenius. Indeed, by Lemma \hyperref[bp11]{13.2} every Weil \(\overline{\mathbf Q}_\ell\)-sheaf over a point is mixed, and Lemma \hyperref[bp14]{13.3} reads the bound of a mixed sheaf off its eigenvalues.

% object: lemma-weil-sheaves-point-mixed (blueprint/chap-Weights/sec-Weil-sheaves-and-weights/lemma-weil-sheaves-point-mixed/_section.md)
\phantomsection\label{bp11}%
\begin{bpstatementit}{Lemma~13.2}{ (Weil sheaves on a point are mixed)}
Every constructible Weil \(\overline{\mathbf Q}_\ell\)-sheaf on \(\operatorname{Spec}k_0\) is mixed. A constructible Weil \(E\)-sheaf on \(\operatorname{Spec}k_0\) becomes mixed after extending scalars to a finite extension of \(E\) containing the eigenvalues of Frobenius.
\end{bpstatementit}

% proof: proof/lemma-weil-sheaves-point-mixed
\begin{proof}[{Proof}]\phantomsection\label{bp12}%
A constructible Weil sheaf on \(\operatorname{Spec}k_0\) is a finite-dimensional vector space \(V\) with an automorphism \(F\), the geometric Frobenius. In both cases we may assume that \(V\) is defined over a finite extension \(E'\) of \(\mathbf Q_\ell\) containing the eigenvalues of \(F\). For \(w\in\mathbf R\) let \(V_w\subset V\) be the sum of the generalized eigenspaces of \(F\) for the eigenvalues \(\alpha\) with \(|\iota(\alpha)|=q^{w/2}\). These subspaces are \(F\)-stable, only finitely many are nonzero, and \(V=\bigoplus_wV_w\). Let \(w_1<\dots<w_m\) be the weights that occur. Then the subspaces \(V_{w_1}\oplus\dots\oplus V_{w_j}\) form a finite filtration by Weil subsheaves with quotients \(V_{w_j}\), and \(V_{w_j}\) is pure of weight \(w_j\).
\end{proof}

% text: prose-13-05
\phantomsection\label{bp13}%
Extending scalars is needed for \(E\)-sheaves. Suppose \(2\) is not a square in \(\mathbf Q_\ell\). Then the rank-two \(\mathbf Q_\ell\)-sheaf on which Frobenius has eigenvalues \(1\pm\sqrt2\) is irreducible. Its two eigenvalues have different weights, so it has no filtration with pure quotients over \(\mathbf Q_\ell\).

% object: lemma-bounding-weights-mixed-sheaf (blueprint/chap-Weights/sec-Weil-sheaves-and-weights/lemma-bounding-weights-mixed-sheaf/_section.md)
\phantomsection\label{bp14}%
\begin{bpstatementit}{Lemma~13.3}{ (Bounding the weights of a mixed sheaf)}
A mixed sheaf has weights at most \(w\) if and only if it is punctually of weights \(\le w\).
\end{bpstatementit}

% proof: proof/lemma-bounding-weights-mixed-sheaf
\begin{proof}[{Proof}]\phantomsection\label{bp15}%
At every closed point \(x\) the eigenvalues of \(\operatorname{Frob}_x\) on the stalk of a mixed sheaf are those on the stalks of the quotients of its filtration. If the quotients have weights at most \(w\), the sheaf is therefore punctually of weights \(\le w\). Conversely, each nonzero quotient has a nonzero stalk at some closed point \(x\), and the eigenvalues of \(\operatorname{Frob}_x\) there are eigenvalues on the stalk of the sheaf. So if the sheaf is punctually of weights \(\le w\), the weight of the quotient is at most \(w\).
\end{proof}

% text: prose-13-06
\phantomsection\label{bp16}%
The lemma does not say that a sheaf punctually of weights \(\le w\) is mixed; the filtration has to be found. Nor need the filtration split: weights see only the characteristic polynomials of Frobenius, so a Jordan block and its semisimplification have the same weights. The following example is a mixed sheaf that is not pure and whose filtration does not split.

% object: example-kummer-extension (blueprint/chap-Weights/sec-Weil-sheaves-and-weights/example-kummer-extension/_section.md)
\phantomsection\label{bp17}%
\begin{bpstatement}{Example~13.4}{ (The Kummer extension)}
On \(\mathbf G_{m,k_0}\) with coordinate \(t\), the \(\mu_{\ell^n}\)-torsors \(u^{\ell^n}=t\) give the Kummer cocycle \(\kappa\), a continuous map from \(\pi_1(\mathbf G_{m,k_0})\) to \(\mathbf Z_\ell(1)\). Let \(\chi\) be the cyclotomic character, by which \(\pi_1\) acts on \(\mathbf Z_\ell(1)\). The matrices

\[
\sigma\longmapsto\begin{pmatrix}\chi(\sigma)&\kappa(\sigma)\\0&1\end{pmatrix}
\]

define a lisse sheaf \(\mathcal E\) on \(\mathbf G_{m,k_0}\), an extension \(0\to\overline{\mathbf Q}_\ell(1)\to\mathcal E\to\overline{\mathbf Q}_\ell\to0\). Geometric Frobenius acts on \(\overline{\mathbf Q}_\ell(1)\) by \(q_x^{-1}\). Hence at a closed point \(x\) the eigenvalues of \(\operatorname{Frob}_x\) are \(q_x^{-1}\) and \(1\), of weights \(-2\) and \(0\), and \(\mathcal E\) is mixed of weights at most \(0\) and not pure. The filtration does not split. Indeed, on \(\pi_1(\mathbf G_{m,k})\) the character \(\chi\) is trivial, and \(\kappa\) is a homomorphism onto \(\mathbf Z_\ell(1)\) because the covers \(u^{\ell^n}=t\) of \(\mathbf G_{m,k}\) are connected. So the line \(\overline{\mathbf Q}_\ell(1)\) is the only stable line. At each closed point, by contrast, the two eigenvalues are distinct and the stalk is the direct sum of its eigenlines.
\end{bpstatement}

% text: prose-13-07
\phantomsection\label{bp18}%
The behaviour of weights under the usual operations comes down to the eigenvalues of an operator on a filtered space, on duals and on tensor and exterior powers.

% object: lemma-filtrations-spectra (blueprint/chap-Weights/sec-Weil-sheaves-and-weights/lemma-filtrations-spectra/_section.md)
\phantomsection\label{bp19}%
\begin{bpstatementit}{Lemma~13.5}{ (Filtrations and spectra)}
Let \(V\) have a finite filtration preserved by an automorphism \(u\). Then

\[
\det(1-Tu\mid V)=\prod_a\det(1-Tu\mid\operatorname{gr}_aV).
\]

In particular, the eigenvalues, with algebraic multiplicity, are the union of the eigenvalues on the graded pieces. If the eigenvalues of \(u\) on \(V\) are \(\alpha_1,\ldots,\alpha_d\), then those on \(V^\vee\) are \(\alpha_i^{-1}\), and those on \(\bigwedge^rV\) are \(\alpha_{i_1}\cdots\alpha_{i_r}\) for \(i_1<\cdots<i_r\). If \(v\) on \(V'\) has eigenvalues \(\beta_j\), those of \(u\otimes v\) are \(\alpha_i\beta_j\).
\end{bpstatementit}

% proof: proof/lemma-filtrations-spectra
\begin{proof}[{Proof}]\phantomsection\label{bp20}%
In a basis adapted to the filtration, the matrix of \(u\) is block upper triangular, so its characteristic polynomial is the product of those of the diagonal blocks. For the other constructions we may choose bases over \(\overline{\mathbf Q}_\ell\) in which \(u\) and \(v\) are upper triangular. Then, in the induced tensor basis and in the ordered wedge basis, the resulting matrices are upper triangular with the stated diagonal entries. Finally, the matrix on the dual is the inverse transpose.
\end{proof}

% text: prose-13-08
\phantomsection\label{bp21}%
At a closed point, each operation in the next lemma and proposition acts on stalks by one of the constructions of the lemma: a subobject, a quotient, a tensor product or a dual, or, for an inverse image, a power of Frobenius. So the behaviour of weights is a matter of eigenvalues. Mixedness needs in addition some bookkeeping to produce the filtrations.

% object: lemma-subobjects-quotients-weil-sheaves (blueprint/chap-Weights/sec-Weil-sheaves-and-weights/lemma-subobjects-quotients-weil-sheaves/_section.md)
\phantomsection\label{bp22}%
\begin{bpstatementit}{Lemma~13.6}{ (Subobjects and quotients of Weil sheaves)}
Subobjects and quotients of a pure Weil sheaf of weight \(w\) are pure of weight \(w\). Subquotients of a mixed sheaf are mixed, and a subquotient of a sheaf that is mixed of weights at most \(w\) is mixed of weights at most \(w\).
\end{bpstatementit}

% proof: proof/lemma-subobjects-quotients-weil-sheaves
\begin{proof}[{Proof}]\phantomsection\label{bp23}%
By Lemma \hyperref[bp19]{13.5}, the eigenvalues on a subobject or a quotient occur among those of the whole. Hence subobjects and quotients of pure sheaves are pure. For mixed sheaves, we intersect a mixed filtration with a subobject, or push it into a quotient. The graded pieces are then subquotients of the original pure pieces, hence pure of the same weights, and so the upper bound persists.
\end{proof}

% object: proposition-weights-under-usual-operations (blueprint/chap-Weights/sec-Weil-sheaves-and-weights/proposition-weights-under-usual-operations/_section.md)
\phantomsection\label{bp24}%
\begin{bpstatementit}{Proposition~13.7}{ (Weights under the usual operations)}
Extensions of mixed sheaves are mixed, and an extension of two sheaves with a common upper bound has the same bound.  Tensor products of pure sheaves of weights \(a\) and \(b\) are pure of weight \(a+b\), and the dual of a pure lisse sheaf of weight \(a\) is pure of weight \(-a\). A tensor product of mixed sheaves of weights at most \(a\) and at most \(b\) is mixed of weights at most \(a+b\). The dual of a mixed lisse sheaf is mixed, of the negated weights; when its filtration has lisse graded pieces, their duals are the graded pieces of the annihilator filtration of the dual.  Inverse image preserves purity, mixedness and their upper bounds.  If a complex \(K\) has weights at most \(w\), then \(K[r]\) has weights at most \(w+r\).

For \(b\in\overline{\mathbf Q}_\ell^\times\), let \(\mathcal B_b\) be the geometrically constant rank-one Weil sheaf on which a degree-one Frobenius acts by \(b\). Tensoring by \(\mathcal B_b\) adds \(\lambda_\iota(b)=2\log_q|\iota(b)|\) to every weight.  The Tate sheaf \(\overline{\mathbf Q}_\ell(1)\) has geometric Frobenius \(q^{-1}\), so twisting by \((r)\) subtracts \(2r\) from every weight.
\end{bpstatementit}

% proof: proof/proposition-weights-under-usual-operations
\begin{proof}[{Proof}]\phantomsection\label{bp25}%
Each assertion has two parts: a statement about eigenvalues at closed points, which follows from Lemma \hyperref[bp19]{13.5}, and, for mixedness, the construction of a filtration with pure quotients.

An extension is filtered by a filtration of its subobject followed by the inverse images of a filtration of its quotient. The graded pieces of the result are those of the two given filtrations, so a common upper bound persists.

For tensor products we use that the coefficients form a field, so that the associated graded of the tensor product of two filtrations is the tensor product of the two associated gradeds. By Lemma \hyperref[bp19]{13.5}, the eigenvalues of tensor products and duals are the products \(\alpha_i\beta_j\) and the inverses \(\alpha_i^{-1}\), whose weights are \(a+b\) and \(-a\).

For the dual of a mixed lisse sheaf \(G\) with mixed filtration \((G_i)\), suppose first that the graded pieces \(G_i/G_{i-1}\) are lisse. Dualizing is exact on lisse sheaves, so the annihilator filtration \(G_i^\perp=\ker(G^\vee\to G_i^\vee)\) has graded pieces \((G_i/G_{i-1})^\vee\), which are pure of the negated weights by the computation above. In general the graded pieces need not be lisse: on a curve, a pure lisse \(G\) has the filtration \(j_!j^*G\subseteq G\), whose quotient \(i_*i^*G\) is a skyscraper with dual zero. We then argue by Noetherian induction on \(X\). Let \(U\subseteq X\) be a dense open, stable under Frobenius, on which all graded pieces are lisse, and let \(i:Z\to X\) be its closed complement. On \(U\) the first case applies to \(j^*G\). Since \(G\) is lisse, \(i^*(G^\vee)\simeq(i^*G)^\vee\), and \(i^*G\) is a mixed lisse sheaf on \(Z\), so \(i^*(G^\vee)\) is mixed by induction, of the negated weights. Extension by zero and direct image along a closed immersion are exact and do not change stalks, so they carry mixed filtrations to mixed filtrations of the same weights, and the exact sequence

\[
0\to j_!j^*(G^\vee)\to G^\vee\to i_*i^*(G^\vee)\to 0
\]

shows that \(G^\vee\) is mixed, of the negated weights.

Inverse images do not change weights, because the weight is normalized by the degree. Indeed, at a closed point of an inverse image the Frobenius is a positive power of the Frobenius at the image point, and the residue field degree is multiplied by the same integer. For shifts, \(\mathcal H^i(K[r])=\mathcal H^{i+r}(K)\) is mixed of weights at most \(w+i+r\), which is the bound \(w+r\) for \(K[r]\).

The local eigenvalue of \(\mathcal B_b\) at \(x\) is \(b^{\deg(x)}\), of weight \(2\log_q|\iota(b)|\), and tensoring by \(\mathcal B_b\) multiplies every eigenvalue by it. For the Tate sheaf, recall that arithmetic Frobenius acts on roots of unity by raising to the \(q\)th power. Geometric Frobenius is its inverse, so it acts on \(\overline{\mathbf Q}_\ell(1)\) by \(q^{-1}\), of weight \(-2\).
\end{proof}


% object: sec-Lisse-sheaves-as-representations (blueprint/chap-Weights/sec-Lisse-sheaves-as-representations/_section.md), content hidden
\subsection*{\texorpdfstring{13.2\quad Lisse sheaves as representations}{13.2 Lisse sheaves as representations}}
\phantomsection\label{bp26}%
\addcontentsline{toc}{subsection}{\texorpdfstring{13.2\quad Lisse sheaves as representations}{13.2 Lisse sheaves as representations}}
% text: prose-13-09
\phantomsection\label{bp27}%
Let \(X\) be a connected normal scheme of finite type over a field and \(L\) a lisse \(E_X\)-module. Lemma \bpout{3.24} gives lattices in \(L\) only on the strata of a finite stratification, and the nodal curve of Section \bpout{3.8} shows that without normality a global lattice need not exist. On a normal \(X\), however, we find a lattice on all of \(X\), and as a consequence lisse sheaves are the same as continuous representations of the fundamental group. Here normality is used to show that the étale \(X\)-scheme representing lattices in \(L\) has a component that is finite étale and surjective. Once we have such a component, we argue as in Lemma \bpout{3.24}.

% object: lemma-geometric-coverings-normal-scheme (blueprint/chap-Weights/sec-Lisse-sheaves-as-representations/lemma-geometric-coverings-normal-scheme/_section.md)
\phantomsection\label{bp28}%
\begin{bpstatementit}{Lemma~13.8}{ (Geometric coverings of a normal scheme)}
Let \(X\) be a nonempty connected normal scheme of finite type over a field. Let \(p:Y\to X\) be separated and étale. Suppose that for every valuation ring \(V\), with fraction field \(K\), every commutative diagram

\[
\begin{CD}
\operatorname{Spec}K @>>> Y\\
@VVV @VVpV\\
\operatorname{Spec}V @>>> X
\end{CD}
\]

has a lift \(\operatorname{Spec}V\to Y\). Then every connected component of \(Y\) is finite étale and surjective over \(X\).
\end{bpstatementit}

% proof: proof/lemma-geometric-coverings-normal-scheme
\begin{proof}[{Proof}]\phantomsection\label{bp29}%
We compare a component \(Y_1\) with the normalization \(N\) of \(X\) in the function field of \(Y_1\). By Zariski's main theorem, \(Y_1\) is an open subscheme of \(N\), and by the valuative hypothesis this open subscheme contains every point of \(N\).

Since \(Y\) is étale over \(X\), it is locally noetherian and normal [Stacks, \href{https://stacks.math.columbia.edu/tag/033C}{033C}]. Hence its connected components are integral, open and closed [Stacks, \href{https://stacks.math.columbia.edu/tag/033N}{033N}]. Fix a nonempty component \(Y_1\). Since a valuative lift whose generic point lies in \(Y_1\) stays in \(Y_1\), the hypothesis also holds for \(Y_1\to X\).

The image of \(Y_1\) is nonempty and open, hence contains the generic point \(\eta\) of \(X\). The generic fibre of \(Y_1\) is integral and étale over \(k(X)\), so it is \(\operatorname{Spec}k(Y_1)\). Thus the function field \(k(Y_1)\) is a finite separable extension of \(k(X)\), and we let \(N\) be the normalization of \(X\) in \(k(Y_1)\). Since \(X\) is Nagata, \(N\) is finite over \(X\) [Stacks, \href{https://stacks.math.columbia.edu/tag/035B}{035B}, \href{https://stacks.math.columbia.edu/tag/0AVK}{0AVK}]. We claim that \(Y_1\to X\) factors through a map \(f:Y_1\to N\). Indeed, since \(Y_1\) is normal, an element of \(k(Y_1)\) that is integral over the coordinate ring of an affine open of \(X\) lies in every local ring of \(Y_1\) over that open. These regular functions agree in \(k(Y_1)\), and hence they glue. The map \(f\) is separated, locally of finite type and locally quasi-finite, because its fibres lie in the zero-dimensional fibres of \(Y_1\to X\).

We show that \(f\) is an open immersion. Choose affine opens \(W\subseteq Y_1\) and \(N_0\subseteq N\) with \(f(W)\subseteq N_0\). By Zariski's main theorem [Stacks, \href{https://stacks.math.columbia.edu/tag/05K0}{05K0}], \(W\) is an open subscheme of a finite \(N_0\)-scheme, and the reduced closure of \(W\) in it is integral, finite and birational over \(N_0\), with function field \(k(Y_1)\). Since \(N_0\) is normal, this closure is \(N_0\) [Stacks, \href{https://stacks.math.columbia.edu/tag/0AB1}{0AB1}], and therefore \(W\to N_0\) is an open immersion. These local inverses glue: on the intersection of two such open images, the inverse sections agree generically, hence everywhere by separatedness, and so they glue to an inverse on the union of the images.

Finally, let \(n\in N\) and choose a valuation ring \(V\subseteq k(Y_1)\) dominating \(\mathcal O_{N,n}\) [Stacks, \href{https://stacks.math.columbia.edu/tag/00IA}{00IA}]. By hypothesis, the generic map \(\operatorname{Spec}k(Y_1)\to Y_1\) extends to a map \(\operatorname{Spec}V\to Y_1\), whose composite with \(f\) agrees with the canonical map \(\operatorname{Spec}V\to N\) generically and over \(X\). By separatedness of \(N\to X\) the two maps agree, so \(n\) lies in the open image of \(Y_1\). We conclude that \(Y_1=N\) is finite étale over \(X\). In particular its image is open and closed, and being nonempty, it is all of \(X\).
\end{proof}

% text: prose-13-10
\phantomsection\label{bp30}%
We apply the lemma to the scheme of lattices in a lisse sheaf. There the valuative hypothesis says that a lattice defined at the generic point of a valuation ring extends over the whole ring. This holds because, locally, lattices form a constant sheaf on a discrete set.

% object: lemma-global-lattice-normal-scheme (blueprint/chap-Weights/sec-Lisse-sheaves-as-representations/lemma-global-lattice-normal-scheme/_section.md)
\phantomsection\label{bp31}%
\begin{bpstatementit}{Lemma~13.9}{ (Global lattices on a normal scheme)}
Let \(X\) be a nonempty connected normal scheme of finite type over a field, and let \(E/\mathbf Q_\ell\) be finite. Every lisse \(E_X\)-module \(L\), in the pro-étale coefficient category of Section \bpout{2.12}, contains a lisse \(\mathcal O_{E,X}\)-submodule \(M\) with \(M[1/\ell]=L\). The reductions \(M_n=M/\pi_E^nM\) are finite locally constant étale sheaves, and \(M\simeq\varprojlim_nM_n\).
\end{bpstatementit}

% proof: proof/lemma-global-lattice-normal-scheme
\begin{proof}[{Proof}]\phantomsection\label{bp32}%
The scheme of lattices satisfies the valuative condition of Lemma \hyperref[bp28]{13.8}, so it has a finite étale surjective component, and we push the universal lattice of this component down to \(X\) by the trace.

We may assume that \(L\) has positive rank \(r\). Let \(Y\) be the \(X\)-scheme representing lattices in \(L\) (Lemma \bpout{3.15}), a separated surjective étale \(X\)-scheme with a universal lattice. Pro-étale locally on \(X\), the scheme \(Y\) is constant: the proof of that lemma shows that, for a quasi-compact pro-étale covering \(U\to X\), with \(U\) a pro-étale affine on which \(L\) is free,

\[
Y\times_XU\simeq
\coprod_{\GL_r(E)/\GL_r(\mathcal O_E)} U,
\]

where \(\GL_r(E)/\GL_r(\mathcal O_E)\) is the discrete set of lattices in \(E^r\).

We check the valuative property of Lemma \hyperref[bp28]{13.8}. Let \(V\) be a valuation ring with fraction field \(K\), and suppose given \(\operatorname{Spec}V\to X\) and a lift \(\xi:\operatorname{Spec}K\to Y\). Since \(U\to X\) is flat and surjective, there are a local extension of valuation rings \(V\to V'\) and a compatible map \(\operatorname{Spec}V'\to U\) [Stacks, \href{https://stacks.math.columbia.edu/tag/089F}{089F}]. This extension is faithfully flat, because it is local and \(V'\) is torsion-free, hence flat, over \(V\) [Stacks, \href{https://stacks.math.columbia.edu/tag/0539}{0539}]. By the displayed description of \(Y\times_XU\), the pullback of \(\xi\) to the spectrum of the fraction field of \(V'\) picks out a lattice in \(E^r\), and the constant section with this value extends it to a map \(\xi':\operatorname{Spec}V'\to Y\).

We descend \(\xi'\) by fpqc descent, which needs its two pullbacks to \(R=V'\otimes_VV'\) to agree. Over the domain \(V'\otimes_VK\), the maps \(\xi'\) and \(\xi\) agree at the generic point, hence everywhere by separatedness, so the two pullbacks agree after base change to \(K\). This suffices: since \(R\) is flat over \(V\), the map \(R\to R\otimes_VK\) is injective, and therefore \(\operatorname{Spec}(R\otimes_VK)\) is schematically dense in \(\operatorname{Spec}R\). As the equalizer of the two pullbacks is closed by separatedness, they agree on \(\operatorname{Spec}R\). Descent then gives a lift \(\operatorname{Spec}V\to Y\) restricting to \(\xi\), which is unique by separatedness and the density of \(\operatorname{Spec}K\).

By Lemma \hyperref[bp28]{13.8}, every nonempty component \(p_1:Y_1\to X\) of \(Y\) is therefore finite étale and surjective. We fix one, and let \(N\subseteq p_1^*L\) be its universal lattice. As in Lemma \bpout{3.24}, set

\[
M=\operatorname{im}\bigl(
p_{1*}N\longrightarrow p_{1*}p_1^*L
\xrightarrow{\operatorname{Tr}}L\bigr).
\]

The sheaf \(M\) is lisse with \(M[1/\ell]=L\). Indeed, after the cover is split into \(d\) sheets, this image is the sum of \(d\) lattices. On a common pro-étale refinement these become constant lattices in \(E^r\), and their sum is a finite free \(\mathcal O_E\)-lattice spanning \(E^r\).

By Theorem \bpout{2.43} the finite reductions \(M_n=M/\pi_E^nM\) are finite locally constant étale sheaves, and \(M=\varprojlim_nM_n\) by Lemma \bpout{2.51}, which we may check where \(M\) is free.
\end{proof}

% text: prose-13-11
\phantomsection\label{bp33}%
Galois theory of finite étale covers turns statements about the finite reductions of such an \(M\) into statements about representations.

% object: proposition-lattices-representations-normal-scheme (blueprint/chap-Weights/sec-Lisse-sheaves-as-representations/proposition-lattices-representations-normal-scheme/_section.md)
\phantomsection\label{bp34}%
\begin{bpstatementit}{Proposition~13.10}{ (Lattices and representations on a normal scheme)}
Let \(X\) be a nonempty connected normal scheme of finite type over a field, let \(E/\mathbf Q_\ell\) be finite, and fix a geometric point \(\bar x\). By Lemma \hyperref[bp31]{13.9}, every lisse \(E_X\)-module \(L\) contains a lisse \(\mathcal O_{E,X}\)-submodule \(M\) with \(M[1/\ell]=L\) and \(M\simeq\varprojlim_nM/\pi_E^nM\), whose reductions are finite locally constant étale sheaves. The fibre functor at \(\bar x\), \(L\mapsto\bigl(\varprojlim_n(M/\pi_E^nM)_{\bar x}\bigr)\otimes_{\mathcal O_E}E\), does not depend on the choice of \(M\) up to canonical isomorphism, and it gives an equivalence between lisse \(E_X\)-modules and finite-dimensional continuous \(E\)-representations of the ordinary profinite group \(\pi_1(X,\bar x)\).
\end{bpstatementit}

% proof: proof/proposition-lattices-representations-normal-scheme
\begin{proof}[{Proof}]\phantomsection\label{bp35}%
Let \(M\subset L\) be as in the statement, with reductions \(M_n=M/\pi_E^nM\). The fibres of these reductions define a continuous action of \(\pi_1(X,\bar x)\) on the finite free completed fibre

\[
M_{\bar x}:=\varprojlim_n(M_n)_{\bar x}.
\]

The fibre functor of the proposition takes \(L\) to the rational fibre \(M_{\bar x}\otimes_{\mathcal O_E}E\); in other words, we apply the fibre functor of finite covers to the finite reductions and complete the result.

Full faithfulness comes from the Galois theory of finite étale covers of the connected scheme \(X\). It says that taking the fibre at \(\bar x\) is an equivalence between finite locally constant sheaves on \(X_{\acute{e}t}\) and finite sets with continuous \(\pi_1(X,\bar x)\)-action, compatible with module structures. To pass from the finite reductions to \(L\) we need to clear denominators. Let \(M\subseteq L\) and \(M'\subseteq L'\) be lattices in two lisse \(E_X\)-modules. Then a morphism \(\phi:L\to L'\) has a uniform denominator: \(\pi_E^c\phi\) maps \(M\) into \(M'\) for some \(c\geq0\). To see this, bound the finitely many matrix entries on a quasi-compact pro-étale covering trivializing both lattices.

For faithfulness, suppose the fibre of \(\phi\) vanishes. Then the fibres of all the reductions of \(\pi_E^c\phi\) vanish, so these reductions vanish by the Galois correspondence. Since \(M'=\varprojlim_nM'_n\), it follows that \(\pi_E^c\phi=0\). For fullness, take an equivariant map between the rational fibres. After multiplication by some \(\pi_E^c\), it maps \(M_{\bar x}\) into \(M'_{\bar x}\). Its reductions are then equivariant maps of finite fibres, so by the Galois correspondence they come from compatible maps \(M_n\to M'_n\). The inverse limit of these is a map \(M\to M'\) whose rationalization, divided by \(\pi_E^c\), has the given fibre.

The same uniform denominator, applied to the identity of \(L\), shows that the rational fibre does not depend on the lattice. Indeed, if \(M_1\) and \(M_2\) are two lattices in \(L\), then \(\pi_E^c\,\mathrm{id}_{L}\) sends \(M_1\) into \(M_2\) for some \(c\geq0\), and symmetrically. On fibres these maps, divided by \(\pi_E^c\), therefore give mutually inverse equivariant isomorphisms between the two rational fibres.

For essential surjectivity we build a stable lattice and glue its reductions. A continuous representation of the profinite group \(\pi_1(X,\bar x)\) on \(E^r\) has compact image. Hence the \(\mathcal O_E\)-span of the translates of a lattice is bounded and contains that lattice. As \(\mathcal O_E\) is noetherian, this span is a stable finite free lattice, and its reductions define compatible finite locally constant étale sheaves. These sheaves become constant simultaneously on the surjective pro-étale cover obtained as the inverse limit of the finite étale covers corresponding to their kernels. Their inverse limit is therefore lisse over \(\mathcal O_E\), and its rationalization is lisse over \(E\). Finally, this construction is inverse to the fibre functor.
\end{proof}


% object: sec-The-Deligne-Laumon-theorem-2 (blueprint/chap-Weights/sec-The-Deligne-Laumon-theorem-2/_section.md), content hidden
\subsection*{\texorpdfstring{13.3\quad The Deligne–Laumon theorem}{13.3 The Deligne–Laumon theorem}}
\phantomsection\label{bp36}%
\addcontentsline{toc}{subsection}{\texorpdfstring{13.3\quad The Deligne–Laumon theorem}{13.3 The Deligne–Laumon theorem}}
% text: prose-13-12
\phantomsection\label{bp37}%
For a separated \(k_0\)-scheme of finite type \(X_0\), we write \(D_{\leq w}(X_0)\) for the category of bounded constructible Weil complexes \(K\) on \(X_0\) of weights at most \(w\) in the sense of Definition \hyperref[bp9]{13.1}, that is, such that \(\mathcal H^i(K)\) is mixed of weights at most \(w+i\) for every \(i\). The Deligne–Laumon theorem says that for a separated morphism of finite type \(f:X_0\to Y_0\), the functor \(Rf_!\) maps \(D_{\leq w}(X_0)\) into \(D_{\leq w}(Y_0)\). The shift by \(i\) makes this one statement for all degrees: for the structure morphism of \(X_0\), one condition on \(R\Gamma_c(X,K)\) bounds all the groups \(H^n_c(X,K)\) at once. Moreover, by Proposition \hyperref[bp24]{13.7}, replacing \(K\) by \(K[r]\) only moves the bound by \(r\). The proof occupies Chapter \hyperref[bp271]{18} and, for a single curve, the sections before it, and Corollary \hyperref[bp326]{19.2} deduces from the theorem the purity of the cohomology of a smooth proper variety.

% object: theorem-weights-rf-shriek (blueprint/chap-Weights/sec-The-Deligne-Laumon-theorem-2/theorem-weights-rf-shriek/_section.md)
\phantomsection\label{bp38}%
\begin{bpstatementit}{Theorem~13.11}{ (Weights of $Rf_!$)}
For every separated morphism of finite type \(f:X_0\to Y_0\),

\[
Rf_!D_{\leq w}(X_0)\subset D_{\leq w}(Y_0).
\]

In particular, if \(K\in D_{\leq w}(X_0)\), then \(H_c^n(X,K)\) has weights at most \(w+n\).
\end{bpstatementit}



% object: chap-Purity-of-real-sheaves (blueprint/chap-Purity-of-real-sheaves/_section.md), content hidden
\section*{\texorpdfstring{14\quad Purity of real sheaves}{14 Purity of real sheaves}}
\phantomsection\label{bp39}%
\addcontentsline{toc}{section}{\texorpdfstring{14\quad Purity of real sheaves}{14 Purity of real sheaves}}
% text: prose-14-01
\phantomsection\label{bp40}%
In Chapter \hyperref[bp223]{17} the Fourier transform of a pure sheaf produces a lisse sheaf \(M'\) whose Frobenius eigenvalues we cannot compute. What we do know is that a twist of \(M'\) is a quotient of a sheaf whose traces are real, and from this alone we must conclude that \(M'\) is pure. In Chapter \hyperref[bp271]{18} the same need returns for the direct images of a family of curves, which we can make real but do not know to be mixed. Both arguments rest on one theorem (Theorem \hyperref[bp120]{14.26}):

\begin{quote}
on a smooth geometrically connected curve, every irreducible constituent of a real lisse Weil sheaf is pure, of weight equal to its determinant weight.
\end{quote}

Before this theorem we know the weights of one kind of sheaf only: a lisse sheaf of rank one on a curve has finite geometric monodromy, and therefore it is pure (Section \hyperref[bp41]{14.1}). In particular the determinant of every lisse sheaf is pure, and its weight divided by the rank, the determinant weight, is the average weight of the eigenvalues at each point. The theorem says that on an irreducible constituent of a real sheaf no eigenvalue deviates from this average.

The tool is the trace formula. On an affine curve \(U_0\) it expresses the \(L\)-function of a lisse sheaf \(\mathcal G\) as

\[
L(U_0,\mathcal G,t)
=\prod_{x\in|U_0|}\det(1-\operatorname{Frob}_xt^{\deg(x)}\mid\mathcal G_{\bar x})^{-1}
=\frac{\det(1-Ft\mid H^1_c(U,\mathcal G))}{\det(1-Ft\mid H^2_c(U,\mathcal G))}.
\]

The poles of the right-hand side are controlled by \(H^2_c\), which is easy to compute (Lemma \hyperref[bp104]{14.20}). On the left, a large eigenvalue at a point \(x\) gives a pole of the local factor at \(x\) close to the origin. But this pole could be cancelled by a zero of the product of the other factors, as in \((1-2t)^{-1}\cdot(1-2t)=1\). Positive factors, however, cannot cancel in this way. Indeed, suppose all factors are power series with nonnegative coefficients and constant term \(1\). Then the coefficients of each factor are bounded by those of the product, and Lemma \hyperref[bp113]{14.24} shows that no factor can converge on a smaller disc than the product.

This idea is Rankin's. In 1939 he improved the bound for Ramanujan's function \(\tau(n)\) to \(O(n^{29/5})\) by studying the Dirichlet series \(\sum_n\tau(n)^2n^{-s}\), whose coefficients are nonnegative. Deligne used the same argument in his first proof [\hyperref[bp350]{Del74}].

For the Euler factors to be positive, the traces must be real. We therefore call a lisse Weil sheaf real if its local characteristic polynomials have real coefficients under the fixed embedding \(\iota\) of Convention \hyperref[bp7]{\bpslug{situation-{}\allowbreak{}weil}} (Definition \hyperref[bp107]{14.21}). For example, every sheaf with rational characteristic polynomials is real, such as a Tate twist of the constant sheaf. On an even tensor power \(\mathcal G^{\otimes2m}\) of a real sheaf the traces are \(2m\)th powers of real numbers, and this makes the Euler factors positive. From this, Proposition \hyperref[bp116]{14.25} bounds each local eigenvalue \(\alpha\) of \(\mathcal G\) by \(|\iota(\alpha)|^2\le q_x^\beta\), where \(\beta\) is the largest determinant weight of the constituents.

This is not yet the theorem, because a real sheaf need not be pure. For example, the sheaf \(\overline{\mathbf Q}_\ell\oplus\overline{\mathbf Q}_\ell(-1)\) is real, with constituents of weights \(0\) and \(2\), but the bound only gives \(2\) for both. To separate the constituents we pass to exterior powers. In \(\bigwedge^2(\overline{\mathbf Q}_\ell\oplus\overline{\mathbf Q}_\ell(-1))=\det\overline{\mathbf Q}_\ell(-1)\otimes\overline{\mathbf Q}_\ell\) the factor \(\overline{\mathbf Q}_\ell\) appears multiplied by the known eigenvalue \(q_x\) of \(\det\overline{\mathbf Q}_\ell(-1)\), so the bound \(2\) for this summand leaves the bound \(0\) for \(\overline{\mathbf Q}_\ell\) itself. In general, Proposition proposition-determinant-weights-tensor-operations finds the exterior power in which a given group of constituents carries the largest determinant weight, and the same argument applies there. Positivity, however, only ever gives upper bounds. The lower bound comes from applying the same argument to the dual, and in Chapter \hyperref[bp319]{19} this role is played by Poincaré duality.

% object: sec-Sheaves-of-rank-one (blueprint/chap-Purity-of-real-sheaves/sec-Sheaves-of-rank-one/_section.md), content hidden
\subsection*{\texorpdfstring{14.1\quad Sheaves of rank one}{14.1 Sheaves of rank one}}
\phantomsection\label{bp41}%
\addcontentsline{toc}{subsection}{\texorpdfstring{14.1\quad Sheaves of rank one}{14.1 Sheaves of rank one}}
% text: prose-14-02
\phantomsection\label{bp42}%
Rankin's argument bounds the local eigenvalues of a sheaf by its determinant weights, so it needs the determinant of every lisse sheaf to be pure (Corollary \hyperref[bp65]{14.8}). Since a determinant has rank one, we prove more generally that every rank-one lisse Weil sheaf on a smooth geometrically connected curve is pure (Corollary \hyperref[bp61]{14.6}), and the later sections Sections \hyperref[bp69]{14.2}–\hyperref[bp118]{Section 14.4} are built on this fact.

In rank one the question is about characters. Let \(U_0\) be a smooth geometrically connected curve and \(s\in W(U_0)\) a lift of Frobenius of degree one (Convention \hyperref[bp7]{\bpslug{situation-{}\allowbreak{}weil}}). Then a rank-one lisse Weil sheaf is a character \(\chi:W(U_0)\to E^\times\), continuous on \(\pi_1(U)\). If \(\chi\) is trivial on \(\pi_1(U)\), the sheaf is one of the sheaves \(\mathcal B_b\) with \(b=\chi(s)\); at a closed point \(x\) its eigenvalue is \(b^{\deg(x)}\), so it is pure of weight \(2\log_q|\iota(b)|\). A Kummer sheaf is the next simplest case. If \(n\mid q-1\), the cover \(t\mapsto t^n\) of \(\mathbf G_m\) is a \(\mu_n\)-torsor, and so a character of \(\mu_n(k_0)\) gives a rank-one sheaf on \(\mathbf G_{m,k_0}\) whose geometric monodromy is cyclic of order dividing \(n\). Its eigenvalues are values of the character on \(\mu_n(k_0)\), hence roots of unity, and it is pure of weight \(0\).

Up to a twist, every rank-one sheaf behaves like these two. Indeed, by Theorem \hyperref[bp58]{14.5} the image \(\chi(\pi_1(U))\) is finite. It follows that \(\chi(\operatorname{Frob}_x)\) is \(\chi(s)^{\deg(x)}\) times a root of unity, so the sheaf is pure of weight \(2\log_q|\iota(\chi(s))|\).

Finiteness is a statement about Weil sheaves, because on the geometric curve there are rank-one local systems with infinite monodromy. For example, the tame fundamental group of \(\mathbf G_{m,k}\) has the quotient \(\mathbf Z_\ell(1)\simeq\mathbf Z_\ell\), and \(a\mapsto(1+\ell)^a\) is a character of infinite order. But it has no Weil structure. Indeed, a Weil structure forces \(\chi(s^{-1}gs)=\chi(g)\) for \(g\in\pi_1(U)\), while conjugation by \(s\) acts on \(\mathbf Z_\ell(1)\) by multiplication by a power \(q^{\pm1}\). A character of \(\mathbf Z_\ell(1)\) that extends to the Weil group therefore satisfies \(\chi(g)^{q-1}=1\), and it has finite order.

The general proof follows the same pattern. Since there is no explicit quotient to compute with, it uses a logarithm instead: a power \(\chi^N\) takes values in \(1+\mathfrak m_E^r\), where the \(\ell\)-adic logarithm converges, and then \(\log\chi^N\) is an additive character of \(\pi_1(U)\) invariant under conjugation by \(s\), that is, a class in \(H^1(U,E)\) fixed by Frobenius (Lemma \hyperref[bp55]{14.4}). Such a class is zero (Proposition \hyperref[bp51]{14.3}), so \(\chi^N\) is trivial on \(\pi_1(U)\).

For a proper curve, Frobenius fixes no class in \(H^1\) because the Jacobian has only finitely many rational points. If \(C_0\) is a proper curve with Jacobian \(J_0\), Kummer theory and Poincaré duality identify \(\varprojlim_nH^1(C,\mathbf Z/\ell^n)\) with the dual of the Tate module \(T_\ell J\) (Proposition \hyperref[bp48]{14.2}). A fixed vector in the dual therefore means that \(1\) is an eigenvalue of Frobenius on \(V_\ell J=T_\ell J\otimes\mathbf Q_\ell\), and after scaling, Frobenius would then fix a nonzero element of \(T_\ell J\). Such an element is a compatible system of \(\ell\)-power torsion points of \(J\) defined over \(k_0\), of unbounded order. These points cannot all lie in the finite group \(J_0(k_0)\), and so there is no such element (Proposition \bpout{B.15}).

For an open curve, the residue sequence contains in addition the boundary term \(H^0(D,E(-1))\). A power of Frobenius acts on it by a power of \(q\ne1\), so it has no fixed vector either. For example, for \(\mathbf G_m\) this is the computation above, since there \(H^1(\mathbf G_{m,k},\mathbf Q_\ell)=\mathbf Q_\ell(-1)\), with Frobenius acting by \(q\).

% object: subsec-Frobenius-invariants (blueprint/chap-Purity-of-real-sheaves/sec-Sheaves-of-rank-one/subsec-Frobenius-invariants/_section.md), content hidden
\subsubsection*{\texorpdfstring{14.1.1\quad Frobenius invariants}{14.1.1 Frobenius invariants}}
\phantomsection\label{bp43}%
\addcontentsline{toc}{subsubsection}{\texorpdfstring{14.1.1\quad Frobenius invariants}{14.1.1 Frobenius invariants}}
% text: prose-14-03
\phantomsection\label{bp44}%
Over \(\mathbf C\) the link between \(H^1\) of a curve and its Jacobian is classical. There the Jacobian is a torus \(V/\Lambda\) with \(\Lambda=H_1(C,\mathbf Z)\), and \(H^1(C,\mathbf Z)=\Hom(\Lambda,\mathbf Z)\). Proposition \hyperref[bp48]{14.2} is the \(\ell\)-adic version, with \(T_\ell J\) in place of \(\Lambda\otimes\mathbf Z_\ell\). Here the link with the Jacobian is Kummer theory, which is a statement about \(\mu_{\ell^n}\) and line bundles. For this reason we work with the coefficients \(\mathbf Z/\ell^n\) and pass to \(E\) only at the end. Apart from this, the only arithmetic input is the finiteness of \(J_0(k_0)\), which holds because \(J_0\) is of finite type over a finite field.

% object: lemma-proper-smooth-model (blueprint/chap-Purity-of-real-sheaves/sec-Sheaves-of-rank-one/subsec-Frobenius-invariants/lemma-proper-smooth-model/_section.md)
\phantomsection\label{bp45}%
\begin{bpstatementit}{Lemma~14.1}{ (A proper smooth model)}
A smooth geometrically connected curve \(U_0/k_0\) has a smooth projective geometrically connected compactification \(j_0:U_0\hookrightarrow C_0\). The reduced boundary \(D_0\) is finite étale over \(k_0\).
\end{bpstatementit}

% proof: proof/lemma-proper-smooth-model
\begin{proof}[{Proof}]\phantomsection\label{bp46}%
We compactify by Nagata's theorem and then normalize. On a curve, normality is the same as regularity, and over a perfect field regularity is smoothness, so the normalization is the model we want.

By Theorem \bpout{A.2}, \(U_0\) is an open subscheme of a proper \(k_0\)-scheme, and we take the reduced closure of \(U_0\) in it and normalize this closure in \(k_0(U_0)\).  Since a scheme of finite type over a field is Nagata, the normalization is finite [Stacks, \href{https://stacks.math.columbia.edu/tag/035B}{035B}, \href{https://stacks.math.columbia.edu/tag/035S}{035S}], and it is an isomorphism over the normal open \(U_0\).  The resulting proper normal curve is regular, because a one-dimensional noetherian normal local domain is a discrete valuation ring. Hence it is smooth over the perfect field \(k_0\) [Stacks, \href{https://stacks.math.columbia.edu/tag/0B8Y}{0B8Y}].  Its boundary is zero-dimensional and proper over \(k_0\), hence finite. Since \(k_0\) is perfect, the reduced boundary is étale.  After base change to \(k\), the curve \(U\) is dense and connected, and therefore its closure \(C\) is connected.  Finally, a proper scheme of dimension one over a field is projective [Stacks, \href{https://stacks.math.columbia.edu/tag/0A26}{0A26}].
\end{proof}

% text: prose-14-04
\phantomsection\label{bp47}%
Once we have a smooth projective model \(C_0\), the Jacobian \(J_0\) of \(C_0\) exists (Appendix \bpout{C}), and we can compare the cohomology of \(C\) with its Tate module. We first make the comparison with finite coefficients, where Kummer theory applies directly. The proposition ends with a matrix description of Frobenius, which is the form used in Proposition \hyperref[bp51]{14.3}.

% object: proposition-kummer-duality-integral-tate-module (blueprint/chap-Purity-of-real-sheaves/sec-Sheaves-of-rank-one/subsec-Frobenius-invariants/proposition-kummer-duality-integral-tate-module/_section.md)
\phantomsection\label{bp48}%
\begin{bpstatementit}{Proposition~14.2}{ (Kummer, duality and the integral Tate module)}
Let \(C_0/k_0\) be a smooth proper geometrically connected curve with Jacobian \(J_0\), and put \(A_n=\mathbf Z/\ell^n\).  There are canonical compatible Galois-equivariant identifications

\[
H^1(C,A_n(1))\simeq J[\ell^n](k),\qquad
\varprojlim_n H^1(C,A_n)
\simeq\Hom_{\mathbf Z_\ell}(T_\ell J,\mathbf Z_\ell).
\]

Here the transition maps are multiplication by \(\ell\) on \(J[\ell^n](k)\) and reduction of coefficients on \(H^1(C,A_n)\).  If arithmetic Frobenius on \(T_\ell J\) has matrix \(B\), then geometric Frobenius on \(\Hom_{\mathbf Z_\ell}(T_\ell J,\mathbf Z_\ell)\) has matrix \(B^t\) in the dual basis.
\end{bpstatementit}

% proof: proof/proposition-kummer-duality-integral-tate-module
\begin{proof}[{Proof}]\phantomsection\label{bp49}%
By the Kummer sequence, \(H^1(C,\mu_{\ell^n})\) is the \(\ell^n\)-torsion of the Picard group, and Poincaré duality turns this into a statement about the untwisted \(H^1\).

The first identification is Proposition \bpout{C.17}, which also gives its Galois equivariance, and it shows that the coefficient map \(z\mapsto z^\ell\) from \(\mu_{\ell^{n+1}}\) to \(\mu_{\ell^n}\) corresponds to \([\ell]\) on the Jacobian.

For the dual Tate module, Poincaré duality at finite level pairs the untwisted \(H^1\) with the twisted one that we have just identified with \(J[\ell^n](k)\). Derived duality gives a statement about the single group \(H^1\) because all the groups involved are free \(A_n\)-modules. Indeed, we have \(H^0(C,A_n)=A_n\) and \(H^2(C,A_n(1))=A_n\) by Proposition \bpout{D.2}, while the group \(H^1(C,A_n(1))\) is free of rank \(2g\) by the identification above and Corollary \bpout{C.15}. Since \(A_n(1)\simeq A_n\) over \(k\), the untwisted \(H^1\) is free of rank \(2g\) as well. Finally, a bounded complex with projective cohomology modules is the sum of its shifted cohomology modules in the derived category, so \(A_n\)-duality acts degree by degree. In this situation derived Poincaré duality, Theorem \bpout{10.12}, applied to the constant sheaf, induces in degree one an isomorphism

\[
H^1(C,A_n)\xrightarrow{\sim}
\Hom_{A_n}(H^1(C,A_n(1)),A_n).
\]

This map is the adjoint of cup product followed by the trace map of the duality theorem, and both cup product and trace commute with reduction of coefficients and with the action of \(\operatorname{Gal}(k/k_0)\).  Because of the twist on the second factor, Frobenius acts trivially on the target \(A_n\).

In the limit, the finite-level dualities combine into a duality with the Tate module. Since the transition maps \([\ell]\) are surjective, we have \(J[\ell^n](k)=T_\ell J/\ell^nT_\ell J\). Moreover \(T_\ell J\) is finite free by Proposition \bpout{B.14}, and a homomorphism of finite free \(\mathbf Z_\ell\)-modules is determined by its reductions modulo all \(\ell^n\), and every compatible system of reductions lifts. Hence the inverse limit of the adjoints at finite level is the dual of \(T_\ell J\). We obtain a perfect pairing

\[
\varprojlim_nH^1(C,A_n)\times T_\ell J\longrightarrow\mathbf Z_\ell,
\]

invariant under Frobenius.  Since geometric Frobenius acts on \(T_\ell J\) by \(B^{-1}\), its matrix on \(\varprojlim_nH^1(C,A_n)\) in the dual basis is \((B^{-1})^{-t}=B^t\).
\end{proof}

% text: prose-14-05
\phantomsection\label{bp50}%
For the elliptic curve \(y^2=x^3-x\) over \(\mathbf F_3\) of the introduction, arithmetic Frobenius \(B\) has eigenvalues \(\pm i\sqrt3\) on \(T_\ell J\). By the proposition, geometric Frobenius acts on \(\varprojlim_nH^1(C,A_n)\) with the same eigenvalues. Since \(\det(B^t-1)=4\ne0\), it has no fixed vector. The next proposition proves the same for every smooth curve, including open ones.

% object: proposition-no-fixed-class-degree-one (blueprint/chap-Purity-of-real-sheaves/sec-Sheaves-of-rank-one/subsec-Frobenius-invariants/proposition-no-fixed-class-degree-one/_section.md)
\phantomsection\label{bp51}%
\begin{bpstatementit}{Proposition~14.3}{ (No fixed class in degree one)}
For a smooth geometrically connected curve \(U_0/k_0\), geometric Frobenius has no fixed vector on \(H^1(U,E)\) for every finite \(E/\mathbf Q_\ell\).  More precisely,

\[
\left(\varprojlim_nH^1(U,\mathbf Z/\ell^n)\right)^{F=1}=0.
\]
\end{bpstatementit}

% proof: proof/proposition-no-fixed-class-degree-one
\begin{proof}[{Proof}]\phantomsection\label{bp52}%
On the compactification the vanishing holds because \(J_0(k_0)\) is finite, and on the boundary it holds because Frobenius acts by a power of \(q\), which is not \(1\). The residue sequence reduces the open case to the proper one, and at the end we pass from integral to rational coefficients.

\emph{The proper case.} Let \(U_0=C_0\) be proper, with Jacobian \(J_0\). Here the cohomology is the dual of the Tate module, and the question is whether \(1\) is an eigenvalue of Frobenius on it.  By Proposition \bpout{B.15}, arithmetic Frobenius has no fixed vector on \(T_\ell J\); indeed, a fixed compatible tuple would consist of \(\ell\)-power torsion points of the finite group \(J_0(k_0)\). This means that \(B-1\) is invertible on \(V_\ell J\).  On the other side, geometric Frobenius acts on the untwisted \(H^1\) by \(B^t\) (Proposition \hyperref[bp48]{14.2}). Since \(\det(B^t-1)=\det(B-1)\ne0\), it fixes neither a nonzero rational vector nor a nonzero integral compatible tuple.  The argument applies to every power of Frobenius, so the same holds over every finite extension of \(k_0\), which we need in the open case.

\emph{The open case.} A fixed class has zero residues at the boundary points, so it extends to the compactification, where the proper case shows that it is zero. To make this precise, choose \(j:U\hookrightarrow C\) as in Lemma \hyperref[bp45]{14.1}, with boundary \(i:D\hookrightarrow C\). At each finite level there is then a canonical residue sequence

\[
0\longrightarrow H^1(C,A_n)\longrightarrow H^1(U,A_n)
\longrightarrow H^0(D,A_n(-1)).
\]

To obtain it with the correct twist and Frobenius action, we use the dualizing complexes.  By transitivity, \(i^!\omega_C=\omega_D=A_n\). Since \(\omega_C=A_n(1)[2]\), it follows that \(i^!A_{n,C}=A_{n,D}(-1)[-2]\).  The localization triangle \(i_*i^!A_{n,C}\to A_{n,C}\to Rj_*A_{n,U}\) now yields the displayed sequence, and all its maps are natural in the coefficients and under automorphisms of the model.

Choose \(m>0\) such that \(F^m\) fixes all geometric boundary points, and put \(Q=q^m\).  On \(\varprojlim_nH^0(D,A_n(-1))\), which is finite free over \(\mathbf Z_\ell\), the operator \(F^m\) acts as \(Q\). Since \(Q-1\ne0\), it has no nonzero fixed element. A fixed compatible class on \(U\) therefore has zero residues at every level, so by exactness it lifts uniquely to \(H^1(C,A_n)\). By uniqueness the lifts are compatible and fixed by \(F^m\), and the proper case over \(\mathbf F_Q\) shows that they vanish.

\emph{Rational coefficients.} To pass from compatible systems of classes to \(H^1(U,E)\), we identify \(H^1(U,\mathbf Q_\ell)\) with the rationalization of the inverse limit and check that rationalizing does not create fixed vectors. By Proposition \bpout{2.55}, applied to the constant sheaves \(A_n\), we have \(R\Gamma(U,\mathbf Z_\ell)=R\varprojlim_nR\Gamma(U,A_n)\), and hence an exact sequence

\[
0\to\varprojlim\nolimits^1_nH^0(U,A_n)\to H^1(U,\mathbf Z_\ell)
\to\varprojlim_nH^1(U,A_n)\to0.
\]

The \(\varprojlim^1\)-term vanishes because the groups \(H^0(U,A_n)=A_n\) have surjective transition maps.  On \(U_{\mathrm{pro}\acute{e}t}\) the sheaf \(\mathbf Q_\ell\) is \(\mathbf Z_\ell[1/\ell]\) (Lemma \bpout{2.51}), a filtered colimit of copies of \(\mathbf Z_\ell\), and cohomology on the quasi-compact and quasi-separated scheme \(U\) commutes with filtered colimits of sheaves (Corollary \bpout{2.29}, applied to sheaves placed in degree zero).  So \(H^1(U,\mathbf Q_\ell)=H^1(U,\mathbf Z_\ell)\otimes\mathbf Q_\ell\) is the rationalization of the inverse limit.

The inverse limit of the residue sequences is exact on the left, and rationalization is exact, so we obtain an exact sequence of finite-dimensional \(\mathbf Q_\ell\)-spaces \(0\to H^1(C,\mathbf Q_\ell)\to H^1(U,\mathbf Q_\ell)\to I\to0\), where \(I\) is the image of the residue map inside the rationalized boundary term. By the proper case, \(F^m-1\) is invertible on \(H^1(C,\mathbf Q_\ell)\). On the boundary term, \(F^m-1\) acts by \(Q-1\), so it is injective on the \(F^m\)-stable subspace \(I\). Now let \(v\) be a vector of \(H^1(U,\mathbf Q_\ell)\) killed by \(F^m-1\). Its image in \(I\) is killed by \(F^m-1\), hence zero. Thus \(v\) comes from \(H^1(C,\mathbf Q_\ell)\), where it is zero by the proper case. Since a vector fixed by \(F\) is fixed by \(F^m\), it follows that Frobenius has no fixed rational vector, and extending scalars to \(E\) gives the assertion over \(E\).
\end{proof}


% object: subsec-Finite-monodromy (blueprint/chap-Purity-of-real-sheaves/sec-Sheaves-of-rank-one/subsec-Finite-monodromy/_section.md), content hidden
\subsubsection*{\texorpdfstring{14.1.2\quad Finite monodromy}{14.1.2 Finite monodromy}}
\phantomsection\label{bp53}%
\addcontentsline{toc}{subsubsection}{\texorpdfstring{14.1.2\quad Finite monodromy}{14.1.2 Finite monodromy}}
% text: prose-14-06
\phantomsection\label{bp54}%
With Proposition \hyperref[bp51]{14.3} in hand, the logarithm argument sketched at the start of Section \hyperref[bp41]{14.1} proves Theorem \hyperref[bp58]{14.5}. Corollaries \hyperref[bp61]{Corollary 14.6} and \hyperref[bp65]{Corollary 14.8} then translate finiteness of monodromy into weights. For a lisse sheaf of higher rank, however, the determinant weight (Definition \hyperref[bp64]{14.7}) is the only weight we can attach until Rankin's argument. Throughout, \(U_0\) is a smooth geometrically connected curve over \(k_0\) and \(s\in W(U_0)\) is a lift of degree one, as in Convention \hyperref[bp7]{\bpslug{situation-{}\allowbreak{}weil}}. In this setting, the first step identifies the additive characters to which the logarithm leads.

% object: lemma-invariant-additive-characters (blueprint/chap-Purity-of-real-sheaves/sec-Sheaves-of-rank-one/subsec-Finite-monodromy/lemma-invariant-additive-characters/_section.md)
\phantomsection\label{bp55}%
\begin{bpstatementit}{Lemma~14.4}{ (Invariant additive characters)}
Every continuous additive homomorphism \(a:\pi_1(U)\to E\) satisfying \(a(s^{-1}gs)=a(g)\) is zero.
\end{bpstatementit}

% proof: proof/lemma-invariant-additive-characters
\begin{proof}[{Proof}]\phantomsection\label{bp56}%
After scaling, the coordinates of \(a\) reduce modulo \(\ell^n\) to a compatible system of Frobenius-fixed classes in \(H^1(U,\mathbf Z/\ell^n)\), which vanish by Proposition \hyperref[bp51]{14.3}. This needs \(H^1\) to be identified with homomorphisms.

Recall that for a finite abelian group \(G\) the group \(H^1(U,G)\) classifies \(G\)-torsors. By [Stacks, \href{https://stacks.math.columbia.edu/tag/0DV5}{0DV5}], finite locally constant sheaves on the connected scheme \(U\) correspond to finite continuous \(\pi_1(U)\)-sets, and under this correspondence a \(G\)-torsor becomes a finite \(\pi_1(U)\)-set with a commuting simply transitive action of \(G\), that is, a continuous homomorphism \(\pi_1(U)\to G\). Hence \(H^1(U,G)=\Hom_{\mathrm{cont}}(\pi_1(U),G)\). In this description Frobenius acts on \(H^1(U,G)\) through the action of \(s\) on \(\pi_1(U)\) by conjugation. The Frobenius-fixed classes are therefore the homomorphisms \(b\) with \(b(s^{-1}gs)=b(g)\).

Choose a \(\mathbf Q_\ell\)-basis of \(E\), with coordinate maps \(t_i\).  Since \(\pi_1(U)\) is compact, each \(t_i a(\pi_1(U))\) is bounded, and after multiplying all coordinates by one power of \(\ell\) we may assume that they take values in \(\mathbf Z_\ell\).  Each coordinate is then a continuous homomorphism \(\pi_1(U)\to\mathbf Z_\ell\) invariant under conjugation by \(s\). Its reductions modulo \(\ell^n\) are compatible classes in \(H^1(U,\mathbf Z/\ell^n)\) fixed by Frobenius. By the proposition they vanish, and hence every coordinate of \(a\) is zero.
\end{proof}

% text: prose-14-07
\phantomsection\label{bp57}%
The invariance hypothesis of the lemma is automatic for characters of the Weil group. If \(\chi:W(U_0)\to E^\times\) is a character, then \(\chi(s^{-1}gs)=\chi(s)^{-1}\chi(g)\chi(s)=\chi(g)\), because \(E^\times\) is commutative. So the lemma applies as soon as \(\chi\) has been made additive on \(\pi_1(U)\).

% object: theorem-rank-one-sheaves-have-finite (blueprint/chap-Purity-of-real-sheaves/sec-Sheaves-of-rank-one/subsec-Finite-monodromy/theorem-rank-one-sheaves-have-finite/_section.md)
\phantomsection\label{bp58}%
\begin{bpstatementit}{Theorem~14.5}{ (Rank-one sheaves have finite geometric monodromy)}
Let \(U_0/k_0\) be a smooth geometrically connected curve, and let \(\chi:W(U_0)\to E^\times\) be the character of a rank-one lisse Weil sheaf.  Then \(\chi(\pi_1(U))\) is finite.
\end{bpstatementit}

% proof: proof/theorem-rank-one-sheaves-have-finite
\begin{proof}[{Proof}]\phantomsection\label{bp59}%
We show that a fixed power of \(\chi\) is trivial on \(\pi_1(U)\), by applying Lemma \hyperref[bp55]{14.4} to its logarithm. For this we need a subgroup \(1+\mathfrak m_E^r\) of finite index in \(\mathcal O_E^\times\) on which the logarithm is an injective homomorphism, and a power of \(\chi\) that lands in it.

\emph{Step 1: the logarithm.} The compact group \(\chi(\pi_1(U))\) lies in \(\mathcal O_E^\times\), because its image under the normalized valuation is a finite subgroup of \(\mathbf Z\) and therefore zero. Now write \(e=v_E(\ell)\) and choose an integer \(r>e/(\ell-1)\). We claim that \(\log\) is injective on \(1+\mathfrak m_E^r\). For \(x\in\mathfrak m_E^r\) and \(n\ge1\) we have \(v_E(n)=e\,v_\ell(n)\) and, by Legendre's formula, \(v_\ell(n!)\le(n-1)/(\ell-1)\), so that

\[
v_E(x^n/n)\ge nr-e\,v_\ell(n),
\qquad
v_E(x^n/n!)\ge r+(n-1)\bigl(r-e/(\ell-1)\bigr).
\]

Both bounds tend to infinity with \(n\) and are at least \(r\): for the second this is clear, and for the first it follows from \(n-1\ge\ell^{v_\ell(n)}-1\ge(\ell-1)v_\ell(n)\) and \((\ell-1)r>e\), which give \((n-1)r\ge e\,v_\ell(n)\). The logarithm and exponential series therefore converge on \(\mathfrak m_E^r\); the logarithm maps \(1+\mathfrak m_E^r\) into \(\mathfrak m_E^r\), and the exponential maps \(\mathfrak m_E^r\) into \(1+\mathfrak m_E^r\). By the usual identities of formal power series, \(\log:1+\mathfrak m_E^r\to\mathfrak m_E^r\) and \(\exp\) are mutually inverse continuous group homomorphisms.

\emph{Step 2: a power of \(\chi\).} The quotient \(\mathcal O_E^\times/(1+\mathfrak m_E^r)\) is finite, being isomorphic to \((\mathcal O_E/\mathfrak m_E^r)^\times\); let \(N>0\) be an integer that kills it. As \(\chi(\pi_1(U))\subset\mathcal O_E^\times\), we have \(\chi(g)^N\in1+\mathfrak m_E^r\) for \(g\in\pi_1(U)\). Hence \(a(g)=\log(\chi(g)^N)\) defines a continuous homomorphism \(\pi_1(U)\to E\), which is invariant under conjugation by \(s\) because \(\chi\) takes values in an abelian group.

\emph{Step 3: conclusion.} Lemma \hyperref[bp55]{14.4} gives \(a=0\), so \(\chi(g)^N=1\) for every \(g\in\pi_1(U)\), the logarithm being injective. The theorem follows, since \(T^N-1\) has only finitely many roots in \(E\).
\end{proof}

% text: prose-14-08
\phantomsection\label{bp60}%
Finiteness gives the sheaf a weight, and since the lift \(s\) is a choice, part of the proof checks that the weight does not depend on it.

% object: corollary-rank-one-weights (blueprint/chap-Purity-of-real-sheaves/sec-Sheaves-of-rank-one/subsec-Finite-monodromy/corollary-rank-one-weights/_section.md)
\phantomsection\label{bp61}%
\begin{bpstatementit}{Corollary~14.6}{ (Rank-one weights)}
Every rank-one lisse Weil sheaf on a smooth geometrically connected curve is punctually \(\iota\)-pure.  For \(c=\chi(s)\) its weight is \(\lambda_\iota(c)=2\log_q|\iota(c)|\).
\end{bpstatementit}

% proof: proof/corollary-rank-one-weights
\begin{proof}[{Proof}]\phantomsection\label{bp62}%
We separate from \(\chi\) the part that depends only on the degree. Put \(\varepsilon(w)=\chi(w)c^{-\deg w}\), a character of \(W(U_0)\) that is trivial on \(s\).  Every element of \(W(U_0)\) can be written as \(w=gs^d\) with \(g\in\pi_1(U)\) and \(d=\deg w\), and then \(\varepsilon(w)=\chi(g)\chi(s)^dc^{-d}=\chi(g)\). So \(\varepsilon\) has the same image as \(\chi|_{\pi_1(U)}\), which is finite by Theorem \hyperref[bp58]{14.5}.  At a closed point \(x\) the eigenvalue is \(\varepsilon(\operatorname{Frob}_x)c^{\deg(x)}\), where the first factor has finite order, hence complex absolute value one. Therefore the eigenvalue has absolute value \(q_x^{\lambda_\iota(c)/2}\).  This number does not depend on \(s\), because changing \(s\) multiplies \(c\) by a value of \(\chi\) on \(\pi_1(U)\), which has finite order.
\end{proof}

% text: prose-14-09
\phantomsection\label{bp63}%
For a Kummer sheaf on \(\mathbf G_{m,k_0}\) the corollary gives weight \(0\), since \(\chi(s)\) is then a root of unity. For \(\mathcal B_b\) it gives \(\lambda_\iota(b)\), the weight found in Proposition \hyperref[bp24]{13.7}.

Before purity is known, the eigenvalues of a lisse sheaf \(\mathcal V\) of rank \(r\) at a closed point are not controlled individually. Their product, however, is controlled, since it is the eigenvalue of the rank-one sheaf \(\det\mathcal V\), whose weight is given by Corollary \hyperref[bp61]{14.6}. Dividing by \(r\) we get the average weight of the eigenvalues, which is the only candidate for a weight of \(\mathcal V\). Indeed, if \(\mathcal V\) is pure of weight \(w\), then \(\det\mathcal V\) is pure of weight \(rw\), and the number defined below is \(w\). For a sheaf that is not pure, on the other hand, the number is only an average. For example, on \(\overline{\mathbf Q}_\ell\oplus\overline{\mathbf Q}_\ell(-1)\) the determinant is \(\overline{\mathbf Q}_\ell(-1)\), of weight \(2\), and the number is \(1\), the weight of no eigenvalue. For this reason we apply it to irreducible constituents: Proposition \hyperref[bp116]{14.25} bounds the eigenvalues of a real sheaf by the largest determinant weight of its constituents, and Theorem \hyperref[bp120]{14.26} then shows that each irreducible constituent is pure of weight equal to its determinant weight.

% object: definition-determinant-weight (blueprint/chap-Purity-of-real-sheaves/sec-Sheaves-of-rank-one/subsec-Finite-monodromy/definition-determinant-weight/_section.md)
\phantomsection\label{bp64}%
\begin{bpstatement}{Definition~14.7}{ (Determinant weight)}
Let \(\mathcal V\) be a nonzero lisse Weil sheaf of rank \(r\) on a smooth geometrically connected curve \(U_0\), let \(\rho\) be the corresponding representation of \(W(U_0)\), and let \(s\in W(U_0)\) be a lift of degree one. The \emph{determinant weight} of \(\mathcal V\) is

\[
\operatorname{dw}_\iota(\mathcal V)
=\frac{\lambda_\iota(\det\rho(s))}{r}
\]

(this does not depend on \(s\), Corollary \hyperref[bp65]{14.8}). For a lisse Weil sheaf \(\mathcal V\) we write \(\operatorname{DW}(\mathcal V)\) for the set of determinant weights of its irreducible constituents.
\end{bpstatement}

% object: corollary-determinant-weights (blueprint/chap-Purity-of-real-sheaves/sec-Sheaves-of-rank-one/subsec-Finite-monodromy/corollary-determinant-weights/_section.md)
\phantomsection\label{bp65}%
\begin{bpstatementit}{Corollary~14.8}{ (Determinant weights)}
Let \(\mathcal V\) be a nonzero lisse Weil sheaf of rank \(r\) on a smooth geometrically connected curve.  Its determinant has finite geometric monodromy and is punctually pure of weight \(r\operatorname{dw}_\iota(\mathcal V)\). The number \(\operatorname{dw}_\iota(\mathcal V)\) does not depend on the choice of \(s\) or of a finite coefficient field over which \(\mathcal V\) is defined.
\end{bpstatementit}

% proof: proof/corollary-determinant-weights
\begin{proof}[{Proof}]\phantomsection\label{bp66}%
The sheaf \(\det\mathcal V\) is a rank-one lisse Weil sheaf defined over a finite extension of \(E\), so Theorem \hyperref[bp58]{14.5} and Corollary \hyperref[bp61]{14.6} apply to it.  By the formula of Corollary \hyperref[bp61]{14.6}, with \(c=\det\rho(s)\), the weight of \(\det\mathcal V\) is \(r\operatorname{dw}_\iota(\mathcal V)\), and it does not depend on \(s\). Finally, a finite extension of scalars changes neither the character nor its absolute value under the fixed embedding.
\end{proof}

% text: prose-14-10
\phantomsection\label{bp67}%
The corollary applies to every irreducible constituent of \(\mathcal V\) and of every tensor or exterior power of \(\mathcal V\). Indeed, a finite-dimensional representation has a composition series, and each factor is again lisse and defined over a finite extension of \(\mathbf Q_\ell\), because after enlarging \(E\) the invariant subspaces and quotients have bases defined over \(E\).

% object: remark-determinant-weight-pure-constituents (blueprint/chap-Purity-of-real-sheaves/sec-Sheaves-of-rank-one/subsec-Finite-monodromy/remark-determinant-weight-pure-constituents/_section.md)
\phantomsection\label{bp68}%
\begin{bpstatement}{Remark~14.9}{ (Determinant weight of a sheaf with pure constituents)}
Let \(\mathcal V\) be a nonzero lisse Weil sheaf of rank \(r\) on a smooth geometrically connected curve. If the irreducible constituents \(\mathcal V_1,\dots,\mathcal V_n\) of \(\mathcal V\) are pure of weights \(w_1,\dots,w_n\), then \(\operatorname{dw}_\iota(\mathcal V)=\sum_a\operatorname{rank}(\mathcal V_a)w_a/r\). In particular a sheaf pure of weight \(w\) has determinant weight \(w\).
\end{bpstatement}



% object: sec-Determinant-weights (blueprint/chap-Purity-of-real-sheaves/sec-Determinant-weights/_section.md), content hidden
\subsection*{\texorpdfstring{14.2\quad Determinant weights}{14.2 Determinant weights}}
\phantomsection\label{bp69}%
\addcontentsline{toc}{subsection}{\texorpdfstring{14.2\quad Determinant weights}{14.2 Determinant weights}}
% text: prose-14-11
\phantomsection\label{bp70}%
Before purity is known, the only weight we can attach to an irreducible lisse Weil sheaf \(Q\) on a curve is its determinant weight \(\operatorname{dw}_\iota(Q)\), the weight of \(\det Q\) divided by the rank (Definition \hyperref[bp64]{14.7} and Corollary \hyperref[bp65]{14.8}). This section proves the two facts about determinant weights that Rankin's argument in Section \hyperref[bp102]{14.3} needs for the tensor and exterior powers of a real sheaf \(\mathcal G\). First, every simple factor of \(\mathcal G^{\otimes2m}\) has determinant weight at most \(2m\) times the largest determinant weight of the constituents of \(\mathcal G\), and this bounds the weights of \(H^2_c\). Second, the determinant weights of the simple factors of \(\bigwedge^r\mathcal G\) are known exactly, and this is what lets exterior powers separate the constituents.

The difficulty is that a determinant only sees an average. For irreducible \(A\) and \(B\) of ranks \(a\) and \(b\), we have \(\det(A\otimes B)=(\det A)^{\otimes b}\otimes(\det B)^{\otimes a}\). Hence the determinant weight of \(A\otimes B\) as a whole is \(\operatorname{dw}_\iota(A)+\operatorname{dw}_\iota(B)\). But \(A\otimes B\) is usually reducible, and which subquotients occur depends on the monodromy. The rank-weighted average of the determinant weights of its simple factors is the number above, but nothing so far prevents one factor from lying above it and another below. If the factors were pure, the determinant weight of each would be its weight, but purity is what we are trying to prove.

The model to keep in mind is a central character. If a group acts on \(V\) and some central element \(z\) acts on \(V\) by the scalar \(c\), then \(z\) acts on every subquotient of \(V^{\otimes a}\otimes(V^\vee)^{\otimes b}\) by \(c^{a-b}\), without any knowledge of how these representations decompose. We therefore look for a central element of positive degree in the monodromy group; its scalars will record weights.

Let \(H\subseteq G\) be the Zariski closures in \(\GL(V)\) of the images of \(\pi_1(U)\) and of \(W(U_0)\), the geometric and arithmetic monodromy groups. If \(H\) were trivial, the image \(\rho(s)\) of a degree-one lift would itself be central. In general we use the structure of the identity component,

\[
G^\circ=Z(G^\circ)^\circ H^\circ
\]

for \(V\) semisimple (Lemma \hyperref[bp87]{14.15}). The reason is that the Lie algebra of \(G^\circ\) is reductive, hence the sum of its center and its derived algebra, and since \(G/H\) is abelian, the derived algebra lies in the Lie algebra of \(H\). Accordingly, in Proposition \hyperref[bp92]{14.17} we write a power of \(\rho(s)\) as a central torus element times an element of \(H\), and then we take the product of its conjugates over \(G/G^\circ\). The result is an element \(z=h\rho(s)^N\), central in \(G\), with \(h\in H\) and \(N>0\).

On a simple factor \(Q\) of any tensor construction, \(z\) acts by a scalar \(c_Q\), and taking determinants gives \(c_Q^{\dim Q}=\det(h\mid Q)\det(s\mid Q)^N\). Since determinants have finite geometric monodromy, \(\det(h\mid Q)\) is a root of unity, and from this Lemma \hyperref[bp96]{14.18} concludes that \(\operatorname{dw}_\iota(Q)=\lambda_\iota(c_Q)/N\), where \(\lambda_\iota(a)=2\log_q|\iota(a)|\) is the additive function of Proposition \hyperref[bp24]{13.7}. Since scalars multiply under tensor products, we obtain Proposition \hyperref[bp99]{14.19}.

The facts on algebraic groups that we use, namely the relation between closed subgroups, their identity components and their Lie algebras, and the structure of a Lie algebra with a faithful semisimple representation, are proved in Appendix \bpout{F}.

% object: subsec-Representations-and-filtrations (blueprint/chap-Purity-of-real-sheaves/sec-Determinant-weights/subsec-Representations-and-filtrations/_section.md), content hidden
\subsubsection*{\texorpdfstring{14.2.1\quad Representations and filtrations}{14.2.1 Representations and filtrations}}
\phantomsection\label{bp71}%
\addcontentsline{toc}{subsubsection}{\texorpdfstring{14.2.1\quad Representations and filtrations}{14.2.1 Representations and filtrations}}
% text: prose-14-12
\phantomsection\label{bp72}%
We treat lisse sheaves as representations of the Weil group. From the geometry of the curve we use only one fact, proved in Section \hyperref[bp41]{14.1}: the geometric fundamental group acts through a finite group on every determinant. Everything else is linear algebra and the structure of algebraic groups.

Throughout, \(U_0\) is a smooth geometrically connected curve over \(k_0\) and \(s\in W(U_0)\) is a lift of degree one, as in Convention \hyperref[bp7]{\bpslug{situation-{}\allowbreak{}weil}}, so that \(W(U_0)=\pi_1(U)\rtimes\langle s\rangle\).

% object: definition-weil-representation (blueprint/chap-Purity-of-real-sheaves/sec-Determinant-weights/subsec-Representations-and-filtrations/definition-weil-representation/_section.md)
\phantomsection\label{bp73}%
\begin{bpstatement}{Definition~14.10}{ (Weil representation)}
A \emph{Weil representation} is a finite-dimensional \(\overline{\mathbf Q}_\ell\)-representation of \(W(U_0)\) whose restriction to \(\pi_1(U)\) is continuous.
\end{bpstatement}

% text: prose-14-13
\phantomsection\label{bp74}%
By Convention \hyperref[bp7]{\bpslug{situation-{}\allowbreak{}weil}} the Weil representations are the lisse Weil \(\overline{\mathbf Q}_\ell\)-sheaves on \(U_0\), and so, by Corollary \hyperref[bp65]{14.8}, \(\pi_1(U)\) acts through a finite group on the determinant of every Weil representation. A Weil representation \(V\) is semisimple if it is a direct sum of simple representations, and the semisimplification \(V^{\mathrm{ss}}\) of \(V\) is the direct sum of its composition factors.

Two elementary facts are needed. First, semisimplicity is preserved by restriction to a normal subgroup, which lets us apply Lie theory to the identity component of the monodromy group below. Second, the composition factors of tensor and exterior powers are computed from those of the given representations, so it suffices to prove the formulas for semisimple representations.

% object: lemma-restriction-normal-subgroup (blueprint/chap-Purity-of-real-sheaves/sec-Determinant-weights/subsec-Representations-and-filtrations/lemma-restriction-normal-subgroup/_section.md)
\phantomsection\label{bp75}%
\begin{bpstatementit}{Lemma~14.11}{ (Restriction to a normal subgroup)}
Let \(N\) be a normal subgroup of a group \(\Gamma\). The restriction to \(N\) of a finite-dimensional semisimple \(\overline{\mathbf Q}_\ell\)-representation of \(\Gamma\) is semisimple.
\end{bpstatementit}

% proof: proof/lemma-restriction-normal-subgroup
\begin{proof}[{Proof}]\phantomsection\label{bp76}%
It suffices to consider a simple \(\Gamma\)-module \(V\). A nonzero \(N\)-submodule \(A\) of least dimension is simple, and since \(N\) is normal, each translate \(gA\) is again a simple \(N\)-submodule. The sum of all these translates is a nonzero \(\Gamma\)-submodule, hence equals \(V\). It remains to see that a finite sum of simple submodules is semisimple. Indeed, the intersection of a simple submodule with the sum of those already chosen is either zero or the whole submodule, so we add it in the first case and omit it in the second, and after finitely many steps we obtain a direct-sum decomposition.
\end{proof}

% text: prose-14-14
\phantomsection\label{bp77}%
The second fact concerns representations that are not semisimple. A nonsplit extension of \(\overline{\mathbf Q}_\ell(-1)\) by \(\overline{\mathbf Q}_\ell\), such as the twist \(\mathcal E(-1)\) of the Kummer extension of Example \hyperref[bp17]{13.4}, does not decompose. Yet the composition factors of its tensor and exterior powers are determined by those of the sheaf, as the next lemma shows by filtering. This is what lets us pass to semisimplifications in Proposition \hyperref[bp99]{14.19}.

% object: lemma-tensor-exterior-powers-filtered-representations (blueprint/chap-Purity-of-real-sheaves/sec-Determinant-weights/subsec-Representations-and-filtrations/lemma-tensor-exterior-powers-filtered-representations/_section.md)
\phantomsection\label{bp78}%
\begin{bpstatementit}{Lemma~14.12}{ (Tensor and exterior powers of filtered representations)}
Suppose \(V\) and \(V'\) have composition factors \(A_1,\ldots,A_t\) and \(B_1,\ldots,B_u\). There is a \(W(U_0)\)-stable filtration of \(V\otimes V'\) with graded pieces \(A_a\otimes B_b\), one for each pair \((a,b)\). There is a \(W(U_0)\)-stable filtration of \(\bigwedge^rV\) with graded pieces

\[
\bigotimes_{a=1}^t\bigwedge^{n_a}A_a,
\qquad 0\le n_a\le\dim A_a,\qquad\sum_an_a=r.
\]

The simple factors of \(V\otimes V'\) and \(\bigwedge^rV\), counted with multiplicity, are those of the respective graded pieces. Every graded piece of \(\bigwedge^rV\) displayed above is nonzero.
\end{bpstatementit}

% proof: proof/lemma-tensor-exterior-powers-filtered-representations
\begin{proof}[{Proof}]\phantomsection\label{bp79}%
We construct the two filtrations and then compare simple factors by the Jordan–Hölder argument. The tensor product is filtered by the filtration on \(V\), and each quotient \(A_a\otimes V'\) by the filtration on \(V'\); since tensoring over a field is exact, the quotients are as asserted.

For the exterior power we cannot simply filter factor by factor, because a wedge product may use several vectors from the same composition factor. Instead we sort wedges of basis vectors by how many vectors they take from each factor, after lifting bases of the \(A_a\) to a basis of \(V\) adapted to its filtration. Thus a wedge of basis vectors determines a tuple \((n_a)\), by counting how many vectors it uses from each lifted basis. Since the basis is adapted, the group \(W(U_0)\) carries a lifted basis vector of \(A_a\) into the span of the lifted bases of \(A_1,\ldots,A_a\). Hence the action is block upper triangular: each term of the image of a basis wedge either has the same tuple or has a tuple obtained from it by moving some of its vectors to smaller indices \(a\). Choose a total order on the tuples compatible with these moves. Then the spans of the initial segments of basis wedges form a stable filtration, and on the quotient indexed by \((n_a)\) the diagonal blocks act as on \(\bigotimes_a\bigwedge^{n_a}A_a\), a space of dimension \(\prod_a\binom{\dim A_a}{n_a}>0\). For example, for \(V\) an extension of \(A_2\) by \(A_1\), both of dimension one, and \(r=1\), the tuples are \((1,0)\) and \((0,1)\), and the filtration is the given one. For \(r=2\) there is the single tuple \((1,1)\), and \(\bigwedge^2V=A_1\otimes A_2\).

The graded pieces of the filtrations just built need not be simple, so to compare simple factors we refine. By the Schreier refinement theorem, any two finite filtrations have refinements with isomorphic graded pieces, matched by the Zassenhaus isomorphism \((V_{i-1}+V_i\cap W_j)/(V_{i-1}+V_i\cap W_{j-1})\simeq
(W_{j-1}+W_j\cap V_i)/(W_{j-1}+W_j\cap V_{i-1})\); for two composition series this is the Jordan–Hölder theorem. Applying the same refinement to a composition series and to the tensor or exterior filtration, we obtain the statement about simple factors. In particular, these multisets do not change when \(V\) and \(V'\) are replaced by their semisimplifications.
\end{proof}


% object: subsec-Central-elements (blueprint/chap-Purity-of-real-sheaves/sec-Determinant-weights/subsec-Central-elements/_section.md), content hidden
\subsubsection*{\texorpdfstring{14.2.2\quad Central elements}{14.2.2 Central elements}}
\phantomsection\label{bp80}%
\addcontentsline{toc}{subsubsection}{\texorpdfstring{14.2.2\quad Central elements}{14.2.2 Central elements}}
% text: prose-14-15
\phantomsection\label{bp81}%
In the arithmetic monodromy group we construct a central element of positive degree (Proposition \hyperref[bp92]{14.17}). It acts by a scalar on each simple factor, and this scalar determines the determinant weight of the factor. As explained above, we construct the element for semisimple representations, and then use stable filtrations to obtain the formulas for arbitrary representations.

The degree must be positive. Indeed, an element of \(H\) that is central in \(G\) also acts by a scalar on each simple factor \(Q\). But, as for every element of \(H\), its determinant on \(Q\) is a root of unity, so under \(\iota\) the scalar has absolute value one, and it records no weight. By contrast, an element \(z=h\rho(s)^N\) with \(N>0\) contains \(N\) factors of Frobenius, and its scalar \(c_Q\) on \(Q\) satisfies \(\lambda_\iota(c_Q)=N\operatorname{dw}_\iota(Q)\) (Lemma \hyperref[bp96]{14.18}).

We produce the central element by Lie theory, which applies to algebraic groups and not to the abstract images of \(\pi_1(U)\) and \(W(U_0)\). For \(\mathcal B_b\) with \(b\) not a root of unity, the image of \(W(U_0)\) in \(\GL_1\) is the infinite cyclic group generated by \(b\), which is not an algebraic subgroup; its Zariski closure is all of \(\GL_1\), since the proper closed subgroups of \(\mathbf G_m\) are finite. Passing to the closure loses nothing that we need. Indeed, stabilizers are closed, so a subspace of \(V\) is stable under a group of matrices if and only if it is stable under its Zariski closure. For this reason Lemma \hyperref[bp87]{14.15} and Proposition \hyperref[bp92]{14.17} are statements about these closures.

% object: definition-monodromy-groups (blueprint/chap-Purity-of-real-sheaves/sec-Determinant-weights/subsec-Central-elements/definition-monodromy-groups/_section.md)
\phantomsection\label{bp82}%
\begin{bpstatement}{Definition~14.13}{ (Monodromy groups)}
Let \(\rho:W(U_0)\to\GL(V)\) be a Weil representation. Its \emph{geometric} and \emph{arithmetic monodromy groups} are the Zariski closures \(H\) of \(\rho(\pi_1(U))\) and \(G\) of \(\rho(W(U_0))\) in \(\GL(V)\).
\end{bpstatement}

% text: prose-14-16
\phantomsection\label{bp83}%
We keep the letters \(H\) and \(G\) for these two groups in the rest of this section.

% object: lemma-normality-geometric-monodromy (blueprint/chap-Purity-of-real-sheaves/sec-Determinant-weights/subsec-Central-elements/lemma-normality-geometric-monodromy/_section.md)
\phantomsection\label{bp84}%
\begin{bpstatementit}{Lemma~14.14}{ (Normality of geometric monodromy)}
The groups \(H\) and \(G\) are reduced closed algebraic subgroups of \(\GL(V)\), and \(H\triangleleft G\). Moreover every commutator of two elements of \(G\) lies in \(H\).
\end{bpstatementit}

% proof: proof/lemma-normality-geometric-monodromy
\begin{proof}[{Proof}]\phantomsection\label{bp85}%
Each assertion holds for the images of \(\pi_1(U)\) and \(W(U_0)\), and we pass to Zariski closures by continuity. First, multiplication and inversion are continuous for the Zariski topology, so the closure of a subgroup is a subgroup. Next, conjugation by \(\rho(W(U_0))\) preserves \(H\) because \(\pi_1(U)\) is normal in \(W(U_0)\), and since the normalizer of a closed subgroup is closed, conjugation by \(G\) preserves \(H\) as well. Finally, since \(W(U_0)/\pi_1(U)\simeq\mathbf Z\) is abelian, commutators of elements of \(\rho(W(U_0))\) lie in \(\rho(\pi_1(U))\). As the commutator map \(G\times G\to\GL(V)\) is continuous, the same holds on the closures.
\end{proof}

% text: prose-14-17
\phantomsection\label{bp86}%
For example, for \(V=\mathcal B_b\) of Proposition \hyperref[bp24]{13.7} the group \(H\) is trivial. Hence \(\rho(s)=b\) is central in \(G\), and the determinant weight of \(\mathcal B_b\) is \(\lambda_\iota(b)\).

% object: lemma-identity-component (blueprint/chap-Purity-of-real-sheaves/sec-Determinant-weights/subsec-Central-elements/lemma-identity-component/_section.md)
\phantomsection\label{bp87}%
\begin{bpstatementit}{Lemma~14.15}{ (The identity component)}
If \(V\) is a semisimple \(W(U_0)\)-representation, then

\[
G^\circ=Z(G^\circ)^\circ H^\circ.
\]
\end{bpstatementit}

% proof: proof/lemma-identity-component
\begin{proof}[{Proof}]\phantomsection\label{bp88}%
A connected algebraic group is determined by its Lie algebra among its closed connected subgroups, so it suffices to show that the Lie algebras of \(Z(G^\circ)^\circ\) and \(H^\circ\) together span that of \(G^\circ\). We show that this Lie algebra is reductive, hence the sum of its center and its derived algebra, and that the derived algebra lies in the Lie algebra of \(H\) because commutators lie in \(H\).

\emph{Reductivity.} Write \(\mathfrak g=\operatorname{Lie}(G^\circ)\). Then

\[
\mathfrak g=Z(\mathfrak g)+[\mathfrak g,\mathfrak g],
\]

by Proposition \bpout{F.8}, once we know that \(V\) is a faithful semisimple \(\mathfrak g\)-module. The \(\rho(W(U_0))\)-stable subspaces are precisely the \(G\)-stable subspaces, because the stabilizer of a subspace is closed. Hence \(V\) is semisimple as a \(G(\overline{\mathbf Q}_\ell)\)-module, and by Lemma \hyperref[bp75]{14.11} it is also semisimple as a \(G^\circ(\overline{\mathbf Q}_\ell)\)-module. Then Lemma \bpout{F.3} makes it a faithful semisimple \(\mathfrak g\)-module.

\emph{The center and the derived algebra.} Put \(T=Z(G^\circ)^\circ\) and \(\mathfrak h=\operatorname{Lie}(H^\circ)\). The comparison between groups and Lie algebras in Lemma \bpout{F.3} gives \(\operatorname{Lie}(T)=Z(\mathfrak g)\). Moreover \([\mathfrak g,\mathfrak g]\subseteq\mathfrak h\), because the commutator morphism \(G\times G\to G\) factors through the closed subgroup \(H\): it does so on \(\overline{\mathbf Q}_\ell\)-points by Lemma \hyperref[bp84]{14.14}, and \(G\times G\) is reduced. To deduce the inclusion, let \(X,Y\in\mathfrak g\) and put \(R=\overline{\mathbf Q}_\ell[\varepsilon,\eta]/(\varepsilon^2,\eta^2)\). The commutator of \(1+\varepsilon X\) and \(1+\eta Y\) in \(G(R)\) is \(1+\varepsilon\eta[X,Y]\), and it lies in \(H(R)\). Here the subring \(\overline{\mathbf Q}_\ell[\varepsilon\eta]\subset R\) is a ring of dual numbers, and the equations of \(H\) vanish on \(1+\varepsilon\eta[X,Y]\) already over this subring. Hence \([X,Y]\in\operatorname{Lie}(H)=\mathfrak h\), and

\[
\mathfrak g=\operatorname{Lie}(T)+\mathfrak h.
\]

\emph{Back to groups.} The multiplication \(T\times H^\circ\to G^\circ\) is a homomorphism, because \(T\) is central. By Lemma \bpout{F.2} its image is a closed connected subgroup, whose Lie algebra contains both summands on the right, hence equals that of \(G^\circ\). By the same lemma the image is therefore \(G^\circ\).
\end{proof}

% text: prose-14-18
\phantomsection\label{bp89}%
Rankin's argument needs the determinant weights of the simple factors of \(\mathcal G^{\otimes2m}\), of the exterior powers of \(\mathcal G\) and of the dual, so we give a name to the representations built in this way; the monodromy group acts on all of them.

% object: definition-tensor-construction (blueprint/chap-Purity-of-real-sheaves/sec-Determinant-weights/subsec-Central-elements/definition-tensor-construction/_section.md)
\phantomsection\label{bp90}%
\begin{bpstatement}{Definition~14.16}{ (Tensor construction)}
A \emph{tensor construction} on a Weil representation \(V\) is a representation obtained from \(V\) and \(V^\vee\) by finitely many tensor products, direct sums, exterior powers and subquotients by \(W(U_0)\)-stable subspaces.
\end{bpstatement}

% text: prose-14-19
\phantomsection\label{bp91}%
Since the stabilizer of a subspace is closed, a \(W(U_0)\)-stable subspace is \(G\)-stable, and so \(G\) acts on every tensor construction through its action on \(V\).

% object: proposition-central-element-arithmetic-monodromy-group (blueprint/chap-Purity-of-real-sheaves/sec-Determinant-weights/subsec-Central-elements/proposition-central-element-arithmetic-monodromy-group/_section.md)
\phantomsection\label{bp92}%
\begin{bpstatementit}{Proposition~14.17}{ (A central element of the arithmetic monodromy group)}
Let \(V_1,\ldots,V_n\) be semisimple Weil representations and form the monodromy groups \(H\subseteq G\) of \(V=\bigoplus_iV_i\), and let \(s\in W(U_0)\) have degree one. There are an integer \(N>0\) and \(h\in H(\overline{\mathbf Q}_\ell)\) such that

\[
z=h\,\rho(s)^N\in Z(G)(\overline{\mathbf Q}_\ell).
\]

The action of \(z\) on each tensor construction on \(V\) preserves every \(W(U_0)\)-stable subquotient and is scalar on each simple subquotient.
\end{bpstatementit}

% proof: proof/proposition-central-element-arithmetic-monodromy-group
\begin{proof}[{Proof}]\phantomsection\label{bp93}%
We write a power of \(\rho(s)\) as a central torus element times an element of \(H\), by Lemma \hyperref[bp87]{14.15}, and average the torus part over the components of \(G\) so that it becomes central in all of \(G\). The factors from \(H\) that accumulate in this process are collected in \(h\); their determinants are roots of unity, so by Lemma \hyperref[bp96]{14.18} they do no harm. Note that the direct sum \(V\) of the \(V_i\) is again semisimple, so Lemma \hyperref[bp87]{14.15} applies to it.

Since the component group of \(G\) is finite (Lemma \bpout{F.2}), we may choose \(m>0\) with \(\rho(s)^m\in G^\circ\). By Lemma \hyperref[bp87]{14.15}, we can write \(\rho(s)^m=ct\) with \(c\in Z(G^\circ)^\circ\) and \(t\in H^\circ\). Choose representatives \(g_1,\ldots,g_b\) of the components of \(G\) and set

\[
z=\prod_{j=1}^bg_jcg_j^{-1},\qquad N=mb.
\]

The element \(c\) is central in \(G^\circ\) but perhaps not in \(G\), since the other components may conjugate it nontrivially. The product \(z\) of its conjugates, one for each component, is however central in \(G\). Indeed, the factors lie in the abelian group \(Z(G^\circ)^\circ\), and since replacing \(g_j\) by \(g_ju\) with \(u\in G^\circ\) does not change the corresponding factor, conjugation by an element of \(G\) permutes the factors. Next, every commutator lies in \(H\), so each factor has the same image as \(c\) in \(G/H\). Since \(c\in H\rho(s)^m\), we get \(z\in H\rho(s)^{mb}\), which gives \(h\) and \(N=mb>0\).

Being central in \(G\), the element \(z\) commutes with the action of \(W(U_0)\) on tensor constructions. Moreover, as noted after the definition of tensor constructions, every \(W(U_0)\)-stable subspace of a tensor construction is \(G\)-stable. Hence \(z\) acts on the \(W(U_0)\)-stable subquotients and commutes with the action of \(W(U_0)\) on them. On a simple subquotient, therefore, an eigenspace of \(z\) is nonzero and \(W(U_0)\)-stable, hence is the whole space; in other words, \(z\) acts by a scalar.
\end{proof}


% object: subsec-Tensor-operations (blueprint/chap-Purity-of-real-sheaves/sec-Determinant-weights/subsec-Tensor-operations/_section.md), content hidden
\subsubsection*{\texorpdfstring{14.2.3\quad Tensor operations}{14.2.3 Tensor operations}}
\phantomsection\label{bp94}%
\addcontentsline{toc}{subsubsection}{\texorpdfstring{14.2.3\quad Tensor operations}{14.2.3 Tensor operations}}
% text: prose-14-20
\phantomsection\label{bp95}%
The central element acts on each simple factor of a tensor construction by a scalar, and from these scalars we read off determinant weights. Since Weil representations are the lisse Weil \(\overline{\mathbf Q}_\ell\)-sheaves, Definition \hyperref[bp64]{14.7} applies to them.

% object: lemma-determinant-weight-from-central-scalar (blueprint/chap-Purity-of-real-sheaves/sec-Determinant-weights/subsec-Tensor-operations/lemma-determinant-weight-from-central-scalar/_section.md)
\phantomsection\label{bp96}%
\begin{bpstatementit}{Lemma~14.18}{ (Determinant weight from a central scalar)}
In the setting of Proposition \hyperref[bp92]{14.17}, let \(Q\) be a nonzero subquotient of a tensor construction on \(V\) on which the element \(z\) acts by a scalar \(c_Q\). Then

\[
\operatorname{dw}_\iota(Q)=\frac{\lambda_\iota(c_Q)}{N}.
\]
\end{bpstatementit}

% proof: proof/lemma-determinant-weight-from-central-scalar
\begin{proof}[{Proof}]\phantomsection\label{bp97}%
The determinant of \(Q\) is an algebraic character of \(H\), and its image on the dense subgroup \(\rho(\pi_1(U))\) is finite, and hence its image on all of \(H\) lies in the same finite group of roots of unity. Taking determinants in \(z=h\rho(s)^N\), we find

\[
c_Q^{\dim Q}=\det(h\mid Q)\det(s\mid Q)^N.
\]

Now apply \(\lambda_\iota\), which vanishes on the first factor, and divide by \(N\dim Q\).
\end{proof}

% text: prose-14-21
\phantomsection\label{bp98}%
The lemma reduces determinant weights of tensor constructions to a computation with the scalars on the composition factors of \(V\). By Proposition \hyperref[bp92]{14.17}, the element \(z\) acts by a scalar on every simple subquotient of a tensor construction on \(V\), so the formula of the lemma applies to each of them. Since the scalars by which a single element acts multiply under tensor products and exterior powers, this gives the following formulas.

% object: proposition-determinant-weights-tensor-operations (blueprint/chap-Purity-of-real-sheaves/sec-Determinant-weights/subsec-Tensor-operations/proposition-determinant-weights-tensor-operations/_section.md)
\phantomsection\label{bp99}%
\begin{bpstatementit}{Proposition~14.19}{ (Determinant weights of tensor operations)}
Let \(V\) and \(V'\) be nonzero Weil representations, let \(A_1,\ldots,A_t\) be the composition factors of \(V\), and put \(\gamma_a=\operatorname{dw}_\iota(A_a)\).

\begin{enumerate}
\item \(\operatorname{DW}(V\otimes V')=\operatorname{DW}(V)+\operatorname{DW}(V')\).
\item \(\operatorname{DW}(\bigwedge^rV)=\bigl\{\sum_an_a\gamma_a :
      0\le n_a\le\dim A_a,\ \sum_an_a=r\bigr\}\) for \(0\le r\le\dim V\).
\item \(\operatorname{DW}(V^\vee)=-\operatorname{DW}(V)\).
\end{enumerate}
\end{bpstatementit}

% proof: proof/proposition-determinant-weights-tensor-operations
\begin{proof}[{Proof}]\phantomsection\label{bp100}%
We use one central element for \(V\) and \(V'\) simultaneously, so that its scalars on the factors of \(V\otimes V'\) are products of its scalars on the factors of \(V\) and of \(V'\). Since Proposition \hyperref[bp92]{14.17} requires semisimple representations, we apply it to \(V^{\mathrm{ss}}\oplus V'^{\mathrm{ss}}\). This changes nothing, because by Lemma \hyperref[bp78]{14.12} the simple factors of \(V\otimes V'\) and \(\bigwedge^rV\) depend only on the composition factors of \(V\) and \(V'\). Now let \(c_a\) and \(c'_b\) be the scalars by which its central element of degree \(N\) acts on the composition factors \(A_a\) of \(V\) and \(A'_b\) of \(V'\). By Lemma \hyperref[bp96]{14.18}, \(\gamma_a=\lambda_\iota(c_a)/N\) and \(\operatorname{dw}_\iota(A'_b)=\lambda_\iota(c'_b)/N\).

(1) The central element acts by \(c_ac'_b\) on \(A_a\otimes A'_b\), and hence on every simple factor of it. By the same lemma, each such factor has determinant weight \(\gamma_a+\operatorname{dw}_\iota(A'_b)\). Moreover, each \(A_a\otimes A'_b\) is nonzero, and Lemma \hyperref[bp78]{14.12} identifies the simple factors of these tensor products with those of \(V\otimes V'\).

(2) The central element acts on \(\bigotimes_a\bigwedge^{n_a}A_a\) by \(\prod_ac_a^{n_a}\). So each simple factor of this nonzero representation has determinant weight \(\sum_an_a\gamma_a\), and Lemma \hyperref[bp78]{14.12} identifies these factors with those of \(\bigwedge^rV\).

(3) We have \(\det(s\mid A^\vee)=\det(s\mid A)^{-1}\), and duality carries a composition series of \(V\) to one of \(V^\vee\) whose factors are the duals \(A_a^\vee\).
\end{proof}

% text: prose-14-22
\phantomsection\label{bp101}%
In particular every simple factor of \(V^{\otimes m}\) has determinant weight at most \(m\max\operatorname{DW}(V)\). In Section \hyperref[bp118]{14.4} we use this bound, and part (2), for the constituents of a real sheaf before we know that they are pure.



% object: sec-Rankins-bound (blueprint/chap-Purity-of-real-sheaves/sec-Rankins-bound/_section.md), content hidden
\subsection*{\texorpdfstring{14.3\quad Rankin's bound}{14.3 Rankin's bound}}
\phantomsection\label{bp102}%
\addcontentsline{toc}{subsection}{\texorpdfstring{14.3\quad Rankin's bound}{14.3 Rankin's bound}}
% text: prose-14-23
\phantomsection\label{bp103}%
We bound the local eigenvalues of a real sheaf by the eigenvalues of Frobenius on \(H^2_c\) of its even tensor powers (Proposition \hyperref[bp116]{14.25}). Only \(H^2_c\) is needed, and Lemma \hyperref[bp104]{14.20} computes it by Poincaré duality as the coinvariants of the geometric monodromy.

Rankin's argument uses one fact from each side of the trace formula. On the cohomological side, an affine curve has \(H^0_c=0\), so only \(H^2_c\) can produce poles of the \(L\)-function, and we can compute this group without knowing anything about \(H^1_c\): its eigenvalues are \(q\) times eigenvalues of Frobenius on a quotient of the generic fibre. On the side of the Euler product, if the local factors have nonnegative coefficients, then each of them is holomorphic wherever the product is.

For example, for the constant sheaf on \(\mathbf G_m\) we have

\[
L(\mathbf G_m,\overline{\mathbf Q}_\ell,t)=\prod_{x\in|\mathbf G_m|}(1-t^{\deg(x)})^{-1}
=\frac{1-t}{1-qt},
\]

where \(1-t\) comes from \(H^1_c\) and the pole at \(t=q^{-1}\) from \(H^2_c=\overline{\mathbf Q}_\ell(-1)\). Each local factor has its poles on the circle \(|t|=1\), outside the disc \(|t|<q^{-1}\), as Lemma \hyperref[bp113]{14.24} predicts. In general we apply the two facts to the even tensor powers of a real sheaf, whose \(H^2_c\) is bounded by the determinant weights of Section \hyperref[bp69]{14.2}, and this gives Proposition \hyperref[bp116]{14.25}.

% object: lemma-degree-zero-top-cohomology (blueprint/chap-Purity-of-real-sheaves/sec-Rankins-bound/lemma-degree-zero-top-cohomology/_section.md)
\phantomsection\label{bp104}%
\begin{bpstatementit}{Lemma~14.20}{ (Degree zero and top cohomology)}
Let \(U_0/\mathbf F_q\) be a smooth geometrically connected curve, let \(\mathcal F\) be a lisse Weil sheaf on \(U_0\), and let \(V\) be its geometric generic fibre, a representation of the geometric fundamental group \(\pi_1(U)\).  There are canonical Frobenius-equivariant identifications

\[
H^0(U,\mathcal F)=V^{\pi_1(U)},\qquad
H_c^2(U,\mathcal F)=V_{\pi_1(U)}(-1).
\]

There is no compactly supported cohomology outside degrees \(0,1,2\). If \(U_0\) is affine, \(H_c^0(U,\mathcal F)=0\). If \(\mathcal F\) is punctually of weights \(\le a\), the weights of \(H^0\) are at most \(a\) and those of \(H_c^2\) at most \(a+2\).  If \(\mathcal F\) is pure of weight \(a\), the two groups, when nonzero, are pure of weights \(a\) and \(a+2\), respectively.
\end{bpstatementit}

% proof: proof/lemma-degree-zero-top-cohomology
\begin{proof}[{Proof}]\phantomsection\label{bp105}%
We deduce both formulas from the description of a lisse sheaf by its representation, using Poincaré duality to reduce \(H^2_c\) to \(H^0\), and then we read off the weights from the stalks.

A global section of a local system is determined by its value at \(\bar\eta\), and a vector of the fibre extends to a global section exactly when it is fixed by the monodromy. This is the first formula.  For the second, by degreewise Poincaré duality (Corollary \bpout{10.13}) we have

\[
H_c^2(U,\mathcal F)^\vee
=H^0(U,\mathcal F^\vee(1))=(V^\vee)^{\pi_1(U)}(1).
\]

Since the dual of \((V^\vee)^{\pi_1(U)}\) is \(V_{\pi_1(U)}\), this gives the second formula, with its twist. Moreover, the identification is induced by the canonical duality pairing and is therefore Frobenius equivariant.

Duality also gives the vanishing outside degrees \(0,1,2\). Indeed, for \(i<0\) there is nothing to prove, and for \(i>2\) the space \(H_c^i(U,\mathcal F)\) is dual to \(H^{2-i}(U,\mathcal F^\vee(1))\), which is zero because \(2-i<0\). For torsion coefficients the vanishing comes instead from the bound on cohomological dimension of Appendix \bpout{D}.

Suppose that \(U_0\) is affine, and choose a smooth proper completion of \(U\). Its boundary is nonempty, because a connected proper curve has only constant global functions and hence is not affine. A compactly supported section extends by zero to the completion, so it vanishes on a neighbourhood of the boundary. Since a section of a lisse sheaf on a connected curve is determined by its generic value, the section is zero, and \(H_c^0(U,\mathcal F)=0\).

For the weights we compare the global Frobenius on these two spaces with the local Frobenius at a closed point. Since \(\pi_1(U)\) acts trivially on \(V^{\pi_1(U)}\) and on \(V_{\pi_1(U)}\), the action of \(W(U_0)\) on them factors through the degree quotient.  If \(\beta\) is an eigenvalue of the operator of degree one, then \(\beta^{\deg(x)}\) is an eigenvalue of \(\operatorname{Frob}_x\) on \(\mathcal F_{\bar x}\) for every closed point \(x\). Indeed, \(V^{\pi_1(U)}\) is a subspace and \(V_{\pi_1(U)}\) a quotient of \(V\simeq\mathcal F_{\bar x}\), both stable under \(\operatorname{Frob}_x\in\pi_1(U)s^{\deg(x)}\), which acts on them as \(s^{\deg(x)}\).  If \(\mathcal F\) is punctually of weights \(\le a\), it follows that \(|\iota(\beta)|\le q^{a/2}\), and if \(\mathcal F\) is pure of weight \(a\), equality holds. Finally, the twist \((-1)\) multiplies each global eigenvalue by \(q\), which gives the stated weights for \(H_c^2\).
\end{proof}

% text: prose-14-24
\phantomsection\label{bp106}%
For the constant sheaf on \(\mathbf G_m\) the lemma gives \(H^2_c=\overline{\mathbf Q}_\ell(-1)\), of weight \(2\), as in the computation at the start of the section. Now take a Kummer sheaf of a nontrivial character on \(\mathbf G_m\). Then \(H^0_c=0\) because \(\mathbf G_m\) is affine, and \(H^2_c=0\) because the coinvariants vanish.

The positivity we want is \(a^{2m}\ge0\) for the traces \(a\) of all powers \(\operatorname{Frob}_x^n\) on \(\mathcal G\), and for this all these traces must be real. It is not enough to ask that the trace of \(\operatorname{Frob}_x\) itself be real. For example, if the eigenvalues at \(x\) are \(i\) and \(1-i\), the trace is \(1\), but the trace of the square is \(-1+(1-i)^2=-1-2i\). By the Newton identities, on the other hand, all the power sums are real exactly when the characteristic polynomial is real, and this is the condition below. Its basic examples are the constant sheaf, its Tate twists and Kummer sheaves of quadratic characters. More generally, in Chapter \hyperref[bp122]{15}, Proposition \hyperref[bp148]{15.9} makes every pure sheaf a summand of a real one after a twist. The condition enters the proof of Theorem \hyperref[bp120]{14.26} through Lemma \hyperref[bp109]{14.22}.

% object: definition-real-sheaves (blueprint/chap-Purity-of-real-sheaves/sec-Rankins-bound/definition-real-sheaves/_section.md)
\phantomsection\label{bp107}%
\begin{bpstatement}{Definition~14.21}{ (Real sheaves)}
A Weil sheaf \(\mathcal G\) is \(\iota\)-real if \(\iota\det(1-T\operatorname{Frob}_x\mid\mathcal G_{\bar x})\) has real coefficients for every closed point; equivalently, the multiset of its complex eigenvalues at each point is stable under complex conjugation.  In particular, the traces of all positive powers of local Frobenius are real.
\end{bpstatement}

% text: prose-14-25
\phantomsection\label{bp108}%
The next lemma is the \(\ell\)-adic form of Rankin's trick. Since the even powers of a real trace are nonnegative, the Euler factors of \(\mathcal G^{\otimes2m}\) have nonnegative coefficients.

% object: lemma-positivity-euler-factors (blueprint/chap-Purity-of-real-sheaves/sec-Rankins-bound/lemma-positivity-euler-factors/_section.md)
\phantomsection\label{bp109}%
\begin{bpstatementit}{Lemma~14.22}{ (Positivity of Euler factors)}
Let \(\mathcal G\) be an \(\iota\)-real Weil sheaf, let \(m\ge1\), and let \(x\) be a closed point. Then the power series \(\iota\det(1-\operatorname{Frob}_xt^{\deg(x)}\mid\mathcal G^{\otimes2m}_{\bar x})^{-1}\) has nonnegative real coefficients and constant coefficient one.
\end{bpstatementit}

% proof: proof/lemma-positivity-euler-factors
\begin{proof}[{Proof}]\phantomsection\label{bp110}%
We write an inverse characteristic polynomial as the exponential of a series in the traces of powers, whose coefficients are then visibly nonnegative. For a square matrix \(A\), triangularizing \(A\) and expanding the logarithm of each factor we get the formal identity

\[
\det(1-tA)^{-1}
=\exp\left(\sum_{n\ge1}\operatorname{Tr}(A^n)\frac{t^n}{n}\right).
\]

On an even tensor power, the trace of the \(n\)th power of Frobenius is \(\operatorname{Tr}(\operatorname{Frob}_x^n\mid\mathcal G)^{2m}\). Since \(\mathcal G\) is real, this is nonnegative after applying \(\iota\).  Every coefficient of the exponential is then a finite sum of products of nonnegative numbers. This remains true after replacing \(t\) by \(t^{\deg(x)}\).
\end{proof}

% object: example-even-tensor-powers (blueprint/chap-Purity-of-real-sheaves/sec-Rankins-bound/example-even-tensor-powers/_section.md)
\phantomsection\label{bp111}%
\begin{bpstatement}{Example~14.23}{ (Why even tensor powers)}
The sheaf \(\mathcal B_{-1}\) is real, since \(\operatorname{Frob}_x\) acts by \((-1)^{\deg(x)}\), but its Euler factor at a point of degree one is \((1+t)^{-1}=1-t+t^2-\cdots\), with negative coefficients. Its square \(\mathcal B_{-1}\otimes\mathcal B_{-1}\) is the constant sheaf, whose Euler factors \((1-t^{\deg(x)})^{-1}\) have nonnegative coefficients, as the lemma predicts for \(m=1\).
\end{bpstatement}

% text: prose-14-26
\phantomsection\label{bp112}%
Without positivity a factor can converge on a smaller disc than the product, as in the cancellation at the start of this chapter. With positivity, however, the coefficients of each factor are bounded by those of the product.

% object: lemma-positivity-convergence (blueprint/chap-Purity-of-real-sheaves/sec-Rankins-bound/lemma-positivity-convergence/_section.md)
\phantomsection\label{bp113}%
\begin{bpstatementit}{Lemma~14.24}{ (Positivity and convergence)}
Let

\[
L(t)=\prod_j f_j(t),\qquad f_j\in1+t\mathbf R_{\ge0}[[t]],
\]

be a coefficientwise locally finite product. If \(L\) is the Taylor series of a function holomorphic on \(|t|<R\), then every \(f_j\) has radius of convergence at least \(R\), and if \(f_j=P_j^{-1}\) with \(P_j(0)=1\), then \(P_j\) has no zero in that disc.
\end{bpstatementit}

% proof: proof/lemma-positivity-convergence
\begin{proof}[{Proof}]\phantomsection\label{bp114}%
We compare coefficients and then apply Cauchy's estimate. To compare, write \(f_j=\sum_na_{j,n}t^n\) and \(L=\sum_nA_nt^n\). Since the product is coefficientwise locally finite, \(A_n\) is already the coefficient of \(t^n\) in a finite subproduct containing \(f_j\), which is \(f_j\) times a series \(g\) with constant coefficient one and nonnegative coefficients. The coefficient of \(t^n\) in \(f_jg\) is \(a_{j,n}\) plus a sum of products of nonnegative numbers, so that \(0\le a_{j,n}\le A_n\); positivity enters only here.  Next, for every \(0<r<R\), Cauchy's estimate gives \(|A_n|\le C_rr^{-n}\). Hence the radius of convergence of \(f_j\) is at least \(r\) for every such \(r\), and so at least \(R\).  A zero of \(P_j\) in the disc would be a pole of \(f_j=1/P_j\), which is impossible.
\end{proof}

% text: prose-14-27
\phantomsection\label{bp115}%
In the proposition below, the twist \((-1)\) on \(H^2_c\) costs a factor \(q\), and this loss disappears when we take \(2m\)th roots and let \(m\) grow.

% object: proposition-bound-largest-determinant-weight (blueprint/chap-Purity-of-real-sheaves/sec-Rankins-bound/proposition-bound-largest-determinant-weight/_section.md)
\phantomsection\label{bp116}%
\begin{bpstatementit}{Proposition~14.25}{ (A bound by the largest determinant weight)}
Let \(U_0/\mathbf F_q\) be a smooth geometrically connected affine curve and let \(\mathcal G\ne0\) be a lisse \(\iota\)-real Weil sheaf.  If \(\beta\) is the largest determinant weight of its irreducible Weil constituents, then

\[
|\iota(\alpha)|^2\le q_x^\beta
\quad(x\in|U_0|,\ \alpha\text{ an eigenvalue on }\mathcal G_{\bar x}).
\]
\end{bpstatementit}

% proof: proof/proposition-bound-largest-determinant-weight
\begin{proof}[{Proof}]\phantomsection\label{bp117}%
We bound the poles of \(L(U_0,\mathcal G^{\otimes2m},t)\) by determinant weights through \(H^2_c\), and positivity then passes the bound to each local factor.

We may assume that \(\mathcal G\) is semisimple. Indeed, by Lemma \hyperref[bp19]{13.5}, replacing \(\mathcal G\) by its semisimplification changes neither its local characteristic polynomials nor the determinant weights of its constituents.

\emph{The top cohomology.} Fix \(m\ge1\) and put \(\mathcal T=\mathcal G^{\otimes2m}\). We show that every eigenvalue \(\mu\) of \(F\) on \(H_c^2(U,\mathcal T)\) satisfies \(|\iota(\mu)|\le q^{m\beta+1}\); these eigenvalues determine the possible poles of the \(L\)-function. First, every constituent of \(\mathcal T\) has determinant weight at most \(2m\beta\), by Proposition proposition-determinant-weights-tensor-operations. Next, Lemma \hyperref[bp104]{14.20} identifies \(H_c^2(U,\mathcal T)\) with \(T_{\pi_1(U)}(-1)\), where \(T\) is the generic fibre of \(\mathcal T\). The geometric fundamental group acts trivially on \(T_{\pi_1(U)}\), so its simple Weil factors are one-dimensional, as one sees by triangularizing the operator induced by a degree-one lift. Let \(b\) be the eigenvalue on such a factor. Since \(T_{\pi_1(U)}\) is a Weil quotient of \(T\), the factor is a constituent of \(\mathcal T\), and its determinant weight is \(\lambda_\iota(b)\). From \(\lambda_\iota(b)=2\log_q|\iota(b)|\le2m\beta\) we get \(|\iota(b)|\le q^{m\beta}\). The twist \((-1)\) multiplies eigenvalues by \(q\), and the claim follows.

\emph{The local factors.} The bound on \(H^2_c\) makes the \(L\)-function holomorphic on a disc, and positivity passes this to each local factor. As \(H_c^0=0\) by Lemma \hyperref[bp104]{14.20}, Theorem \bpout{12.30} gives

\[
L(U_0,\mathcal T,t)
=\prod_{x\in|U_0|}\det(1-\operatorname{Frob}_x t^{\deg(x)}\mid\mathcal T_{\bar x})^{-1}
=\frac{\det(1-Ft\mid H_c^1(U,\mathcal T))}
{\det(1-Ft\mid H_c^2(U,\mathcal T))}.
\]

This is an identity of formal power series, and the product is coefficientwise locally finite because there are only finitely many closed points of bounded degree. Now put \(R_m=q^{-m\beta-1}\). By the bound on \(H_c^2\), the denominator has no zero on the disc \(|t|<R_m\) after applying \(\iota\), so the rational function, and with it its Taylor series, is holomorphic on this disc.

Each local factor \(\det(1-\operatorname{Frob}_x t^{\deg(x)}\mid\mathcal T_{\bar x})\) has no zero on the same disc. Indeed, after applying \(\iota\) the local factors of \(L(U_0,\mathcal T,t)\) are power series with nonnegative coefficients and constant coefficient one by Lemma \hyperref[bp109]{14.22}, so that Lemma \hyperref[bp113]{14.24} applies.

Now let \(\alpha\) be a local eigenvalue of \(\mathcal G\). Then \(\alpha^{2m}\) is an eigenvalue on \(\mathcal T\) by Lemma \hyperref[bp19]{13.5}, and \(1-\iota(\alpha)^{2m}t^{\deg(x)}\) divides the local factor. Its roots have absolute value \(|\iota(\alpha)|^{-2m/\deg(x)}\), so by the previous paragraph this number must be at least \(R_m\), that is,

\[
|\iota(\alpha)|^2\le q_x^{\beta+1/m}.
\]

Letting \(m\) tend to infinity proves the bound.
\end{proof}


% object: sec-Exterior-powers-and-purity (blueprint/chap-Purity-of-real-sheaves/sec-Exterior-powers-and-purity/_section.md), content hidden
\subsection*{\texorpdfstring{14.4\quad Exterior powers and purity}{14.4 Exterior powers and purity}}
\phantomsection\label{bp118}%
\addcontentsline{toc}{subsection}{\texorpdfstring{14.4\quad Exterior powers and purity}{14.4 Exterior powers and purity}}
% text: prose-14-28
\phantomsection\label{bp119}%
Proposition \hyperref[bp116]{14.25} only bounds the local eigenvalues by the largest determinant weight. To bound the eigenvalues of each constituent by its own determinant weight, we apply the proposition to exterior powers, which separate the determinant weights of the simple factors.

Tensor powers cannot do this. For example, suppose \(\mathcal G=\mathcal G_1\oplus\mathcal G_2\) with determinant weights \(\gamma_1>\gamma_2\). In \(\mathcal G^{\otimes m}\) the largest determinant weight \(m\gamma_1\) belongs to the factors of \(\mathcal G_1^{\otimes m}\), while every factor involving \(\mathcal G_2\) lies strictly below it. So the bound of Proposition \hyperref[bp116]{14.25} says nothing sharp about \(\mathcal G_2\).

An exterior power, by contrast, uses each basis vector of \(\mathcal G_1\) at most once. Thus if \(r_1\) is the rank of \(\mathcal G_1\), every graded piece of \(\bigwedge^{r_1+1}\mathcal G\) involves \(\mathcal G_2\), and the summand \(\det\mathcal G_1\otimes\mathcal G_2\) has the largest determinant weight \(r_1\gamma_1+\gamma_2\). Since Corollary \hyperref[bp65]{14.8} gives the eigenvalues of \(\det\mathcal G_1\), a bound for this summand is a bound for \(\mathcal G_2\). We argue on an affine curve, where the proposition applies, and then pass to an arbitrary curve by restricting to affine opens, which preserves irreducibility.

% object: theorem-purity-real-sheaves-curves (blueprint/chap-Purity-of-real-sheaves/sec-Exterior-powers-and-purity/theorem-purity-real-sheaves-curves/_section.md)
\phantomsection\label{bp120}%
\begin{bpstatementit}{Theorem~14.26}{ (Purity of real sheaves on curves)}
On a smooth geometrically connected curve over \(\mathbf F_q\), every irreducible constituent of an \(\iota\)-real lisse Weil sheaf is punctually \(\iota\)-pure of weight equal to its determinant weight. In particular, every \(\iota\)-real lisse Weil sheaf on the curve is mixed.
\end{bpstatementit}

% proof: proof/theorem-purity-real-sheaves-curves
\begin{proof}[{Proof}]\phantomsection\label{bp121}%
We group the constituents by determinant weight and place each group, tensored with the determinants of the groups of larger weight, in an exterior power where it carries the largest determinant weight. There Proposition \hyperref[bp116]{14.25} bounds the eigenvalues of the group by its own determinant weight, and the same argument applied to the dual gives the reverse bound.

Let \(\mathcal G\) be an \(\iota\)-real lisse Weil sheaf on the curve. We first prove the assertion when the curve is affine. There, replacing \(\mathcal G\) by its semisimplification changes neither its local characteristic polynomials nor its irreducible constituents, and it keeps the sheaf real, so we may assume that \(\mathcal G\) is semisimple and write

\[
\mathcal G=\bigoplus_{i=1}^t\mathcal G_i,
\qquad \gamma_1>\cdots>\gamma_t,
\qquad r_i=\operatorname{rank}\mathcal G_i,
\]

where \(\mathcal G_i\) is the sum of the irreducible constituents with determinant weight \(\gamma_i\).

Fix \(i\), and let us find the exterior power in which \(\mathcal G_i\), together with the determinants of \(\mathcal G_1,\dots,\mathcal G_{i-1}\), carries the largest determinant weight. To this end set \(m=\sum_{j<i}r_j\), and put \(B=\sum_{j<i}r_j\gamma_j\).  By Proposition proposition-determinant-weights-tensor-operations, the largest determinant weight of a constituent of \(\bigwedge^{m+1}\mathcal G\) is \(B+\gamma_i\).  Indeed, it is obtained by taking \(n_j=r_j\) for every \(j<i\) and \(n_i=1\), and every other choice gives a smaller sum, because the \(\gamma_j\) are strictly decreasing.  Moreover, exterior powers of a real sheaf are real. Indeed, their eigenvalues are the products of eigenvalues with distinct indices, and complex conjugation permutes this multiset.  Hence Proposition \hyperref[bp116]{14.25} applies to \(\bigwedge^{m+1}\mathcal G\).

The eigenvalues of \(\mathcal G_i\) appear in \(\bigwedge^{m+1}\mathcal G\) not by themselves but multiplied by the eigenvalues of the determinants of the groups of larger weight, and those eigenvalues are known. To make this precise, let \(x\) be a closed point of the curve and \(\alpha\) an eigenvalue of \(\operatorname{Frob}_x\) on \((\mathcal G_i)_{\bar x}\). The decomposition of the exterior power of a direct sum shows that \(\bigwedge^{m+1}\mathcal G\) contains the summand

\[
(\det\mathcal G_1\otimes\cdots\otimes\det\mathcal G_{i-1})
\otimes\mathcal G_i.
\]

If \(\beta_j\) is the eigenvalue of \(\operatorname{Frob}_x\) on \((\det\mathcal G_j)_{\bar x}\), then the product \(\alpha\beta_1\cdots\beta_{i-1}\) is an eigenvalue of \(\operatorname{Frob}_x\) on the stalk of this summand at \(\bar x\). Since \(\det\mathcal G_j\) is pure of weight \(r_j\gamma_j\), we have \(|\iota(\beta_1\cdots\beta_{i-1})|^2=q_x^B\). Now Proposition \hyperref[bp116]{14.25} gives

\[
|\iota(\alpha)|^2q_x^B\le q_x^{B+\gamma_i},
\qquad |\iota(\alpha)|^2\le q_x^{\gamma_i}.
\]

For the reverse bound, note that the dual \(\mathcal G^\vee\) is real as well, with determinant weights \(-\gamma_t>\cdots>-\gamma_1\). Applying the same argument to the summand \(\mathcal G_i^\vee\), we get \(|\iota(\alpha^{-1})|^2\le q_x^{-\gamma_i}\), and hence \(|\iota(\alpha)|^2=q_x^{\gamma_i}\). Since every irreducible constituent of \(\mathcal G\) is a summand of exactly one \(\mathcal G_i\), this proves the theorem on an affine curve.

On an arbitrary curve we restrict to an affine neighbourhood of a given closed point \(x\), which suffices because purity is a condition at each closed point. It remains to check that the restriction changes neither the constituents nor their determinant weights. The neighbourhood is a nonempty open subscheme of a smooth geometrically integral curve, so it is geometrically connected, and restriction to it preserves irreducibility. Indeed, a connected finite étale cover of a normal connected curve is integral, and hence remains connected after restriction to a dense open. So the maps of geometric and of arithmetic fundamental groups are surjective, and since restriction preserves the degree map, the map of Weil groups is surjective too. Hence a subspace of the fibre is stable on the open if and only if it is stable on the curve.

The restriction of a real sheaf is real, and by Corollary \hyperref[bp61]{14.6} the determinant keeps its weight. The affine case therefore applies at \(x\), and since \(x\) was arbitrary, we obtain purity on the whole curve. Finally, the sheaf is mixed because a composition series is a finite filtration with pure quotients.
\end{proof}



% object: chap-Pure-sheaves-on-curves (blueprint/chap-Pure-sheaves-on-curves/_section.md), content hidden
\section*{\texorpdfstring{15\quad Pure sheaves on curves}{15 Pure sheaves on curves}}
\phantomsection\label{bp122}%
\addcontentsline{toc}{section}{\texorpdfstring{15\quad Pure sheaves on curves}{15 Pure sheaves on curves}}
% text: prose-15-01
\phantomsection\label{bp123}%
How much does the trace function of a sheaf know about its Frobenius eigenvalues? For a sheaf of rank one, pure of weight \(w\), on a dense open \(U_0\) of the line, it knows the weight. Indeed, every trace over \(\mathbf F_{q^n}\) has absolute value \(q^{nw/2}\), so the sum \(S_n\) of the \(|t(x)|^2\) over \(U_0(\mathbf F_{q^n})\) is \(q^{nw}\#U_0(\mathbf F_{q^n})\), about \(q^{n(w+1)}\), and the growth rate \(q^{w+1}\) of these sums therefore determines \(w\). In higher rank, however, a trace is a sum of eigenvalues, and the eigenvalues can cancel. For example, on \(\mathcal B_b\oplus\mathcal B_{-b}\) the trace at a point of degree \(n\) is \(b^n+(-b)^n\), which is zero for every odd \(n\). We show that the cancellation cannot persist: for a nonzero pure lisse sheaf of weight \(w\) on an affine curve, \(\limsup_n\frac1n\log_qS_n=w+1\) (Theorem \hyperref[bp162]{15.14}).

The reason is positivity once more. Recall that a sheaf \(\mathcal G\) is real if the characteristic polynomials of Frobenius on its stalks have real coefficients under \(\iota\) (Definition \hyperref[bp107]{14.21}). For such \(\mathcal G\), pure of weight \(w\), \(|t_{\mathcal G}|^2\) is the trace function of \(\mathcal G\otimes\mathcal G\), and hence the generating series of the \(S_n\) is the logarithmic derivative of \(L(\mathcal G\otimes\mathcal G,t)\). Positivity of the Euler factors (Lemma \hyperref[bp109]{14.22}) shows that \(L(\mathcal G\otimes\mathcal G,t)\) is not a polynomial, and Lemma \hyperref[bp104]{14.20} puts its poles on \(|t|=q^{-w-1}\). Together these force the sums \(S_n\) to reach the size \(q^{n(w+1)}\) infinitely often. The restriction to real sheaves costs nothing, since an arbitrary pure sheaf becomes real after a twist to weight zero and the addition of its dual (Proposition \hyperref[bp148]{15.9}).

This is what the Fourier argument of Chapter \hyperref[bp223]{17} needs. That argument starts from a pure lisse sheaf \(M\) of weight \(w\) on a dense open \(j:U_0\hookrightarrow\mathbf A^1\) and produces a pure lisse sheaf \(M'\) on a dense open \(j'\) of the dual line. It knows the trace function of \(M'\) exactly, as the discrete Fourier transform of that of \(M\), but it cannot see the weight of \(M'\) directly. Since Plancherel's formula multiplies the sums \(S_n\) by \(q^n\), the theorem shows that \(M'\) has weight \(w+1\).

The group we finally want to bound, \(H^1_c(\mathbf A^1,j_*M)\), is the stalk of \(j'_*M'\) at the origin. But the origin may be one of the points where \(M'\) is not lisse, and there we need to know that \(j'_*\) does not create larger weights (Proposition \hyperref[bp134]{15.5}).

% object: sec-Weights-of-middle-extensions (blueprint/chap-Pure-sheaves-on-curves/sec-Weights-of-middle-extensions/_section.md), content hidden
\subsection*{\texorpdfstring{15.1\quad Weights of middle extensions}{15.1 Weights of middle extensions}}
\phantomsection\label{bp124}%
\addcontentsline{toc}{subsection}{\texorpdfstring{15.1\quad Weights of middle extensions}{15.1 Weights of middle extensions}}
% text: prose-15-02
\phantomsection\label{bp125}%
Let \(j\) be the inclusion of a dense open in a smooth curve, and let \(\mathcal F\) be a lisse sheaf on the open. The Fourier argument needs an extension of \(\mathcal F\) to the whole curve that loses nothing at the boundary points. The ordinary direct image \(j_*\mathcal F\) is such an extension, whereas extension by zero loses something. For example, for the constant sheaf on \(\mathbf G_m\subset\mathbf A^1\), \(j_!\overline{\mathbf Q}_\ell\) has stalk \(0\) at the origin, while the constant sheaf on \(\mathbf A^1\) is the natural extension. In general, by Lemma \hyperref[bp127]{15.2}, \(j_*\mathcal F\) recovers a lisse extension whenever one exists, and its stalk at a boundary point is the space of inertia invariants.

If \(\mathcal F\) is punctually of weights \(\le a\), then so is \(j_*\mathcal F\), at the boundary points as well (Proposition \hyperref[bp134]{15.5}). In Chapter \hyperref[bp223]{17} this bounds the weights of the boundary terms in the Fourier transform of a middle extension.

% object: definition-middle-extension (blueprint/chap-Pure-sheaves-on-curves/sec-Weights-of-middle-extensions/definition-middle-extension/_section.md)
\phantomsection\label{bp126}%
\begin{bpstatement}{Definition~15.1}{ (Middle extension)}
Let \(j:U\hookrightarrow X\) be a dense open in a smooth curve and let \(\mathcal F\) be a lisse sheaf on \(U\). The \emph{middle extension} of \(\mathcal F\) is the ordinary direct image \(j_*\mathcal F\) of sheaves.
\end{bpstatement}

% object: lemma-boundary-stalk (blueprint/chap-Pure-sheaves-on-curves/sec-Weights-of-middle-extensions/lemma-boundary-stalk/_section.md)
\phantomsection\label{bp127}%
\begin{bpstatementit}{Lemma~15.2}{ (The boundary stalk)}
Let \(j:U_0\hookrightarrow X_0\) be a dense open in a smooth curve, let \(\mathcal F\) be a lisse Weil sheaf on \(U_0\), and let \(j_*\) be the ordinary direct image of sheaves. Let \(\bar z\) be a geometric point over a point \(z\) of \(X_0-U_0\), and let \(I_z\) be the inertia group at \(z\), acting on the local generic fibre \(V\) of \(\mathcal F\). There is a canonical Frobenius-equivariant identification

\[
(j_*\mathcal F)_{\bar z}=V^{I_z}.
\]

If \(\mathcal F\) extends to a lisse sheaf on \(X_0\), the adjunction map from that extension to \(j_*\mathcal F\) is an isomorphism.
\end{bpstatementit}

% proof: proof/lemma-boundary-stalk
\begin{proof}[{Proof}]\phantomsection\label{bp128}%
We prove the formula for finite coefficients, where it is Galois theory of the punctured henselian neighbourhood, and pass to the limit.

The stalk of \(j_*\mathcal F\) at \(\bar z\) is the colimit of the sections of \(\mathcal F\) over the inverse images of \(U\) in étale neighbourhoods of \(\bar z\), that is, the group of sections over the punctured spectrum of the strict henselization at \(\bar z\). This punctured spectrum is the spectrum of a field, whose absolute Galois group is the inertia group \(I_z\), so finite locally constant sheaves there correspond to finite continuous representations of the inertia group, and their sections are the invariant vectors. Applying this to the finite reductions of a stable lattice and passing to the inverse limit, we obtain the formula with integral coefficients.

Two compatibilities take us to rational coefficients and to \(\overline{\mathbf Q}_\ell\). Rationalization commutes with taking invariants, since a power of the uniformizer carries an invariant rational vector into the lattice.  A finite extension of scalars also commutes with invariants. Indeed, it is exact, hence commutes with finite intersections of kernels, and in a finite-dimensional space the intersection of the kernels of all \(g-1\) is the intersection of finitely many of them. This gives the formula with rational coefficients, after finite extension of scalars, and hence also over \(\overline{\mathbf Q}_\ell\). Since all these maps commute with the action of the decomposition group, the identification is Frobenius equivariant. Finally, if \(\mathcal F\) extends to a lisse sheaf on \(X_0\), the local inertia acts trivially, and so the adjunction map is an isomorphism on every geometric stalk.
\end{proof}

% object: example-boundary-stalks-g-m (blueprint/chap-Pure-sheaves-on-curves/sec-Weights-of-middle-extensions/example-boundary-stalks-g-m/_section.md)
\phantomsection\label{bp129}%
\begin{bpstatement}{Example~15.3}{ (Boundary stalks on $\mathbf G_m\subset\mathbf A^1$)}
Let \(j:\mathbf G_{m,k_0}\hookrightarrow\mathbf A^1_{k_0}\) and \(z=0\). The covers \(u^n=t\) are totally ramified at \(0\), so the inertia \(I_0\) maps onto the monodromy of a Kummer sheaf. For the Kummer sheaf of a nontrivial character, \(I_0\) acts on the line \(V\) through that character. Hence \(V^{I_0}=0\), and the stalk of \(j_*\) of this sheaf at \(\bar0\) is zero, so here \(j_*\) and \(j_!\) agree. For the Kummer extension \(\mathcal E\) of Example \hyperref[bp17]{13.4}, on the other hand, \(I_0\) acts by the matrices \(\bigl(\begin{smallmatrix}1&\kappa\\0&1\end{smallmatrix}\bigr)\) with \(\kappa\) running over \(\mathbf Z_\ell(1)\), whose common fixed vectors form the line \(\overline{\mathbf Q}_\ell(1)\). So the stalk of \(j_*\mathcal E\) at \(0\) is \(\overline{\mathbf Q}_\ell(1)\), of weight \(-2\), while \(\mathcal E\) has weights \(-2\) and \(0\); the stalk at the boundary keeps only part of the weights.
\end{bpstatement}

% text: prose-15-03
\phantomsection\label{bp130}%
To bound \(j_*\mathcal F\) we compare the \(L\)-functions of \(\mathcal F\) and of \(j_*\mathcal F\) on a disc. To see that the Euler product converges there, we need a crude bound for the number of points of a curve.

% object: lemma-point-count-bound (blueprint/chap-Pure-sheaves-on-curves/sec-Weights-of-middle-extensions/lemma-point-count-bound/_section.md)
\phantomsection\label{bp131}%
\begin{bpstatementit}{Lemma~15.4}{ (A point-count bound)}
For a smooth affine curve \(U_0/k_0\) there is a constant \(D\) such that \(\#U_0(k_n)\le Dq^n\) for all \(n\ge1\).  For a finite \(k_0\)-scheme, the number of \(k_n\)-points is bounded independently of \(n\).
\end{bpstatementit}

% proof: proof/lemma-point-count-bound
\begin{proof}[{Proof}]\phantomsection\label{bp132}%
We map \(U_0\) finitely and flatly onto the affine line, whose number of \(k_n\)-points is \(q^n\), and bound the size of the fibres.

By Noether normalization [Stacks, \href{https://stacks.math.columbia.edu/tag/00OY}{00OY}] there is a finite map \(U_0\to\mathbf A^1_{k_0}\), and the coordinate ring of \(U_0\) is a finite torsion-free module over \(k_0[t]\). To see this, let \(f\ne0\) in \(k_0[t]\) and \(fa=0\). On each irreducible component of \(U_0\) the image of \(f\) is nonzero, since otherwise the component would be finite over \(k_0[t]/(f)\) and hence of dimension zero. Since the coordinate ring of each component is a domain, \(fa=0\) forces \(a\) to vanish on every component, and then \(a=0\) because \(U_0\) is reduced.  As \(k_0[t]\) is a principal ideal domain, it follows that the coordinate ring of \(U_0\) is free of some rank \(D\).  Hence, after any extension of the base field, every fibre has a coordinate ring of dimension \(D\), and by the Chinese remainder theorem it has at most \(D\) rational points.  Summing over \(\mathbf A^1(k_n)\), we obtain the first bound.  The second holds because the number of rational points of a finite scheme is at most its number of geometric points.
\end{proof}

% text: prose-15-04
\phantomsection\label{bp133}%
In the next proof we use the lemma to show that the logarithm of the Euler product of \(\mathcal F\) converges on a disc of radius \(q^{-a/2-1}\), and on this disc we compare it with \(j_*\mathcal F\).

% object: proposition-weights-j-boundary-points (blueprint/chap-Pure-sheaves-on-curves/sec-Weights-of-middle-extensions/proposition-weights-j-boundary-points/_section.md)
\phantomsection\label{bp134}%
\begin{bpstatementit}{Proposition~15.5}{ (Weights of $j_*$ at the boundary points)}
Let \(j:U_0\hookrightarrow X_0\) be a dense open of a smooth geometrically connected curve.  If a lisse Weil sheaf \(\mathcal F\) on \(U_0\) is punctually of weights \(\le a\), so is \(j_*\mathcal F\) on \(X_0\).
\end{bpstatementit}

% proof: proof/proposition-weights-j-boundary-points
\begin{proof}[{Proof}]\phantomsection\label{bp135}%
Comparing the \(L\)-functions of \(\mathcal F\) and of \(j_*\mathcal F\) gives the weaker bound \(a+2\) at the boundary, and tensor powers remove the extra \(2\).

The assertion is local at a point of \(X_0-U_0\), so we may replace \(X_0\) by an affine neighbourhood of that point. The complement of \(U_0\) is an effective Cartier divisor, so after shrinking this neighbourhood it is cut out by one equation \(f\), and then \(U_0=D(f)\) is affine as well. Write \(i:Z_0=X_0-U_0\hookrightarrow X_0\). The hypotheses are preserved by this replacement. Indeed, forming \(j_*\) commutes with restriction to open subschemes, and the restriction of \(\mathcal F\) is still punctually of weights \(\le a\). Moreover, a nonempty open of \(X_0\) is still geometrically connected.

The boundary Euler factors of \(j_*\mathcal F\) are the ratio \(L_X/L_U\) of the two \(L\)-functions defined below. We show that \(L_U\) has no zero and \(L_X\) no pole on the disc \(|t|<q^{-a/2-1}\), so that the boundary factors have no pole there, and this bounds their eigenvalues.

\emph{The \(L\)-function on \(U_0\).} Put \(L_U(t)=L(U_0,\mathcal F,t)\). We show that \(\iota L_U\) is holomorphic and nonvanishing on the disc, by bounding the coefficients of its logarithm. The logarithm of the Euler product has coefficients

\[
\log\iota L_U(t)=\sum_{n\ge1}\frac{c_n}{n}t^n,
\qquad
c_n=\sum_{d\mid n}d\sum_{\deg(x)=d}
\iota\operatorname{Tr}(\operatorname{Frob}_x^{n/d}\mid\mathcal F_{\bar x}).
\]

Let \(r=\operatorname{rank}\mathcal F\). For a closed point \(x\) of degree \(d\mid n\), the bound on the weights gives \(|\iota\operatorname{Tr}(\operatorname{Frob}_x^{n/d}\mid\mathcal F_{\bar x})|
\le rq_x^{(n/d)a/2}=rq^{na/2}\). Since \(\sum_{d\mid n}d\cdot\#\{x\in|U_0|:\deg(x)=d\}=\#U_0(k_n)\), Lemma \hyperref[bp131]{15.4} gives \(|c_n|\le rDq^{n(a/2+1)}\). Hence the logarithm converges on the disc \(|t|<R=q^{-a/2-1}\), and its exponential is a nonvanishing holomorphic function there. By the trace formula, its Taylor series is that of the rational function \(\iota L_U(t)\), so the two agree on the whole disc.

\emph{The \(L\)-function on \(X_0\).} Let \(L_X(t)=L(X_0,j_*\mathcal F,t)\). We show that \(\iota L_X\) is holomorphic on the same disc. By the trace formula its poles come from \(H^0_c\) and \(H^2_c\) of \(j_*\mathcal F\), and we compute both groups.

The group \(H^0_c\) vanishes. Indeed, the support of a compactly supported section of \(j_*\mathcal F\) on \(X\) is closed and proper over \(k\), and since \(X\) is an affine curve, the support is finite. So the section vanishes on a dense open subset of \(U\), and as \(\mathcal F\) is lisse and \(U\) is connected, it vanishes on all of \(U\). Since a section of \(j_*\mathcal F\) is determined by its restriction to \(U\), the section is zero. For \(H^2_c\), note that the quotient sheaf in

\[
0\longrightarrow j_!\mathcal F\longrightarrow j_*\mathcal F
\longrightarrow i_*i^*j_*\mathcal F\longrightarrow0
\]

is supported on the finite boundary, whose cohomology is concentrated in degree zero.  From the long exact sequence we obtain \(H_c^2(X,j_*\mathcal F)=H_c^2(U,\mathcal F)\), and by Lemma \hyperref[bp104]{14.20} its eigenvalues have absolute value at most \(q^{a/2+1}\).  By the trace formula, every pole of \(\iota L_X\) is the inverse of a Frobenius eigenvalue on \(H^0_c\) or \(H^2_c\). Since \(H^0_c=0\), the poles have absolute value at least \(q^{-a/2-1}=R\). Hence \(\iota L_X\) is holomorphic on \(|t|<R\).

\emph{The boundary factors.} On \(U_0\) the sheaves \(\mathcal F\) and \(j_*\mathcal F\) have the same stalks, so the two Euler products differ only by the factors at the points of \(Z_0\), and we get

\[
L_X(t)=L_U(t)\prod_{z\in Z_0}
\det(1-\operatorname{Frob}_z t^{\deg(z)}\mid(j_*\mathcal F)_{\bar z})^{-1}.
\]

Since \(L_U\) has no zero on the disc, the finite product on the right is holomorphic there. As it is the reciprocal of a product of polynomials with constant term one, none of these polynomials has a zero on the disc.  Now, for an eigenvalue \(\alpha\) at \(z\), the roots of \(1-\iota(\alpha)t^{\deg(z)}\) have absolute value \(|\iota(\alpha)|^{-1/\deg(z)}\). This number is therefore at least \(R\), which is the preliminary bound

\[
|\iota(\alpha)|^2\le q_z^{a+2}.
\]

\emph{Tensor powers.} The error \(2\) in the exponent comes from the twist \((-1)\) on \(H^2_c\). Since it does not grow with the weight, it becomes negligible for high tensor powers. We therefore apply the preliminary bound to \(\mathcal F^{\otimes m}\) for every \(m\ge1\), a sheaf punctually of weights \(\le ma\).  Let \(V\) be the local generic fibre of \(\mathcal F\) at \(z\). By Lemma \hyperref[bp127]{15.2}, the eigenvalue \(\alpha\) of \(\operatorname{Frob}_z\) on \((j_*\mathcal F)_{\bar z}\) has an eigenvector \(v\in V^{I_z}\), and then \(v^{\otimes m}\) is a nonzero invariant eigenvector in \(V^{\otimes m}\) with eigenvalue \(\alpha^m\). By the same lemma, \(\alpha^m\) is therefore an eigenvalue on the stalk of \(j_*(\mathcal F^{\otimes m})\) at \(\bar z\), and the preliminary bound for \(\mathcal F^{\otimes m}\) gives \(|\iota(\alpha)|^{2m}\le q_z^{ma+2}\).  Taking \(m\)th roots and letting \(m\) tend to infinity, we obtain \(|\iota(\alpha)|^2\le q_z^a\) at every boundary point. On \(U_0\), finally, the bound is the hypothesis.
\end{proof}

% text: prose-15-05
\phantomsection\label{bp136}%
The proposition gives only an upper bound at the boundary, and in general nothing better is true. For example, for a family of elliptic curves with multiplicative reduction at \(z\), the first cohomology of the fibres is pure of weight \(1\), while its inertia invariants at \(z\) form a line of weight \(0\). When the inertia acts trivially, however, we get the lower bound by duality.

% object: corollary-unramified-boundary-point (blueprint/chap-Pure-sheaves-on-curves/sec-Weights-of-middle-extensions/corollary-unramified-boundary-point/_section.md)
\phantomsection\label{bp137}%
\begin{bpstatementit}{Corollary~15.6}{ (An unramified boundary point)}
Let \(\mathcal F\) be a lisse sheaf, pure of weight \(a\), on a smooth geometrically connected curve with a point removed. If the inertia at the removed point acts trivially, then the stalk at that point of the direct image of \(\mathcal F\) to the curve is pure of weight \(a\). In particular, a lisse Weil sheaf on a smooth geometrically connected curve that is pure on a dense open is pure on the whole curve.
\end{bpstatementit}

% proof: proof/corollary-unramified-boundary-point
\begin{proof}[{Proof}]\phantomsection\label{bp138}%
Purity follows from the upper bounds for the sheaf and for its dual, provided the stalk of the dual is the dual of the stalk, and trivial inertia guarantees this. So we apply Proposition \hyperref[bp134]{15.5} to the sheaf and to its dual, which is pure of weight \(-a\).  Since the inertia acts trivially, the stalk of the extended dual is the dual of the stalk of the extension, and so every eigenvalue \(\alpha\) satisfies \(|\iota\alpha|^2\le q_z^a\) and \(|\iota\alpha^{-1}|^2\le q_z^{-a}\), and equality holds.

For a sheaf that is already lisse on the curve, the inertia is trivial, and Lemma \hyperref[bp127]{15.2} identifies the stalk of \(j_*\) with the original stalk. The first assertion, applied at each of the finitely many points outside the dense open, then gives the second.
\end{proof}

% text: prose-15-06
\phantomsection\label{bp139}%
The proof of Proposition \hyperref[bp134]{15.5}, together with Lemma \hyperref[bp104]{14.20}, also gives the extreme degrees of the cohomology of a middle extension on an affine curve.

% object: lemma-extreme-degrees-affine-curve (blueprint/chap-Pure-sheaves-on-curves/sec-Weights-of-middle-extensions/lemma-extreme-degrees-affine-curve/_section.md)
\phantomsection\label{bp140}%
\begin{bpstatementit}{Lemma~15.7}{ (The extreme degrees on an affine curve)}
Let \(j:U_0\hookrightarrow C_0\) be a dense open in a smooth geometrically connected affine curve, let \(F\) be lisse on \(U_0\), and put \(V=F_{\bar\eta}\). Then \(H_c^i(C,j_*F)=0\) for \(i\ne1,2\), and there is a canonical identification

\[
H_c^2(C,j_*F)=V_{\pi_1(U)}(-1).
\]

If \(F\) is pure of weight \(w\), the group in degree two, when nonzero, is pure of weight \(w+2\).
\end{bpstatementit}

% proof: proof/lemma-extreme-degrees-affine-curve
\begin{proof}[{Proof}]\phantomsection\label{bp141}%
The vanishing of \(H^0_c\) was shown in the proof of Proposition \hyperref[bp134]{15.5}. For the higher degrees, let \(i\) be the inclusion of \(C_0-U_0\). Since sheaves with finite support have compact-support cohomology only in degree zero, the localization sequence for \(0\to j_!F\to j_*F\to i_*i^*j_*F\to0\) identifies \(H^r_c(C,j_*F)\) with \(H_c^r(U,F)\) for \(r\ge2\). Lemma \hyperref[bp104]{14.20} gives the value and the weight of this group for \(r=2\), and its vanishing for \(r>2\).
\end{proof}

% text: prose-15-07
\phantomsection\label{bp142}%
For the perverse t-structure in Chapter \hyperref[bp167]{16} we need the same range of degrees for arbitrary constructible sheaves.

% object: lemma-cohomological-range-curve (blueprint/chap-Pure-sheaves-on-curves/sec-Weights-of-middle-extensions/lemma-cohomological-range-curve/_section.md)
\phantomsection\label{bp143}%
\begin{bpstatementit}{Lemma~15.8}{ (Cohomological range on a curve)}
Let \(C\) be a smooth curve over an algebraically closed field and let \(F\) be a constructible sheaf on \(C\) with coefficients in \(E\) or \(\overline{\mathbf Q}_\ell\). Then \(H^i_c(C,F)=0\) for \(i\notin\{0,1,2\}\).
\end{bpstatementit}

% proof: proof/lemma-cohomological-range-curve
\begin{proof}[{Proof}]\phantomsection\label{bp144}%
Choose a dense open \(j:U\hookrightarrow C\) on which \(F\) is lisse, and let \(i\) be the inclusion of the finite complement. The sequence \(0\to j_!j^*F\to F\to i_*i^*F\to0\) gives a long exact sequence in which \(H^r_c(C,i_*i^*F)\) vanishes for \(r\ne0\), because the support is finite. It therefore suffices to show \(H^r_c(U,j^*F)=0\) for \(r\notin\{0,1,2\}\). For \(r<0\) there is nothing to prove, and for \(r>2\), degreewise Poincaré duality (Corollary \bpout{10.13}, with \(d=1\)) shows that \(H^r_c(U,j^*F)\) is dual to \(H^{2-r}(U,(j^*F)^\vee(1))\), which is zero because \((j^*F)^\vee\) is a sheaf and \(2-r<0\).
\end{proof}

% text: prose-15-08
\phantomsection\label{bp145}%
A constructible Weil sheaf on a curve \(C_0\) over \(k_0\) gives a constructible sheaf on \(C=C_0\times_{k_0}k\), so the lemma applies to it as well.


% object: sec-Growth-of-squared-traces (blueprint/chap-Pure-sheaves-on-curves/sec-Growth-of-squared-traces/_section.md), content hidden
\subsection*{\texorpdfstring{15.2\quad Growth of squared traces}{15.2 Growth of squared traces}}
\phantomsection\label{bp146}%
\addcontentsline{toc}{subsection}{\texorpdfstring{15.2\quad Growth of squared traces}{15.2 Growth of squared traces}}
% text: prose-15-09
\phantomsection\label{bp147}%
In Chapter \hyperref[bp223]{17} we know the trace function of the Fourier transform \(M'\) exactly, but not its weight. For a pure sheaf the trace function determines the weight: it is the growth rate of the sums of squared traces (Theorem \hyperref[bp162]{15.14}).

% object: proposition-twisting-weight-zero-adding-dual (blueprint/chap-Pure-sheaves-on-curves/sec-Growth-of-squared-traces/proposition-twisting-weight-zero-adding-dual/_section.md)
\phantomsection\label{bp148}%
\begin{bpstatementit}{Proposition~15.9}{ (Twisting to weight zero and adding the dual)}
Let \(\mathcal F\ne0\) be lisse and pure of weight \(w\) on a nonempty finite-type scheme over \(k_0\).  There is \(b\in\overline{\mathbf Q}_\ell^\times\) with \(|\iota(b)|=q^{-w/2}\).  Put

\[
\mathcal F^0=\mathcal F\otimes\mathcal B_b,
\qquad \mathcal E=\mathcal F^0\oplus(\mathcal F^0)^\vee.
\]

Then \(\mathcal E\) is real and pure of weight zero, contains \(\mathcal F^0\) as a direct summand, and at points over \(k_n\) its trace is twice the real part of the trace of \(\mathcal F^0\).
\end{bpstatementit}

% proof: proof/proposition-twisting-weight-zero-adding-dual
\begin{proof}[{Proof}]\phantomsection\label{bp149}%
A twist to weight zero puts all local eigenvalues on the unit circle, where inversion is complex conjugation. Adding the dual therefore makes each multiset of eigenvalues stable under conjugation.

To find the twist, choose a closed point \(x\), of degree \(d\), at which the stalk is nonzero, and a Frobenius eigenvalue \(\alpha\) at \(x\).  Let \(b\in\overline{\mathbf Q}_\ell\) be a \(d\)th root of \(\alpha^{-1}\).  By purity, \(|\iota(b)|=|\iota(\alpha)|^{-1/d}=q^{-w/2}\), and Proposition \hyperref[bp24]{13.7} shows that \(\mathcal F^0\) is pure of weight zero.

Its local eigenvalues \(\beta_j\) therefore have absolute value one after applying \(\iota\). The local eigenvalues of its dual are the \(\beta_j^{-1}\), and \(\iota(\beta_j^{-1})=\overline{\iota(\beta_j)}\). So the union of the two multisets is stable under complex conjugation, and \(\mathcal E\) is real and pure. In the same way, at a point \(x\) over \(k_n\) the eigenvalues of \(\operatorname{Frob}_{x,n}\) on \((\mathcal F^0)^\vee\) become, after \(\iota\), the complex conjugates of those on \(\mathcal F^0\). Hence \(\iota\) of the trace of \((\mathcal F^0)^\vee\) is the complex conjugate of \(\iota\) of the trace of \(\mathcal F^0\), and the trace of \(\mathcal E\) is twice its real part.
\end{proof}

% text: prose-15-10
\phantomsection\label{bp150}%
The number \(\nu\) below is the growth exponent of the sums of squared traces, shifted by \(1\) so that it equals the weight. We take a \(\limsup\) because the sums can vanish for infinitely many \(n\), as the example of \(\mathcal B_b\oplus\mathcal B_{-b}\) at the start of this chapter shows.

% object: definition-growth-squared-traces (blueprint/chap-Pure-sheaves-on-curves/sec-Growth-of-squared-traces/definition-growth-squared-traces/_section.md)
\phantomsection\label{bp151}%
\begin{bpstatement}{Definition~15.10}{ (Growth of squared traces)}
Let \(K\) be a constructible Weil complex on an affine curve \(X_0\). For \(n\ge1\) and \(x\in X_0(k_n)\), let \(\operatorname{Frob}_{x,n}\) be the Frobenius acting on the stalk \(K_{\bar x}\) (Theorem \bpout{12.24}), and set

\[
t_{K,n}(x)=\sum_i(-1)^i\operatorname{Tr}
(\operatorname{Frob}_{x,n}\mid\mathcal H^i(K)_{\bar x}),\qquad
S_n(K)=\sum_{x\in X_0(k_n)}|\iota t_{K,n}(x)|^2.
\]

Put

\[
\nu(K)=\limsup_{n\to\infty}\frac{\log_q S_n(K)}n-1,
\qquad \log_q0=-\infty.
\]
\end{bpstatement}

% text: prose-15-11
\phantomsection\label{bp152}%
A cohomological shift changes the weight of a complex in the sense of Definition \hyperref[bp9]{13.1}. But it only multiplies the trace function by a sign, and so it leaves \(\nu\) unchanged.

In the theorem below we find \(\nu\) from the radius of convergence of a generating series, which is a logarithmic derivative of an \(L\)-function.

% object: lemma-radius-squared-traces (blueprint/chap-Pure-sheaves-on-curves/sec-Growth-of-squared-traces/lemma-radius-squared-traces/_section.md)
\phantomsection\label{bp153}%
\begin{bpstatementit}{Lemma~15.11}{ (Radius of convergence)}
Let \(K\) be a constructible Weil complex on an affine curve \(X_0\). The series \(\varphi_K(t)=\sum_{n\ge1}S_n(K)t^{n-1}\) has radius of convergence \(q^{-\nu(K)-1}\), where this number is read as \(+\infty\) if \(\nu(K)=-\infty\) and as \(0\) if \(\nu(K)=+\infty\).
\end{bpstatementit}

% proof: proof/lemma-radius-squared-traces
\begin{proof}[{Proof}]\phantomsection\label{bp154}%
By the Cauchy–Hadamard formula the radius is the inverse of \(\limsup_nS_n(K)^{1/n}=q^{\nu(K)+1}\); the shift of the index by one does not change it.
\end{proof}

% text: prose-15-12
\phantomsection\label{bp155}%
To compute \(\nu\) we also need upper bounds for lisse sheaves and for complexes with finite support, as well as the effect of twisting by \(\mathcal B_b\).

% object: lemma-bounds-squared-traces (blueprint/chap-Pure-sheaves-on-curves/sec-Growth-of-squared-traces/lemma-bounds-squared-traces/_section.md)
\phantomsection\label{bp156}%
\begin{bpstatementit}{Lemma~15.12}{ (Bounds for squared traces)}
Let \(X_0\) be a smooth affine curve over \(k_0\), and let \(a\) be a real number.

\begin{enumerate}
\item If \(F\) is a lisse Weil sheaf of rank \(r\) on \(X_0\), punctually of weights \(\le a\), then \(S_n(F)\le Dr^2q^{n(a+1)}\) for all \(n\), with \(D\) as in Lemma \hyperref[bp131]{15.4}, and hence \(\nu(F)\le a\).
\item If \(K\) is a constructible Weil complex on \(X_0\) supported at finitely many closed points, and every eigenvalue \(\alpha\) of \(\operatorname{Frob}_x\) on every \(\mathcal H^i(K)_{\bar x}\) satisfies \(|\iota(\alpha)|\le q_x^{a/2}\), then there is a constant \(C\) with \(S_n(K)\le Cq^{na}\) for all \(n\), and hence \(\nu(K)\le a-1\).
\item For every constructible Weil complex \(K\) on \(X_0\) and every \(b\in\overline{\mathbf Q}_\ell^\times\), we have \(\nu(K\otimes\mathcal B_b)=\nu(K)+\lambda_\iota(b)\), with \(\mathcal B_b\) as in Proposition \hyperref[bp24]{13.7}.
\end{enumerate}
\end{bpstatementit}

% proof: proof/lemma-bounds-squared-traces
\begin{proof}[{Proof}]\phantomsection\label{bp157}%
For (1), the trace at each point of \(X_0(k_n)\) has absolute value at most \(rq^{na/2}\), and Lemma \hyperref[bp131]{15.4} gives \(S_n\le Dr^2q^{n(a+1)}\). For (2), Lemma \hyperref[bp131]{15.4} bounds the number of points of the support with values in \(k_n\) independently of \(n\), so it suffices to bound the trace at each of them. Since the support consists of finitely many points with finite-dimensional stalks, the sum \(c\) of the ranks of the stalks of the \(\mathcal H^i(K)\) is finite, and at each such point the trace has absolute value at most \(cq^{na/2}\). For (3), the trace of \(\mathcal B_b\) at a \(k_n\)-point is \(b^n\), so tensoring by \(\mathcal B_b\) multiplies \(S_n\) by \(|\iota(b)|^{2n}=q^{n\lambda_\iota(b)}\).
\end{proof}

% text: prose-15-13
\phantomsection\label{bp158}%
A correction at finitely many points of small weight does not change \(\nu\). We use this to discard the boundary terms in Corollary \hyperref[bp165]{15.15}.

% object: lemma-growth-finite-corrections (blueprint/chap-Pure-sheaves-on-curves/sec-Growth-of-squared-traces/lemma-growth-finite-corrections/_section.md)
\phantomsection\label{bp159}%
\begin{bpstatementit}{Lemma~15.13}{ (Finite corrections)}
Let \(K\) and \(L\) be constructible Weil complexes on an affine curve \(X_0\), and let \(a\) be a real number. Suppose that there is a constant \(C\) with \(\sum_{x\in X_0(k_n)}|\iota(t_{K,n}-t_{L,n})(x)|^2\le Cq^{na}\) for all \(n\). If \(\max(\nu(K),\nu(L))>a-1\), then \(\nu(K)=\nu(L)\).
\end{bpstatementit}

% proof: proof/lemma-growth-finite-corrections
\begin{proof}[{Proof}]\phantomsection\label{bp160}%
The inequality \(|z+u|^2\le2|z|^2+2|u|^2\), applied to \(t_{L,n}\) and the difference \(t_{K,n}-t_{L,n}\), gives \(S_n(K)\le2S_n(L)+2Cq^{na}\), so that \(\nu(K)\le\max(\nu(L),a-1)\), and by symmetry \(\nu(L)\le\max(\nu(K),a-1)\).  If one of \(\nu(K)\), \(\nu(L)\) exceeds \(a-1\), these two inequalities force equality.
\end{proof}

% text: prose-15-14
\phantomsection\label{bp161}%
By counting points, Lemma lemma-bounds-squared-traces gives \(\nu\le w\) for a sheaf pure of weight \(w\). A priori, however, the squared traces could be much smaller than this trivial bound, through cancellation among the eigenvalues at each point. The next theorem shows that on average they are not.

% object: theorem-weights-sums-squared-traces (blueprint/chap-Pure-sheaves-on-curves/sec-Growth-of-squared-traces/theorem-weights-sums-squared-traces/_section.md)
\phantomsection\label{bp162}%
\begin{bpstatementit}{Theorem~15.14}{ (Weights and sums of squared traces)}
Let \(X_0/\mathbf F_q\) be a smooth geometrically connected affine curve. For every nonzero pure lisse Weil sheaf \(\mathcal G\) of weight \(w\),

\[
\nu(\mathcal G)=w.
\]
\end{bpstatementit}

% proof: proof/theorem-weights-sums-squared-traces
\begin{proof}[{Proof}]\phantomsection\label{bp163}%
Lemma lemma-bounds-squared-traces gives \(\nu(\mathcal G)\le w\), so, by Lemma \hyperref[bp153]{15.11}, we must show that the radius of \(\varphi_{\mathcal G}\) is not larger than \(q^{-w-1}\). For a real sheaf we find a pole of the logarithmic derivative of an \(L\)-function on that circle, and Proposition \hyperref[bp148]{15.9} reduces the general case to this one by twisting to weight zero and adding the dual.

\emph{The real case.} Suppose first that \(\mathcal G\) is real, and put \(\mathcal T=\mathcal G\otimes\mathcal G\). The squared traces of \(\mathcal G\) are then the traces of \(\mathcal T\), and we express their generating series through the \(L\)-function of \(\mathcal T\). Since \(\iota t_{\mathcal G,n}(x)\) is real, we have \(\iota t_{\mathcal T,n}(x)=|\iota t_{\mathcal G,n}(x)|^2\) for \(x\in X_0(k_n)\). The Euler product and the trace formula then give

\[
\varphi_{\mathcal G}(t)=\frac{L'(t)}{L(t)},\qquad
L(t)=\frac{\iota\det(1-Ft\mid H_c^1(X,\mathcal T))}
{\iota\det(1-Ft\mid H_c^2(X,\mathcal T))}.
\]

For the first equality, group the local logarithms by degree. Then the coefficient of \(t^{n-1}\) in \(L'/L\) is the sum of the traces over \(X_0(k_n)\), which is \(S_n(\mathcal G)\). As \(\mathcal T\) is pure of weight \(2w\), Lemma \hyperref[bp104]{14.20} shows that all roots of the denominator have absolute value \(R=q^{-w-1}\).

We show that the function \(L\) has a pole on the circle \(|t|=R\). This is not automatic, since a priori all the eigenvalues on \(H^2_c\) could cancel against eigenvalues on \(H^1_c\), but positivity rules this out. Indeed, by Lemma \hyperref[bp109]{14.22}, each local Euler factor of \(\mathcal T\) has nonnegative Taylor coefficients, so the coefficients of the full product dominate those of any one factor. Now choose a closed point of the curve. Since \(\mathcal G\) has positive rank and Frobenius is invertible, the Euler factor of \(\mathcal T\) at this point is the reciprocal of a nonconstant polynomial. It has infinitely many nonzero Taylor coefficients, because a polynomial cannot be inverse to a nonconstant polynomial. Hence the full Euler product is not a polynomial, and the rational function \(L\) has a pole, which by the bound on the denominator lies on \(|t|=R\).

A pole of \(L\) is also a pole of its logarithmic derivative: if \(L(t)=(t-t_0)^{-m}u(t)\) near \(t_0\), with \(m>0\) and \(u(t_0)\ne0\), then \(L'/L=-m/(t-t_0)+u'/u\). So the radius of convergence of \(\varphi_{\mathcal G}\) is exactly \(R\), and Lemma \hyperref[bp153]{15.11} gives \(\nu(\mathcal G)=w\).

\emph{The general case.} Let \(\mathcal G\) be an arbitrary pure sheaf, and let \(\mathcal G^0=\mathcal G\otimes\mathcal B_b\) and \(\mathcal E\) be as in Proposition \hyperref[bp148]{15.9}. The sheaf \(\mathcal G^0\) is pure of weight zero, and by Lemma lemma-bounds-squared-traces we have \(\nu(\mathcal G^0)=\nu(\mathcal G)+\lambda_\iota(b)=\nu(\mathcal G)-w\), so it suffices to show that \(\nu(\mathcal G^0)=0\). The sheaf \(\mathcal E=\mathcal G^0\oplus(\mathcal G^0)^\vee\) is nonzero, real and pure of weight zero, so by the case already treated \(\nu(\mathcal E)=0\).

This lower bound passes to \(\mathcal G^0\), because the trace of \(\mathcal E\) is twice the real part of that of \(\mathcal G^0\). In other words, for every \(n\) and \(x\in X_0(k_n)\) we have \(\iota t_{\mathcal E,n}(x)=2\operatorname{Re}(\iota t_{\mathcal G^0,n}(x))\). Hence \(S_n(\mathcal E)\le4S_n(\mathcal G^0)\), and so \(\nu(\mathcal G^0)\ge0\). Since the opposite inequality comes from Lemma lemma-bounds-squared-traces, we obtain \(\nu(\mathcal G^0)=0\), that is, \(\nu(\mathcal G)=w\).
\end{proof}

% text: prose-15-15
\phantomsection\label{bp164}%
Under the hypothesis of Proposition \hyperref[bp218]{16.17}, the Fourier transform of \((j_*M)[1]\) is again a middle extension. But Plancherel applies only to complexes on the whole line, boundary points included, so we need the theorem for \(j_*\) and \(j_!\) on the larger curve. This costs nothing: the boundary points do not change \(\nu\), because their weights are small.

% object: corollary-sums-squared-traces-j-shriek (blueprint/chap-Pure-sheaves-on-curves/sec-Growth-of-squared-traces/corollary-sums-squared-traces-j-shriek/_section.md)
\phantomsection\label{bp165}%
\begin{bpstatementit}{Corollary~15.15}{ (Sums of squared traces of $j_!$ and $j_*$)}
Let \(j:U_0\hookrightarrow X_0\) be a dense open of a smooth geometrically connected affine curve, and let \(\mathcal G\) be a nonzero pure lisse Weil sheaf of weight \(w\) on \(U_0\). Then

\[
\nu(j_!\mathcal G)=\nu(j_*\mathcal G)=w.
\]
\end{bpstatementit}

% proof: proof/corollary-sums-squared-traces-j-shriek
\begin{proof}[{Proof}]\phantomsection\label{bp166}%
For \(j_!\) the trace function is that of \(\mathcal G\) extended by zero, and for \(j_*\) it differs from this at finitely many points of small weight, which Lemma \hyperref[bp159]{15.13} allows us to ignore.

The open curve \(U_0\) is again smooth, geometrically connected and affine, and Theorem \hyperref[bp162]{15.14} on \(U_0\) gives \(\nu(\mathcal G)=w\). Extension by zero only adds zero traces, so \(\nu(j_!\mathcal G)=w\). The trace functions of \(j_*\) and of the extension by zero differ only at the points of \(X_0-U_0\), by the trace function of the boundary stalks. These stalks have weights at most \(w\) by Proposition \hyperref[bp134]{15.5}, so if \(i\) denotes the inclusion of \(X_0-U_0\), then Lemma lemma-bounds-squared-traces, applied to the sheaf \(i_*i^*j_*\mathcal G\), bounds the squared sum of the difference by \(Cq^{nw}\). Since \(\nu(j_!\mathcal G)=w>w-1\), Lemma \hyperref[bp159]{15.13} shows that \(\nu(j_*\mathcal G)\) is still \(w\).
\end{proof}



% object: chap-The-Fourier-transform (blueprint/chap-The-Fourier-transform/_section.md), content hidden
\section*{\texorpdfstring{16\quad The Fourier transform}{16 The Fourier transform}}
\phantomsection\label{bp167}%
\addcontentsline{toc}{section}{\texorpdfstring{16\quad The Fourier transform}{16 The Fourier transform}}
% text: prose-16-01
\phantomsection\label{bp168}%
Theorem \hyperref[bp120]{14.26} controls the Frobenius eigenvalues of a real sheaf at the points of a curve. The bound we need next, however, is of a different kind. Let \(M\) be an irreducible lisse sheaf, pure of weight \(w\), on a dense open \(j:U_0\hookrightarrow\mathbf A^1\). The curve case of the Deligne–Laumon theorem comes down to showing that \(H^1_c(\mathbf A^1,j_*M)\) has weights at most \(w+1\). This is a cohomology group, not the stalk of a sheaf at a point, and Rankin's argument does not reach it. Is there a sheaf on another curve whose stalk at one point is this group, and whose weight we can compute from that of \(M\)? Up to a shift, the Fourier transform of \((j_*M)[1]\) is such a sheaf.

For functions on \(\mathbf F_q\) the discrete Fourier transform already does what we want. Fix a nontrivial additive character \(\psi\) and put \(\hat t(y)=\sum_xt(x)\psi(xy)\). Two properties matter. First, the value \(\hat t(0)\) is the sum of all values of \(t\), so a sum over the line becomes a single value. Second, Plancherel's formula \(\sum_y|\hat t(y)|^2=q\sum_x|t(x)|^2\) holds, by orthogonality of characters.

To see Plancherel compute an absolute value, let \(\chi\) be a nontrivial character of \(\mathbf F_q^\times\), extended by \(\chi(0)=0\), and let \(g=\sum_x\chi(x)\psi(x)\) be its Gauss sum. Substituting \(x\mapsto x/y\) shows that \(\hat\chi(y)=\overline\chi(y)g\) for \(y\ne0\), so the transform has the same absolute value \(|g|\) at all \(q-1\) nonzero points, while \(\hat\chi(0)=0\). Plancherel therefore reads \((q-1)|g|^2=q(q-1)\), and we find \(|g|=\sqrt q\). In the language of weights, this means that the transform carries a function of weight \(0\) to one of weight \(1\). For the quadratic character of \(\mathbf F_3\), moreover, Example \hyperref[bp174]{16.2} finds this Gauss sum, up to sign, as a Frobenius eigenvalue on a cohomology group.

The proof in Chapter \hyperref[bp223]{17} follows the same two steps for \(M\): first we show that the transform is pure, and then Plancherel finds its weight.

Grothendieck's function–sheaf dictionary (Proposition \bpout{12.28}) turns the transform into a functor. Let \(\mathcal K_\psi=\mu^*\mathcal L_\psi\) be the pullback of the Artin–Schreier sheaf (Lemma \hyperref[bp171]{16.1}) along the multiplication \(\mu(x,y)=xy\). Its trace function on \(\mathbf F_q\)-points is \(\psi(xy)\), and we put

\[
\operatorname{FT}_\psi K=Rp_{2!}(p_1^*K\otimes\mathcal K_\psi)[1].
\]

On trace functions it is the discrete transform up to sign: for \(y\in\mathbf F_q\),

\[
t_{\operatorname{FT}_\psi K}(y)=-\sum_{x\in\mathbf F_q}t_K(x)\,\psi(xy),
\]

because under the dictionary (Section \bpout{12.6}) pullback, tensor product and \(Rp_{2!}\) become composition, multiplication and summation over the fibres, the three operations from which \(\hat t\) is built, while the shift \([1]\) contributes the sign. Both properties of the discrete transform carry over. The first becomes the statement that the stalk at the origin of the dual line is \(R\Gamma_c(\mathbf A^1,K)[1]\) (Proposition \hyperref[bp183]{16.5}), so for \(K=j_*M\) this complex contains the group we want. Plancherel becomes an identity of trace functions (Corollary \hyperref[bp192]{16.8}), and it is what will compute the weight of the transform. In addition, Fourier inversion, \(\operatorname{FT}_{\psi^{-1}}\operatorname{FT}_\psi K\simeq K(-1)\), holds because \(R\Gamma_c(\mathbf A^1,\mathcal L_\psi)=0\), the sheaf version of \(\sum_x\psi(x)=0\). The \(\ell\)-adic transform is due to Deligne, and Laumon's proof [\hyperref[bp352]{Lau87}] of the Deligne–Laumon theorem is built on it.

To speak of the weight of the transform we need it to be a lisse sheaf on a dense open. But the transform of a sheaf need not be a sheaf. For example, the transform of the constant sheaf is a skyscraper at the origin in degree one, \(\operatorname{FT}_\psi(E)=\delta_0(-1)[-1]\). This is the sheaf version of \(\hat1=q\,\delta_0\), up to the sign coming from the shift. What the transform does preserve, however, is the category of perverse sheaves on the line (Proposition \hyperref[bp214]{16.16}, proved in Section \hyperref[bp212]{16.5} below). Since the transform is an exact equivalence of perverse sheaves, it sends the simple object \((j_*M)[1]\) to a simple object, which is either punctual or a middle extension. Moreover, only character sheaves \(\mathcal L_{\psi(ax)}\) are sent to skyscrapers. Hence, if no geometric constituent of the irreducible sheaf \(M\) is such a character, we have

\[
\operatorname{FT}_\psi((j_*M)[1])\simeq(j'_*M')[1]
\]

for an irreducible lisse \(M'\) on a dense open \(j'\) of the dual line (Proposition \hyperref[bp218]{16.17} below). Comparing stalks at the origin then identifies \(H^1_c(\mathbf A^1,j_*M)\) with a stalk of \(j'_*M'\), and this is the sheaf we asked for.

We use the six operations on the pro-étale site. At first the coefficients are a finite extension \(E/\mathbf Q_\ell\) that contains the \(p\)th roots of unity, where \(p\) is the characteristic of \(k_0\), so that \(E\) contains the values of the additive characters of \(k_0\). Later the coefficients are \(\overline{\mathbf Q}_\ell\).

% object: sec-The-Artin-Schreier-sheaf (blueprint/chap-The-Fourier-transform/sec-The-Artin-Schreier-sheaf/_section.md), content hidden
\subsection*{\texorpdfstring{16.1\quad The Artin–Schreier sheaf}{16.1 The Artin–Schreier sheaf}}
\phantomsection\label{bp169}%
\addcontentsline{toc}{subsection}{\texorpdfstring{16.1\quad The Artin–Schreier sheaf}{16.1 The Artin–Schreier sheaf}}
% text: prose-16-02
\phantomsection\label{bp170}%
The kernel of the transform is a lisse sheaf whose trace function is the additive character \(\psi\), just as \(\psi(xy)\) is the kernel of the discrete transform. It is cut out of the Artin–Schreier cover of the line as an eigenspace of the translations. This cover is Galois with group the additive group \(\mathbf F_p\), acting by \(u\mapsto u+a\). Let \(u^p-u=x\) with \(x\in\mathbf F_{q^n}\). Raising to the powers \(p,p^2,\ldots\) and adding gives \(u^{q^n}=u+\operatorname{Tr}(x)\). In other words, Frobenius acts on the fibre over \(x\) by the translation \(\operatorname{Tr}(x)\in\mathbf F_p\). A character \(\psi_0\) of \(\mathbf F_p\) then turns this translation into the number \(\psi_0(\operatorname{Tr}x)\). Similarly, the Kummer cover \(u^{q-1}=x\) of \(\mathbf G_m\) produces the multiplicative characters in the same way.

Besides the sheaf itself (Lemma \hyperref[bp171]{16.1}), we need the sheaf version of \(\sum_x\psi(x)=0\), namely the vanishing of \(R\Gamma_c(\mathbf A^r,\lambda^*\mathcal L_\psi)\) for a nonzero linear form \(\lambda\) (Lemma \hyperref[bp179]{16.4}). This is stronger than the statement on trace functions, since the trace formula only says that the alternating sum of the traces vanishes over every \(\mathbf F_{q^n}\). Fourier inversion and the stalk computations below, however, need each cohomology group to vanish. This vanishing comes from the eigenspace decomposition of the cover: the total space of the cover is again an affine line, and the translations act trivially on the compactly supported cohomology of a line. Hence the eigenspace of a nontrivial character is zero.

% object: lemma-artin-schreier-character-sheaf (blueprint/chap-The-Fourier-transform/sec-The-Artin-Schreier-sheaf/lemma-artin-schreier-character-sheaf/_section.md)
\phantomsection\label{bp171}%
\begin{bpstatementit}{Lemma~16.1}{ (The Artin–Schreier character sheaf)}
Let \(\psi:k_0\to E^\times\) be a nontrivial additive character. There is a lisse rank-one \(E\)-sheaf \(\mathcal L_\psi\) on \(\mathbf A^1_{k_0}\), with a specified trivialization at zero, such that

\[
m^*\mathcal L_\psi\simeq
\operatorname{pr}_1^*\mathcal L_\psi\otimes
\operatorname{pr}_2^*\mathcal L_\psi,
\qquad
\mathcal L_\psi^\vee\simeq\mathcal L_{\psi^{-1}}.
\]

Here \(m\) is the addition map, and the isomorphisms respect the trivializations at zero. For \(x\in\mathbf F_{q^n}\), the Frobenius trace is

\[
t_{\mathcal L_\psi,n}(x)
=\psi\bigl(\operatorname{Tr}_{\mathbf F_{q^n}/k_0}(x)\bigr).
\]

In particular \(\mathcal L_\psi\) is pure of weight zero.
\end{bpstatementit}

% proof: proof/lemma-artin-schreier-character-sheaf
\begin{proof}[{Proof}]\phantomsection\label{bp172}%
We define \(\mathcal L_\psi\) as an eigensheaf in the direct image of the constant sheaf under an Artin–Schreier cover, and read off its properties from the corresponding properties of the cover: addition of roots for the first isomorphism, a fibrewise pairing for the second, and the action of Frobenius on roots for the trace.

Fix a nontrivial character \(\psi_0:\mathbf F_p\to E^\times\). Since the trace pairing on \(k_0\) is nondegenerate, there is \(a\in k_0^\times\) such that \(\psi(z)=\psi_0(\operatorname{Tr}_{k_0/\mathbf F_p}(az))\). The cover \(u^p-u=ax\) of the affine line is finite étale, with group of translations \(\mathbf F_p\). Let \(\mathcal L_\psi\) be the image of

\[
e_{\psi_0}=\frac1p\sum_{c\in\mathbf F_p}\psi_0(c)^{-1}\tau_c^*
\]

on the direct image of the constant sheaf \(E\) under this cover. It is lisse of rank one, since its stalk is the \(\psi_0\)-eigenspace in the regular representation, and the root \(u=0\) over \(x=0\) specifies the trivialization at zero. For the first isomorphism, adding roots of two Artin–Schreier equations identifies the contracted product of the two torsors with the torsor for the sum, and passing to the character eigenspaces we obtain the isomorphism. For the second, the pairing on the direct image given by summing products over the fibres is perfect and invariant under the translations. Since, for an invariant pairing, the \(\chi\)- and \(\chi'\)-eigenspaces of two characters are orthogonal unless \(\chi\chi'=1\), our pairing pairs the \(\psi_0\)-eigenspace perfectly with the \(\psi_0^{-1}\)-eigenspace. Since \(\psi^{-1}(z)=\psi_0^{-1}(\operatorname{Tr}_{k_0/\mathbf F_p}(az))\), the \(\psi_0^{-1}\)-eigenspace of the same cover is \(\mathcal L_{\psi^{-1}}\), and we obtain \(\mathcal L_\psi^\vee\simeq\mathcal L_{\psi^{-1}}\).

For a root \(u\) above \(x\in\mathbf F_{q^n}\) we have

\[
u^{q^n}=u+
\operatorname{Tr}_{\mathbf F_{q^n}/\mathbf F_p}(ax).
\]

Geometric Frobenius acts on functions on the torsor by pullback along this map, and hence on the chosen eigenspace by \(\psi_0(\operatorname{Tr}(ax))\). By transitivity of field traces, this is the asserted trace. Since these scalars are roots of unity, their images under any field embedding have absolute value one.
\end{proof}

% text: prose-16-03
\phantomsection\label{bp173}%
We check the lemma by hand at two points. Over \(\mathbf F_p\) with \(a=1\), the fibre of the cover over \(x=1\) is \(u^p-u=1\), which has no root in \(\mathbf F_p\). Indeed, Frobenius acts on the \(p\) roots by \(u\mapsto u^p=u+1\), a translation of order \(p\), and the trace of \(\mathcal L_\psi\) at \(x=1\) is therefore \(\psi(1)=\psi_0(1)\ne1\). Over \(x=0\), by contrast, the roots are the elements of \(\mathbf F_p\), so they are all fixed and the trace is \(1\).

% object: example-gauss-sum-over-f-3 (blueprint/chap-The-Fourier-transform/sec-The-Artin-Schreier-sheaf/example-gauss-sum-over-f-3/_section.md)
\phantomsection\label{bp174}%
\begin{bpstatement}{Example~16.2}{ (A Gauss sum over $\mathbf F_3$)}
Take \(q=3\), a primitive cube root of unity \(\omega\in\overline{\mathbf Q}_\ell\) and \(\psi(x)=\omega^x\) on \(\mathbf F_3\). Pull \(\mathcal L_\psi\) back along \(h:\mathbf A^1\to\mathbf A^1\), \(x\mapsto x^2\). By the lemma, the trace of Frobenius on the stalk of \(h^*\mathcal L_\psi\) at \(x\in\mathbf F_3\) is \(\psi(x^2)\), so the trace formula gives

\[
\sum_r(-1)^r\operatorname{Tr}\bigl(F\mid H^r_c(\mathbf A^1_k,h^*\mathcal L_\psi)\bigr)
=\sum_{x\in\mathbf F_3}\omega^{x^2}=1+2\omega ,
\]

and \((1+2\omega)^2=-3\). The sheaf \(h^*\mathcal L_\psi\) is lisse of rank one. It is geometrically nontrivial, since the cover \(u^3-u=x^2\) is geometrically irreducible. Hence its geometric coinvariants vanish, and \(H^0_c=H^2_c=0\) (Lemma \hyperref[bp140]{15.7}). The sum is therefore \(-\operatorname{Tr}(F\mid H^1_c)\).

To see that \(H^1_c\) is a line, we use the running example. The cover \(u^3-u=x^2\) is the affine part \(Y_0\) of the elliptic curve \(C_0\colon y^2=x^3-x\) of the introduction, with the coordinates renamed (\(x^2=u^3-u\)). It lies over \(\mathbf A^1\) by \(\pi(u,x)=x\). Since \(\pi\) is the pullback along \(h\) of the Artin–Schreier cover, \(\pi_*E\) is the direct sum of \(h^*\mathcal L_{\psi'}\) over the three characters \(\psi'\) of \(\mathbf F_3\). Over \(k=\overline{\mathbf F}_3\) we obtain

\[
H^1_c(Y,E)=H^1_c(\mathbf A^1,E)\oplus H^1_c(\mathbf A^1,h^*\mathcal L_\psi)
\oplus H^1_c(\mathbf A^1,h^*\mathcal L_{\psi^{-1}}).
\]

The left side is \(H^1(C,E)\), because \(C\) minus \(Y\) is the single point at infinity, and it has dimension \(2\). Since the first summand on the right is zero, it remains to see that the last two summands have the same dimension. For this, note that the automorphism \((u,x)\mapsto(-u,ix)\) of \(Y\), with \(i^2=-1\) in \(\mathbf F_9\), conjugates the translation by \(c\) into the translation by \(-c\). It therefore exchanges the last two summands, and each of them is a line. This automorphism is defined only over \(\mathbf F_9\), but we use it only to count dimensions. The decomposition itself is stable under Frobenius, because the translations are defined over \(\mathbf F_3\).

Thus Frobenius acts on the line \(H^1_c(\mathbf A^1,h^*\mathcal L_\psi)\) by \(-(1+2\omega)\). This is a square root of \(-3\), and it is one of the two reciprocal roots \(\pm i\sqrt3\) of \(P_1(T)=1+3T^2\) found in the introduction. Moreover, the number \(1+2\omega=\omega-\omega^2\) is the Gauss sum \(g\) of the quadratic character of \(\mathbf F_3\). The classical fact \(|g|=\sqrt q\) is thus the statement that this \(H^1_c\) is pure of weight one. Chapter \hyperref[bp223]{17} proves the analogue for every such sheaf.
\end{bpstatement}

% text: prose-16-04
\phantomsection\label{bp175}%
Fourier inversion integrates the kernel over a whole fibre \(\mathbf A^r\) of a vector bundle, so we need the vanishing on \(\mathbf A^r\) and not only on the line. In coordinates in which the linear form is the first coordinate, the integral over \(\mathbf A^r\) is a product of an integral over the line and one over \(\mathbf A^{r-1}\), and the Künneth formula makes this precise.

Let \(f:X\to S\) and \(g:Y\to S\) be separated morphisms of finite type, let \(p_1:X\times_SY\to X\) and \(p_2:X\times_SY\to Y\) be the projections, and let \(h=f\circ p_1:X\times_SY\to S\) be the structure map. For complexes \(K\) on \(X\) and \(L\) on \(Y\) we put \(K\boxtimes_SL=p_1^*K\otimes p_2^*L\).

% object: lemma-kunneth-formula-compact-supports (blueprint/chap-The-Fourier-transform/sec-The-Artin-Schreier-sheaf/lemma-kunneth-formula-compact-supports/_section.md)
\phantomsection\label{bp176}%
\begin{bpstatementit}{Lemma~16.3}{ (Künneth formula for compact supports)}
In this situation, let \(K,L\) be bounded constructible complexes with coefficients in a field. Then the maps in

\[
Rh_!(p_1^*K\otimes p_2^*L)
\simeq Rf_!(K\otimes Rp_{1!}p_2^*L)
\simeq Rf_!(K\otimes f^*Rg_!L)
\simeq Rf_!K\otimes Rg_!L
\]

of the projection formula for \(p_1\), the proper base change isomorphism \(Rp_{1!}p_2^*L\to f^*Rg_!L\) and the projection formula for \(f\) are isomorphisms. We call their composite \(Rh_!(K\boxtimes_SL)\simeq Rf_!K\otimes Rg_!L\) the \emph{Künneth isomorphism}.
\end{bpstatementit}

% proof: proof/lemma-kunneth-formula-compact-supports
\begin{proof}[{Proof}]\phantomsection\label{bp177}%
We integrate first along the fibres of \(p_1\), which are copies of the fibres of \(g\). This is possible because \(h=f\circ p_1\), so that \(Rh_!=Rf_!Rp_{1!}\). Since the square formed by \(p_1\), \(p_2\), \(f\) and \(g\) is cartesian, proper base change gives \(Rp_{1!}p_2^*L\simeq f^*Rg_!L\). The two projection formulas are isomorphisms as well, and hence so is the composite.
\end{proof}

% text: prose-16-05
\phantomsection\label{bp178}%
Over a point, this is the Künneth isomorphism for cohomology with compact support. All three isomorphisms in the composite are natural, so when \(X\), \(Y\), \(S\), \(f\), \(g\), \(K\) and \(L\) are defined over \(k_0\), the Künneth isomorphism commutes with the Weil structures.

The case \(\lambda=0\) of the next lemma is needed as well, since Fourier inversion integrates the kernel along the diagonal, where the character is trivial.

% object: lemma-integral-linear-character (blueprint/chap-The-Fourier-transform/sec-The-Artin-Schreier-sheaf/lemma-integral-linear-character/_section.md)
\phantomsection\label{bp179}%
\begin{bpstatementit}{Lemma~16.4}{ (The integral of a linear character)}
Let \(\lambda\) be a linear form on \(\mathbf A^r\) over an algebraically closed field of characteristic \(p\). Then

\[
R\Gamma_c(\mathbf A^r,\lambda^*\mathcal L_\psi)
\simeq
\begin{cases}
E(-r)[-2r],&\lambda=0,\\
0,&\lambda\ne0.
\end{cases}
\]

For \(\lambda=0\), the isomorphism is the trace map for cohomology with compact support.
\end{bpstatementit}

% proof: proof/lemma-integral-linear-character
\begin{proof}[{Proof}]\phantomsection\label{bp180}%
On the line, from the cohomology of \(\mathbf P^1\) and localization at the point at infinity, we get \(R\Gamma_c(\mathbf A^1,E)=E(-1)[-2]\), identified by the trace map (Appendix \bpout{D}). Every translation acts trivially on it, since it preserves the trace and the cohomology is concentrated in degree two. Now let \(\pi:\mathbf A^1\to\mathbf A^1\), \(u\mapsto u^p-u\), be the Artin–Schreier cover, which is finite étale and whose total space is again an affine line. Let \(\mathcal L_{\psi_0}\) be the \(\psi_0\)-eigenspace of \(\pi_*E\), that is, the sheaf of the proof of Lemma \hyperref[bp171]{16.1} for \(a=1\). The projector \(e_{\psi_0}\) of that proof cuts out \(R\Gamma_c(\mathbf A^1,\mathcal L_{\psi_0})\) from \(R\Gamma_c(\mathbf A^1,\pi_*E)=R\Gamma_c(\mathbf A^1,E)\) as the \(\psi_0\)-eigenspace of the deck translations. As we have just seen, these translations act trivially on the cohomology of the upper line, so, since \(\psi_0\) is nontrivial, the eigenspace is zero, that is, \(R\Gamma_c(\mathbf A^1,\mathcal L_{\psi_0})=0\). Finally, by its construction, \(\mathcal L_\psi\) is the pullback of \(\mathcal L_{\psi_0}\) by the automorphism \(x\mapsto ax\) of the line, and hence \(R\Gamma_c(\mathbf A^1,\mathcal L_\psi)=0\) as well.

On \(\mathbf A^r\) with \(\lambda\ne0\), we choose linear coordinates in which \(\lambda\) is the first coordinate. Then \(\lambda^*\mathcal L_\psi\) is the external product of \(\mathcal L_\psi\) on the first factor with the constant sheaf on \(\mathbf A^{r-1}\), and the Künneth isomorphism of Lemma \hyperref[bp176]{16.3} writes its cohomology as a tensor product with the first factor \(R\Gamma_c(\mathbf A^1,\mathcal L_\psi)=0\). For \(\lambda=0\), on the other hand, the sheaf is constant, so the trace map of Definition \bpout{9.24} for the projection \(\mathbf A^r\to\operatorname{Spec}k\) is an isomorphism \(R\Gamma_c(\mathbf A^r,E)(r)[2r]\simeq E\). On constant coefficients this map is the trace map of affine space (9.11), obtained by integrating along one coordinate line at a time. Finally, all these calculations take place with coefficients in \(E\) and commute with extension of scalars.
\end{proof}


% object: sec-Fourier-inversion (blueprint/chap-The-Fourier-transform/sec-Fourier-inversion/_section.md), content hidden
\subsection*{\texorpdfstring{16.2\quad Fourier inversion}{16.2 Fourier inversion}}
\phantomsection\label{bp181}%
\addcontentsline{toc}{subsection}{\texorpdfstring{16.2\quad Fourier inversion}{16.2 Fourier inversion}}
% text: prose-16-06
\phantomsection\label{bp182}%
We integrate with compact supports, \(Rp_{2!}\) rather than \(Rp_{2*}\), for two reasons. First, proper base change computes the stalks of \(Rp_{2!}\) as compactly supported cohomology of the fibres, and this puts \(H^1_c(\mathbf A^1,j_*M)\) at the origin. Second, the trace formula holds for \(Rp_{2!}\), so that on trace functions the transform becomes the discrete Fourier transform (16.3) (proved in Section \hyperref[bp190]{16.3}).

The shift \([r]\) is the one for which the transform preserves the perverse heart. For example, on the line we will find \(\operatorname{FT}_\psi(E[1])=\delta_0(-1)\), whereas without the shift the perverse sheaf \(E[1]\) would go to \(\delta_0(-1)[-1]\), which is not perverse.

We work with a vector bundle \(V\to S\) rather than the line alone, because the proofs move along the base. For instance, the stalk formula is base change to a point of the dual bundle, and Fourier inversion is a computation on \(V\times_SV^\vee\times_SV\), which reduces, fibre by fibre, to Lemma \hyperref[bp179]{16.4} on \(\mathbf A^r\). We prove these two properties in Propositions \hyperref[bp183]{16.5} and \hyperref[bp185]{Proposition 16.6}, while the preservation of the heart is Proposition \hyperref[bp214]{16.16}.

Let \(S\) be a scheme of finite type over \(k_0\) or over \(k\), let \(V\to S\) be a vector bundle of rank \(r\), and let \(V^\vee\) be its dual bundle. Write \(p_1,p_2\) for the projections of \(V\times_SV^\vee\), and set

\[
\mathcal K_\psi=\langle\ ,\ \rangle^*\mathcal L_\psi,
\qquad
\operatorname{FT}_\psi(K)
=Rp_{2!}(p_1^*K\otimes\mathcal K_\psi)[r].
\tag{16.1}
\]

The kernel \(\mathcal K_\psi\) is lisse of rank one, and since the six operations preserve constructibility, this defines a triangulated functor \(\operatorname{FT}_\psi:\Dbc(V,E)\to\Dbc(V^\vee,E)\), normalized by the shift \([r]\).

% object: proposition-base-change-fourier-inversion (blueprint/chap-The-Fourier-transform/sec-Fourier-inversion/proposition-base-change-fourier-inversion/_section.md)
\phantomsection\label{bp183}%
\begin{bpstatementit}{Proposition~16.5}{ (Base change and stalks)}
Let \(g:S'\to S\) be a morphism of schemes with \(S'\) of finite type over \(k_0\) or over \(k\), let \(V'=V\times_SS'\), and let \(g_V:V'\to V\) and \(g_{V^\vee}:V'^\vee\to V^\vee\) be the projections.

\begin{enumerate}
\item The proper base change isomorphism for \(p_2\) induces an isomorphism \(g_{V^\vee}^*\operatorname{FT}_\psi(K)\simeq
  \operatorname{FT}_\psi(g_V^*K)\).
\item For \(N\in\Dbc(S,E)\), the projection formula for \(p_2\) induces an isomorphism \(\operatorname{FT}_\psi(K\otimes\pi^*N)\simeq
  \operatorname{FT}_\psi(K)\otimes\pi^{\vee*}N\), where \(\pi:V\to S\) and \(\pi^\vee:V^\vee\to S\) are the projections.
\item For a geometric point \(\bar y\) of \(V^\vee\), the proper base change map gives \[
   (\operatorname{FT}_\psi K)_{\bar y}
   \simeq
   R\Gamma_c\bigl(V_{\bar y},
   K|_{V_{\bar y}}\otimes
   \mathcal L_{\psi\langle x,y\rangle}\bigr)[r].
   \tag{16.2}
   \]
\end{enumerate}
\end{bpstatementit}

% proof: proof/proposition-base-change-fourier-inversion
\begin{proof}[{Proof}]\phantomsection\label{bp184}%
For (1), the formation of \(V\times_SV^\vee\), of the pairing and of \(\mathcal K_\psi\) commutes with base change along \(g\), and proper base change for \(Rp_{2!}\) gives the isomorphism. For (2), we have \(p_1^*\pi^*N=p_2^*\pi^{\vee*}N\), and the projection formula for \(p_2\) gives the isomorphism. For (3), the fibre of \(p_2\) over \(\bar y\) is \(V_{\bar y}\), on which \(\mathcal K_\psi\) restricts to \(\mathcal L_{\psi\langle x,y\rangle}\), and proper base change for \(Rp_{2!}\) along \(\bar y\to V^\vee\) gives the formula.
\end{proof}

% object: proposition-fourier-inversion (blueprint/chap-The-Fourier-transform/sec-Fourier-inversion/proposition-fourier-inversion/_section.md)
\phantomsection\label{bp185}%
\begin{bpstatementit}{Proposition~16.6}{ (Fourier inversion)}
Under \(V^{\vee\vee}=V\) there are natural isomorphisms of functors \(\Dbc(V,E)\to\Dbc(V,E)\)

\[
\operatorname{FT}_{\psi^{-1}}\operatorname{FT}_\psi(K)
\simeq K(-r),
\qquad
\operatorname{FT}_{\psi}\operatorname{FT}_\psi(K)
\simeq [-1]_V^*K(-r).
\]
\end{bpstatementit}

% proof: proof/proposition-fourier-inversion
\begin{proof}[{Proof}]\phantomsection\label{bp186}%
The composite of two transforms is again an integral transform, and we compute its kernel. It will turn out to be the constant sheaf on the diagonal, twisted by \((-r)\), which up to this twist is the kernel of the identity functor. Let \(r_{ij}\) be the projections of \(V\times_SV^\vee\times_SV\) onto pairs of factors. By base change and the projection formula, the composite of two transforms is the transform with the convolution kernel

\[
Rr_{13!}(r_{12}^*\mathcal K_\psi\otimes
r_{23}^*\mathcal K_{\psi^{-1}})[2r].
\]

At a point \((x,z)\), the unshifted integral over the middle factor is the cohomology of the character sheaf attached to \(y\mapsto\psi\langle x-z,y\rangle\), and hence it vanishes for \(x\ne z\) by Lemma \hyperref[bp179]{16.4}. Thus the unshifted convolution is supported on the diagonal, so the localization triangle identifies it canonically with the direct image of its restriction to the diagonal. There the unit trivialization of the character sheaf and the trace map for the vector bundle identify the restriction with \(E(-r)[-2r]\). Since the shift \([2r]\) cancels \([-2r]\), we obtain the first inversion formula. If instead the two kernels use the same character, then \(x-z\) is replaced by \(x+z\), the support is the antidiagonal, and we obtain the second formula. Since the traces and unit trivializations used here are defined over \(k_0\), these isomorphisms commute with the Weil structures.
\end{proof}

% text: prose-16-07
\phantomsection\label{bp187}%
On the line, the transform exchanges the two simplest perverse sheaves, the shifted constant sheaf and the skyscraper at the origin, up to a twist. This is the sheaf version of \(\hat 1=q\,\delta_0\) and \(\hat\delta_0=1\).

% object: lemma-fourier-transform-constant-sheaf (blueprint/chap-The-Fourier-transform/sec-Fourier-inversion/lemma-fourier-transform-constant-sheaf/_section.md)
\phantomsection\label{bp188}%
\begin{bpstatementit}{Lemma~16.7}{ (Transform of the constant sheaf)}
Let \(V=\mathbf A^1_{k_0}\) over \(S=\operatorname{Spec}k_0\), and let \(\delta_0\) be the constant skyscraper sheaf \(E\) at the origin. There are isomorphisms of Weil complexes

\[
\operatorname{FT}_\psi(E)\simeq\delta_0(-1)[-1],
\qquad \operatorname{FT}_\psi(\delta_0)\simeq E[1].
\]
\end{bpstatementit}

% proof: proof/lemma-fourier-transform-constant-sheaf
\begin{proof}[{Proof}]\phantomsection\label{bp189}%
For \(K=E\), the stalk formula (16.2) at \(y\) is the integral of the character sheaf \(\mathcal L_{\psi(xy)}\) over the line, shifted by one. By Lemma \hyperref[bp179]{16.4} it vanishes for \(y\ne0\) and equals \(E(-1)[-2]\) at \(y=0\). Hence \(\operatorname{FT}_\psi(E)\) is supported at the origin, and with the shift \([1]\) it is \(\delta_0(-1)[-1]\). For \(K=\delta_0\), the pullback by the first projection is supported on \(\{0\}\times\mathbf A^1\). Since the kernel is trivial there and the second projection is an isomorphism, this gives \(E[1]\). Moreover, the trace map and the unit trivialization are defined over \(k_0\), so both identifications are Frobenius-equivariant.
\end{proof}


% object: sec-The-Plancherel-formula (blueprint/chap-The-Fourier-transform/sec-The-Plancherel-formula/_section.md), content hidden
\subsection*{\texorpdfstring{16.3\quad The Plancherel formula}{16.3 The Plancherel formula}}
\phantomsection\label{bp190}%
\addcontentsline{toc}{subsection}{\texorpdfstring{16.3\quad The Plancherel formula}{16.3 The Plancherel formula}}
% text: prose-16-08
\phantomsection\label{bp191}%
On trace functions, the Fourier transform is, up to the sign \((-1)^r\), the discrete Fourier transform over \(\mathbf F_{q^n}\) for the character \(\psi\circ\operatorname{Tr}_{\mathbf F_{q^n}/k_0}\) (equation (16.3) below). The Plancherel formula then says how it changes the growth of squared trace sums, and for a pure sheaf this growth is the weight (Theorem \hyperref[bp162]{15.14}). This is how we find the weight of the transform in Chapter \hyperref[bp223]{17}: on the line, Plancherel multiplies the sums of squared traces by \(q^n\), and so it raises the weight read off from their growth by one.

In this section \(S=\operatorname{Spec}k_0\) and \(V=\mathbf A^r_{k_0}\). For a Weil complex \(K\) on a \(k_0\)-scheme \(X_0\), we write \(t_{K,n}(x)\) for the alternating sum of the Frobenius traces at \(x\in X_0(\mathbf F_{q^n})\), as in Definition \hyperref[bp151]{15.10}. With this notation, the dictionary (Proposition \bpout{12.28}) computes the trace function of the transform step by step. For \(y\in\mathbf F_{q^n}^r\), the shift contributes \((-1)^r\) (chain formula of Proposition \bpout{11.2}), part (3) sums over the fibre of \(p_2\), and part (2) splits the trace of the tensor product:

\[
t_{\operatorname{FT}_\psi K,n}(y)
=(-1)^r\,t_{Rp_{2!}(p_1^*K\otimes\mathcal K_\psi),n}(y)
=(-1)^r\sum_{x\in\mathbf F_{q^n}^r}t_{p_1^*K,n}(x,y)\,
t_{\mathcal K_\psi,n}(x,y).
\]

By part (1), \(t_{p_1^*K,n}(x,y)=t_{K,n}(x)\), and by part (1) and Lemma \hyperref[bp171]{16.1}, \(t_{\mathcal K_\psi,n}(x,y)=\psi(\operatorname{Tr}_{\mathbf F_{q^n}/k_0}\langle x,y\rangle)\). Hence

\[
t_{\operatorname{FT}_\psi K,n}(y)
=(-1)^r\sum_{x\in\mathbf F_{q^n}^r}
t_{K,n}(x)\,
\psi\bigl(\operatorname{Tr}_{\mathbf F_{q^n}/k_0}
\langle x,y\rangle\bigr).
\tag{16.3}
\]

% object: corollary-plancherel (blueprint/chap-The-Fourier-transform/sec-The-Plancherel-formula/corollary-plancherel/_section.md)
\phantomsection\label{bp192}%
\begin{bpstatementit}{Corollary~16.8}{ (Plancherel)}
For every \(n\ge1\),

\[
\sum_y|\iota t_{\operatorname{FT}_\psi K,n}(y)|^2
=q^{nr}\sum_x|\iota t_{K,n}(x)|^2.
\]

Consequently, for \(r=1\), the number \(\nu\) of Definition \hyperref[bp151]{15.10} satisfies \(\nu(\operatorname{FT}_\psi K)=\nu(K)+1\) whenever these numbers are finite.
\end{bpstatementit}

% proof: proof/corollary-plancherel
\begin{proof}[{Proof}]\phantomsection\label{bp193}%
As for the classical transform, the proof is by orthogonality of characters. We expand the square of the sum in (16.3), using that the complex conjugate of a character value is its inverse. This writes the left side as a double sum over \(x,z\in\mathbf F_{q^n}^r\) of \(\iota t_{K,n}(x)\overline{\iota t_{K,n}(z)}\) times a character sum over \(y\), which we now compute. For \(x,z\in\mathbf F_{q^n}^r\), the orthogonality of characters gives

\[
\sum_y\iota\psi\bigl(\operatorname{Tr}\langle x-z,y\rangle\bigr)
=\begin{cases}q^{nr},&x=z,\\0,&x\ne z.\end{cases}
\]

Indeed, for \(x\ne z\) the linear form \(y\mapsto\operatorname{Tr}_{\mathbf F_{q^n}/k_0}\langle x-z,y\rangle\) is nonzero, since the trace pairing is nondegenerate, and hence surjective onto \(k_0\). As \(\psi\) is nontrivial, there is \(y_0\) with \(c=\psi(\operatorname{Tr}\langle x-z,y_0\rangle)\ne1\), and substituting \(y+y_0\) for \(y\) multiplies the sum by \(\iota(c)\), so the sum is zero. Therefore only the diagonal terms remain, and they give the asserted identity. On a line, finally, the factor \(q^n\) raises \(\nu\) by one.
\end{proof}

% text: prose-16-09
\phantomsection\label{bp194}%
For the constant sheaf \(E\) on the line and \(n=1\), formula (16.3) gives the trace \(-\sum_x\psi(xy)\), which is \(-q\) at \(y=0\) and \(0\) elsewhere, in agreement with \(\operatorname{FT}_\psi(E)=\delta_0(-1)[-1]\). Both sides of Plancherel are then \(q^2\), and over \(\mathbf F_{q^n}\) they are \(q^{2n}\). This gives \(\nu(E)=0\) and \(\nu(\operatorname{FT}_\psi E)=1\), as the corollary predicts.


% object: sec-Perverse-sheaves-on-a-curve (blueprint/chap-The-Fourier-transform/sec-Perverse-sheaves-on-a-curve/_section.md), content hidden
\subsection*{\texorpdfstring{16.4\quad Perverse sheaves on a curve}{16.4 Perverse sheaves on a curve}}
\phantomsection\label{bp195}%
\addcontentsline{toc}{subsection}{\texorpdfstring{16.4\quad Perverse sheaves on a curve}{16.4 Perverse sheaves on a curve}}
% text: prose-16-10
\phantomsection\label{bp196}%
To apply Theorem \hyperref[bp162]{15.14} to the transform of a pure lisse sheaf \(M\), we need to know that the transform of \((j_*M)[1]\) is again a lisse sheaf on a dense open, placed in degree \(-1\). Although the transform sends the constant sheaf to a skyscraper (Lemma \hyperref[bp188]{16.7}), it preserves the perverse heart, which on a smooth curve has an elementary description in terms of ordinary cohomology sheaves and punctual support.

We can read off the shape of this description from the stalk formula. The stalk of the transform of \(G[1]\) at \(y\) is \(R\Gamma_c(\mathbf A^1,G\otimes\mathcal L_{\psi(yx)})[2]\). For a sheaf \(G\) without punctual sections \(H^0_c\) vanishes, and the proof of Proposition \hyperref[bp214]{16.16} below shows that \(H^2_c\) is nonzero only at finitely many \(y\). So the transform lives in degrees \(-1\) and \(0\), with a punctual sheaf in degree \(0\), and this suggests allowing skyscrapers in degree \(0\) in \({}^pD^{\le0}\). Orthogonality then forces the condition on the other side: a sheaf in degree \(0\) lies in \({}^pD^{\ge1}\) only if it receives no nonzero map from a skyscraper, that is, only if it has no punctual subsheaf. The perverse t-structure on a curve is the t-structure built from these two conditions.

What we finally use about it is the list of its simple objects (Lemma \hyperref[bp209]{16.14}). By this list, an exact equivalence must send an irreducible middle extension to an irreducible middle extension or to an irreducible punctual sheaf.

Recall that a \emph{t-structure} on a triangulated category \(D\) is a pair of full subcategories \(D^{\le0}\supset D^{\le-1}\) and \(D^{\ge0}\supset D^{\ge1}\), where \(D^{\le n}=D^{\le0}[-n]\) and \(D^{\ge n}=D^{\ge0}[-n]\), such that \(\Hom(A,B)=0\) for \(A\in D^{\le0}\) and \(B\in D^{\ge1}\), and every object sits in a triangle \(A\to K\to
B\to\) with \(A\in D^{\le0}\) and \(B\in D^{\ge1}\). Its \emph{heart} is \(D^{\le0}\cap D^{\ge0}\); for the ordinary t-structure it is the category of sheaves. We write \(D^{[a,b]}=D^{\ge a}\cap D^{\le b}\). The triangle is unique up to unique isomorphism, so, as for the ordinary t-structure, it gives truncation functors and cohomology objects of \(K\) in the heart. Finally, a functor is \emph{t-exact} if it preserves both subcategories.

Let \(C\) be a smooth curve over a perfect field. On \(C\) the perverse t-structure is built from punctual sheaves.

% object: definition-punctual-sheaf (blueprint/chap-The-Fourier-transform/sec-Perverse-sheaves-on-a-curve/definition-punctual-sheaf/_section.md)
\phantomsection\label{bp197}%
\begin{bpstatement}{Definition~16.9}{ (Punctual sheaf)}
A constructible sheaf on \(C\) is \emph{punctual} if its support is a finite set of closed points.
\end{bpstatement}

% text: prose-16-11
\phantomsection\label{bp198}%
The truncation triangle of the perverse t-structure splits a sheaf into a punctual part and a part without punctual subsheaves, so we first construct the largest punctual subsheaf. In the following lemma “sheaf” means a constructible sheaf with coefficients \(E\) or \(\overline{\mathbf Q}_\ell\).

% object: lemma-largest-punctual-subsheaf (blueprint/chap-The-Fourier-transform/sec-Perverse-sheaves-on-a-curve/lemma-largest-punctual-subsheaf/_section.md)
\phantomsection\label{bp199}%
\begin{bpstatementit}{Lemma~16.10}{ (The largest punctual subsheaf)}
A lisse sheaf on \(C\) has no nonzero punctual subsheaf. Let \(F\) be a constructible sheaf on \(C\), let \(j:U\hookrightarrow C\) be a dense open on which \(F\) is lisse, and put

\[
F_{\rm p}=\ker(F\longrightarrow j_*j^*F).
\]

Then \(F_{\rm p}\) is the largest punctual subsheaf of \(F\), and \(F/F_{\rm p}\) has no nonzero punctual subsheaf. In particular, \(F_{\rm p}\) does not depend on \(U\).
\end{bpstatementit}

% proof: proof/lemma-largest-punctual-subsheaf
\begin{proof}[{Proof}]\phantomsection\label{bp200}%
Let \(L\) be a lisse sheaf. For a geometric point \(\bar x\) and a geometric generic point \(\bar\eta\) specializing to it, the generization map \(L_{\bar x}\to L_{\bar\eta}\) is injective. Indeed, for a sheaf defined over a finite extension \(E\) of \(\mathbf Q_\ell\), this holds for the finite locally constant reductions of the lattice in Proposition \hyperref[bp34]{13.10}, and we pass to the inverse limit, invert the uniformizer and extend scalars. Now a punctual subsheaf \(G\subset L\) has \(G_{\bar\eta}=0\), and since the generization maps of \(G\) and \(L\) are compatible, \(G_{\bar x}=0\).

Now let \(F\) and \(U\) be as in the statement. The sheaf \(F_{\rm p}\) is supported on the finite set \(C\setminus U\), hence punctual. Conversely, a punctual subsheaf \(G\) of \(F\) restricts to a punctual subsheaf of the lisse sheaf \(j^*F\), hence to zero. Since its map to \(j_*j^*F\) is adjoint to \(j^*G\to j^*F\), this map vanishes, and \(G\subset F_{\rm p}\). Finally, \(j_*j^*F\) has no nonzero punctual subsheaf. Indeed, if \(H\subset j_*j^*F\) is punctual, then \(j^*H=0\), and the inclusion is zero since \(\Hom(H,j_*j^*F)=\Hom(j^*H,j^*F)=0\). As \(F/F_{\rm p}\) embeds in \(j_*j^*F\), it has no nonzero punctual subsheaf either.
\end{proof}

% object: definition-perverse-sheaves-curve (blueprint/chap-The-Fourier-transform/sec-Perverse-sheaves-on-a-curve/definition-perverse-sheaves-curve/_section.md)
\phantomsection\label{bp201}%
\begin{bpstatement}{Definition~16.11}{ (Perverse sheaves on a curve)}
Let \(\Lambda\) be \(E\) or \(\overline{\mathbf Q}_\ell\). Define two classes of objects of \(D^b_c(C_{\mathrm{pro}\acute{e}t},\Lambda)\) by

\[
\begin{split}
{}^pD^{\le0}&=\{K:\mathcal H^aK=0\ (a>0),\
\mathcal H^0K\text{ is punctual}\},\\
{}^pD^{\ge0}&=\{K:\mathcal H^aK=0\ (a<-1),\
\mathcal H^{-1}K\text{ has no nonzero punctual subsheaf}\},
\end{split}
\]

and \({}^pD^{\le n}={}^pD^{\le0}[-n]\), \({}^pD^{\ge n}={}^pD^{\ge0}[-n]\). A \emph{perverse sheaf} on \(C\) is an object of \({}^pD^{\le0}\cap{}^pD^{\ge0}\). For \(C=C_0\times_{k_0}k\), a \emph{perverse Weil sheaf} on \(C_0\) is a perverse sheaf on \(C\) with a Weil structure, with morphisms as in Definition \bpout{10.18}.
\end{bpstatement}

% text: prose-16-12
\phantomsection\label{bp202}%
By Lemma \hyperref[bp203]{16.12} below these classes form a t-structure, the perverse t-structure. We write \({}^pD^{[a,b]}\) and \({}^p\mathcal H^n\) for its subcategories and cohomology objects.

On the line, the two basic examples are \(E[1]\) and the skyscraper \(\delta_0\), and by Lemma \hyperref[bp188]{16.7} the Fourier transform exchanges them up to a Tate twist. This is the simplest instance of the exactness proved in Proposition \hyperref[bp214]{16.16}.

% object: lemma-perverse-sheaves-curve (blueprint/chap-The-Fourier-transform/sec-Perverse-sheaves-on-a-curve/lemma-perverse-sheaves-curve/_section.md)
\phantomsection\label{bp203}%
\begin{bpstatementit}{Lemma~16.12}{ (Perverse sheaves on a curve)}
The classes \({}^pD^{\le0}\) and \({}^pD^{\ge0}\) of Definition \hyperref[bp201]{16.11} form a bounded t-structure on \(D^b_c(C_{\mathrm{pro}\acute{e}t},\overline{\mathbf Q}_\ell)\), and its heart, the category of perverse sheaves on \(C\), is abelian. A punctual sheaf in degree zero is perverse, and so is \((j_*M)[1]\) for a lisse sheaf \(M\) on a dense open \(j:U\hookrightarrow C\).
\end{bpstatementit}

% proof: proof/lemma-perverse-sheaves-curve
\begin{proof}[{Proof}]\phantomsection\label{bp204}%
The truncation triangle of \(K\) splits \(\mathcal H^0K\) into its largest punctual subsheaf, which goes into \({}^pD^{\le0}\), and the quotient, which goes into \({}^pD^{\ge1}\), using Lemma \hyperref[bp199]{16.10}.

The shift axioms are immediate from the definitions. For orthogonality, if \(A\in{}^pD^{\le0}\) and \(B\in{}^pD^{\ge1}\), then ordinary truncation identifies \(\Hom(A,B)\) with \(\Hom(\mathcal H^0A,\mathcal H^0B)\), which is zero because the source is punctual and the target has no punctual subsheaf. To construct the truncation triangle for \(K\), put \(K'=\tau^{\le0}K\) and \(F=\mathcal H^0K\), and let \(A\) be the fibre of \(K'\to F/F_{\rm p}\). Then the cohomology of \(A\) is that of \(K\) below zero, and \(\mathcal H^0A=F_{\rm p}\), so that \(A\in{}^pD^{\le0}\). By the octahedral axiom, applied to \(A\to K'\to K\), we obtain a triangle \(A\to K\to B\to\), in which the cone \(B\) is an extension of \(\tau^{\ge1}K\) by \(F/F_{\rm p}\) in degree zero, so it has no negative ordinary cohomology and \(\mathcal H^0B=F/F_{\rm p}\). Since \(F/F_{\rm p}\) has no punctual subsheaf by Lemma \hyperref[bp199]{16.10}, \(B\in{}^pD^{\ge1}\). Finally, the t-structure is bounded because the ordinary one is.

The heart of any t-structure is abelian, with kernels and cokernels computed from the perverse cohomology of a cone [\hyperref[bp353]{BBD82}, Théorème 1.3.6]; the proof uses only the octahedral axiom and orthogonality. Punctual sheaves in degree zero clearly lie in the heart. So does \((j_*M)[1]\), since \(j_*M\) has no punctual subsheaf, as follows from Lemma \hyperref[bp199]{16.10} applied to \(F=j_*M\), for which \(F_{\rm p}=0\).
\end{proof}

% text: prose-16-13
\phantomsection\label{bp205}%
Unwinding the definition, a perverse sheaf on \(C\) is a complex in ordinary degrees \(-1,0\) whose \(\mathcal H^0\) is punctual and whose \(\mathcal H^{-1}\) has no nonzero punctual subsheaf.

% object: lemma-perverse-weil-sheaves (blueprint/chap-The-Fourier-transform/sec-Perverse-sheaves-on-a-curve/lemma-perverse-weil-sheaves/_section.md)
\phantomsection\label{bp206}%
\begin{bpstatementit}{Lemma~16.13}{ (Perverse Weil sheaves)}
The perverse Weil sheaves on \(C_0\) form an abelian category, and the functor forgetting the Weil structure, to perverse sheaves on \(C\), is exact and reflects exactness.
\end{bpstatementit}

% proof: proof/lemma-perverse-weil-sheaves
\begin{proof}[{Proof}]\phantomsection\label{bp207}%
Frobenius pullback is an exact equivalence for the ordinary t-structure, and it preserves punctual support and the absence of punctual subsheaves. It is therefore t-exact for the perverse t-structure, which means that it carries the kernel and cokernel of a compatible morphism of perverse sheaves to the kernel and cokernel of the pulled-back morphism. The Weil isomorphisms and their inverses therefore induce inverse isomorphisms on the kernel and cokernel, and by naturality the coimage–image map is compatible with them. So perverse Weil sheaves form an abelian category, with kernels and cokernels computed in perverse sheaves on \(C\), and this is the assertion about the forgetful functor.
\end{proof}

% text: prose-16-14
\phantomsection\label{bp208}%
An exact equivalence of abelian categories sends simple objects to simple objects. Once we know that the transform is exact on the heart, the following list of simple objects will tell us what the transform of a middle extension can be.

% object: lemma-simple-objects (blueprint/chap-The-Fourier-transform/sec-Perverse-sheaves-on-a-curve/lemma-simple-objects/_section.md)
\phantomsection\label{bp209}%
\begin{bpstatementit}{Lemma~16.14}{ (The simple objects)}
In the category of perverse sheaves on \(C\), and in that of perverse Weil sheaves on \(C_0\), the simple objects are the irreducible punctual objects and the objects \((j_*M)[1]\) with \(M\) nonzero irreducible lisse on a dense open, where irreducibility is taken in the geometric or in the Weil category, respectively. These middle extensions have no nonzero punctual subobject or quotient in the perverse heart.
\end{bpstatementit}

% proof: proof/lemma-simple-objects
\begin{proof}[{Proof}]\phantomsection\label{bp210}%
Since a middle extension has no punctual subobject or quotient and \(M\) is irreducible on the open, the middle extension is simple. Conversely, ordinary truncation splits an arbitrary perverse sheaf into a punctual part and a part of the form \(F[1]\), and simplicity forces one of them to vanish.

We first show that \(P=(j_*M)[1]\) has no punctual subobject or quotient. The adjunction calculations below are geometric, so they apply equally to compatible Weil morphisms. Let \(i\) be the inclusion of a finite set of closed points and \(A\) a sheaf on it. A map \(P\to i_*A\) is zero, because by adjunction it corresponds to a map from \(i^*P=(i^*j_*M)[1]\) to the sheaf \(A\), that is, to a negative Ext. A map in the other direction is also zero, and to see this we first enlarge the boundary. Since \(M\) is lisse, \(j_*M\) does not change when we replace \(U\) by a smaller dense open and \(M\) by its restriction, so we may assume that the support of \(A\) lies in \(Z=C\setminus U\), and we let \(i\) now denote the inclusion of \(Z\). Then the adjunction \(j^*\dashv Rj_*\) shows that \(i^!Rj_*M=0\). In the triangle

\[
j_*M\longrightarrow Rj_*M\longrightarrow
\tau^{\ge1}Rj_*M\longrightarrow
\]

the last term is supported on \(Z\), so that \(i^!\) agrees with \(i^*\) on it. Applying \(i^!\) and using \(i^!Rj_*M=0\), we find \(i^!j_*M\simeq(i^*\tau^{\ge1}Rj_*M)[-1]\), which has ordinary cohomology only in degrees at least two. Hence \(i^!P\) has ordinary cohomology only in degrees at least one. Therefore \(\Hom(i_*A,P)=\Hom(A,i^!P)=0\).

To see that \(P\) is simple, note that on a dense open a perverse subobject of \(P\) is a shifted lisse subsheaf of \(M\). By irreducibility it is \(0\) or everything there, and the difference is then punctual. In detail, consider a perverse exact sequence \(0\to Q\to P\to R\to0\), restrict it to a smaller dense open on which all ordinary cohomology sheaves are lisse. There \(P=M[1]\), and the sequence is a shifted exact sequence of lisse sheaves. Moreover, restriction to a dense open preserves irreducibility, because the fundamental group of a dense open of a normal connected curve surjects onto that of the curve, as in the proof of Theorem \hyperref[bp120]{14.26}, and the same holds for the Weil groups. Hence \(Q\) or \(R\) has zero generic restriction. Such a perverse object is punctual. Indeed, its \(\mathcal H^{-1}\) is then supported on a finite set, and a sheaf without punctual subsheaves with this property is zero, while its \(\mathcal H^0\) is punctual by definition. Thus \(Q\) or \(R\) is a punctual subobject or quotient of \(P\), which by the vanishing just proved is zero, and so \(P\) is simple.

Conversely, let \(P\) be simple. We show that \(P\) is punctual or a middle extension of an irreducible lisse sheaf. Ordinary truncation gives a perverse exact sequence

\[
0\longrightarrow(\mathcal H^{-1}P)[1]
\longrightarrow P\longrightarrow\mathcal H^0P\longrightarrow0.
\]

Hence a simple \(P\) is either punctual or of the form \(F[1]\), where \(F\) has no punctual subsheaf. In the second case, choose a dense open \(U\) on which \(F\) is lisse and put \(M=F|_U\). Then the map \(F\to j_*M\) is injective, since its kernel is \(F_{\rm p}=0\) (Lemma \hyperref[bp199]{16.10}), and its cokernel \(B\) is punctual. Rotating the corresponding triangle of sheaves we obtain the perverse exact sequence \(0\to B\to F[1]\to(j_*M)[1]\to0\), whose last term is nonzero, so \(B=0\) by simplicity. It remains to see that \(M\) is irreducible. Suppose that \(M\) had a proper nonzero lisse subsheaf \(M_1\). Then the object \(j_!M_1[1]\) is perverse, because \(j_!M_1\subset j_*M_1\) has no punctual subsheaf, and the image of \(j_!M_1[1]\to P\) in the abelian heart is proper and nonzero, as one sees on \(U\). This contradicts simplicity. Finally, if \(P\) is punctual, then by the long exact cohomology sequence its subobjects in the heart are punctual sheaves in degree zero, so \(P\) is simple exactly when it is irreducible as a punctual sheaf.
\end{proof}

% object: example-extension-by-zero-not-simple (blueprint/chap-The-Fourier-transform/sec-Perverse-sheaves-on-a-curve/example-extension-by-zero-not-simple/_section.md)
\phantomsection\label{bp211}%
\begin{bpstatement}{Example~16.15}{ (Extension by zero is not simple)}
Let \(j:\mathbf G_m\hookrightarrow\mathbf A^1\) and \(i:0\hookrightarrow\mathbf A^1\) be the inclusions over \(k\), and put \(\delta_0=i_*E\). The sheaf \(j_!E\) has no punctual subsheaf, as its stalk at the origin is zero, so \(j_!E[1]\) is perverse. It is not simple. Rotating the triangle of the exact sequence \(0\to j_!E\to E\to\delta_0\to0\) gives the triangle \(\delta_0\to j_!E[1]\to E[1]\to\). Its three terms are perverse, and so

\[
0\longrightarrow\delta_0\longrightarrow j_!E[1]
\longrightarrow E[1]\longrightarrow0
\]

is exact in the perverse heart. Thus, while among sheaves the skyscraper is a quotient of \(E\), in the perverse heart it becomes a subobject of \(j_!E[1]\). By contrast, the middle extension of \(E|_{\mathbf G_m}\) is \(E\) itself (Lemma \hyperref[bp127]{15.2}), and \(E[1]\) is simple. On the other hand \((E\oplus\delta_0)[1]\) is not perverse, because \(\mathcal H^{-1}\) has the punctual subsheaf \(\delta_0\), which is the subsheaf \(F_{\rm p}\) of Lemma \hyperref[bp199]{16.10}, the kernel of \(E\oplus\delta_0\to j_*j^*(E\oplus\delta_0)=E\).
\end{bpstatement}


% object: sec-Perverse-t-exactness (blueprint/chap-The-Fourier-transform/sec-Perverse-t-exactness/_section.md), content hidden
\subsection*{\texorpdfstring{16.5\quad Perverse t-exactness}{16.5 Perverse t-exactness}}
\phantomsection\label{bp212}%
\addcontentsline{toc}{subsection}{\texorpdfstring{16.5\quad Perverse t-exactness}{16.5 Perverse t-exactness}}
% text: prose-16-15
\phantomsection\label{bp213}%
The transform \(\operatorname{FT}_\psi\) is t-exact for the perverse t-structure on the line (Proposition \hyperref[bp214]{16.16}). By the stalk formula (16.2), this comes down to the compactly supported cohomology of sheaves on an affine curve in the extreme degrees zero and two.

% object: proposition-fourier-transform-perverse-heart (blueprint/chap-The-Fourier-transform/sec-Perverse-t-exactness/proposition-fourier-transform-perverse-heart/_section.md)
\phantomsection\label{bp214}%
\begin{bpstatementit}{Proposition~16.16}{ (Fourier transform and the perverse heart)}
On the geometric affine line, \(\operatorname{FT}_\psi\) is a t-exact equivalence of bounded constructible derived categories for the perverse t-structure of Lemma \hyperref[bp203]{16.12}, with inverse \(\operatorname{FT}_{\psi^{-1}}(-)(1)\). Hence it is an exact equivalence of the perverse hearts, and its canonical Frobenius compatibilities induce an exact equivalence of the categories of perverse Weil sheaves.
\end{bpstatementit}

% proof: proof/proposition-fourier-transform-perverse-heart
\begin{proof}[{Proof}]\phantomsection\label{bp215}%
We show that the geometric transform and its inverse preserve \({}^pD^{\le0}\). Being mutually inverse, they then preserve \({}^pD^{\ge0}\) by orthogonality, and the Weil statement follows by Lemma \hyperref[bp206]{16.13}, since the forgetful functor on perverse Weil sheaves is exact and reflects exactness. For the first step, the key point is that the stalks of the transform are compactly supported cohomology groups on the affine line, in degrees \(0,1,2\), and that the group in degree two is nonzero only at the finitely many points \(y\) for which the twisted sheaf has a geometrically constant quotient.

\emph{Right exactness.} Since the transform is triangulated and \({}^pD^{\le0}\) is stable under extensions and positive shifts, ordinary truncation reduces this to punctual sheaves in degree zero and to objects \(G[1]\) with \(G\) constructible. Indeed, an object of \({}^pD^{\le0}\) is an iterated extension of its shifted ordinary cohomology sheaves \(\mathcal H^a(K)[-a]\) with \(a\le0\); the one with \(a=0\) is punctual, and the others are of the form \(G[1]\) shifted by a nonnegative amount. Punctual sheaves map into \({}^pD^{\le0}\), since by the stalk formula the transform of a geometric skyscraper is a lisse sheaf shifted by one.

For \(G[1]\) we show that the transform lies in ordinary degrees \(-2,-1,0\) and that its degree-zero sheaf has finite support. The first assertion follows from the stalk formula and Lemma \hyperref[bp143]{15.8}. For the second, note that the degree-zero stalk at a closed geometric point \(y\) is

\[
H_c^2(\mathbf A^1,G\otimes\mathcal L_{\psi(yx)}).
\]

On a dense open where \(G\) is lisse with generic representation \(V\), write \(\chi_y\) for the character of \(\pi_1\) given by the restriction of \(\mathcal L_{\psi(yx)}\). By localization we may discard the finitely many remaining points, which contribute only in degree zero, and Lemma \hyperref[bp140]{15.7} (whose proof of the formula for \(H^2_c\) uses only the geometric curve) shows that this stalk is \((V\otimes\chi_y)_{\pi_1}(-1)\). In particular, it is nonzero only if \(V\) has the rank-one quotient \(\chi_y^{-1}\), and we show that this happens for only finitely many \(y\). First, the characters \(\chi_y\) are pairwise distinct. Indeed, if \(\chi_y=\chi_{y'}\) with \(y\ne y'\), then \(\mathcal L_{\psi((y-y')x)}\) is trivial on the dense open, hence on the line by the surjectivity of fundamental groups. This is impossible, because its compactly supported cohomology vanishes by Lemma \hyperref[bp179]{16.4}, whereas the constant sheaf has nonzero top cohomology with compact support. Second, if \(V\) has pairwise nonisomorphic rank-one quotients \(\chi_1,\dots,\chi_n\), then the induced map \(V\to\chi_1\oplus\dots\oplus\chi_n\) is surjective. To see this, note that its image is a subrepresentation of a sum of pairwise nonisomorphic simple representations, hence the sum of some of the \(\chi_i\), and it maps onto each \(\chi_i\). So \(n\le\dim V\), and \(\chi_y^{-1}\) is a quotient of \(V\) for at most \(\dim V\) points \(y\). The degree-zero sheaf therefore has only finitely many nonzero closed geometric stalks. Since it is constructible, it has finite support, and the transform of \(G[1]\) belongs to \({}^pD^{\le0}\).

\emph{Exactness.} The argument for \(\operatorname{FT}_\psi\) applies equally to \(\psi^{-1}\), and a Tate twist does not affect supports or degrees. Both \(\operatorname{FT}_\psi\) and its quasi-inverse \(\operatorname{FT}_{\psi^{-1}}(-)(1)\) of Proposition \hyperref[bp185]{16.6} therefore preserve \({}^pD^{\le0}\). Since \({}^pD^{\ge0}\) is the right orthogonal of \({}^pD^{\le-1}\), the two mutually inverse equivalences also preserve \({}^pD^{\ge0}\). Both functors are therefore t-exact, and they induce exact equivalences of the geometric hearts. With their canonical Frobenius compatibilities and the equivariant inversion isomorphisms of Proposition \hyperref[bp185]{16.6}, they then give quasi-inverse functors on perverse Weil sheaves, which are exact because the forgetful functors are exact and reflect exactness.
\end{proof}


% object: sec-Middle-extensions (blueprint/chap-The-Fourier-transform/sec-Middle-extensions/_section.md), content hidden
\subsection*{\texorpdfstring{16.6\quad Middle extensions}{16.6 Middle extensions}}
\phantomsection\label{bp216}%
\addcontentsline{toc}{subsection}{\texorpdfstring{16.6\quad Middle extensions}{16.6 Middle extensions}}
% text: prose-16-16
\phantomsection\label{bp217}%
Let \(j:U_0\hookrightarrow\mathbf A^1_{k_0}\) be a dense open, and let \(M\) be an irreducible lisse Weil sheaf on \(U_0\). The argument of Chapter \hyperref[bp223]{17} needs the transform of the middle extension of \(M\) to be again a middle extension, that is, to come from a lisse sheaf \(M'\) on a dense open of the dual line. This holds unless \(M\) has a geometric constituent of the form \(\mathcal L_{\psi(ax)}\), and the hypothesis of the next proposition excludes these. They behave differently: for \(a\ne0\) the sheaf \(\mathcal L_{\psi(ax)}\) has no cohomology with compact support on the line by Lemma \hyperref[bp179]{16.4}, and its transform is a skyscraper rather than a lisse sheaf of weight one. Excluding them, however, costs nothing in the application, because a sheaf unramified at infinity has no such constituent other than the constant one, and constant constituents contribute nothing, since \(H^1_c(\mathbf A^1,\overline{\mathbf Q}_\ell)=0\).

% object: proposition-fourier-transform-middle-extension (blueprint/chap-The-Fourier-transform/sec-Middle-extensions/proposition-fourier-transform-middle-extension/_section.md)
\phantomsection\label{bp218}%
\begin{bpstatementit}{Proposition~16.17}{ (The Fourier transform of a middle extension)}
Suppose that \(M\) is nonzero and irreducible, and that none of its geometric irreducible constituents is the restriction of \(\mathcal L_{\psi(ax)}\) for some \(a\in\bar k_0\) (the case \(a=0\) being the constant sheaf). Then there are a dense open \(j':U'_0\hookrightarrow\mathbf A^{1\vee}_{k_0}\) and a nonzero irreducible lisse Weil sheaf \(M'\) on it such that

\[
\operatorname{FT}_\psi((j_*M)[1])\simeq(j'_*M')[1].
\]
\end{bpstatementit}

% proof: proof/proposition-fourier-transform-middle-extension
\begin{proof}[{Proof}]\phantomsection\label{bp219}%
By Fourier inversion, a punctual transform would force every geometric constituent of \(M\) to be a linear character sheaf, and the hypothesis excludes this.

By Proposition \hyperref[bp214]{16.16}, the transform of the simple object \((j_*M)[1]\) is simple, and so by Lemma \hyperref[bp209]{16.14} it is either punctual or of the displayed form. Restricting the kernel to \(\{a\}\times\mathbf A^1\) shows that the inverse transform of a geometric skyscraper at \(a\) is \(\mathcal L_{\psi(-ax)}(1)[1]\). Now, if the transform were punctual, it would geometrically be an iterated extension of skyscrapers at finitely many points \(a\), so by Fourier inversion and exactness \((j_*M)[1]\) would be geometrically an iterated extension of the objects \(\mathcal L_{\psi(-ax)}(1)[1]\). Restricting to \(U\), we see that every geometric constituent of \(M\) would be one of these linear characters, which the hypothesis excludes. Finally, the open \(U'_0\) and the sheaf \(M'\) are defined over \(k_0\), since we may choose a Weil-stable open on which the constructible transform is lisse.
\end{proof}

% text: prose-16-17
\phantomsection\label{bp220}%
The sheaf \(M'\) is determined by \(M\) up to isomorphism and shrinking of \(U'_0\), since it is the restriction to \(U'_0\) of \(\mathcal H^{-1}\operatorname{FT}_\psi((j_*M)[1])\). In the application we check the hypothesis of the proposition with the following criterion, whose proof is a computation of ramification at infinity.

% object: lemma-unramified-at-infinity-no-linear-characters (blueprint/chap-The-Fourier-transform/sec-Middle-extensions/lemma-unramified-at-infinity-no-linear-characters/_section.md)
\phantomsection\label{bp221}%
\begin{bpstatementit}{Lemma~16.18}{ (Sheaves unramified at infinity)}
Let \(M\) be an irreducible lisse Weil sheaf on a dense open \(U_0\subset\mathbf A^1_{k_0}\). If \(M\) is unramified at infinity and is not geometrically constant, then no geometric constituent of \(M\) is of the form \(\mathcal L_{\psi(ax)}\) with \(a\in\bar k_0\).
\end{bpstatementit}

% proof: proof/lemma-unramified-at-infinity-no-linear-characters
\begin{proof}[{Proof}]\phantomsection\label{bp222}%
For \(a\ne0\) the sheaf \(\mathcal L_{\psi(ax)}\) is ramified at infinity, and the trivial character is excluded because \(M\) is irreducible and not geometrically constant, its geometric restriction being semisimple. For the first point, write \(t=1/x\) and express the character by the Artin–Schreier equation \(u^p-u=c/t\) with \(c\ne0\). Let \(v\) be the normalized valuation of the field generated by a root \(u\) over the completion \(k((t))\) at infinity, and let \(e\) be the ramification index, so that \(v(t)=e\) and \(e\le p\). Since \(v(c/t)=-e<0\), we must have \(v(u)<0\), so that \(v(u^p)=pv(u)<v(u)\) and \(v(u^p-u)=pv(u)\). Now \(pv(u)=-e\) forces \(p\mid e\), and so \(e=p\): the cover is totally ramified. Inertia at infinity therefore maps onto the Galois group \(\mathbf F_p\) of this cover, on which \(\psi\) is nontrivial. Since \(M\) is unramified at infinity, this proves the first point, because a character that is nontrivial on inertia cannot be a constituent of a representation on which inertia acts trivially.

For the trivial character, we apply Lemma \hyperref[bp75]{14.11} to the geometric fundamental group, which is normal in the Weil group: the restriction of \(M\) to it is semisimple, and so a trivial geometric constituent gives nonzero geometric invariants. Since this space of invariants is Weil-stable, the irreducibility of \(M\) shows that \(M\) is geometrically constant.
\end{proof}



% object: chap-The-Deligne-Laumon-theorem-for-curves (blueprint/chap-The-Deligne-Laumon-theorem-for-curves/_section.md), content hidden
\section*{\texorpdfstring{17\quad The Deligne–Laumon theorem for curves}{17 The Deligne–Laumon theorem for curves}}
\phantomsection\label{bp223}%
\addcontentsline{toc}{section}{\texorpdfstring{17\quad The Deligne–Laumon theorem for curves}{17 The Deligne–Laumon theorem for curves}}
% text: prose-17-01
\phantomsection\label{bp224}%
Let \(F\) be a pure lisse Weil sheaf of weight \(w\) on a smooth curve \(C\). Why should \(H^1_c(C,F)\) have weights at most \(w+1\)? For the constant sheaf on a smooth projective curve this is Weil's theorem, proved with the Jacobian, and for an arbitrary \(F\) there is no such object to work with. Yet the Deligne–Laumon theorem comes down to exactly this bound. The theorem says that \(H^r_c(C,F)\) has weights at most \(w+r\) (Theorem \hyperref[bp269]{17.15}), and only \(r=1\) requires work. Indeed, on an affine curve \(H^0_c\) vanishes, while by Lemma \hyperref[bp140]{15.7} \(H^2_c\) is a twist of the coinvariants of geometric monodromy, pure of weight \(w+2\). For \(r=1\) the bound can be strict: on \(\mathbf G_m\) the group \(H^1_c\) comes from the boundary and has weight \(0\) (see the example of Chapter \hyperref[bp3]{13}).

For a Gauss sum the answer takes two steps. The transform \(\hat\chi(y)=\overline\chi(y)g\) has the same absolute value for all \(y\ne0\), and Plancherel finds the common value \(\sqrt q\). On an open of the affine line, Chapter \hyperref[bp167]{16} makes the same two steps available for sheaves. More precisely, let \(M\) be irreducible on a dense open \(j:U_0\hookrightarrow\mathbf A^1\), with no geometric constituent of the form \(\mathcal L_{\psi(ax)}\), and let \(M'\) be the lisse sheaf with \(\operatorname{FT}_\psi((j_*M)[1])\simeq(j'_*M')[1]\) (Proposition \hyperref[bp218]{16.17}). Then

\[
H^1_c(\mathbf A^1,j_*M)\simeq(j'_*M')_{0'}
\]

(Corollary \hyperref[bp241]{17.6}). Since a stalk of a middle extension has weights at most the weight of the lisse sheaf (Proposition \hyperref[bp134]{15.5}), it suffices to show that \(M'\) is pure of weight \(w+1\). Here the second step is the easy one, because once \(M'\) is pure, the sums of squared traces of Chapter \hyperref[bp122]{15} give its weight \(w+1\) (Theorem \hyperref[bp238]{17.5}).

The first step, purity of \(M'\), is the deep one. The only purity theorem we have is Theorem \hyperref[bp120]{14.26}, for real sheaves. But the transform of a real sheaf is not real, since \(\psi\) takes nonreal values. For example, over \(\mathbf F_3\) let \(M\) be the Kummer sheaf of the quadratic character \(\chi\) on \(\mathbf G_m\), which is real. By the computation of Chapter \hyperref[bp167]{16} and the signs in (16.3), the trace of \(M'\) at \(y\in\mathbf F_3^\times\) is \(-\chi(y)g\). Here \(g=\sum_x\chi(x)\psi(x)=\omega-\omega^2=1+2\omega\) is a square root of \(-3\), which is not real. In this case, admittedly, a twist by \(\mathcal B_{1/g}\) makes \(M'\) real, but in general no twist is known to do this.

The plan is to realize \(M'\), up to a twist, as a quotient of a real sheaf on the dual line, to which that theorem applies. To build it we add a conjugate. First we twist \(M\) to weight zero, \(M^0=M\otimes\mathcal B_b\). Then we transform \(j_!M^0\) with the kernel \(\mathcal K_\psi\) and \(j_!(M^0)^\vee\) with the kernel \(\mathcal K_{\psi^{-1}}\). At each point the traces of the two transforms are complex conjugate, so their sum is a real sheaf \(H\) on the dual line. Since the twist \(M'\otimes\mathcal B_b\) is a quotient of \(H\) (Proposition \hyperref[bp234]{17.4}), purity of real sheaves then makes \(M'\) pure. We use \(j_!\) rather than \(j_*\) because its transform is a sheaf, and by Lemma \hyperref[bp227]{17.1} this transform differs from the transform of \(j_*M\) by a lisse sheaf of weights at most \(w\).

Passing from one irreducible \(M\) to every pure sheaf needs care, because \(j_*\) is not exact in general. For example, the extension \(0\to\overline{\mathbf Q}_\ell(1)\to F\to\overline{\mathbf Q}_\ell\to0\) on \(\mathbf G_m\) given by the Kummer class of the coordinate is lisse, but \(j_*F\to j_*\overline{\mathbf Q}_\ell\) is not surjective at \(0\), since the monodromy of \(F\) around \(0\) is a nontrivial unipotent matrix. For sheaves of one weight \(w\), however, exactness holds by a weight gap (Proposition \hyperref[bp249]{17.9}): by local duality at the boundary points (Lemma \hyperref[bp245]{17.7}), the connecting map \(j_*F_2\to R^1j_*F_1\) goes from weights at most \(w\) to weights at least \(w+2\), and so it vanishes.

% object: sec-Purity-of-the-lisse-part (blueprint/chap-The-Deligne-Laumon-theorem-for-curves/sec-Purity-of-the-lisse-part/_section.md), content hidden
\subsection*{\texorpdfstring{17.1\quad Purity of the lisse part}{17.1 Purity of the lisse part}}
\phantomsection\label{bp225}%
\addcontentsline{toc}{subsection}{\texorpdfstring{17.1\quad Purity of the lisse part}{17.1 Purity of the lisse part}}
% text: prose-17-02
\phantomsection\label{bp226}%
Throughout this section and the next, \(j:U_0\hookrightarrow\mathbf A^1_{k_0}\) is a dense open and \(M\) is a nonzero irreducible lisse Weil sheaf on \(U_0\), pure of weight \(w\). We assume that no geometric irreducible constituent of \(M\) is the restriction of \(\mathcal L_{\psi(ax)}\) for some \(a\in\bar k_0\), and we let \(M'\) be the lisse Weil sheaf on a dense open \(j':U'_0\hookrightarrow\mathbf A^{1\vee}_{k_0}\) with \(\operatorname{FT}_\psi((j_*M)[1])\simeq(j'_*M')[1]\) (Proposition \hyperref[bp218]{16.17}).

We build the real sheaf from extensions by zero, because for them we can check reality on the trace function. Indeed, at a point of \(U_0\) the eigenvalues of \(M^0\) and of its dual have absolute value one, and hence they are complex conjugate. At a boundary point, however, the stalk of \(j_*M^0\) is the space of inertia invariants, and the stalk of \(j_*(M^0)^\vee\) is the space of invariants of the dual representation, but nothing we know at this stage makes their Frobenius eigenvalues conjugate. With \(j_!\), by contrast, both stalks vanish there, and the trace function of the sum of the two transforms is clearly real.

The price is that \(j_!M\) is not a simple perverse sheaf, so its transform \(G\) is not \(j'_*M'\) but an extension of \(j'_*M'\) by the transform \(B'\) of the boundary. This \(B'\) is a lisse sheaf of weights at most \(w\) (Lemma \hyperref[bp227]{17.1}), and in the proof of Proposition \hyperref[bp234]{17.4} we use only that \(G\) is a sheaf and that it surjects onto \(j'_*M'\).

% object: lemma-fourier-transform-boundary-sequence (blueprint/chap-The-Deligne-Laumon-theorem-for-curves/sec-Purity-of-the-lisse-part/lemma-fourier-transform-boundary-sequence/_section.md)
\phantomsection\label{bp227}%
\begin{bpstatementit}{Lemma~17.1}{ (Fourier transform of the boundary sequence)}
Put \(B=j_*M/j_!M\). There are sheaves \(G\) and \(B'\) on the dual line such that \(\operatorname{FT}_\psi(j_!M)=G\), \(\operatorname{FT}_\psi(B)=B'[1]\), and the sequence

\[
0\longrightarrow B'\longrightarrow G
\longrightarrow j'_*M'\longrightarrow0
\]

is exact. The sheaf \(B'\) is lisse on the entire dual line and has weights at most \(w\).
\end{bpstatementit}

% proof: proof/lemma-fourier-transform-boundary-sequence
\begin{proof}[{Proof}]\phantomsection\label{bp228}%
The three assertions are of different kinds. That the transform of \(j_!M\) is a single sheaf is a vanishing statement for the cohomology of the twisted sheaves \(M\otimes\mathcal L_{\psi(yx)}\) on \(U\), whereas the exact sequence comes from transforming the boundary sequence \(0\to j_!M\to j_*M\to B\to0\) term by term. Finally, the bound on \(B'\) comes from the stalks of \(B\), which are known from Proposition \hyperref[bp134]{15.5}.

For \(G\), the stalk formula of Proposition \hyperref[bp183]{16.5} says that the stalk of \(\operatorname{FT}_\psi(j_!M)\) at \(y\) is \(R\Gamma_c(U,M\otimes\mathcal L_{\psi(yx)})[1]\), so we must show that only \(H^1_c\) is nonzero. We apply Lemma \hyperref[bp140]{15.7} to the lisse sheaf \(M\otimes\mathcal L_{\psi(yx)}\) on the affine curve \(U\), with \(j\) the identity. It shows that the group in degree zero vanishes and that the group in degree two is a twist of the coinvariants of geometric monodromy on \(M\otimes\mathcal L_{\psi(yx)}\). But a nonzero coinvariant would be a geometric quotient of \(M\) isomorphic to \(\mathcal L_{\psi(-yx)}\), and our assumption on \(M\) excludes such a constituent, so the group in degree two vanishes as well. The shift then places the remaining group \(H^1_c\) in degree zero, and the transform is a sheaf \(G\).

The transform of \(B\) is computed point by point. After geometric base change, the punctual sheaf \(B\) is a finite sum of skyscrapers with coefficient modules, and the transform of each skyscraper is a linear character sheaf tensored with that module and shifted by one. Since a linear character sheaf is lisse on the whole dual line, \(\operatorname{FT}_\psi(B)=B'[1]\) with \(B'\) lisse.

Applying the Fourier transform to the triangle obtained from \(0\to j_!M\to j_*M\to B\to0\), and using Proposition \hyperref[bp218]{16.17} for the middle term, we obtain a triangle \(G\to j'_*M'\to B'[1]\to\). Here \(G\) and \(j'_*M'\) are sheaves in degree zero, and \(B'[1]\) is a sheaf in degree \(-1\). Hence the long exact sequence of cohomology sheaves has only three nonzero terms, and it is the stated exact sequence.

The weights of \(B'\) come from those of \(B\), because a character sheaf has weight zero, and the weights of \(B\) are at most \(w\) by Proposition \hyperref[bp134]{15.5}. Only at boundary points that are not rational does this need an argument. So let \(s\) be a boundary point of degree \(d\) and \(y\) a rational point of the dual line. The geometric points over \(s\) contribute \(d\) summands to \(B'_{\bar y}\), which \(\operatorname{Frob}_y\) permutes cyclically, and its \(d\)th power acts on one summand as \(\operatorname{Frob}_s\) on \(B_{\bar s}\) times a root of unity coming from the character sheaf. Every eigenvalue \(\beta\) of \(\operatorname{Frob}_y\) on these summands therefore satisfies \(\beta^d=\zeta\alpha\), with \(\zeta\) a root of unity and \(\alpha\) an eigenvalue of \(\operatorname{Frob}_s\). Since \(|\iota\alpha|\le q^{dw/2}\), we get \(|\iota\beta|\le q^{w/2}\). At a point \(y\) of higher degree we replace \(q\) by the cardinality of \(\kappa(y)\).

For a filtration of \(B'\) with pure quotients, choose for each boundary point, over its residue field, a Frobenius-stable full flag of the coefficient module, which exists because \(\overline{\mathbf Q}_\ell\) is algebraically closed. Each successive quotient is then a line, on which Frobenius acts by one of its eigenvalues. Tensoring with the linear character sheaves and summing over the Frobenius orbit, we obtain a finite Weil filtration of \(B'\) whose quotients are pure of the weights of the coefficient eigenvalues, and these weights are at most \(w\).
\end{proof}

% object: example-kummer-no-boundary (blueprint/chap-The-Deligne-Laumon-theorem-for-curves/sec-Purity-of-the-lisse-part/example-kummer-no-boundary/_section.md)
\phantomsection\label{bp229}%
\begin{bpstatement}{Example~17.2}{ (A Kummer sheaf has no boundary term)}
Let \(M\) be the Kummer sheaf of a nontrivial character \(\chi\) on \(U_0=\mathbf G_m\). It is irreducible and pure of weight \(0\). Since it is tame at infinity, it has no geometric constituent \(\mathcal L_{\psi(ax)}\), and so it satisfies the hypotheses of this section. Inertia at the origin acts on \(M\) through \(\chi\), nontrivially, and hence \((j_*M)_{\bar0}=0\) by Lemma \hyperref[bp127]{15.2}. This means that \(j_!M=j_*M\) and \(B=0\), and the lemma gives \(G\simeq j'_*M'\). On \(\mathbf F_q\)-points this is visible: by (16.3) the trace of \(G\) at \(y\ne0\) is \(-\sum_x\chi(x)\psi(xy)=-\overline\chi(y)g\), with \(g=\sum_x\chi(x)\psi(x)\), and for the quadratic character this is the trace of \(M'\) computed at the start of Chapter \hyperref[bp223]{17}. At \(y=0\), on the other hand, the trace of \(G\) is \(-\sum_x\chi(x)=0\), in agreement with \(H^1_c(\mathbf G_m,M)=0\).
\end{bpstatement}

% text: prose-17-03
\phantomsection\label{bp230}%
The argument of Proposition \hyperref[bp148]{15.9} also applies to twists by the Fourier kernel, and we need this in the proof of Proposition \hyperref[bp234]{17.4}.

% object: lemma-twisted-envelope (blueprint/chap-The-Deligne-Laumon-theorem-for-curves/sec-Purity-of-the-lisse-part/lemma-twisted-envelope/_section.md)
\phantomsection\label{bp231}%
\begin{bpstatementit}{Lemma~17.3}{ (Twisting and adding the dual)}
Let \(\mathcal A\) be lisse and pure of weight zero, and let \(\mathcal L\) be lisse of rank one with local eigenvalues of finite order. Then

\[
(\mathcal A\otimes\mathcal L)
\oplus(\mathcal A^\vee\otimes\mathcal L^{-1})
\]

is real and pure of weight zero, and its extension by zero along an open immersion is real as well.
\end{bpstatementit}

% proof: proof/lemma-twisted-envelope
\begin{proof}[{Proof}]\phantomsection\label{bp232}%
As in the proof of Proposition \hyperref[bp148]{15.9}, the local eigenvalues of \(\mathcal A\otimes\mathcal L\) have absolute value one after applying \(\iota\), and those of the second summand are their inverses, hence their complex conjugates. The extension by zero is real as well, because its stalks outside the open are zero, and so its local characteristic polynomials there are equal to \(1\).
\end{proof}

% text: prose-17-04
\phantomsection\label{bp233}%
In the next proof we use Lemma \hyperref[bp227]{17.1} twice. Applying its vanishing argument to \(M^0\) and to its dual, we see that the transform of the real sheaf is a single sheaf. Then, using the surjection \(G\to j'_*M'\) of the lemma, we exhibit a twist of \(M'\) as a quotient of that real sheaf.

% object: proposition-purity-lisse-part-fourier-transform (blueprint/chap-The-Deligne-Laumon-theorem-for-curves/sec-Purity-of-the-lisse-part/proposition-purity-lisse-part-fourier-transform/_section.md)
\phantomsection\label{bp234}%
\begin{bpstatementit}{Proposition~17.4}{ (Purity of the lisse part of the Fourier transform)}
The sheaf \(M'\) is pure of some real weight \(v\) on the whole open where it is lisse.
\end{bpstatementit}

% proof: proof/proposition-purity-lisse-part-fourier-transform
\begin{proof}[{Proof}]\phantomsection\label{bp235}%
We twist to weight zero and add the dual, as in Proposition \hyperref[bp148]{15.9}, before taking the transform, so that a twist of \(M'\) becomes a quotient of a real sheaf, to which Theorem \hyperref[bp120]{14.26} applies.

Choose \(b\) as in Proposition \hyperref[bp148]{15.9}, so that \(M^0=M\otimes\mathcal B_b\) is pure of weight zero.

On the product of the line and the dual line, consider the sheaf

\[
F=
(p_1^*j_!M^0\otimes\mathcal K_\psi)
\ \oplus\
(p_1^*j_!(M^0)^\vee\otimes\mathcal K_{\psi^{-1}}).
\]

At each closed point, the Frobenius eigenvalues of the second summand are the complex conjugates of those of the first, because the eigenvalues of \(M^0\) have absolute value one and those of the character sheaves are roots of unity, while at the boundary both zero extensions have zero stalks. Hence \(F\) is real, as in Lemma \hyperref[bp231]{17.3}. Moreover, since duality replaces \(a\) by \(-a\), the dual \((M^0)^\vee\) also has no geometric constituent of the form \(\mathcal L_{\psi(ax)}\), so the argument of Lemma \hyperref[bp227]{17.1} shows that the direct image with compact support of \(F\) along \(p_2\) is concentrated in degree one. We set \(H=R^1p_{2!}F\).

Let \(y\) be a point with values in a finite field. Since \(Rp_{2!}F=H[-1]\), the trace formula for \(Rp_{2!}\) (Corollary \bpout{12.26}) gives \(-t_H(y)=\sum_x t_F(x,y)\). For each \(x\) the traces of the two summands of \(F\) at \((x,y)\) are complex conjugate, so \(t_H(y)\) is real. Over the finite extensions of \(\kappa(y)\) the same argument shows that the trace of every power of the local Frobenius is real. These traces are the power sums of the local eigenvalues, and the Newton identities express the coefficients of the characteristic polynomial through the power sums with rational coefficients. Hence the local characteristic polynomials of \(H\) have real coefficients, that is, \(H\) is real.

The first summand of \(H\) is \(R^1p_{2!}(p_1^*j_!M\otimes\mathcal K_\psi)\otimes\mathcal B_b\) by the projection formula for the constant twist, that is, \(G\otimes\mathcal B_b\) with \(G=\operatorname{FT}_\psi(j_!M)\) as in Lemma \hyperref[bp227]{17.1}. Tensoring the surjection \(G\to j'_*M'\) of that lemma with \(\mathcal B_b\), we see that \(M'\otimes\mathcal B_b\) is a quotient of a direct summand of \(H\). To apply the purity theorem we restrict to a dense open on which both \(H\) and \(M'\) are lisse. There \(M'\otimes\mathcal B_b\) is a lisse quotient of the real lisse sheaf \(H\), and it is irreducible by Lemma \hyperref[bp209]{16.14}, so by Theorem \hyperref[bp120]{14.26} it is pure of some real weight. Removing the twist by \(\mathcal B_b\), we find that \(M'\) is pure of some weight \(v\) on that open. Since the open where \(M'\) is lisse is a smooth geometrically connected curve, the last assertion of Corollary \hyperref[bp137]{15.6} extends the purity to all of it.
\end{proof}


% object: sec-The-weight-of-the-Fourier-transform (blueprint/chap-The-Deligne-Laumon-theorem-for-curves/sec-The-weight-of-the-Fourier-transform/_section.md), content hidden
\subsection*{\texorpdfstring{17.2\quad The weight of the Fourier transform}{17.2 The weight of the Fourier transform}}
\phantomsection\label{bp236}%
\addcontentsline{toc}{subsection}{\texorpdfstring{17.2\quad The weight of the Fourier transform}{17.2 The weight of the Fourier transform}}
% text: prose-17-05
\phantomsection\label{bp237}%
By Proposition \hyperref[bp234]{17.4}, \(M'\) is pure of some weight \(v\), and Plancherel now finds \(v\). Purity had to come first, since the growth exponent \(\nu\) of Definition \hyperref[bp151]{15.10} is the weight only for pure sheaves.

% object: theorem-weight-fourier-transform (blueprint/chap-The-Deligne-Laumon-theorem-for-curves/sec-The-weight-of-the-Fourier-transform/theorem-weight-fourier-transform/_section.md)
\phantomsection\label{bp238}%
\begin{bpstatementit}{Theorem~17.5}{ (The weight of the Fourier transform)}
The sheaf \(M'\) is pure of weight exactly \(w+1\).
\end{bpstatementit}

% proof: proof/theorem-weight-fourier-transform
\begin{proof}[{Proof}]\phantomsection\label{bp239}%
Let \(v\) be the weight of \(M'\) given by Proposition \hyperref[bp234]{17.4}. Shifts do not change the squared absolute values of traces, and, if a nonzero lisse sheaf on a dense open of the line is pure of weight \(u\), then its direct image \(j_*\) has \(\nu=u\) (Corollary \hyperref[bp165]{15.15}). Applying this to \(M'\) and to \(M\), and using the Plancherel formula (Corollary \hyperref[bp192]{16.8}), we get

\[
v=\nu(j'_*M')
=\nu(\operatorname{FT}_\psi((j_*M)[1]))
=\nu((j_*M)[1])+1=w+1.
\]

The first equality holds because \(M'\) is pure of weight \(v\), the second is the defining isomorphism of \(M'\), the third is Plancherel, and the last holds because \(M\) is pure of weight \(w\).
\end{proof}

% text: prose-17-06
\phantomsection\label{bp240}%
In the example of the chapter opening, \(M\) is the Kummer sheaf of the quadratic character over \(\mathbf F_3\), of weight \(0\), and the trace of \(M'\) at \(y\in\mathbf F_3^\times\) is \(-\chi(y)g\) with \(g^2=-3\). Its absolute value is \(\sqrt3=3^{1/2}\), which is weight \(1=w+1\), as the theorem predicts.

The weight of \(M'\) bounds \(H^1_c\), because the stalk of the transform at the origin of the dual line is the compactly supported cohomology of \(j_*M\).

% object: corollary-h-1-c-stalk-fourier (blueprint/chap-The-Deligne-Laumon-theorem-for-curves/sec-The-weight-of-the-Fourier-transform/corollary-h-1-c-stalk-fourier/_section.md)
\phantomsection\label{bp241}%
\begin{bpstatementit}{Corollary~17.6}{ ($H^1_c$ as a stalk of the Fourier transform)}
We have \(H_c^2(\mathbf A^1,j_*M)=0\), and there is a Frobenius-equivariant isomorphism \(H_c^1(\mathbf A^1,j_*M)\simeq(j'_*M')_{0'}\) with the stalk at the origin \(0'\) of the dual line. In particular \(H_c^1(\mathbf A^1,j_*M)\) has weights at most \(w+1\).
\end{bpstatementit}

% proof: proof/corollary-h-1-c-stalk-fourier
\begin{proof}[{Proof}]\phantomsection\label{bp242}%
We compute the stalk at \(0'\) of \(\operatorname{FT}_\psi((j_*M)[1])\) in two ways: by the stalk formula, which gives the cohomology of \(j_*M\), and through the isomorphism with \((j'_*M')[1]\), which gives a stalk of a middle extension.

At the rational origin \(0'\) of the dual line, the kernel \(\mathcal K_\psi\) restricts to the pullback of \(\mathcal L_\psi\) along the constant map to \(0\), which is canonically trivial (on traces this is \(\psi(0)=1\)). The stalk formula of Proposition \hyperref[bp183]{16.5}, which is compatible with Frobenius, therefore gives

\[
\bigl(\operatorname{FT}_\psi((j_*M)[1])\bigr)_{0'}
=R\Gamma_c(\mathbf A^1,j_*M)[2].
\]

We compare this with the stalk of \((j'_*M')[1]\), which is concentrated in degree \(-1\). In degree \(0\) we therefore see that \(H_c^2(\mathbf A^1,j_*M)=0\), and in degree \(-1\) we obtain

\[
H_c^1(\mathbf A^1,j_*M)
\simeq(j'_*M')_{0'}
=(M'_{\bar\eta})^{I_{0'}},
\]

where \(I_{0'}\) is the inertia group at the origin \(0'\) of the dual line acting on the generic fibre of \(M'\), and the last equality is the formula of Lemma \hyperref[bp127]{15.2} for the stalk of a middle extension. By Proposition \hyperref[bp134]{15.5} the weights of this stalk are at most the weight of \(M'\), which is \(w+1\) by Theorem \hyperref[bp238]{17.5}.
\end{proof}


% object: sec-The-affine-line (blueprint/chap-The-Deligne-Laumon-theorem-for-curves/sec-The-affine-line/_section.md), content hidden
\subsection*{\texorpdfstring{17.3\quad The affine line}{17.3 The affine line}}
\phantomsection\label{bp243}%
\addcontentsline{toc}{subsection}{\texorpdfstring{17.3\quad The affine line}{17.3 The affine line}}
% text: prose-17-07
\phantomsection\label{bp244}%
To pass from irreducible sheaves to arbitrary ones we show that \(j_*\) is exact on sequences of pure sheaves of one weight. The reason is a gap in weights. Let \(0\to F_1\to F\to F_2\to0\) be an exact sequence of lisse sheaves on a dense open \(U\) of a curve \(C\). Applying \(j_*\) gives an exact sequence \(0\to j_*F_1\to j_*F\to j_*F_2\to R^1j_*F_1\), and \(R^1j_*F_1\) is supported on \(C\setminus U\). Now suppose that \(F_1\) and \(F_2\) are pure of weight \(w\). Then, by Proposition \hyperref[bp134]{15.5}, the source of the last map has weights at most \(w\) at a point of \(C\setminus U\). On the other hand, by local duality, which we prove first, the stalk of \(R^1j_*F_1\) at such a point is the dual of the invariants of \(F_1^\vee\), twisted by \((-1)\), and this space has weights at least \(w+2\). Since a Frobenius-equivariant map between spaces with disjoint weights is zero, \(j_*\) is exact on such sequences (Proposition \hyperref[bp249]{17.9}).

The hypothesis of equal weights is needed. In the Kummer extension \(0\to\overline{\mathbf Q}_\ell(1)\to F\to\overline{\mathbf Q}_\ell\to0\) of the chapter opening, \(F_1=\overline{\mathbf Q}_\ell(1)\) has weight \(-2\) and \(F_2=\overline{\mathbf Q}_\ell\) has weight \(0\), so at the origin the stalks of \(j_*\overline{\mathbf Q}_\ell\) and of \(R^1j_*\overline{\mathbf Q}_\ell(1)\) both have weight \(0\), and the connecting map is indeed nonzero.

Exactness lets us extend Corollary \hyperref[bp241]{17.6} from one irreducible sheaf with no geometric constituent of the form \(\mathcal L_{\psi(ax)}\) to a pure sheaf unramified at infinity (Theorem \hyperref[bp252]{17.10}) by induction along a composition series. Each factor of the series is either geometrically constant, and then \(H^1_c=0\), or it satisfies the hypotheses of the corollary.

% object: lemma-local-duality-point-c-u (blueprint/chap-The-Deligne-Laumon-theorem-for-curves/sec-The-affine-line/lemma-local-duality-point-c-u/_section.md)
\phantomsection\label{bp245}%
\begin{bpstatementit}{Lemma~17.7}{ (Local duality at a point of $C\setminus U$)}
Let \(j_0:U_0\hookrightarrow C_0\) be a dense open in a smooth curve over \(k_0\), and let \(F\) be a lisse Weil \(\overline{\mathbf Q}_\ell\)-sheaf on \(U_0\). For a geometric point \(\bar s\) over a closed point \(s\in C_0\setminus U_0\), there is a canonical perfect pairing, equivariant for the local geometric Frobenius at \(s\),

\[
(j_*F^\vee)_{\bar s}\otimes(R^1j_*F)_{\bar s}
\longrightarrow\overline{\mathbf Q}_\ell(-1).
\]
\end{bpstatementit}

% proof: proof/lemma-local-duality-point-c-u
\begin{proof}[{Proof}]\phantomsection\label{bp246}%
The two stalks are cohomology sheaves of the complexes \(i^*Rj_*F^\vee\) (in degree zero) and \(i^*Rj_*F\) (in degree one) on the boundary. We compute the Verdier dual of the second complex by duality on the curve, and we find the first one, shifted and twisted.

Let \(i:Z\hookrightarrow C\) be the inclusion of the whole finite complement, and put \(K_F=i^*Rj_*F\). For a lisse sheaf \(G\) on \(U\), the adjunction \(j^*\dashv Rj_*\) shows that \(i^!Rj_*G=0\), and applying \(i^!\) to the localization triangle for \(Rj_*G\) we obtain \(i^!j_!G\simeq i^*Rj_*G[-1]\). Since duality exchanges \(Rj_*\) with \(j_!\) and \(i^*\) with \(i^!\), this formula lets us write the result again in the form \(i^*Rj_*\). Indeed, by Verdier duality on the smooth curve, \(\mathbb D_UF=F^\vee(1)[2]\) for the lisse sheaf \(F\), and we find

\[
\begin{split}
\mathbb D_Z K_F
&\simeq i^!\mathbb D_C Rj_*F
\simeq i^!j_!\mathbb D_UF\\
&\simeq i^*Rj_*(F^\vee(1)[2])[-1]
=K_{F^\vee}(1)[1].
\end{split}
\]

On the finite geometric scheme \(Z\), Verdier duality is the exact duality of coefficient vector spaces, so taking degree \(-1\) at \(\bar s\) we get \((R^1j_*F)_{\bar s}^\vee=(j_*F^\vee)_{\bar s}(1)\), which is the pairing. It is Frobenius-equivariant because the adjunction, localization and duality maps are natural in \(F\) and respect its Weil structure.
\end{proof}

% text: prose-17-08
\phantomsection\label{bp247}%
For \(F\) pure lisse of weight \(w\), the pairing shows that \((R^1j_*F)_{\bar s}\) has weights at least \(w+2\). Indeed, it puts this stalk in duality with a stalk of \(j_*F^\vee\), which has weights at most \(-w\), and the pairing takes values in \(\overline{\mathbf Q}_\ell(-1)\), which has weight \(2\). The next proposition uses this gap between the source and the target of the connecting map.

% object: example-local-duality-origin (blueprint/chap-The-Deligne-Laumon-theorem-for-curves/sec-The-affine-line/example-local-duality-origin/_section.md)
\phantomsection\label{bp248}%
\begin{bpstatement}{Example~17.8}{ (Local duality at the origin of the line)}
Let \(j:\mathbf G_m\hookrightarrow\mathbf A^1\) over \(k_0\), let \(F=\overline{\mathbf Q}_\ell\), which is pure of weight \(0\), and let \(s=0\). Since inertia at the origin acts trivially on the constant sheaf, \((j_*\overline{\mathbf Q}_\ell)_{\bar 0}=\overline{\mathbf Q}_\ell\) by Lemma \hyperref[bp127]{15.2}, with Frobenius acting by \(1\), of weight \(0\). On the other side, the Kummer covers \(t\mapsto t^{\ell^n}\) of a punctured neighbourhood of the origin give \((R^1j_*\overline{\mathbf Q}_\ell)_{\bar 0}=\overline{\mathbf Q}_\ell(-1)\) (Chapter \bpout{6}), on which geometric Frobenius acts by \(q\), of weight \(2\). The pairing of Lemma \hyperref[bp245]{17.7} is the multiplication \(\overline{\mathbf Q}_\ell\otimes\overline{\mathbf Q}_\ell(-1)\to\overline{\mathbf Q}_\ell(-1)\). Thus the two stalks have weights \(0\) and \(2\), and the gap of two comes entirely from the Tate twist. Indeed, with coefficients \(\overline{\mathbf Q}_\ell(1)\), \(H^1\) of a punctured disc is spanned by the Kummer class of the coordinate, which Frobenius fixes, so with coefficients \(\overline{\mathbf Q}_\ell=\overline{\mathbf Q}_\ell(1)\otimes\overline{\mathbf Q}_\ell(-1)\) Frobenius acts by \(q\).
\end{bpstatement}

% object: proposition-exactness-fixed-pure-weight (blueprint/chap-The-Deligne-Laumon-theorem-for-curves/sec-The-affine-line/proposition-exactness-fixed-pure-weight/_section.md)
\phantomsection\label{bp249}%
\begin{bpstatementit}{Proposition~17.9}{ (Exactness at a fixed pure weight)}
Let \(j_0:U_0\hookrightarrow C_0\) be as in Lemma \hyperref[bp245]{17.7}. If \(0\to F_1\to F\to F_2\to0\) is a short exact sequence of lisse Weil sheaves on \(U_0\), with \(F_1\) and \(F_2\) pure of the same weight \(w\), then the sequence

\[
0\longrightarrow j_*F_1\longrightarrow j_*F
\longrightarrow j_*F_2\longrightarrow0
\]

is exact.
\end{bpstatementit}

% proof: proof/proposition-exactness-fixed-pure-weight
\begin{proof}[{Proof}]\phantomsection\label{bp250}%
Applying \(j_*\) to the short exact sequence gives the long exact sequence

\[
0\to j_*F_1\to j_*F\to j_*F_2\xrightarrow{\ \delta\ }R^1j_*F_1,
\]

so it suffices to show that the connecting map \(\delta\) vanishes. It is Frobenius-equivariant, and we show that its source and target have disjoint weights. Since \(R^1j_*F_1\) is zero on \(U\), we look at a point of \(C\setminus U\). There the stalk of the source has weights at most \(w\) by Proposition \hyperref[bp134]{15.5}. For the target, note that the dual \(F_1^\vee\) is pure of weight \(-w\), so the same result shows that the stalk of \(j_*F_1^\vee\) has weights at most \(-w\). By Lemma \hyperref[bp245]{17.7} the stalk of the target is the dual of this stalk twisted by \((-1)\), and hence the target has weights at least \(w+2\). Now the image of \(\delta\) is a quotient of the source and a submodule of the target, so its eigenvalues would have weights both at most \(w\) and at least \(w+2\). Hence \(\delta\) vanishes at each point of \(C\setminus U\). As \(R^1j_*F_1\) is zero on \(U\), it follows that \(\delta=0\).
\end{proof}

% text: prose-17-09
\phantomsection\label{bp251}%
All factors of a composition series of a pure lisse sheaf are pure of the same weight, so the proposition applies to every step of such a series. In this way we pass from the irreducible sheaves of Corollary \hyperref[bp241]{17.6} to arbitrary ones.

% object: theorem-weight-bound-affine-line (blueprint/chap-The-Deligne-Laumon-theorem-for-curves/sec-The-affine-line/theorem-weight-bound-affine-line/_section.md)
\phantomsection\label{bp252}%
\begin{bpstatementit}{Theorem~17.10}{ (The weight bound on the affine line)}
Let \(j:U_0\hookrightarrow\mathbf A^1_{k_0}\) be a dense open, and let \(F\) be a lisse sheaf on \(U_0\), pure of weight \(w\) and unramified at infinity. Then \(H_c^1(\mathbf A^1,j_*F)\) has weights at most \(w+1\).
\end{bpstatementit}

% proof: proof/theorem-weight-bound-affine-line
\begin{proof}[{Proof}]\phantomsection\label{bp253}%
We prove the bound for the irreducible factors of \(F\) and pass to \(F\) by induction. The irreducible factors are of two kinds: geometrically constant ones, for which \(H^1_c\) vanishes, and the others, to which Corollary \hyperref[bp241]{17.6} applies.

Take a finite composition series of the finite-dimensional Weil representation of \(F\). Each irreducible lisse factor is pure of weight \(w\), since its local eigenvalues are eigenvalues of a subquotient. It is again unramified at infinity, since inertia at infinity acts trivially on every subquotient of a representation on which it acts trivially. An irreducible factor that is geometrically constant is a sheaf \(\mathcal B_b\) of Proposition \hyperref[bp24]{13.7}, because the geometric action is trivial and an invertible linear operator over \(\overline{\mathbf Q}_\ell\) has an eigenline. Its middle extension is then constant on the affine line, and so its \(H_c^1\) vanishes by the trivial case of Lemma \hyperref[bp179]{16.4}. Every other factor \(F_i\) is irreducible, unramified at infinity and not geometrically constant, so by Lemma \hyperref[bp221]{16.18} it has no geometric constituent of the form \(\mathcal L_{\psi(ax)}\), and we may apply Corollary \hyperref[bp241]{17.6} with \(M=F_i\).

The induction needs \(j_*\) to turn the composition series into a filtration of \(j_*F\) whose quotients are the \(j_*F_i\). Since all factors are pure of the same weight \(w\), Proposition \hyperref[bp249]{17.9} shows that the short exact sequences of the series remain exact after applying \(j_*\). So consider a step \(0\to F'\to F''\to F_i\to0\) of the series, where we may assume that the bound holds for \(F'\). In the resulting long exact sequence of compactly supported cohomology, the image of \(H^1_c(\mathbf A^1,j_*F')\) in \(H^1_c(\mathbf A^1,j_*F'')\) is a quotient of the former, and the cokernel embeds in \(H^1_c(\mathbf A^1,j_*F_i)\). Since the bound \(w+1\) is preserved under subquotients and extensions, we conclude by induction that it holds for \(F\).
\end{proof}

% text: prose-17-10
\phantomsection\label{bp254}%
The condition at infinity enters only through Lemma \hyperref[bp221]{16.18}, to verify the hypothesis of Proposition \hyperref[bp218]{16.17}: it excludes the constituents \(\mathcal L_{\psi(ax)}\) with \(a\ne0\), whose transforms are skyscrapers rather than middle extensions. In the last section we remove this hypothesis, together with the restriction to opens of the line.


% object: sec-Arbitrary-curves (blueprint/chap-The-Deligne-Laumon-theorem-for-curves/sec-Arbitrary-curves/_section.md), content hidden
\subsection*{\texorpdfstring{17.4\quad Arbitrary curves}{17.4 Arbitrary curves}}
\phantomsection\label{bp255}%
\addcontentsline{toc}{subsection}{\texorpdfstring{17.4\quad Arbitrary curves}{17.4 Arbitrary curves}}
% text: prose-17-11
\phantomsection\label{bp256}%
Theorem \hyperref[bp252]{17.10} assumes that the curve is an open of the affine line and that the sheaf is unramified at infinity. We remove the second hypothesis by moving a point where the sheaf is lisse to infinity, and the first by taking the direct image along a finite map to the line.

For an upper bound it is enough to replace \(H^1_c\) by a group that surjects onto it, and each move does this. For the first move, let \(s\) be a point where the sheaf is lisse. Since a point has no \(H^1_c\), deleting \(s\) gives a surjection \(H^1_c(W\setminus\{s\},F)\to H^1_c(W,F)\), and after a finite extension of the constant field and an automorphism of \(\mathbf P^1\) sending \(s\) to infinity, the sheaf is unramified there. For the second, a finite map \(\pi:C\to\mathbf A^1\) does not change compactly supported cohomology at all, since \(R\Gamma_c(\mathbf A^1,\pi_*j_!F)=R\Gamma_c(C,F)\). However, \(\pi_*j_!F\) is lisse and pure only off the branch locus, so we remove the bad points by one more localization. Such a map \(\pi\) exists by Riemann–Roch, because a function with a single pole of order prime to \(p\) defines a finite generically étale map (Lemma \hyperref[bp263]{17.13}). Curves that are not affine, and the degrees \(r\ne1\), are treated afterwards in Theorem \hyperref[bp269]{17.15}.

% object: lemma-weight-bound-constant-extension (blueprint/chap-The-Deligne-Laumon-theorem-for-curves/sec-Arbitrary-curves/lemma-weight-bound-constant-extension/_section.md)
\phantomsection\label{bp257}%
\begin{bpstatementit}{Lemma~17.11}{ (Weight bounds and finite constant extensions)}
Let \(X_0\) be a separated scheme of finite type over \(k_0\), let \(F\) be a constructible Weil sheaf on \(X_0\), let \(k_m\supset k_0\) be the extension of degree \(m\), and let \(u\in\mathbf R\). The group \(H^r_c(X,F)\) is the same for \(X_0\) and for \(X_0\otimes_{k_0}k_m\), and the geometric Frobenius of \(k_m\) acts on it as the \(m\)th power of that of \(k_0\). Every eigenvalue \(\alpha\) of the Frobenius of \(k_0\) on \(H^r_c(X,F)\) satisfies \(|\iota\alpha|\le q^{u/2}\) if and only if every eigenvalue \(\beta\) of the Frobenius of \(k_m\) satisfies \(|\iota\beta|\le q^{mu/2}\).
\end{bpstatementit}

% proof: proof/lemma-weight-bound-constant-extension
\begin{proof}[{Proof}]\phantomsection\label{bp258}%
An extension of degree \(m\) of the constant field does not change the geometric scheme \(X\), the sheaf on it, or its cohomology, and it replaces the Frobenius of \(k_0\) by its \(m\)th power. So the eigenvalues \(\beta\) are the \(\alpha^m\), and \(|\iota\alpha^m|\le q^{mu/2}\) holds if and only if \(|\iota\alpha|\le q^{u/2}\).
\end{proof}

% text: prose-17-12
\phantomsection\label{bp259}%
To move a point to infinity we need a rational point, and hence a finite extension of the constant field. By the lemma this does not affect a weight bound.

% object: lemma-finite-extension-base-field-moving (blueprint/chap-The-Deligne-Laumon-theorem-for-curves/sec-Arbitrary-curves/lemma-finite-extension-base-field-moving/_section.md)
\phantomsection\label{bp260}%
\begin{bpstatementit}{Lemma~17.12}{ (Moving a lisse point to infinity)}
Let \(W_0\) be a dense open of \(\mathbf P^1_{k_0}\), and let \(F\) be a lisse Weil sheaf on \(W_0\), pure of weight \(w\). There are a finite extension \(k_m\) of \(k_0\), a dense open \(V_0\) of \(\mathbf A^1_{k_m}\), a lisse Weil sheaf \(G\) on \(V_0\), pure of weight \(w\) and unramified at infinity, and a surjection

\[
H_c^1(V,G)\twoheadrightarrow H_c^1(W,F)
\]

that commutes with the geometric Frobenius of \(k_m\).
\end{bpstatementit}

% proof: proof/lemma-finite-extension-base-field-moving
\begin{proof}[{Proof}]\phantomsection\label{bp261}%
We delete a rational point where \(F\) is lisse, which replaces \(H^1_c\) by a group surjecting onto it, and move that point to infinity by an automorphism of \(\mathbf P^1\). First choose a closed point of \(W_0\) and a finite extension \(k_m\) of the constant field over which it has a rational point \(s\). The restriction of \(F\) to \(W_0\otimes_{k_0}k_m\) is again pure of weight \(w\), because a residue extension raises an eigenvalue and the residue cardinality to the same power. Next, some element of \(\mathrm{PGL}_2(k_m)\) sends \(s\) to infinity, and we transport \(F|_{W\setminus\{s\}}\) by this automorphism. As \(F\) is lisse at \(s\), its local inertia there is trivial, so the transported sheaf \(G\) extends to a lisse sheaf across the new infinity. It is defined on the image \(V_0\) of \(W_0\setminus\{s\}\), a dense open of the affine line, and \(H^1_c(V,G)=H^1_c(W\setminus\{s\},F)\) because the automorphism identifies the two pairs. The surjection follows from the localization sequence

\[
H_c^1(W\setminus\{s\},F)\to H_c^1(W,F)\to H_c^1(\{s\},F_{\bar s})=0,
\]

in which the last group vanishes because a point has no cohomology in positive degrees. Since every construction is defined over \(k_m\), the map commutes with the Frobenius of \(k_m\).
\end{proof}

% text: prose-17-13
\phantomsection\label{bp262}%
By the lemma, for any open of the projective line and any pure lisse sheaf on it, we may bound \(H^1_c\) with Theorem \hyperref[bp252]{17.10}. To reach an arbitrary affine curve we map it to the line, and we have to choose the map so that the direct image is again pure and lisse on a dense open.

% object: lemma-finite-generically-etale-map-line (blueprint/chap-The-Deligne-Laumon-theorem-for-curves/sec-Arbitrary-curves/lemma-finite-generically-etale-map-line/_section.md)
\phantomsection\label{bp263}%
\begin{bpstatementit}{Lemma~17.13}{ (A finite generically étale map to the line)}
Let \(C_0\) be a smooth geometrically connected affine curve. There are a smooth curve \(C'_0\) containing \(C_0\) with finite complement and a finite generically étale map \(\pi:C'_0\to\mathbf A^1_{k_0}\). If \(F\) is a pure lisse sheaf of weight \(w\) on \(C_0\), then the sheaf \(K=\pi_*j_!F\), where \(j:C_0\hookrightarrow C'_0\), is lisse and pure of weight \(w\) outside the branch locus and the image of \(C'_0\setminus C_0\). Moreover, \(R\Gamma_c(\mathbf A^1,K)=R\Gamma_c(C,F)\) canonically.
\end{bpstatementit}

% proof: proof/lemma-finite-generically-etale-map-line
\begin{proof}[{Proof}]\phantomsection\label{bp264}%
The map is given by a function with a single pole, at a point at infinity, and we choose the order of this pole prime to \(p\) so that the map is generically étale. We find the function by Riemann–Roch, we show by a local computation at the pole that its differential is nonzero, and we check the assertions about \(K\) stalk by stalk.

Let \(\overline C_0\) be the smooth projective completion, choose a closed point \(s\) of \(\overline C_0\) outside \(C_0\), and choose \(m\) prime to \(p\) with \((m-1)\deg(s)>2g(\overline C_0)-2\). Both divisors \(ms\) and \((m-1)s\) then have degree \(>2g-2\), so Riemann–Roch (Theorem \bpout{C.3}) gives \(h^0(\mathcal O(ns))=n\deg(s)+1-g\) for \(n=m,m-1\), and hence

\[
h^0(\mathcal O(ms))-h^0(\mathcal O((m-1)s))=\deg(s)>0.
\]

Therefore there is a function \(f\) whose only pole is at \(s\), of exact order \(m\). The sections \(1,f\) define a nonconstant map \(\overline C_0\to\mathbf P^1\) whose inverse image of infinity is \(\{s\}\), and this map is proper with finite fibres, hence finite. To see that it is generically étale, choose a uniformizer \(t\) at \(s\) and write \(f=t^{-m}u\) with \(u\) a unit. Since \(\kappa(s)/k_0\) is separable and the curve is smooth, \(dt\) generates the module of differentials near \(s\). Writing \(du=u'\,dt\), we get

\[
df=t^{-m-1}(-mu+tu')\,dt\ne0,
\]

since \(-mu+tu'\) is congruent to the unit \(-mu\) modulo \(t\), and \(mu\) is a unit because \(p\nmid m\). The differential of this finite map between smooth curves is therefore invertible at the generic point, and the map is generically étale. We set \(C'_0=\overline C_0\setminus\{s\}\) and let \(\pi\) be the restriction of the map to \(C'_0\).

Over the open in the statement, \(\pi\) is finite étale and its inverse image lies in \(C_0\), where \(j_!F=F\). Moreover, the direct image of a lisse sheaf under a finite étale map is lisse, because étale locally the map is a disjoint union of copies of the base. Hence \(K\) is lisse there. For purity, note that at a closed point \(y\) of this open, the stalk of \(K\) is a sum over the geometric points of the fibre, which Frobenius permutes. Consider such a cycle of length \(e\), corresponding to a point over \(y\) with residue extension of degree \(e\). On it the \(e\)th power of the operator is the local Frobenius upstairs, so an eigenvalue \(\beta\) satisfies \(|\iota\beta|^e=q_y^{ew/2}\), and the weight is still \(w\). Since \(R\pi_!=\pi_*\) for the finite map \(\pi\), the identification of cohomology follows by composing direct images with compact support. It commutes with Frobenius because the construction is defined over \(k_0\).
\end{proof}

% text: prose-17-14
\phantomsection\label{bp265}%
For example, take for \(C_0\) the affine curve \(y^2=x^3-x\) over \(\mathbf F_3\), the complement of the point at infinity of the projective curve. There \(g=1\), the point \(s=\infty\) is rational, and \(m=2\) satisfies \((m-1)\deg(s)>0\). Accordingly, the function \(x\) has a single pole, of order two, which is prime to \(3\), and the map is the double cover \((x,y)\mapsto x\) of the line, whose extension to the projective curve is branched over \(0\), \(\pm1\) and \(\infty\).

With both lemmas the bound on the line extends to every affine curve.

% object: theorem-weight-bound-affine-curve (blueprint/chap-The-Deligne-Laumon-theorem-for-curves/sec-Arbitrary-curves/theorem-weight-bound-affine-curve/_section.md)
\phantomsection\label{bp266}%
\begin{bpstatementit}{Theorem~17.14}{ (The weight bound on an affine curve)}
For a pure lisse Weil sheaf \(F\) of weight \(w\) on a smooth geometrically connected affine curve \(C_0\), the group \(H_c^1(C,F)\) has weights at most \(w+1\).
\end{bpstatementit}

% proof: proof/theorem-weight-bound-affine-curve
\begin{proof}[{Proof}]\phantomsection\label{bp267}%
We push the sheaf forward to the line and realize \(H^1_c\) as a quotient of a group bounded by Theorem \hyperref[bp252]{17.10}. This uses two surjections: one removes the points where the direct image fails to be lisse and pure, and the other moves a lisse point to infinity.

Choose \(\pi\) as in Lemma \hyperref[bp263]{17.13} and put \(K=\pi_*j_!F\). Let \(W_0\) be the dense open on which \(K\) is lisse and pure. Since its complement is finite, localization gives

\[
H_c^1(W,K|_W)\twoheadrightarrow
H_c^1(\mathbf A^1,K)=H_c^1(C,F).
\]

By Lemma \hyperref[bp260]{17.12}, after a finite extension of the constant field there is a further surjection from \(H_c^1(V,G)\) onto \(H_c^1(W,K|_W)\), with \(G\) pure of weight \(w\) and unramified at infinity. Let \(v:V\hookrightarrow\mathbf A^1\) be the inclusion and \(i:Z\hookrightarrow\mathbf A^1\) its finite complement. The localization sequence for \(v_*G\) reads

\[
H_c^0(Z,i^*v_*G)\longrightarrow H_c^1(V,G)
\longrightarrow H_c^1(\mathbf A^1,v_*G)\longrightarrow0.
\]

The first group has weights at most \(w\). Indeed, the stalks of \(i^*v_*G\) at the closed points of \(Z_0\) have weights at most \(w\) by Proposition \hyperref[bp134]{15.5}, and \(H_c^0(Z,i^*v_*G)\) is the direct sum, over the closed points \(z\) of \(Z_0\), of the stalks at the \(e=[\kappa(z):k_m]\) geometric points over \(z\), where \(k_m\) is the constant field of \(V_0\). The geometric Frobenius of \(k_m\) permutes these \(e\) stalks cyclically, and its \(e\)th power on each of them is conjugate to the local Frobenius \(F_z\). An eigenvalue \(\beta\) on such a cycle therefore satisfies \(\beta^e=\alpha\) for an eigenvalue \(\alpha\) of \(F_z\), so that \(|\iota\beta|^e=|\iota\alpha|\le q_z^{w/2}=(q^m)^{ew/2}\), and \(|\iota\beta|\le(q^m)^{w/2}\), which is weight at most \(w\) over \(k_m\). This is the case of the finite map \(Z_0\to\operatorname{Spec}k_m\) in Lemma \hyperref[bp275]{18.1}, whose proof is the same computation. The last group has weights at most \(w+1\) by Theorem \hyperref[bp252]{17.10}. As an extension of a subgroup of the last by a quotient of the first, the middle group therefore has weights at most \(w+1\). This bound passes to quotients under the two surjections, and by Lemma \hyperref[bp257]{17.11} it holds over the original constant field.
\end{proof}

% text: prose-17-15
\phantomsection\label{bp268}%
To reach a curve that is not affine, we pass to a dense affine open and view the sheaf on the whole curve as a middle extension.

% object: theorem-weight-bound-curve (blueprint/chap-The-Deligne-Laumon-theorem-for-curves/sec-Arbitrary-curves/theorem-weight-bound-curve/_section.md)
\phantomsection\label{bp269}%
\begin{bpstatementit}{Theorem~17.15}{ (The weight bound on a curve)}
Let \(j:U_0\hookrightarrow C_0\) be a dense open in a smooth curve over \(k_0\), and let \(F\) be a pure lisse Weil sheaf of weight \(w\) on \(U_0\). Then for every \(r\in\mathbf Z\) the group \(H_c^r(C,j_*F)\) has weights at most \(w+r\). In particular, for \(U_0=C_0\), the group \(H_c^r(C,F)\) has weights at most \(w+r\).
\end{bpstatementit}

% proof: proof/theorem-weight-bound-curve
\begin{proof}[{Proof}]\phantomsection\label{bp270}%
On a geometrically connected curve we compare \(j_*F\) with its restriction to a dense affine open, where the affine case applies, and we deduce the bounds from the localization sequence. The general case reduces to this one, because a curve that is not geometrically connected becomes a disjoint union of geometrically connected ones after a finite constant extension.

Suppose first that \(C_0\) is geometrically connected, and choose a dense affine open \(v:V_0\hookrightarrow U_0\). Since \(F\) is lisse on \(U_0\), the unit \(F\to v_*F|_V\) is an isomorphism by the last assertion of Lemma \hyperref[bp127]{15.2}. Hence \(j_*F\) is also the middle extension of \(F|_V\) from \(V_0\) to \(C_0\). On \(V\), by Theorem \hyperref[bp266]{17.14} and Lemma \hyperref[bp140]{15.7}, \(H_c^r(V,F)\) has weights at most \(w+r\) for every \(r\), and \(H_c^0(V,F)=0\). On the complement, write \(Z_0=C_0\setminus V_0\) and \(i\) for its inclusion; then the restriction of \(j_*F\) to the finite set \(Z_0\) has weights at most \(w\) by Proposition \hyperref[bp134]{15.5}, and its cohomology is concentrated in degree zero. These two pieces are joined by the localization sequence for \(V\subset C\), which reads

\[
0=H_c^0(V,F)\to H_c^0(C,j_*F)\to H^0(Z,i^*j_*F)\to H_c^1(V,F)
\to H_c^1(C,j_*F)\to0,
\]

followed by an isomorphism \(H_c^2(V,F)\simeq H_c^2(C,j_*F)\), and all higher groups vanish. Over \(\operatorname{Spec}k_0\) the bounds are conditions on eigenvalues (Lemmas \hyperref[bp11]{13.2} and \hyperref[bp14]{Lemma 13.3}), and we read them off degree by degree. In degree zero, the group \(H_c^0(C,j_*F)\) injects into \(H^0(Z,i^*j_*F)\), which has weights at most \(w\): Frobenius permutes the stalks at the geometric points over each closed point of \(Z_0\) cyclically, with the local Frobenius as the power that returns to a stalk, as in the proof of Theorem \hyperref[bp266]{17.14}. In degree one, \(H_c^1(C,j_*F)\) is a quotient of \(H_c^1(V,F)\), which has weights at most \(w+1\). In degree two the groups are isomorphic, and the bound \(w+2\) holds for \(V\).

In general, the connected components of \(C\) are defined over a finite extension \(k_m\) of \(k_0\), and over \(k_m\) each of them is a smooth geometrically connected curve. For example, \(\mathbf P^1_{\mathbf F_{q^2}}\), viewed as a curve over \(\mathbf F_q\), is connected, but \(C\) has two components, which Frobenius exchanges and which are defined over \(\mathbf F_{q^2}\). As Frobenius permutes the components, only its \(m\)th power acts on the cohomology of each of them. On such a component, \(j_*F\) is the middle extension of the restriction of \(F\) to the dense open where the component meets \(U\), and this restriction is pure of weight \(w\). Moreover, the group \(H^r_c(C,j_*F)\) is the direct sum of the groups of these components, compatibly with the \(m\)th power of Frobenius. Hence the bound holds over \(k_m\) by the first case, and by Lemma \hyperref[bp257]{17.11} it holds over \(k_0\).
\end{proof}



% object: chap-The-Deligne-Laumon-theorem-in-general (blueprint/chap-The-Deligne-Laumon-theorem-in-general/_section.md), content hidden
\section*{\texorpdfstring{18\quad The Deligne–Laumon theorem in general}{18 The Deligne–Laumon theorem in general}}
\phantomsection\label{bp271}%
\addcontentsline{toc}{section}{\texorpdfstring{18\quad The Deligne–Laumon theorem in general}{18 The Deligne–Laumon theorem in general}}
% text: prose-18-01
\phantomsection\label{bp272}%
At this point we know the Deligne–Laumon theorem on curves: a pure lisse sheaf of weight \(w\) on a curve has \(H^r_c\) of weights at most \(w+r\). For the Weil conjectures we need the same bound for the cohomology of a variety of any dimension. In this chapter we prove Theorem \hyperref[bp38]{13.11}, which says that for a separated morphism of finite type \(f:X_0\to Y_0\), the functor \(Rf_!\) maps \(D_{\leq w}(X_0)\) into \(D_{\leq w}(Y_0)\).

Take a surface fibred in curves over a curve, \(f:X_0\to Y_0\). Here the Leray spectral sequence \(H^a_c(Y,R^bf_!\overline{\mathbf Q}_\ell)\Rightarrow H^{a+b}_c(X,\overline{\mathbf Q}_\ell)\) suggests applying the curve case twice, once on the fibres and once on the base. On the fibres there is no difficulty: at a closed point \(y\) the stalk of \(R^bf_!\overline{\mathbf Q}_\ell\) is \(H^b_c(X_{\bar y},\overline{\mathbf Q}_\ell)\), and the curve case says that it has weights at most \(b\). On the base, however, the curve case applies only to pure sheaves. So we need a finite filtration of \(R^bf_!\overline{\mathbf Q}_\ell\) with pure quotients, that is, we need \(R^bf_!\overline{\mathbf Q}_\ell\) to be mixed in the sense of Definition \hyperref[bp9]{13.1}, and bounds at each point separately do not produce such a filtration. In his first proof [\hyperref[bp350]{Del74}] Deligne avoided the problem with a Lefschetz pencil, where the Picard–Lefschetz formula describes the direct image in the middle degree, but for an arbitrary fibration the sheaves on the base are not explicit. Once a mixed filtration exists, on the other hand, the pointwise bound also bounds the weights of its quotients (Lemma \hyperref[bp14]{13.3}). The question is therefore whether the direct images are mixed.

The answer comes once more from positivity (Theorem \hyperref[bp295]{18.8}):

\begin{quote}
every constructible real Weil sheaf on a separated scheme of finite type over \(k_0\) is mixed.
\end{quote}

To see why, let \(S\) be an irreducible constituent of a real lisse sheaf on a smooth affine \(X_0\). In order to compare the eigenvalues of \(S\) at a fixed point \(x_*\) with those at an arbitrary point \(x\), we restrict to a curve through both. On the curve each irreducible constituent is pure by Theorem \hyperref[bp120]{14.26}, so if the restriction of \(S\) stays irreducible, its eigenvalues at \(x\) have the weight read off at \(x_*\). In order to find curves on which \(S\) stays irreducible, we first show that irreducibility can be detected in one finite quotient of the fundamental group (Lemma \hyperref[bp278]{18.2}), and then Bertini's theorem gives a curve through the given points that stays connected in the corresponding finite étale cover (Lemma \hyperref[bp281]{18.3} and Corollary \hyperref[bp284]{18.4}). For constructible sheaves, finally, we stratify.

By Noether normalization over strata (Lemma \hyperref[bp312]{18.14}), an arbitrary \(f\) is built from quasi-finite maps and projections of relative affine lines. The essential case is a smooth family of affine curves \(h:U_0\to Y_0\) with a pure lisse sheaf \(F\) on \(U_0\). To apply the theorem on real sheaves we need direct images that are real, and here the curve case helps a second time. After a twist to weight zero and the addition of the dual, \(F\) becomes a real sheaf \(G\), and the curve case on the fibres makes the direct images of \(G\) real, hence mixed. A twist of \(F\) is a direct summand of \(G\), so the direct images of \(F\) are mixed as well (Theorem \hyperref[bp299]{18.9}).

% object: sec-Quasi-finite-morphisms-and-curves (blueprint/chap-The-Deligne-Laumon-theorem-in-general/sec-Quasi-finite-morphisms-and-curves/_section.md), content hidden
\subsection*{\texorpdfstring{18.1\quad Quasi-finite morphisms and curves}{18.1 Quasi-finite morphisms and curves}}
\phantomsection\label{bp273}%
\addcontentsline{toc}{subsection}{\texorpdfstring{18.1\quad Quasi-finite morphisms and curves}{18.1 Quasi-finite morphisms and curves}}
% text: prose-18-02
\phantomsection\label{bp274}%
To prove Theorem \hyperref[bp295]{18.8} we restrict a sheaf to curves, reassemble it over strata and descend it along constant extensions. For the reassembly and the descent we need that quasi-finite direct images preserve weights. For the restriction we need curves on which given lisse sheaves stay geometrically irreducible.

The first fact is elementary, because a finite direct image only permutes stalks. Let \(y\) be a point downstairs and \(x\) a point over \(y\) with residue extension of degree \(e\). Frobenius at \(y\) cycles the \(e\) copies of the stalk at \(x\), and its eigenvalues are \(e\)th roots of the eigenvalues upstairs. Measured against the residue field, absolute values are therefore unchanged (Lemma \hyperref[bp275]{18.1}).

The second fact is geometric, since restricting to a curve can lose monodromy. On \(\mathbf G_m\times\mathbf G_m\) with coordinates \(x,y\), in odd characteristic, the double cover \(u^2=x\) is connected, but over the curve \(x=1\) it splits into two copies. Accordingly, the Kummer sheaf of the quadratic character in \(x\) is geometrically nontrivial on the surface but constant on that curve. In the same way an irreducible sheaf of higher rank can become reducible on a badly chosen curve. On the other hand, a representation stays absolutely irreducible on a curve if the image of the fundamental group of the curve is large enough, and we show that this can be checked in one finite quotient (Lemma \hyperref[bp278]{18.2}). It then suffices that the curve stays connected in one finite étale cover, and we find such curves with Bertini's irreducibility theorem (Lemma \hyperref[bp281]{18.3}).

% object: lemma-finite-quasi-finite-direct-images (blueprint/chap-The-Deligne-Laumon-theorem-in-general/sec-Quasi-finite-morphisms-and-curves/lemma-finite-quasi-finite-direct-images/_section.md)
\phantomsection\label{bp275}%
\begin{bpstatementit}{Lemma~18.1}{ (Finite and quasi-finite direct images)}
Let \(g:Z_0\to Y_0\) be separated and quasi-finite. Then \(Rg_!\) is exact on sheaves and carries mixed sheaves of weights at most \(w\) to mixed sheaves of weights at most \(w\). It also carries \(D_{\leq w}(Z_0)\) into \(D_{\leq w}(Y_0)\). For a finite map, direct image carries a pure sheaf of weight \(w\) to a pure sheaf of weight \(w\).
\end{bpstatementit}

% proof: proof/lemma-finite-quasi-finite-direct-images
\begin{proof}[{Proof}]\phantomsection\label{bp276}%
We first treat a finite map and then reduce to it by Zariski's main theorem, which writes a quasi-finite map as an open immersion followed by a finite map. All the computation is in the finite case, since extension by zero does not change any nonzero stalk.

For a finite map, exactness follows from Lemma \bpout{4.27} and the comparison of coefficients. The same lemma shows that the stalk at a geometric point \(\bar y\) over a closed point \(y\) is the direct sum of the stalks at the geometric points of \(Z_0\) over \(\bar y\), and the local Frobenius at \(y\) permutes these summands. Let \(z\) be a point of residue degree \(e=[\kappa(z):\kappa(y)]\). Frobenius cycles through the \(e\) summands over \(z\), and its \(e\)th power on each of them is conjugate to the local Frobenius at \(z\). Each eigenvalue \(\beta\) on this cycle therefore satisfies \(\beta^e=\alpha\) for an eigenvalue \(\alpha\) at \(z\). If the input has weight \(w\), then

\[
|\iota(\beta)|^e=|\iota(\alpha)|=q_z^{w/2}=q_y^{ew/2},
\]

and we find \(|\iota(\beta)|=q_y^{w/2}\). This proves that finite direct image preserves purity. For example, for the constant sheaf on \(\operatorname{Spec}\mathbf F_{q^2}\) over \(\operatorname{Spec}\mathbf F_q\), the direct image has rank two and Frobenius swaps the two summands. Its eigenvalues \(\pm1\) have weight \(0\), as the constant sheaf does.

In general, \(Y_0\) is Noetherian, hence quasi-compact and quasi-separated, and Zariski's main theorem in the form [Stacks, \href{https://stacks.math.columbia.edu/tag/05K0}{05K0}] applies. It factors \(g\) as an open immersion \(j:Z_0\hookrightarrow\overline Z_0\) followed by a finite map \(\bar g:\overline Z_0\to Y_0\). By Proposition \bpout{5.28} we have \(Rg_!=\bar g_*j_!\), and both functors are exact. Extension by zero preserves purity, since its other stalks are zero. Now take a mixed sheaf of weights at most \(w\) on \(Z_0\), with a finite filtration by Weil subsheaves whose quotients are pure of weights at most \(w\). Since \(j\) and \(\bar g\) are defined over \(k_0\), the functors \(j_!\) and \(\bar g_*\) take Weil subsheaves to Weil subsheaves. Being exact, they take the filtration to a finite filtration of the direct image whose quotients are the direct images of the pure quotients. By the first part these are pure of the same weights, which proves the assertion on sheaves. Since \(Rg_!\) is exact, it commutes with taking cohomology sheaves, and the assertion on \(D_{\leq w}\) follows.
\end{proof}

% text: prose-18-03
\phantomsection\label{bp277}%
A lisse sheaf stays geometrically irreducible on a curve \(C\) in \(X\) as soon as the image of \(\pi_1(C)\) is large enough. The next lemma shows that it suffices for this image to surject onto one finite quotient of \(\pi_1(X)\), and in Lemma \hyperref[bp281]{18.3} we construct curves with this property.

% object: lemma-detecting-irreducibility-finite-quotient (blueprint/chap-The-Deligne-Laumon-theorem-in-general/sec-Quasi-finite-morphisms-and-curves/lemma-detecting-irreducibility-finite-quotient/_section.md)
\phantomsection\label{bp278}%
\begin{bpstatementit}{Lemma~18.2}{ (Detecting irreducibility in a finite quotient)}
Let \(\Pi\) be profinite and let \(\rho:\Pi\to\GL_E(V)\) be continuous and absolutely irreducible, with \(\dim_E V=n\). There are a \(\Pi\)-stable lattice \(L\subset V\) and an integer \(N\) such that every subgroup \(\Pi'\subseteq\Pi\) with the same image as \(\Pi\) in \(\GL(L/\pi_E^NL)\) acts absolutely irreducibly on \(V\).
\end{bpstatementit}

% proof: proof/lemma-detecting-irreducibility-finite-quotient
\begin{proof}[{Proof}]\phantomsection\label{bp279}%
Recall Burnside's theorem: an algebra of endomorphisms of a finite-dimensional vector space over an algebraically closed field that acts irreducibly is the full endomorphism algebra. Absolute irreducibility thus means that the group spans the full matrix algebra. We detect spanning by a nonzero determinant of \(n^2\) group elements. Such a determinant stays nonzero when we change the elements by a congruence modulo a high power of \(\pi_E\).

The compact group \(\rho(\Pi)\) stabilizes a lattice \(L\). Indeed, the stabilizer of any lattice \(L_0\) is open, hence of finite index in \(\rho(\Pi)\). So \(L_0\) has only finitely many translates, and their sum is a \(\rho(\Pi)\)-stable lattice. By absolute irreducibility and Burnside's theorem, there are elements \(g_1,\ldots,g_{n^2}\) whose matrices span \(\operatorname{End}_E(V)\). In a basis of \(L\), the determinant \(\Delta\) of their matrices, written as column vectors, is nonzero and integral. Choose \(N>v_{\pi_E}(\Delta)\). Suppose that \(g'_1,\ldots,g'_{n^2}\) have matrices congruent to those of the \(g_i\) modulo \(\pi_E^N\). Since the determinant is a polynomial with integer coefficients in the entries, their determinant \(\Delta'\) satisfies \(\Delta'\equiv\Delta\pmod{\pi_E^N}\). Hence \(v_{\pi_E}(\Delta')=v_{\pi_E}(\Delta)\), and in particular \(\Delta'\ne0\). A subgroup as in the statement contains such elements. Its matrices therefore span the whole matrix algebra, also after any field extension, and by Burnside's theorem it acts absolutely irreducibly. Finally, the image of \(\Pi\) in \(\GL(L/\pi_E^NL)\) is a finite quotient of \(\Pi\), since by continuity the kernel is open.
\end{proof}

% text: prose-18-04
\phantomsection\label{bp280}%
For a lisse sheaf on \(X\) the finite quotient corresponds to a connected finite Galois cover \(T\to X\). If \(C\subset X\) is a curve with \(T\times_XC\) connected, then the fundamental group of \(C\) surjects onto the Galois group of \(T\). What we need, then, are curves through given points whose inverse image in a given connected cover stays connected.

% object: lemma-curve-preserving-finite-cover (blueprint/chap-The-Deligne-Laumon-theorem-in-general/sec-Quasi-finite-morphisms-and-curves/lemma-curve-preserving-finite-cover/_section.md)
\phantomsection\label{bp281}%
\begin{bpstatementit}{Lemma~18.3}{ (A curve preserving a finite cover)}
Let \(X\) be a smooth connected affine scheme over \(\overline{\mathbf F}_q\), of dimension at least two. Let \(T\to X\) be a connected finite étale cover and let \(S\subset X\) be a nonempty finite set of closed geometric points. There is a smooth connected closed affine curve \(C\subset X\) containing \(S\), with \(T\times_XC\) connected. The curve and its inclusion are defined over a finite field.
\end{bpstatementit}

% proof: proof/lemma-curve-preserving-finite-cover
\begin{proof}[{Proof}]\phantomsection\label{bp282}%
We cut \(X\) by a general hypersurface through \(S\), show that it is smooth and that both it and its inverse image in \(T\) are connected, and then repeat.

\emph{The linear system.} Embed \(X\) in a projective closure \(\overline X\). For large \(n\), the sections of \(\mathcal O_{\overline X}(n)\) vanishing on \(S\) have base locus exactly \(S\) on \(X\) and define an immersion off \(S\) ([Stacks, \href{https://stacks.math.columbia.edu/tag/0FD5}{0FD5}]). They also surject onto the first-order jets \(\bigoplus_{s\in S}\mathfrak m_s/\mathfrak m_s^2\otimes
\mathcal O_{\overline X}(n)\) at \(S\). Indeed, if \(\mathcal I\) is the ideal sheaf of \(S\), then \(H^1(\overline X,\mathcal I^2(n))=0\) for large \(n\) by Serre vanishing ([Stacks, \href{https://stacks.math.columbia.edu/tag/0B5U}{0B5U}]), and the cohomology sequence of \(0\to\mathcal I^2(n)\to\mathcal I(n)\to(\mathcal I/\mathcal I^2)(n)\to0\) shows that \(H^0(\overline X,\mathcal I(n))\) surjects onto the jets.

\emph{Smoothness.} A general divisor of these sections is smooth on \(X\). The immersion off \(S\) says nothing about the points of \(S\) themselves, and there smoothness is controlled by the jets. A divisor through \(s\) is smooth at \(s\) exactly when its equation has nonzero image in \(\mathfrak m_s/\mathfrak m_s^2\). Since the sections surject onto the jets, the condition that a section have nonzero derivative at each \(s\) is nonempty and open. On \(X\setminus S\) the divisors are the pullbacks of hyperplanes under the immersion of \(X\setminus S\) into a projective space that these sections define. By Bertini's theorem ([Stacks, \href{https://stacks.math.columbia.edu/tag/0FD6}{0FD6}]) a general one is smooth there.

\emph{Connectedness.} A general divisor and its inverse image in \(T\) are irreducible. We use Bertini's irreducibility theorem [Stacks, \href{https://stacks.math.columbia.edu/tag/0G4F}{0G4F}]: if \(Z\) is an irreducible scheme of finite type over an algebraically closed field and \(Z\to\mathbf P^N\) is a morphism whose image has closure of dimension at least two, then the inverse image of a general hyperplane is irreducible. We apply it to the immersion of \(X\setminus S\) above and to its composition with the restriction of \(T\to X\) over \(X\setminus S\). In both cases the hyperplanes correspond to the sections vanishing on \(S\) and their pullbacks to \(T\). Both sources are smooth and connected, hence irreducible. Both images have closure of dimension \(\dim X\ge2\), because \(T\to X\) is finite and surjective. So for a general divisor \(D\), both \(D\cap(X\setminus S)\) and its inverse image in \(T\) are irreducible. The finitely many points over \(S\) cannot be components, since every component of a divisor has dimension \(\dim X-1\ge1\), and the claim follows.

\emph{Induction.} Intersecting these nonempty open conditions in the space of sections, we obtain a smooth connected affine hypersurface \(H\) containing \(S\) whose inverse image in \(T\) is connected. It is smooth, connected and affine of dimension \(\dim X-1\), it contains \(S\), and \(T\times_XH\to H\) is a connected finite étale cover. Thus \(H\) again satisfies the hypotheses of the lemma, and we repeat the construction until the dimension is one. The curve is defined over a finite field, because it has finitely many equations, whose coefficients lie in a finite extension of \(\mathbf F_q\).
\end{proof}

% text: prose-18-05
\phantomsection\label{bp283}%
Combining the two lemmas, we obtain the curves we need for real sheaves. We can treat several sheaves at once, because a single cover dominates the finite quotients attached to all of them.

% object: corollary-curve-preserving-geometric-irreducibility (blueprint/chap-The-Deligne-Laumon-theorem-in-general/sec-Quasi-finite-morphisms-and-curves/corollary-curve-preserving-geometric-irreducibility/_section.md)
\phantomsection\label{bp284}%
\begin{bpstatementit}{Corollary~18.4}{ (A curve preserving geometric irreducibility)}
Let \(X_0\) be a smooth geometrically connected affine scheme of positive dimension. Given finitely many geometrically irreducible lisse Weil sheaves on \(X_0\) and finitely many closed geometric points, there is, after a finite extension of the constant field, a smooth geometrically connected affine curve through those points on which all the sheaves remain geometrically irreducible.
\end{bpstatementit}

% proof: proof/corollary-curve-preserving-geometric-irreducibility
\begin{proof}[{Proof}]\phantomsection\label{bp285}%
We translate irreducibility into the surjectivity of \(\pi_1(C)\) onto one finite group, and then find \(C\) with Bertini's theorem applied to the corresponding cover.

Choose a finite extension \(E\) of \(\mathbf Q_\ell\) over which all the sheaves are defined, and apply Lemma \hyperref[bp278]{18.2} to their geometric representations. Let \(T\to X\) be the connected Galois cover corresponding to the intersection of the kernels of the resulting finite quotients of \(\pi_1(X)\), and let \(Q\) be its group. In dimension one, we take \(C=X\). Otherwise \(\dim X\ge2\), and we apply Lemma \hyperref[bp281]{18.3} to \(T\) and the given points. If no points are given we add one, since the lemma needs \(S\) nonempty. As \(T\times_XC\) is connected, \(\pi_1(C)\) acts transitively on the fibre of the \(Q\)-torsor \(T\times_XC\to C\), so its image in \(Q\) is all of \(Q\). Each of the finite quotients above is a quotient of \(Q\). Hence \(\pi_1(C)\) has the same image as \(\pi_1(X)\) in each of them, and Lemma \hyperref[bp278]{18.2}, applied with \(\Pi=\pi_1(X)\) and \(\Pi'\) the image of \(\pi_1(C)\), shows that the restrictions remain geometrically irreducible. By Lemma \hyperref[bp281]{18.3} the curve and the points are defined over a finite extension of the constant field. The curve \(C\) is geometrically connected because it is connected over \(\overline{\mathbf F}_q\).
\end{proof}


% object: sec-Mixedness-of-real-sheaves (blueprint/chap-The-Deligne-Laumon-theorem-in-general/sec-Mixedness-of-real-sheaves/_section.md), content hidden
\subsection*{\texorpdfstring{18.2\quad Mixedness of real sheaves}{18.2 Mixedness of real sheaves}}
\phantomsection\label{bp286}%
\addcontentsline{toc}{subsection}{\texorpdfstring{18.2\quad Mixedness of real sheaves}{18.2 Mixedness of real sheaves}}
% text: prose-18-06
\phantomsection\label{bp287}%
With these curves we pass from the purity theorem for real sheaves on curves (Theorem \hyperref[bp120]{14.26}) to the mixedness of real sheaves on any base. The curves preserve geometric irreducibility, so we first split a lisse sheaf into geometrically irreducible pieces. These pieces need not be Weil sheaves, because an irreducible Weil sheaf need not be geometrically irreducible. For example, let \(\chi\) be a character of \(\mathbf F_{q^2}^\times\) with \(\chi^q\ne\chi\), and let \(g\) be the map from \(\mathbf G_m\) over \(\mathbf F_{q^2}\) to \(\mathbf G_m\) over \(\mathbf F_q\). Geometrically, the direct image under \(g\) of the Kummer sheaf of \(\chi\) splits into the Kummer sheaves of \(\chi\) and \(\chi^q\), which are distinct. Frobenius exchanges them, and so the direct image is an irreducible Weil sheaf of rank two. After the quadratic constant extension the two pieces are Weil sheaves themselves. In general, the geometric constituents become Weil subquotients after a finite constant extension (Lemma \hyperref[bp288]{18.5}). In the proof of Theorem \hyperref[bp295]{18.8} we descend mixedness back along this extension by a trace retraction. We also reduce constructible sheaves to lisse sheaves on smooth affine strata, and glue the strata using Lemma \hyperref[bp275]{18.1}.

% object: lemma-geometric-constituents-after-constant-extension (blueprint/chap-The-Deligne-Laumon-theorem-in-general/sec-Mixedness-of-real-sheaves/lemma-geometric-constituents-after-constant-extension/_section.md)
\phantomsection\label{bp288}%
\begin{bpstatementit}{Lemma~18.5}{ (Geometric constituents after a constant extension)}
A lisse Weil sheaf on a smooth geometrically connected scheme admits, after one finite constant extension, a finite Weil filtration with geometrically irreducible quotients.
\end{bpstatementit}

% proof: proof/lemma-geometric-constituents-after-constant-extension
\begin{proof}[{Proof}]\phantomsection\label{bp289}%
An irreducible Weil representation is geometrically semisimple. A power of Frobenius fixes each isotypic component, and on it we can triangularize this power. The flags obtained in this way are stable under the Weil group of the constant extension, and by them we refine a composition series to the required filtration.

In detail, take a composition series for the Weil representation. By Lemma \hyperref[bp75]{14.11}, applied to the geometric fundamental group, which is normal in the Weil group, the geometric restriction of each irreducible factor is semisimple. Frobenius permutes the finitely many geometric isotypic components of the factor, and we choose a lift \(s\) of geometric Frobenius to the Weil group and a power \(s^m\) fixing all of them. On a component \(S\otimes M\), geometric monodromy acts irreducibly on \(S\) and trivially on \(M\). Since \(s^m\) fixes the isomorphism type of \(S\), there is an isomorphism \(u\) from \(S\) to its conjugate by \(s^m\). We claim that the action of \(s^m\) on the component has the form \(u\otimes v\) with \(v\) invertible on \(M\). Indeed, the composite of the action of \(s^m\) with \(u^{-1}\otimes1\) commutes with geometric monodromy, and by Schur's lemma the endomorphisms of \(S\otimes M\) commuting with it are \(1\otimes\operatorname{End}(M)\). Triangularizing \(v\), we obtain a stable flag whose geometric factors are \(S\).

We take a common \(m\) for the finitely many factors. Then the subgroup generated by geometric monodromy and \(s^m\) is the Weil group after the constant extension of degree \(m\). Choose a finite extension \(E'\) of the coefficient field over which these flags are defined. Since a subspace defined over \(E'\) is stable under this group if and only if its extension to \(\overline{\mathbf Q}_\ell\) is, the flags define lisse Weil subsheaves, and refining the composition series by them gives the required filtration.
\end{proof}

% text: prose-18-07
\phantomsection\label{bp290}%
In the example above, after the quadratic extension the sheaf \(g_*\mathcal K_\chi\) becomes the direct sum of the Kummer sheaves of \(\chi\) and \(\chi^q\), and either summand gives a filtration as in the lemma. Beyond this example, we use the lemma, together with the curves of Corollary \hyperref[bp284]{18.4}, to prove that real sheaves are mixed (Theorem \hyperref[bp295]{18.8}). In that proof, however, the filtrations with pure quotients exist only after a constant extension and only on the pieces of a stratification, so the next two lemmas recover mixedness of the original sheaf.

% object: lemma-mixedness-descends-constant-extension (blueprint/chap-The-Deligne-Laumon-theorem-in-general/sec-Mixedness-of-real-sheaves/lemma-mixedness-descends-constant-extension/_section.md)
\phantomsection\label{bp291}%
\begin{bpstatementit}{Lemma~18.6}{ (Mixedness and finite constant extensions)}
Let \(X_0\) be a separated scheme of finite type over \(k_0\), let \(g:X'_0=X_0\otimes_{k_0}k_m\to X_0\) be a finite constant extension, and let \(G\) be a constructible Weil sheaf on \(X_0\). If \(g^*G\) is mixed, then \(G\) is mixed.
\end{bpstatementit}

% proof: proof/lemma-mixedness-descends-constant-extension
\begin{proof}[{Proof}]\phantomsection\label{bp292}%
Since \(g^*G\) is mixed, so is \(g_*g^*G\) by Lemma \hyperref[bp275]{18.1}. The unit \(G\to g_*g^*G\) has the Weil-equivariant retraction \((\deg g)^{-1}\operatorname{Tr}\). Indeed, étale locally the unit is the diagonal and the trace is the sum over the copies. These local maps glue because the sum is invariant under permutations of the copies. So \(G\) is a Weil subsheaf of \(g_*g^*G\). A Weil subobject of a mixed sheaf is mixed (Proposition \hyperref[bp24]{13.7}), and therefore \(G\) is mixed.
\end{proof}

% object: lemma-mixedness-open-closed (blueprint/chap-The-Deligne-Laumon-theorem-in-general/sec-Mixedness-of-real-sheaves/lemma-mixedness-open-closed/_section.md)
\phantomsection\label{bp293}%
\begin{bpstatementit}{Lemma~18.7}{ (Mixedness on an open and its complement)}
Let \(j:U_0\to X_0\) be an open immersion with closed complement \(i:Z_0\to X_0\), and let \(G\) be a constructible Weil sheaf on \(X_0\). If \(j^*G\) and \(i^*G\) are mixed, then \(G\) is mixed.
\end{bpstatementit}

% proof: proof/lemma-mixedness-open-closed
\begin{proof}[{Proof}]\phantomsection\label{bp294}%
The sheaf \(G\) is an extension

\[
0\longrightarrow j_!j^*G\longrightarrow G
\longrightarrow i_*i^*G\longrightarrow0,
\]

whose outer terms are mixed by Lemma \hyperref[bp275]{18.1}, since \(j\) and \(i\) are quasi-finite. An extension of mixed sheaves is mixed (Proposition \hyperref[bp24]{13.7}).
\end{proof}

% object: theorem-real-constructible-sheaves-are-mixed (blueprint/chap-The-Deligne-Laumon-theorem-in-general/sec-Mixedness-of-real-sheaves/theorem-real-constructible-sheaves-are-mixed/_section.md)
\phantomsection\label{bp295}%
\begin{bpstatementit}{Theorem~18.8}{ (Real constructible sheaves are mixed)}
Every constructible \(\iota\)-real Weil sheaf on a separated finite-type \(k_0\)-scheme is \(\iota\)-mixed.
\end{bpstatementit}

% proof: proof/theorem-real-constructible-sheaves-are-mixed
\begin{proof}[{Proof}]\phantomsection\label{bp296}%
We join a fixed point \(x_*\) to an arbitrary point \(x\) by a curve on which the geometrically irreducible constituents stay irreducible. On that curve each constituent is pure by Theorem \hyperref[bp120]{14.26}. Its weight can be read off at \(x_*\), so it does not depend on the curve. This handles a lisse sheaf on a smooth affine scheme, and stratifying handles an arbitrary constructible sheaf. Descent along constant extensions and gluing of strata preserve mixedness (Lemmas \hyperref[bp291]{18.6} and \hyperref[bp293]{Lemma 18.7}), so nothing about curves is used beyond Theorem \hyperref[bp120]{14.26}.

\emph{The lisse case.} Suppose that \(X_0\) is smooth, geometrically connected and affine, and that the sheaf \(G\) is lisse. Since the curves of Corollary \hyperref[bp284]{18.4} preserve geometric irreducibility, we first split \(G\) into geometrically irreducible pieces. To do so, we make the constant extension of Lemma \hyperref[bp288]{18.5} and write \(S_1,\ldots,S_r\) for the quotients of its Weil filtration. Note that the sheaf \(G\) remains real after this extension, because the powers of a conjugation-stable multiset of eigenvalues again form a conjugation-stable multiset.

Fix a closed point \(x_*\) of \(X_0\), and let \(x\) be any other closed point. By Corollary \hyperref[bp284]{18.4}, there is a curve through chosen geometric points above \(x_*\) and \(x\) on which every \(S_j\) remains irreducible. The restriction of \(G\) to this curve is real, and Theorem \hyperref[bp120]{14.26} says that the irreducible constituents of a real sheaf on a curve are pure. Since the restriction of \(S_j\) is irreducible, it is one of them, and so it is pure. Let \(\alpha\) be an eigenvalue of \(S_j\) at \(x_*\), and let \(e\) be the residue degree over \(\kappa(x_*)\) of the point of the curve above \(x_*\). At that point, then, the corresponding eigenvalue is \(\alpha^e\) and the residue cardinality is \(q_{x_*}^e\). If the restriction is pure of weight \(u\), then \(|\iota(\alpha)|^e=q_{x_*}^{eu/2}\), so the weight of the restriction of \(S_j\) must be

\[
w_j=2\log_{q_{x_*}}|\iota(\alpha)|,
\]

which does not depend on the curve. In particular, all eigenvalues of \(S_j\) at \(x_*\) have this modulus. At \(x\), the same calculation with powers shows that \(S_j\) is pure of weight \(w_j\). As \(x\) was arbitrary, each \(S_j\) is pure of weight \(w_j\) on \(X_0\). The chosen Weil filtration then shows that \(G\) is mixed after the constant extension, and by Lemma \hyperref[bp291]{18.6}, \(G\) is mixed.

\emph{The general case.} Now let \(G\) be an arbitrary constructible real sheaf. Since nilpotents do not change the étale site, we may assume that \(X_0\) is reduced. In each irreducible component, choose a dense smooth affine open on which \(G\) is lisse and which avoids the other components. Since the complement has smaller dimension, iterating we obtain a finite chain of closed subschemes whose differences are disjoint unions of smooth connected affine pieces on which \(G\) is lisse. After one finite constant extension these pieces are geometrically connected, and such an extension is allowed by Lemma \hyperref[bp291]{18.6}, and it keeps \(G\) real. The lisse case then applies to the positive-dimensional pieces, while at a zero-dimensional piece a Frobenius-stable full flag of each stalk has quotients of rank one, which are pure. The restriction of \(G\) to every piece is therefore mixed, and Lemma \hyperref[bp293]{18.7}, applied successively along the chain of closed subschemes, shows that \(G\) is mixed.
\end{proof}


% object: sec-Relative-affine-curves (blueprint/chap-The-Deligne-Laumon-theorem-in-general/sec-Relative-affine-curves/_section.md), content hidden
\subsection*{\texorpdfstring{18.3\quad Relative affine curves}{18.3 Relative affine curves}}
\phantomsection\label{bp297}%
\addcontentsline{toc}{subsection}{\texorpdfstring{18.3\quad Relative affine curves}{18.3 Relative affine curves}}
% text: prose-18-08
\phantomsection\label{bp298}%
Let \(h:U_0\to Y_0\) be a smooth family of affine curves and \(F\) a pure lisse sheaf on \(U_0\). On each stalk of \(R^ih_!F\), Theorem \hyperref[bp269]{17.15} bounds the eigenvalues by \(w+i\). Suppose that \(R^ih_!F\) is mixed. Each pure quotient of its filtration has a nonzero stalk somewhere, and its eigenvalues there are among those of \(R^ih_!F\). By Lemma \hyperref[bp14]{13.3} the quotients therefore have weights at most \(w+i\). Most of the proof of Theorem \hyperref[bp299]{18.9} is devoted to showing that \(R^ih_!F\) is mixed.

% object: theorem-pure-lisse-sheaves-over-relative (blueprint/chap-The-Deligne-Laumon-theorem-in-general/sec-Relative-affine-curves/theorem-pure-lisse-sheaves-over-relative/_section.md)
\phantomsection\label{bp299}%
\begin{bpstatementit}{Theorem~18.9}{ (Pure lisse sheaves over relative affine curves)}
Let \(h:U_0\to Y_0\) be smooth separated of relative dimension one, with affine geometric fibres, and let \(F\) be a lisse Weil sheaf on \(U_0\), pure of weight \(w\). Then \(R^ih_!F\) is mixed of weights at most \(w+i\).
\end{bpstatementit}

% proof: proof/theorem-pure-lisse-sheaves-over-relative
\begin{proof}[{Proof}]\phantomsection\label{bp300}%
We bound the stalks of \(R^ih_!F\) by the curve case, we show that these sheaves are mixed by twisting to weight zero and adding the dual, and we combine the two facts using Lemma \hyperref[bp14]{13.3}.

\emph{The stalks.} Fix a closed point \(y\) of \(Y_0\). Base change for direct images with compact support (Proposition \bpout{5.28}) gives an identification

\[
(R^ih_!F)_{\bar y}=H_c^i(U_{\bar y},F)
\]

compatible with the local Frobenius operators. The restriction of \(F\) to the fibre is again pure of weight \(w\), because a residue extension raises both an eigenvalue and the residue cardinality to the same power. The fibre is a smooth affine curve over \(\kappa(y)\), so Theorem \hyperref[bp269]{17.15} shows that \(H_c^i(U_{\bar y},F)\) has weights at most \(w+i\). After a finite extension of \(\kappa(y)\) the components of the fibre are geometrically connected. Lemma \hyperref[bp140]{15.7}, applied to each of them with \(j\) the identity, shows that these groups vanish outside degrees \(1\) and \(2\).

\emph{Mixedness.} We prove that \(R^ih_!F\) is mixed. This is trivial for \(F=0\), so we may assume \(F\ne0\). Let \(G=F^0\oplus(F^0)^\vee\), with \(F^0=F\otimes\mathcal B_b\), be the sheaf of Proposition \hyperref[bp148]{15.9}. It is lisse, real and pure of weight zero. At a closed point \(y\) write

\[
P_{i,y}(T)=\det(1-T\operatorname{Frob}_y\mid H_c^i(U_{\bar y},G)).
\]

Applying Lemma \hyperref[bp140]{15.7} to the components of the fibre after a finite extension of \(\kappa(y)\), as above, we see that \(P_{0,y}=1\) and that every reciprocal root of \(P_{2,y}\) has modulus \(q_y\). By the bound just proved, every reciprocal root of \(P_{1,y}\) has modulus at most \(q_y^{1/2}\).

We claim that \(\iota P_{1,y}\) and \(\iota P_{2,y}\) have real coefficients. By the Euler product and the determinant identity of Theorem \bpout{12.30}, after applying \(\iota\) we have

\[
L_\iota(U_y,G,T)=\frac{\iota P_{1,y}(T)}{\iota P_{2,y}(T)}.
\]

Since the Euler factors are real, this rational function has real Taylor coefficients. Both it and its complex conjugate are regular at zero, so they are equal as rational functions. Moreover, the numerator and the denominator are coprime: a common root would give an eigenvalue having both permitted moduli, whereas \(q_y^{1/2}<q_y\). Since a reduced fraction whose numerator and denominator have constant coefficient one is unique, comparing it with its conjugate we find \(\iota P_{1,y},\iota P_{2,y}\in\mathbf R[T]\), which proves the claim. Without the curve case this step would fail, because a common factor of numerator and denominator could cancel, and we would only know that the quotient is real, not each polynomial.

The sheaves \(R^ih_!G\) are constructible by the finiteness theorems, and we have shown that they are real at every closed point. By Theorem \hyperref[bp295]{18.8} each of them has a finite mixed filtration.

To return from \(G\) to \(F\), note that by additivity \(R^ih_!F^0\) is a direct summand of \(R^ih_!G\). A direct summand is a Weil subsheaf, and so it is mixed by Proposition \hyperref[bp24]{13.7}. By the projection formula for the constant twist, it is isomorphic to \(R^ih_!F\otimes\mathcal B_b\). Removing the twist preserves the finite Weil filtration and shifts the weight of every pure quotient by \(w\). Hence \(R^ih_!F\) is mixed.

\emph{Conclusion.} The sheaf \(R^ih_!F\) is mixed, and by the bound on the stalks all its local eigenvalues have modulus at most \(q_y^{(w+i)/2}\). Lemma \hyperref[bp14]{13.3}, with \(w+i\) in place of \(w\), shows that every nonzero quotient in the filtration has weight at most \(w+i\).
\end{proof}


% object: sec-Proof-of-the-Deligne-Laumon-theorem (blueprint/chap-The-Deligne-Laumon-theorem-in-general/sec-Proof-of-the-Deligne-Laumon-theorem/_section.md), content hidden
\subsection*{\texorpdfstring{18.4\quad Proof of the Deligne–Laumon theorem}{18.4 Proof of the Deligne–Laumon theorem}}
\phantomsection\label{bp301}%
\addcontentsline{toc}{subsection}{\texorpdfstring{18.4\quad Proof of the Deligne–Laumon theorem}{18.4 Proof of the Deligne–Laumon theorem}}
% text: prose-18-09
\phantomsection\label{bp302}%
An arbitrary \(f\) reduces to relative affine curves and quasi-finite maps. Over a field, Noether normalization makes an affine variety of dimension \(d\) finite over \(\mathbf A^d\), and \(\mathbf A^d\to\operatorname{Spec}k\) is a composite of \(d\) projections of affine lines. For a separated morphism \(f\) of finite type the same works over a dense open of the target, and we treat the complement by induction on its dimension (Lemma \hyperref[bp312]{18.14}). So we must check the weight bound for two kinds of maps and show that it survives the gluing.

Quasi-finite maps were treated in Lemma \hyperref[bp275]{18.1}. For the projection \(\mathbf A^1_Y\to Y\) (Proposition \hyperref[bp308]{18.12}), consider a sheaf on the relative line. Over a dense open of the base it is lisse off a closed subscheme finite over the base. It is therefore glued from a relative affine curve and a quasi-finite piece. The bound is stable under composition by the Leray spectral sequence (Theorem \hyperref[bp315]{18.15}). It is stable under partitions of the source and of the target by the localization sequences below. These gluing steps use only that mixed sheaves of weights at most \(u\) are closed under subquotients and extensions.

% object: lemma-weights-exact-sequence (blueprint/chap-The-Deligne-Laumon-theorem-in-general/sec-Proof-of-the-Deligne-Laumon-theorem/lemma-weights-exact-sequence/_section.md)
\phantomsection\label{bp303}%
\begin{bpstatementit}{Lemma~18.10}{ (Weights in an exact sequence)}
Let \(F'\to F\to F''\) be an exact sequence of Weil sheaves. If \(F'\) and \(F''\) are mixed of weights at most \(u\), then so is \(F\).
\end{bpstatementit}

% proof: proof/lemma-weights-exact-sequence
\begin{proof}[{Proof}]\phantomsection\label{bp304}%
Let \(a:F'\to F\) and \(b:F\to F''\) be the maps. Exactness at \(F\) gives a short exact sequence

\[
0\to\operatorname{im}(a)\to F\to\operatorname{im}(b)\to0,
\]

where \(\operatorname{im}(a)\) is a quotient of \(F'\) and \(\operatorname{im}(b)\) a subobject of \(F''\). Since mixed sheaves of weights at most \(u\) are closed under subquotients by Lemma \hyperref[bp22]{13.6} and under extensions by Proposition \hyperref[bp24]{13.7}, \(F\) is mixed of weights at most \(u\).
\end{proof}

% text: prose-18-10
\phantomsection\label{bp305}%
By the next lemma, to prove a weight bound for \(R^if_!\) of a morphism \(f\), it suffices to prove it over the pieces of finite partitions of the source and of the target.

% object: lemma-weight-bounds-partitions (blueprint/chap-The-Deligne-Laumon-theorem-in-general/sec-Proof-of-the-Deligne-Laumon-theorem/lemma-weight-bounds-partitions/_section.md)
\phantomsection\label{bp306}%
\begin{bpstatementit}{Lemma~18.11}{ (Weight bounds and partitions)}
Let \(f:X_0\to Y_0\) be separated of finite type, and let \(K\) be a constructible Weil sheaf on \(X_0\), mixed of weights at most \(w\).

\begin{enumerate}
\item Let \(j:U_0\to X_0\) be an open immersion with closed complement \(i:Z_0\to X_0\). If \(R^q(fj)_!j^*K\) and \(R^q(fi)_!i^*K\) are mixed of weights at most \(w+q\) for every \(q\), then so is \(R^qf_!K\).
\item Let \(V_0\subset Y_0\) be open with closed complement \(W_0\), and let \(f_V\) and \(f_W\) be the base changes of \(f\) to \(V_0\) and \(W_0\). If \(R^qf_{V!}(K|_{f^{-1}V_0})\) and \(R^qf_{W!}(K|_{f^{-1}W_0})\) are mixed of weights at most \(w+q\) for every \(q\), then so is \(R^qf_!K\).
\end{enumerate}
\end{bpstatementit}

% proof: proof/lemma-weight-bounds-partitions
\begin{proof}[{Proof}]\phantomsection\label{bp307}%
In both cases the restrictions of \(K\) are again mixed of weights at most \(w\) (Proposition \hyperref[bp24]{13.7}), and we place \(R^qf_!K\) in the middle of an exact sequence of Weil sheaves whose outer terms are bounded.

For (1), applying \(Rf_!\) and the composition isomorphism to the triangle \(j_!j^*K\to K\to i_*i^*K\to\) gives an exact sequence

\[
R^q(fj)_!j^*K\longrightarrow R^qf_!K\longrightarrow R^q(fi)_!i^*K,
\]

which is Frobenius-equivariant by naturality, and Lemma \hyperref[bp303]{18.10} gives the bound for \(R^qf_!K\).

For (2), let \(j':V_0\to Y_0\) and \(i':W_0\to Y_0\) be the inclusions. Base change for direct images with compact support (Proposition \bpout{5.28}) identifies \(j'^*R^qf_!K\) with \(R^qf_{V!}(K|_{f^{-1}V_0})\) and \(i'^*R^qf_!K\) with \(R^qf_{W!}(K|_{f^{-1}W_0})\). By Lemma \hyperref[bp275]{18.1}, the outer terms of

\[
0\longrightarrow j'_!j'^*R^qf_!K\longrightarrow R^qf_!K
\longrightarrow i'_*i'^*R^qf_!K\longrightarrow0
\]

are mixed of weights at most \(w+q\), and Lemma \hyperref[bp303]{18.10} gives the bound for the middle term.
\end{proof}

% object: proposition-relative-affine-line (blueprint/chap-The-Deligne-Laumon-theorem-in-general/sec-Proof-of-the-Deligne-Laumon-theorem/proposition-relative-affine-line/_section.md)
\phantomsection\label{bp308}%
\begin{bpstatementit}{Proposition~18.12}{ (A relative affine line)}
Let \(p:\mathbf A^1_{Y_0}\to Y_0\) be the projection, and let \(G\) be a constructible Weil sheaf on \(\mathbf A^1_{Y_0}\), mixed of weights at most \(w\). Then \(R^ip_!G\) is mixed of weights at most \(w+i\) for every \(i\).
\end{bpstatementit}

% proof: proof/proposition-relative-affine-line
\begin{proof}[{Proof}]\phantomsection\label{bp309}%
We stratify the target so that over each stratum, \(G\) is lisse off a subscheme that is finite over the base. On the lisse part the projection is a relative affine curve, and we apply Theorem \hyperref[bp299]{18.9}. On the finite part it is quasi-finite, and we apply Lemma \hyperref[bp275]{18.1}.

Take a finite mixed filtration of \(G\). In the long exact sequences of \(R^ip_!\) for its steps, the bound \(w+i\) passes from the pure quotients to \(G\) by Lemma \hyperref[bp303]{18.10}. We may therefore assume that \(G\) is pure and constructible of weight \(v\leq w\).

Suppose first that the target \(Y_0=\operatorname{Spec}R\) is integral and affine, so that \(\mathbf A^1_{Y_0}\) is integral. We construct the strata as follows. Choose a dense open \(U\) on which \(G\) is lisse. Since the locus where \(G\) is lisse is stable under Frobenius, we may replace \(U\) by the intersection of its finitely many Frobenius conjugates. Then \(U\) is defined over \(k_0\), and so are the ideal \(I\subset R[T]\) of its reduced complement and the choices below. If the complement is nonempty, choose a nonzero \(g\in I\) and invert its nonzero leading coefficient \(b\). Then over \(R_b\) the polynomial \(b^{-1}g\) is monic and vanishes in \(R_b[T]/I\). Hence the image of \(T\) is integral over \(R_b\), and \(R_b[T]/I\) is generated as an \(R_b\)-module by the powers \(T^k\) with \(k<\deg g\). In particular \(R_b[T]/I\) is finite over \(R_b\); if \(g\) has degree zero, the complement is empty over \(D(b)\).

Over \(D(b)\) the complement is therefore finite, hence quasi-finite. The open \(U\) is smooth of relative dimension one over \(D(b)\), and each of its fibres is an affine line minus finitely many points. So \(U\to D(b)\) satisfies the hypotheses of Theorem \hyperref[bp299]{18.9}. Applying that theorem on \(U\) and Lemma \hyperref[bp275]{18.1} on the complement, we obtain the bound over \(D(b)\) from Lemma lemma-weight-bounds-partitions.

For a general target we argue by induction on \(\dim Y_0\). Since nilpotents do not change the étale site, we may replace \(Y_0\) by its reduction. Take disjoint integral affine opens containing the generic points of its components, and perform this construction on each of them. Their union is open, and its complement has smaller dimension. The case of dimension zero is contained in the construction, because a reduced zero-dimensional scheme is a finite disjoint union of integral affine pieces. In general, the restriction of \(G\) over the closed complement is again mixed of weights at most \(w\), and the inductive hypothesis gives the bound there. Lemma lemma-weight-bounds-partitions, applied to the open union and its closed complement in the target, gives the bound over all of \(Y_0\).
\end{proof}

% object: example-leading-coefficient-stratification (blueprint/chap-The-Deligne-Laumon-theorem-in-general/sec-Proof-of-the-Deligne-Laumon-theorem/example-leading-coefficient-stratification/_section.md)
\phantomsection\label{bp310}%
\begin{bpstatement}{Example~18.13}{ (Why the leading coefficient is inverted)}
Let \(Y_0=\mathbf A^1_{k_0}=\operatorname{Spec}k_0[y]\), and suppose that \(G\) is lisse on the relative line \(\mathbf A^1_{Y_0}=\operatorname{Spec}k_0[y,T]\) off the curve \(Z=V(yT-1)\). This curve is not finite over the base. Indeed, over a point \(y\ne0\) the fibre of \(Z\) is the single point \(T=1/y\), and over \(y=0\) it is empty, since the equation becomes \(-1=0\). So \(Z\to Y_0\) is quasi-finite, but its image is \(D(y)\), which is not closed, whereas a finite map is closed; geometrically, as \(y\) tends to \(0\) the point \(1/y\) runs off to infinity. To repair this, the construction in the proof takes \(g=yT-1\) with leading coefficient \(b=y\). Over \(D(y)\) the polynomial \(T-y^{-1}\) is monic, and indeed \(k_0[y,y^{-1},T]/(yT-1)=k_0[y,y^{-1}]\) is finite over the base. The closed stratum \(y=0\), which has smaller dimension, is treated by induction, and on it the complement of the lisse locus happens to be empty.
\end{bpstatement}

% text: prose-18-11
\phantomsection\label{bp311}%
To apply the proposition to an arbitrary morphism, we factor it, over strata, as a quasi-finite map followed by projections of relative affine lines.

% object: lemma-relative-noether-normalization-over-stratification (blueprint/chap-The-Deligne-Laumon-theorem-in-general/sec-Proof-of-the-Deligne-Laumon-theorem/lemma-relative-noether-normalization-over-stratification/_section.md)
\phantomsection\label{bp312}%
\begin{bpstatementit}{Lemma~18.14}{ (Relative Noether normalization over a stratification)}
Let \(f:X_0\to Y_0\) be a morphism of finite type. There are finite partitions of \(X_0\) and of \(Y_0\) into locally closed pieces such that each piece of \(X_0\) lies over a piece \(V\) of \(Y_0\), and its map to \(V\) factors through a finite morphism to \(\mathbf A^d_V\) for some \(d\).
\end{bpstatementit}

% proof: proof/lemma-relative-noether-normalization-over-stratification
\begin{proof}[{Proof}]\phantomsection\label{bp313}%
We spread out Noether normalization from the generic points of the target, and then argue by induction on the dimension of the complement. Spreading out is possible because a normalization over the fraction field involves finitely many equations with finitely many denominators. Recall Noether normalization over a field \(k'\) [Stacks, \href{https://stacks.math.columbia.edu/tag/00OY}{00OY}]: a nonzero finitely generated \(k'\)-algebra is finite over a polynomial subalgebra.

Let \(R\) be an integral domain with fraction field \(k'\), let \(\operatorname{Spec}R\) be an affine open of the target, and let \(B\) be a finitely generated \(R\)-algebra. If \(B\otimes_Rk'=0\), then \(B_s=0\) for some nonzero \(s\in R\). Otherwise, Noether normalization over \(k'\) provides elements \(t_1,\ldots,t_d\) and monic equations for each of a fixed set of algebra generators of \(B\) over \(k'[t_1,\ldots,t_d]\). Choose a nonzero \(s\in R\) divisible by all denominators of the \(t_i\) and of the coefficients of these equations. Each equation, read in \(B_s\), has a left-hand side that vanishes in \(B\otimes_Rk'\) and is therefore killed by a nonzero element of \(R\). Multiplying \(s\) by these finitely many elements, we may assume that the equations hold in \(B_s\). Then bounded powers of the generators span \(B_s\) over \(R_s[t_1,\ldots,t_d]\), and the resulting map to \(\mathbf A^d_{R_s}\) is finite.

To obtain the partitions, choose disjoint integral affine opens containing the generic points of the reduced target. Partition their inverse images into finitely many reduced affine locally closed pieces and apply the preceding construction to each piece. On each target open, a single product of the finitely many denominators then works for all pieces. Since the complement of the resulting open has smaller dimension, by induction on this complement we obtain finite partitions of the whole source and target. Moreover, all choices are made over \(k_0\).
\end{proof}

% text: prose-18-12
\phantomsection\label{bp314}%
On each piece the map to \(V\) is thus a finite map followed by \(d\) projections of relative affine lines. The bound holds for the finite map by Lemma \hyperref[bp275]{18.1} and for each projection by Proposition \hyperref[bp308]{18.12}, and the next theorem puts the pieces together.

% object: theorem-direct-image-mixed-sheaf (blueprint/chap-The-Deligne-Laumon-theorem-in-general/sec-Proof-of-the-Deligne-Laumon-theorem/theorem-direct-image-mixed-sheaf/_section.md)
\phantomsection\label{bp315}%
\begin{bpstatementit}{Theorem~18.15}{ (Direct image of a mixed sheaf)}
Let \(f:X_0\to Y_0\) be separated of finite type and let \(F\) be mixed of weights at most \(w\). Then \(R^if_!F\) is mixed of weights at most \(w+i\) for every \(i\).
\end{bpstatementit}

% proof: proof/theorem-direct-image-mixed-sheaf
\begin{proof}[{Proof}]\phantomsection\label{bp316}%
Say that a morphism has the property if the conclusion holds for it, for every \(w\) and every mixed \(F\) of weights at most \(w\). Quasi-finite maps and projections of relative affine lines have the property, and by Lemma \hyperref[bp312]{18.14} every \(f\) is built from these two kinds of maps. So it suffices to show that the property is stable under composition and under finite partitions.

For composition, let \(X_0\xrightarrow{h}Z_0\xrightarrow{g}Y_0\) and consider the Leray spectral sequence for compact supports,

\[
E_2^{a,b}=R^ag_!(R^bh_!F)\ \Longrightarrow\ R^{a+b}(gh)_!F.
\]

It is Frobenius-equivariant by the naturality of the truncation tower and of the composition isomorphism for direct image with compact support. If both maps have the property, then \(E_2^{a,b}\) is mixed of weights at most \(w+a+b\). Since each later page is a subquotient of the previous one, the same holds for \(E_\infty^{a,b}\). The abutment filtration of \(R^n(gh)_!F\) is finite, because the operations preserve bounded constructible complexes. Its quotients are the \(E_\infty^{a,n-a}\), and by Lemma \hyperref[bp303]{18.10} \(R^n(gh)_!F\) is mixed of weights at most \(w+n\).

For partitions, we apply Lemma \hyperref[bp312]{18.14} and refine the partitions to filtrations by open and closed pieces. After this refinement we can order the pieces so that each is open in the union of itself and the later ones. This lets us apply the localization sequences one piece at a time. On the pieces, the restriction of \(F\) to each source piece is again mixed of weights at most \(w\), and on every refined source piece the finite map to \(\mathbf A^d_V\) restricts to a separated quasi-finite map, which has the property by Lemma \hyperref[bp275]{18.1}. Moreover, the projection \(\mathbf A^d_V\to V\) is a composition of \(d\) projections of relative affine lines. Each of them has the property by Proposition \hyperref[bp308]{18.12}, and by the composition argument the map from the piece to \(V\) has the property. Finally, we assemble the pieces one at a time in the chosen order. By Lemma lemma-weight-bounds-partitions on the source and Lemma lemma-weight-bounds-partitions on the target, the bound for the pieces implies the bound for \(R^if_!F\) on \(Y_0\).
\end{proof}

% text: prose-18-13
\phantomsection\label{bp317}%
For Theorem \hyperref[bp38]{13.11}, stated in Section \hyperref[bp36]{13.3}, the case of sheaves is Theorem \hyperref[bp315]{18.15}, and the spectral sequence of the truncation tower passes to complexes.

% proof: proof/theorem-weights-rf-shriek
\begin{proof}[{Proof of Theorem 13.11}]\phantomsection\label{bp318}%
Let \(K\in D_{\le w}(X_0)\), so that each \(\mathcal H^b(K)\) is mixed of weights at most \(w+b\). Applying \(Rf_!\) to the ordinary truncation tower of \(K\), which is bounded and functorial, we obtain the spectral sequence

\[
E_2^{a,b}=R^af_!\mathcal H^b(K)
\ \Longrightarrow\ \mathcal H^{a+b}(Rf_!K),
\]

By naturality and the Frobenius compatibility of \(Rf_!\), its differentials and its finite abutment filtration are compatible with the Weil structures. By Theorem \hyperref[bp315]{18.15}, \(E_2^{a,b}\) is mixed of weights at most \(w+a+b\), and so is its subquotient \(E_\infty^{a,b}\). The quotients of the filtration of \(\mathcal H^n(Rf_!K)\) are the \(E_\infty^{a,n-a}\). By Lemma \hyperref[bp303]{18.10} and induction on the length of the filtration, \(\mathcal H^n(Rf_!K)\) is mixed of weights at most \(w+n\). In other words, \(Rf_!K\in D_{\le w}(Y_0)\). The last assertion is the case of the structure map of \(X_0\), for which the geometric stalk of \(R^nf_!K\) is \(H_c^n(X,K)\) with its Frobenius operator.
\end{proof}



% object: chap-Proof-of-the-Weil-conjectures (blueprint/chap-Proof-of-the-Weil-conjectures/_section.md), content hidden
\section*{\texorpdfstring{19\quad Proof of the Weil conjectures}{19 Proof of the Weil conjectures}}
\phantomsection\label{bp319}%
\addcontentsline{toc}{section}{\texorpdfstring{19\quad Proof of the Weil conjectures}{19 Proof of the Weil conjectures}}
% text: prose-19-01
\phantomsection\label{bp320}%
The Deligne–Laumon theorem is an upper bound for weights, stated with \(\ell\)-adic coefficients and one embedding \(\iota:\overline{\mathbf Q}_\ell\to\mathbf C\). Weil's conjectures ask for more: an equality \(|\alpha|=q^{i/2}\), for polynomials \(P_i\) with integer coefficients, under every complex embedding. This chapter closes both gaps.

The lower bound comes from Poincaré duality, and a curve shows how. Let \(X_0\) be a smooth projective curve and \(\alpha\) an eigenvalue of Frobenius on \(H^1(X,\overline{\mathbf Q}_\ell)\). Then Theorem \hyperref[bp38]{13.11}, applied to the constant sheaf, gives \(|\iota(\alpha)|\le q^{1/2}\). On the other hand, cup product is a perfect pairing \(H^1\times H^1\to H^2=\overline{\mathbf Q}_\ell(-1)\), which Frobenius multiplies by \(q\). Hence \(q/\alpha\) is an eigenvalue as well, and the same bound for \(q/\alpha\) says that \(|\iota(\alpha)|\ge q^{1/2}\). In dimension \(d\) the argument is the same: an eigenvalue \(\alpha\) on \(H^i(X,\overline{\mathbf Q}_\ell)\) satisfies \(|\iota(\alpha)|\le q^{i/2}\), while duality makes \(q^d/\alpha\) an eigenvalue on \(H^{2d-i}\), of absolute value at most \(q^{(2d-i)/2}\), and we find \(|\iota(\alpha)|=q^{i/2}\). This equality is Weil's axiom (d), and we prove it in Section \hyperref[bp321]{19.1}, more generally for the higher direct images of a pure lisse sheaf under a smooth proper morphism (Theorem \hyperref[bp323]{19.1}); the case of the constant sheaf over a point is Corollary \hyperref[bp326]{19.2}. Since the Deligne–Laumon theorem holds for every \(\iota\), purity holds for all complex embeddings at once. Indeed, every embedding of a number field into \(\mathbf C\) extends to an isomorphism \(\overline{\mathbf Q}_\ell\simeq\mathbf C\).

The polynomials \(P_{i,\ell}=\det(1-TF\mid H^i(X,\mathbf Q_\ell))\) are defined with \(\ell\)-adic coefficients, and a priori they could depend on \(\ell\). The zeta function does not mention \(\ell\). It has a unique expression \(Z(X_0,T)=A/B\) with \(A,B\in\mathbf Z[T]\) coprime and \(A(0)=B(0)=1\) (Lemma \hyperref[bp333]{19.4}; Sections \hyperref[bp328]{19.2}–\hyperref[bp336]{Section 19.3}). Can each \(P_{i,\ell}\) be read off from \(Z(X_0,T)\)? The trace formula writes \(Z\) as the alternating product of the \(P_{i,\ell}\), but an alternating product can cancel. For example, \(X_0=\mathbf G_m\sqcup
\operatorname{Spec}\mathbf F_q\) is not proper, and it has \(q^m\) points over \(\mathbf F_{q^m}\), as many as \(\mathbf A^1\). Hence \(Z(X_0,T)=1/(1-qT)\). Yet \(H^0_c\) and \(H^1_c\) are both one-dimensional with eigenvalue \(1\), and the factor \(1-T\) cancels between them. For a smooth proper \(X_0\) purity rules this out. The reciprocal roots of \(P_{i,\ell}\) have absolute value \(q^{i/2}\), so no root of \(P_{i,\ell}\) is a root of \(P_{j,\ell}\) for \(i\neq j\). It follows that the product of the \(P_{i,\ell}\) with \(i\) odd is \(A\) and the product with \(i\) even is \(B\). Grouping the reciprocal roots of \(A\) and \(B\) by absolute value then recovers each \(P_{i,\ell}\) from \(Z(X_0,T)\) alone (Theorem \hyperref[bp342]{19.7}). In particular \(P_{i,\ell}\) has integer coefficients and does not depend on \(\ell\).

In terms of point counts, purity gives the classical estimates. For a smooth proper geometrically connected \(X_0\) we obtain \(|\#X_0(\mathbf F_{q^m})-(1+q^{dm})|\le\sum_{i=1}^{2d-1}b_iq^{mi/2}\) (Corollary \hyperref[bp348]{19.10}). In particular, for a curve of genus \(g\) this is Weil's bound \(2gq^{m/2}\), and for elliptic curves it is Hasse's theorem. As an illustration, Example \hyperref[bp346]{19.9} works this out for genus-one curves over \(\mathbf F_3\).

% object: sec-The-Riemann-hypothesis (blueprint/chap-Proof-of-the-Weil-conjectures/sec-The-Riemann-hypothesis/_section.md), content hidden
\subsection*{\texorpdfstring{19.1\quad The Riemann hypothesis}{19.1 The Riemann hypothesis}}
\phantomsection\label{bp321}%
\addcontentsline{toc}{subsection}{\texorpdfstring{19.1\quad The Riemann hypothesis}{19.1 The Riemann hypothesis}}
% text: prose-19-02
\phantomsection\label{bp322}%
For a proper map we have \(Rf_!=Rf_*\), so the upper bound of Theorem \hyperref[bp315]{18.15} applies to ordinary cohomology. If the map is also smooth, Poincaré duality on each fibre pairs \(H^i(X_{\bar y},F)\) perfectly with \(H^{2d-i}(X_{\bar y},F^\vee(d))\). The eigenvalues of the second group are the inverses of those of the first. The second group has weights at most \((-w-2d)+(2d-i)=-w-i\), and this upper bound becomes the lower bound \(w+i\) for the first.

Both hypotheses are needed. Without properness, \(H^1_c\) of \(\mathbf G_m\) has weight \(0\) instead of \(1\). Without smoothness duality fails. For example, the nodal cubic is proper, and its \(H^1\) is a line on which Frobenius acts by \(\pm1\), of weight \(0\). The relative statement is no harder than the absolute one, because we check the bound at one closed point \(y\) at a time, on the smooth proper fibre over \(\kappa(y)\). In the proof of Theorem \hyperref[bp342]{19.7} we use only the case of a point as base and the constant sheaf (Corollary \hyperref[bp326]{19.2}).

% object: theorem-purity-higher-direct-images-under (blueprint/chap-Proof-of-the-Weil-conjectures/sec-The-Riemann-hypothesis/theorem-purity-higher-direct-images-under/_section.md)
\phantomsection\label{bp323}%
\begin{bpstatementit}{Theorem~19.1}{ (Purity of higher direct images under a smooth proper morphism)}
Let \(f_0:X_0\to Y_0\) be smooth proper of pure relative dimension \(d\), and let \(F\) be a lisse Weil sheaf, punctually \(\iota\)-pure of weight \(w\). For every \(i\), the canonical identification

\[
R^if_{0!}F=R^if_{0*}F
\]

is Weil-equivariant, and this constructible sheaf is punctually \(\iota\)-pure of weight \(w+i\). In particular, it is mixed of the single weight \(w+i\).
\end{bpstatementit}

% proof: proof/theorem-purity-higher-direct-images-under
\begin{proof}[{Proof}]\phantomsection\label{bp324}%
Since \(f_0\) is proper, the definition of \(Rf_{0!}\) by compactification identifies it with \(Rf_{0*}\). This identification commutes with the base change maps that define the Weil structures, which proves the first assertion.

Fix a closed point \(y\). By proper base change, the stalk at \(\bar y\) with its local Frobenius is \(H^i(X_{\bar y},F)\). By Theorem \hyperref[bp315]{18.15}, \(R^if_{0*}F\) is mixed of weights at most \(w+i\), so every eigenvalue \(\alpha\) there satisfies

\[
|\iota(\alpha)|\leq q_y^{(w+i)/2}.
\]

The opposite inequality comes from duality. The dual \(F^\vee\) is lisse of weight \(-w\), so \(F^\vee(d)\) has weight \(-w-2d\). On the smooth proper geometric fibre, Corollary \bpout{10.13} and Theorem \bpout{10.21} give a Frobenius-equivariant perfect pairing

\[
H^i(X_{\bar y},F)\times
H^{2d-i}(X_{\bar y},F^\vee(d))\longrightarrow\overline{\mathbf Q}_\ell.
\]

Since the pairing is perfect and takes values in the trivial Weil line, it identifies the second factor with the dual of the first as a Frobenius module, so the eigenvalues on the second factor are the \(\alpha^{-1}\), with multiplicities. Applying Theorem \hyperref[bp315]{18.15} once more, now to the structure map of the fibre over \(\kappa(y)\) and to \(F^\vee(d)\), of weight \(-w-2d\), and to its cohomology in degree \(2d-i\), we get

\[
|\iota(\alpha)|^{-1}
\leq q_y^{(-w-2d+2d-i)/2}=q_y^{(-w-i)/2}.
\]

The two inequalities give \(|\iota(\alpha)|=q_y^{(w+i)/2}\) at every closed point \(y\). Every nonzero pure quotient in a mixed filtration of \(R^if_{0*}F\) has a nonzero stalk at some closed point, so its weight is \(w+i\). This proves the last assertion.
\end{proof}

% text: prose-19-03
\phantomsection\label{bp325}%
For the Weil conjectures we need only the case where the base is a point and the sheaf is constant. Then the fibre is \(X\) itself, and the stalk of \(R^if_*\overline{\mathbf Q}_\ell\) is \(H^i(X,\overline{\mathbf Q}_\ell)\).

% object: corollary-purity-cohomology-smooth-proper-variety (blueprint/chap-Proof-of-the-Weil-conjectures/sec-The-Riemann-hypothesis/corollary-purity-cohomology-smooth-proper-variety/_section.md)
\phantomsection\label{bp326}%
\begin{bpstatementit}{Corollary~19.2}{ (Purity of the cohomology of a smooth proper variety)}
Let \(X_0/\mathbf F_q\) be smooth proper of pure dimension \(d\), put \(X=X_0\times_{\mathbf F_q}\overline{\mathbf F}_q\), let \(\ell\nmid q\) be a prime, and let \(\iota:\overline{\mathbf Q}_\ell\hookrightarrow\mathbf C\) be a field embedding. Then every geometric Frobenius eigenvalue \(\alpha\) on \(H^i(X_{\mathrm{pro}\acute{e}t},\overline{\mathbf Q}_\ell)\) satisfies

\[
|\iota(\alpha)|=q^{i/2}.
\]
\end{bpstatementit}

% proof: proof/corollary-purity-cohomology-smooth-proper-variety
\begin{proof}[{Proof}]\phantomsection\label{bp327}%
The constant sheaf has weight zero, since its local Frobenius is the identity. So Theorem \hyperref[bp323]{19.1} applies to the structure map \(f_0:X_0\to\operatorname{Spec}\mathbf F_q\) with \(w=0\). The base has a single closed point, with \(q_y=q\). The stalk of \(R^if_{0*}\overline{\mathbf Q}_\ell\) there is \(H^i(X,\overline{\mathbf Q}_\ell)\) with its geometric Frobenius, and its eigenvalues therefore satisfy \(|\iota(\alpha)|=q^{i/2}\).
\end{proof}


% object: sec-Rational-power-series (blueprint/chap-Proof-of-the-Weil-conjectures/sec-Rational-power-series/_section.md), content hidden
\subsection*{\texorpdfstring{19.2\quad Rational power series}{19.2 Rational power series}}
\phantomsection\label{bp328}%
\addcontentsline{toc}{subsection}{\texorpdfstring{19.2\quad Rational power series}{19.2 Rational power series}}
% text: prose-19-04
\phantomsection\label{bp329}%
The zeta function has integer coefficients by its definition, but the trace formula proves its rationality only over \(\mathbf Q_\ell\), as an alternating product of \(\ell\)-adic characteristic polynomials. To compare these polynomials for different primes \(\ell\), we need a form of the zeta function that does not mention \(\ell\). We show that a rational power series with integer coefficients is a quotient of integral polynomials in a unique reduced form.

Lemma \hyperref[bp330]{19.3} shows that rationality is a condition of linear algebra on the coefficients, namely that the Hankel matrix \((c_{r+s})\) has finite rank. Such a condition descends from \(\mathbf Q_\ell\) to \(\mathbf Q\). For integrality of the reduced fraction, Lemma \hyperref[bp333]{19.4} uses a lattice in the row span of the Hankel matrix that is stable under the shift. The characteristic polynomial of the shift on this lattice is then integral. For example, the coefficients \(c_n=2^{n+1}-1\) of \(1/((1-2T)(1-T))\) satisfy \(c_{n+2}=3c_{n+1}-2c_n\). The shift matrix \(S\) on the row span has \(\det(1-TS)=(1-2T)(1-T)\), which is the reduced denominator. In the next section we compare the \(\ell\)-adic factorizations with the reduced fraction \(A/B\).

% object: lemma-hankel-rank-rationality (blueprint/chap-Proof-of-the-Weil-conjectures/sec-Rational-power-series/lemma-hankel-rank-rationality/_section.md)
\phantomsection\label{bp330}%
\begin{bpstatementit}{Lemma~19.3}{ (Hankel rank and rationality)}
For a field \(K\) and \(C(T)=\sum_{n\geq0}c_nT^n\in K[[T]]\), the following conditions are equivalent:

\begin{enumerate}
\item \(C(T)\) is rational and regular at zero;
\item the finite submatrices of \((c_{r+s})_{r,s\geq0}\) have bounded rank;
\item the span of the rows \(R_j=(c_{j+s})_{s\geq0}\) is finite-dimensional.
\end{enumerate}

If \(K\subset K'\) and \(C\) has coefficients in \(K\), rationality over \(K'\) implies rationality over \(K\).
\end{bpstatementit}

% proof: proof/lemma-hankel-rank-rationality
\begin{proof}[{Proof}]\phantomsection\label{bp331}%
The equivalence of (2) and (3) is linear algebra. For the equivalence of (1) and (3), note that a denominator \(Q\) gives a linear recurrence among the rows. Conversely, the characteristic polynomial of the shift on a finite-dimensional row span is a denominator. The descent statement holds because (2) does not change when we extend the field.

For example, \(C=1/(1-T)\) has all \(c_n=1\), every row is \((1,1,\ldots)\), and the row span is a line. For \(C=1/(1-T)^2\) we have \(c_n=n+1\), and \(R_j=R_0+j(R_1-R_0)\), so the row span is a plane.

For (3)\(\Rightarrow\)(2), a nonzero \(k\times k\) minor of a finite submatrix gives \(k\) linearly independent rows. So every finite submatrix has rank at most the dimension of the row span. For (2)\(\Rightarrow\)(3), let \(R_{j_1},\ldots,R_{j_m}\) be linearly independent. On their span the coordinate functionals have common kernel zero. They therefore span its dual, and \(m\) of them already do. The corresponding \(m\times m\) minor is nonzero, and so \(m\) is at most the bound in (2).

For (1)\(\Rightarrow\)(3), suppose that \(QC=P\) with \(Q(0)=1\). Comparing coefficients gives a fixed linear recurrence for all sufficiently large indices, and the same recurrence expresses every sufficiently late row in terms of a fixed finite list of earlier rows.

For (3)\(\Rightarrow\)(1), the left shift \(S(v)_n=v_{n+1}\) preserves the finite-dimensional row span \(V\), since \(S(R_j)=R_{j+1}\). Let \(e\) be evaluation at coordinate zero and \(v=R_0\). In any basis of \(V\), with \(M\) the matrix of the shift, we have

\[
C(T)=e(I-TM)^{-1}v
=\frac{e\operatorname{adj}(I-TM)v}{\det(I-TM)}.
\]

The first equality holds coefficient by coefficient. Indeed, \((I-TM)^{-1}=\sum_nT^nM^n\), and \(eM^nv\) is the coordinate zero of \(S^n(R_0)=R_n\), that is, \(c_n\). As the denominator has constant coefficient one, this proves (1).

For the descent, note that condition (2) involves only the coefficients \(c_n\in K\), and the rank of a matrix with entries in \(K\) does not change when we pass to \(K'\). So (2) holds over \(K\) if and only if it holds over \(K'\). Applying (1)\(\Rightarrow\)(2) over \(K'\) and (2)\(\Rightarrow\)(1) over \(K\) gives the descent of rationality.
\end{proof}

% text: prose-19-05
\phantomsection\label{bp332}%
Applied with \(K=\mathbf Q\) and \(K'=\mathbf Q_\ell\), the lemma shows that the zeta function, which the trace formula writes as a quotient of polynomials over \(\mathbf Q_\ell\), is a quotient of polynomials over \(\mathbf Q\). The next lemma shows that the integrality of its coefficients passes to the reduced fraction.

% object: lemma-reduced-fractions-integer-coefficients (blueprint/chap-Proof-of-the-Weil-conjectures/sec-Rational-power-series/lemma-reduced-fractions-integer-coefficients/_section.md)
\phantomsection\label{bp333}%
\begin{bpstatementit}{Lemma~19.4}{ (Reduced fractions with integer coefficients)}
If \(C(T)\in1+T\mathbf Z[[T]]\) is rational over \(\mathbf Q\), then there are unique coprime \(A,B\in\mathbf Z[T]\) such that

\[
C(T)=A(T)/B(T),\qquad A(0)=B(0)=1.
\]
\end{bpstatementit}

% proof: proof/lemma-reduced-fractions-integer-coefficients
\begin{proof}[{Proof}]\phantomsection\label{bp334}%
We repeat the construction in the proof of Lemma \hyperref[bp330]{19.3} with an integral basis of the row span. This gives an integral fraction, which we then reduce with Gauss' lemma.

Let \(V\) be the rational row span of Lemma \hyperref[bp330]{19.3}, and put \(L=V\cap\mathbf Z^{\mathbf Z_{\geq0}}\). Then \(L\) contains all rows, spans \(V\), and is preserved by the left shift. Since \(V\) is finite-dimensional, some \(r\) of the coordinate functionals span its dual, as in the proof of (2)\(\Rightarrow\)(3) in Lemma \hyperref[bp330]{19.3}, and so they are jointly injective on \(V\). These coordinates embed \(L\) into \(\mathbf Z^r\). Since a subgroup of \(\mathbf Z^r\) is free of finite rank, \(L\) is a shift-stable lattice in \(V\), free of rank \(\dim V\).

In a basis of \(L\), the shift matrix \(M\), the row \(e\) and the column \(v=R_0\) are integral. The matrix identity in the proof of Lemma \hyperref[bp330]{19.3} yields integral polynomials

\[
A_0=e\operatorname{adj}(I-TM)v,\qquad B_0=\det(I-TM),
\qquad C=A_0/B_0.
\]

Both constant coefficients are one. Let \(G\) be a primitive integral polynomial that is a greatest common divisor of \(A_0\) and \(B_0\) in \(\mathbf Q[T]\). By Gauss' lemma the quotients \(A_0/G\) and \(B_0/G\) are integral. Since \(G(0)\) divides \(B_0(0)=1\), we may choose the sign of \(G\) so that \(G(0)=1\). Dividing by \(G\), we obtain polynomials \(A\) and \(B\) that are coprime and have constant coefficients one.

For uniqueness, let \(A/B=A'/B'\) be two such fractions, so that \(AB'=A'B\) in \(\mathbf Q[T]\). Since \(B\) is coprime to \(A\), it divides \(B'\), and in the same way \(B'\) divides \(B\). So \(B'=cB\) for a scalar \(c\). Comparing constant coefficients gives \(c=1\), and then \(AB=A'B\) forces \(A=A'\).
\end{proof}

% text: prose-19-06
\phantomsection\label{bp335}%
By Corollary \bpout{12.31}, \(Z(X_0,T)\in\mathbf Q(T)\cap(1+T\mathbf Z[[T]])\) for every separated \(X_0\) of finite type over \(\mathbf F_q\). So Lemma \hyperref[bp333]{19.4} applies. We write \(Z(X_0,T)=A/B\) for the resulting fraction, with \(A\) and \(B\) coprime in \(\mathbf Z[T]\) and \(A(0)=B(0)=1\).


% object: sec-Integrality-and-independence-of-ell (blueprint/chap-Proof-of-the-Weil-conjectures/sec-Integrality-and-independence-of-ell/_section.md), content hidden
\subsection*{\texorpdfstring{19.3\quad Integrality and independence of $\ell$}{19.3 Integrality and independence of ell}}
\phantomsection\label{bp336}%
\addcontentsline{toc}{subsection}{\texorpdfstring{19.3\quad Integrality and independence of $\ell$}{19.3 Integrality and independence of ell}}
% text: prose-19-07
\phantomsection\label{bp337}%
Let \(X_0\) be smooth and proper over \(\mathbf F_q\), of pure dimension \(d\), and let \(A/B\) be the reduced fraction of \(Z(X_0,T)\) (Lemma \hyperref[bp333]{19.4}). We now carry out the argument sketched at the start of the chapter. The point is that purity makes the characteristic polynomials of Frobenius in different degrees pairwise coprime, so no cancellation can occur, and the polynomials can be read off from \(A\) and \(B\). For a prime \(\ell\nmid q\) put

\[
H^i_\ell=H^i(X_{\mathrm{pro}\acute{e}t},\mathbf Q_\ell),\qquad
P_{i,\ell}(T)=\det(1-TF\mid H^i_\ell),\qquad
b_{i,\ell}=\dim_{\mathbf Q_\ell}H^i_\ell.
\]

Since \(X_0\) is proper, \(H^i_c=H^i\), and the determinant identity of Corollary \bpout{12.31} for \(\Lambda=\mathbf Q_\ell\) reads

\[
Z(X_0,T)=\prod_{i=0}^{2d}P_{i,\ell}(T)^{(-1)^{i+1}}.
\tag{19.1}
\]

Since Frobenius is invertible, \(\deg P_{i,\ell}=b_{i,\ell}\).

% object: definition-reciprocal-root (blueprint/chap-Proof-of-the-Weil-conjectures/sec-Integrality-and-independence-of-ell/definition-reciprocal-root/_section.md)
\phantomsection\label{bp338}%
\begin{bpstatement}{Definition~19.5}{ (Reciprocal roots)}
A \emph{reciprocal root} of a polynomial \(P\) over a field is an element \(\alpha\) of an algebraic closure such that \(1-\alpha T\) divides \(P\).
\end{bpstatement}

% text: prose-19-08
\phantomsection\label{bp339}%
Thus the reciprocal roots of \(P_{i,\ell}\) are the Frobenius eigenvalues on \(H^i_\ell\), counted with multiplicity. The next lemma shows that \(A\) and \(B\) sort these eigenvalues by degree, independently of \(\ell\) and of the complex embedding.

% object: lemma-reciprocal-roots-zeta (blueprint/chap-Proof-of-the-Weil-conjectures/sec-Integrality-and-independence-of-ell/lemma-reciprocal-roots-zeta/_section.md)
\phantomsection\label{bp340}%
\begin{bpstatementit}{Lemma~19.6}{ (Reciprocal roots of the zeta function)}
Let \(X_0/\mathbf F_q\) be smooth proper of pure dimension \(d\), let \(Z(X_0,T)=A/B\) be the reduced fraction fixed above, and let \(\ell\nmid q\) be a prime. Then

\[
A=\prod_{i\ \mathrm{odd}}P_{i,\ell},\qquad
B=\prod_{i\ \mathrm{even}}P_{i,\ell}
\quad\hbox{in }\mathbf Q_\ell[T].
\tag{19.2}
\]

Every reciprocal root \(a\in\overline{\mathbf Q}\) of \(A\) or of \(B\) is a nonzero algebraic integer, and there is an integer \(i\), odd for \(A\) and even for \(B\), such that \(|\tau(a)|=q^{i/2}\) for every field embedding \(\tau:\mathbf Q(a)\hookrightarrow\mathbf C\) and such that \(\theta(a)\) is a Frobenius eigenvalue on \(H^i_\ell\) for every field embedding \(\theta:\overline{\mathbf Q}\hookrightarrow\overline{\mathbf Q}_\ell\).
\end{bpstatementit}

% proof: proof/lemma-reciprocal-roots-zeta
\begin{proof}[{Proof}]\phantomsection\label{bp341}%
The identity (19.2) holds because both sides are reduced fractions for \(Z(X_0,T)\) over \(\mathbf Q_\ell\). By Corollary \hyperref[bp326]{19.2}, the reciprocal roots of \(P_{i,\ell}\) have absolute value \(q^{i/2}\) under every embedding \(\iota:\overline{\mathbf Q}_\ell\hookrightarrow\mathbf C\). Since \(q>1\), distinct degrees give distinct absolute values, and so the polynomials \(P_{i,\ell}\) are pairwise coprime. Their products over odd and over even degrees are therefore coprime, and they have constant coefficients one. By (19.1) their quotient is \(Z(X_0,T)\). On the other hand, \(A\) and \(B\) remain coprime in \(\mathbf Q_\ell[T]\), since a Bézout relation over \(\mathbf Q\) remains valid over \(\mathbf Q_\ell\). The two reduced fractions agree by the uniqueness argument in the proof of Lemma \hyperref[bp333]{19.4}, which works over any field.

Write \(A(T)=\prod_a(1-aT)\) and \(B(T)=\prod_b(1-bT)\) over \(\overline{\mathbf Q}\). The reciprocal roots \(a\) and \(b\) are roots of the monic integral polynomials \(U^{\deg A}A(U^{-1})\) and \(U^{\deg B}B(U^{-1})\). Hence they are nonzero algebraic integers. Let \(\theta:\overline{\mathbf Q}\hookrightarrow\overline{\mathbf Q}_\ell\) be an embedding, and let \(a\) be a reciprocal root of \(A\) or \(B\). By (19.2) and the pairwise coprimality of the \(P_{i,\ell}\), there is exactly one \(i\) such that \(\theta(a)\) is a Frobenius eigenvalue on \(H^i_\ell\). It is odd for \(A\) and even for \(B\).

Corollary \hyperref[bp326]{19.2} bounds \(|\iota(\theta(a))|\) for embeddings \(\iota\) of \(\overline{\mathbf Q}_\ell\). The lemma, however, concerns the embeddings of the number field \(\mathbf Q(a)\). We pass from one to the other by extending field embeddings. Let \(\tau:\mathbf Q(a)\hookrightarrow\mathbf C\) be any embedding. The map \(\tau\circ\theta^{-1}\), defined on \(\theta(\mathbf Q(a))\), extends to a field isomorphism \(\iota':\overline{\mathbf Q}_\ell\xrightarrow{\sim}\mathbf C\), because both are algebraically closed fields of characteristic zero with transcendence degree \(2^{\aleph_0}\) over \(\mathbf Q\). We have \(\iota'(\theta(a))=\tau(a)\). Applying Corollary \hyperref[bp326]{19.2} with this \(\iota'\) to the eigenvalue \(\theta(a)\) on \(H^i_\ell\), we find

\[
|\tau(a)|=q^{i/2}.
\]

Since \(q>1\), this absolute value determines \(i\), and so \(i\) does not depend on \(\theta\).
\end{proof}

% object: theorem-weil-conjectures-smooth-proper-schemes (blueprint/chap-Proof-of-the-Weil-conjectures/sec-Integrality-and-independence-of-ell/theorem-weil-conjectures-smooth-proper-schemes/_section.md)
\phantomsection\label{bp342}%
\begin{bpstatementit}{Theorem~19.7}{ (The Weil conjectures for smooth proper schemes)}
Let \(X_0/\mathbf F_q\) be smooth proper of pure dimension \(d\), and let \(P_{i,\ell}\) be as above.

\begin{enumerate}
\item There are unique polynomials \(P_i(T)\in\mathbf Z[T]\) with \(P_i(0)=1\) whose images in \(\mathbf Q_\ell[T]\) are \(P_{i,\ell}(T)\) for every prime \(\ell\nmid q\), and \(Z(X_0,T)=\prod_{i=0}^{2d}P_i(T)^{(-1)^{i+1}}\).
\item The reciprocal roots of \(P_i\) are algebraic integers, and each of them has absolute value \(q^{i/2}\) under every complex embedding.
\item The degrees \(b_i=\deg P_i=\dim_{\mathbf Q_\ell}H^i_\ell\) do not depend on \(\ell\), and \(b_i=b_{2d-i}\). If \(X_0\) is geometrically connected, then \(P_0=1-T\), \(P_{2d}=1-q^dT\), and \(b_0=b_{2d}=1\).
\item Put \(\chi=\sum_i(-1)^ib_i\). Then \(d\chi\) is even and there is \(\varepsilon\in\{1,-1\}\) such that \[
   Z\!\left(X_0,\frac1{q^dT}\right)
   =\varepsilon q^{d\chi/2}T^\chi Z(X_0,T)
   \quad\hbox{in }\mathbf Q(T).
   \]
\end{enumerate}
\end{bpstatementit}

% proof: proof/theorem-weil-conjectures-smooth-proper-schemes
\begin{proof}[{Proof}]\phantomsection\label{bp343}%
For (1) and (2), we define \(P_i\) using only the reduced fraction \(A/B\) of Lemma \hyperref[bp333]{19.4}, by sorting the reciprocal roots of \(A\) and \(B\) by absolute value, and then we compare \(P_i\) with \(P_{i,\ell}\) by Lemma \hyperref[bp340]{19.6}.

For odd \(i\), let \(R_i\) be the submultiset of the reciprocal roots of \(A\) whose absolute value under every complex embedding is \(q^{i/2}\), and for even \(i\) define \(R_i\) in the same way using the reciprocal roots of \(B\). By Lemma \hyperref[bp340]{19.6}, every reciprocal root of \(A\) or \(B\) lies in exactly one \(R_i\). The multisets of reciprocal roots of \(A\) and \(B\) are Galois-stable, with multiplicities. The condition on all complex embeddings is Galois-invariant, and so each \(R_i\) is Galois-stable. Put

\[
P_i(T)=\prod_{a\in R_i}(1-aT).
\]

The coefficients of \(P_i\) are integers, being rational by Galois invariance and algebraic integers because the roots satisfy monic integral equations.

Now let \(\ell\nmid q\) be any prime and \(\theta:\overline{\mathbf Q}\hookrightarrow\overline{\mathbf Q}_\ell\) an embedding. By (19.2) and Lemma \hyperref[bp340]{19.6}, the embedding \(\theta\) maps \(R_i\) onto the multiset of Frobenius eigenvalues on \(H^i_\ell\). Since both polynomials have constant coefficient one, \(P_i=P_{i,\ell}\) in \(\overline{\mathbf Q}_\ell[T]\), hence in \(\mathbf Q_\ell[T]\). As \(\ell\) was arbitrary, this proves the existence in (1), and uniqueness follows by comparing with a single \(P_{i,\ell}\). The product formula for \(Z(X_0,T)\) in (1) is then (19.1), and assertion (2) is Lemma \hyperref[bp340]{19.6}.

The remaining assertions follow from Poincaré duality and the functional equation of Chapter \bpout{12}. For (3), since \(P_i=P_{i,\ell}\), we have \(b_i=\deg P_i=\dim_{\mathbf Q_\ell}H^i_\ell\), which does not depend on \(\ell\), and \(b_i=b_{2d-i}\) by Corollary \bpout{12.33}.

For the second part of (3), assume that \(X_0\) is geometrically connected, so that \(X\) is connected. By Theorem \bpout{10.15}, \(H^0(X,\mathbf Q_\ell)=\mathbf Q_\ell\), with Frobenius acting as the identity, and the trace map \(\operatorname{tr}_X\) identifies \(H^{2d}(X,\mathbf Q_\ell)\) with \(\mathbf Q_\ell(-d)\), on which Frobenius acts by \(q^d\). So \(b_0=b_{2d}=1\), \(P_0=1-T\) and \(P_{2d}=1-q^dT\).

For (4), Poincaré duality pairs the eigenvalues on \(H^i_\ell\) with those on \(H^{2d-i}_\ell\) by \(\alpha\mapsto q^d/\alpha\). Substituting \(T\mapsto1/(q^dT)\) in the product formula therefore exchanges the corresponding factors up to a monomial, and Corollary \bpout{12.33} carries out this computation and gives (4) in \(\mathbf Q_\ell(T)\). Since the zeta function is rational over \(\mathbf Q\), and the scalar \(\varepsilon q^{d\chi/2}\) is a rational number, the identity holds in \(\mathbf Q(T)\).
\end{proof}

% text: prose-19-09
\phantomsection\label{bp344}%
In the following examples we read the theorem backwards. For projective space the point count determines \(Z\), and the grouping of the reciprocal roots by absolute value recovers the \(P_i\). For an elliptic curve the theorem turns one point count into an inequality.

% object: example-zeta-projective-space-elliptic-curve (blueprint/chap-Proof-of-the-Weil-conjectures/sec-Integrality-and-independence-of-ell/example-zeta-projective-space-elliptic-curve/_section.md)
\phantomsection\label{bp345}%
\begin{bpstatement}{Example~19.8}{ (Projective space and elliptic curves)}
We have \(\#\mathbf P^d(\mathbf F_{q^m})=1+q^m+\cdots+q^{dm}\), so that \(Z(\mathbf P^d_{\mathbf F_q},T)=\prod_{i=0}^d(1-q^iT)^{-1}\). Grouping the reciprocal roots by absolute value as in (2), we find \(P_{2i}=1-q^iT\) for \(0\le i\le d\) and \(P_i=1\) for odd \(i\). For an elliptic curve \(C_0\), we have \(P_0=1-T\) and \(P_2=1-qT\) by (3). Since \(b_1=2\) (Appendix \bpout{D}), the polynomial \(P_1\) has two reciprocal roots \(\alpha,\beta\). Poincaré duality matches them as \(\beta=q/\alpha\), so \(P_1=1-aT+qT^2\) with \(a=\alpha+\beta\). The point-counting formula for \(m=1\) gives \(\#C_0(\mathbf F_q)=1-a+q\), that is, \(a=1+q-\#C_0(\mathbf F_q)\in\mathbf Z\). By (2) both \(\alpha\) and \(\beta\) have absolute value \(q^{1/2}\), and we get Hasse's theorem \(|a|\le2q^{1/2}\).
\end{bpstatement}

% object: example-running-example-purity (blueprint/chap-Proof-of-the-Weil-conjectures/sec-Integrality-and-independence-of-ell/example-running-example-purity/_section.md)
\phantomsection\label{bp346}%
\begin{bpstatement}{Example~19.9}{ (The running example: purity)}
For the curve \(C_0\colon y^2=x^3-x\) over \(\mathbf F_3\) of the introduction, the reciprocal roots \(\pm i\sqrt3\) of \(P_1=1+3T^2\) (Example \bpout{12.32}) have absolute value \(\sqrt3\), as (2) predicts. What purity adds to the earlier examples is a bound: for every genus-one curve over \(\mathbf F_3\) we get \(|4-N_1|\le2\sqrt3\), that is, \(1\le N_1\le7\). In the proof, the upper bound \(|\alpha|\le\sqrt3\) comes from the weight bound for \(Rf_!\) (Theorem \hyperref[bp38]{13.11}) applied to the constant sheaf, and \(\alpha\beta=3\) (Example \bpout{10.23}) turns it into the lower bound.
\end{bpstatement}

% text: prose-19-10
\phantomsection\label{bp347}%
The argument of Hasse's theorem works in every dimension. The point-counting formula expresses \(\#X_0(\mathbf F_{q^m})\) through the traces of \(F^m\), and purity bounds each trace by the dimension of the group times the common absolute value of its eigenvalues.

% object: corollary-point-count-estimates (blueprint/chap-Proof-of-the-Weil-conjectures/sec-Integrality-and-independence-of-ell/corollary-point-count-estimates/_section.md)
\phantomsection\label{bp348}%
\begin{bpstatementit}{Corollary~19.10}{ (Point-count estimates)}
Let \(X_0/\mathbf F_q\) be smooth, proper and geometrically connected of dimension \(d\geq1\), and let \(b_i\) be as in Theorem \hyperref[bp342]{19.7}. Then for every \(m\geq1\),

\[
\left|\#X_0(\mathbf F_{q^m})-(1+q^{dm})\right|
\leq\sum_{i=1}^{2d-1}b_iq^{mi/2}.
\]

For a smooth proper geometrically connected curve of genus \(g\), this is

\[
\left|\#X_0(\mathbf F_{q^m})-(1+q^m)\right|
\leq 2gq^{m/2}.
\]
\end{bpstatementit}

% proof: proof/corollary-point-count-estimates
\begin{proof}[{Proof}]\phantomsection\label{bp349}%
Since \(X_0\) is proper, the point-counting formula of Corollary \bpout{12.31} reads

\[
\#X_0(\mathbf F_{q^m})
=\sum_{i=0}^{2d}(-1)^i\operatorname{Tr}(F^m\mid H^i_\ell).
\]

The terms in degrees \(0\) and \(2d\) are \(1\) and \(q^{dm}\) by Theorem theorem-weil-conjectures-smooth-proper-schemes. The left side of the estimate is therefore the absolute value of the sum of the remaining terms. In degree \(i\), the trace is a sum of \(b_i\) eigenvalues of \(F^m\), each of complex modulus \(q^{mi/2}\) by Theorem theorem-weil-conjectures-smooth-proper-schemes. Its absolute value is thus at most \(b_iq^{mi/2}\), and the triangle inequality proves the first bound. For a curve, \(b_1=2g\) by the cohomology calculation of Appendix \bpout{D}.
\end{proof}




\clearpage\phantomsection\hypertarget{bp-outside}{}%
\section*{Outside this export}
\addcontentsline{toc}{section}{Outside this export}
References marked \ensuremath{{}^\ast} are to parts of the blueprint this document leaves out:
\begin{description}
\item[Corollary 2.29] Continuity of cohomology, in Ⅰ The six operations (\bpslug{corollary-{}\allowbreak{}continuity-{}\allowbreak{}cohomology})
\item[Theorem 2.43] Comparison for constructible sheaves, in Ⅰ The six operations (\bpslug{theorem-{}\allowbreak{}comparison-{}\allowbreak{}constructible-{}\allowbreak{}sheaves})
\item[Section 2.12] Coefficient rings, in Ⅰ The six operations (\bpslug{sec-{}\allowbreak{}Coefficient-{}\allowbreak{}rings})
\item[Lemma 2.51] Comparison with completed coefficients, in Ⅰ The six operations (\bpslug{lemma-{}\allowbreak{}comparison-{}\allowbreak{}completed-{}\allowbreak{}coefficients})
\item[Proposition 2.55] Continuous cohomology as pro-étale cohomology, in Ⅰ The six operations (\bpslug{proposition-{}\allowbreak{}continuous-{}\allowbreak{}cohomology-{}\allowbreak{}pro-{}\allowbreak{}etale-{}\allowbreak{}cohomology})
\item[Lemma 3.15] Lattices in lisse $E$-sheaves, in Ⅰ The six operations (\bpslug{lemma-{}\allowbreak{}lattices-{}\allowbreak{}lisse-{}\allowbreak{}e-{}\allowbreak{}sheaves})
\item[Section 3.8] Lattices after stratification, in Ⅰ The six operations (\bpslug{sec-{}\allowbreak{}Lattices-{}\allowbreak{}after-{}\allowbreak{}stratification})
\item[Lemma 3.24] Lattices after a finite stratification, in Ⅰ The six operations (\bpslug{lemma-{}\allowbreak{}lattices-{}\allowbreak{}after-{}\allowbreak{}finite-{}\allowbreak{}stratification})
\item[Lemma 4.27] Finite morphisms, in Ⅰ The six operations (\bpslug{lemma-{}\allowbreak{}finite-{}\allowbreak{}morphisms})
\item[Proposition 5.28] Composition and base change for $Rf_!$, in Ⅰ The six operations (\bpslug{proposition-{}\allowbreak{}composition-{}\allowbreak{}base-{}\allowbreak{}change-{}\allowbreak{}rf-{}\allowbreak{}shriek})
\item[Chapter 6] Constructibility of $f^!$, in Ⅰ The six operations (\bpslug{chap-{}\allowbreak{}Constructibility-{}\allowbreak{}of-{}\allowbreak{}f})
\item[Definition 9.24] Trace map, in Ⅱ Duality and the trace formula (\bpslug{definition-{}\allowbreak{}trace-{}\allowbreak{}map})
\item[Theorem 10.12] Poincaré duality, in Ⅱ Duality and the trace formula (\bpslug{theorem-{}\allowbreak{}poincare-{}\allowbreak{}duality})
\item[Corollary 10.13] Degreewise duality over fields, in Ⅱ Duality and the trace formula (\bpslug{corollary-{}\allowbreak{}degreewise-{}\allowbreak{}duality-{}\allowbreak{}over-{}\allowbreak{}fields})
\item[Theorem 10.15] The trace isomorphism in top degree, in Ⅱ Duality and the trace formula (\bpslug{theorem-{}\allowbreak{}trace-{}\allowbreak{}isomorphism-{}\allowbreak{}top-{}\allowbreak{}degree})
\item[Section 10.5] Weil complexes, in Ⅱ Duality and the trace formula (\bpslug{sec-{}\allowbreak{}Weil-{}\allowbreak{}complexes})
\item[Definition 10.18] Weil structure, in Ⅱ Duality and the trace formula (\bpslug{definition-{}\allowbreak{}weil-{}\allowbreak{}structure})
\item[Theorem 10.21] Poincaré duality and Frobenius, in Ⅱ Duality and the trace formula (\bpslug{theorem-{}\allowbreak{}poincare-{}\allowbreak{}duality-{}\allowbreak{}frobenius})
\item[Example 10.23] The running example: duality, in Ⅱ Duality and the trace formula (\bpslug{example-{}\allowbreak{}running-{}\allowbreak{}example-{}\allowbreak{}duality})
\item[Proposition 11.2] Properties of the trace, in Ⅱ Duality and the trace formula (\bpslug{proposition-{}\allowbreak{}properties-{}\allowbreak{}trace})
\item[Chapter 12] The trace formula, in Ⅱ Duality and the trace formula (\bpslug{chap-{}\allowbreak{}The-{}\allowbreak{}trace-{}\allowbreak{}formula})
\item[Theorem 12.24] Grothendieck–Lefschetz trace formula, in Ⅱ Duality and the trace formula (\bpslug{theorem-{}\allowbreak{}grothendieck-{}\allowbreak{}lefschetz-{}\allowbreak{}trace-{}\allowbreak{}formula})
\item[Section 12.6] The function–sheaf dictionary, in Ⅱ Duality and the trace formula (\bpslug{sec-{}\allowbreak{}The-{}\allowbreak{}function-{}\allowbreak{}sheaf-{}\allowbreak{}dictionary-{}\allowbreak{}2})
\item[Corollary 12.26] Trace functions and direct image with compact support, in Ⅱ Duality and the trace formula (\bpslug{corollary-{}\allowbreak{}trace-{}\allowbreak{}functions-{}\allowbreak{}direct-{}\allowbreak{}image-{}\allowbreak{}compact})
\item[Proposition 12.28] The function–sheaf dictionary, in Ⅱ Duality and the trace formula (\bpslug{proposition-{}\allowbreak{}function-{}\allowbreak{}sheaf-{}\allowbreak{}dictionary})
\item[Theorem 12.30] The trace formula as an Euler product, in Ⅱ Duality and the trace formula (\bpslug{theorem-{}\allowbreak{}trace-{}\allowbreak{}formula-{}\allowbreak{}euler-{}\allowbreak{}product})
\item[Corollary 12.31] Point counting and the zeta function, in Ⅱ Duality and the trace formula (\bpslug{corollary-{}\allowbreak{}point-{}\allowbreak{}counting-{}\allowbreak{}zeta-{}\allowbreak{}function})
\item[Example 12.32] The running example: the trace formula, in Ⅱ Duality and the trace formula (\bpslug{example-{}\allowbreak{}running-{}\allowbreak{}example-{}\allowbreak{}trace-{}\allowbreak{}formula})
\item[Corollary 12.33] Functional equation, in Ⅱ Duality and the trace formula (\bpslug{corollary-{}\allowbreak{}functional-{}\allowbreak{}equation})
\item[Theorem A.2] Nagata's compactification theorem, in Appendices (\bpslug{theorem-{}\allowbreak{}nagata-{}\allowbreak{}s-{}\allowbreak{}compactification-{}\allowbreak{}theorem})
\item[Proposition B.14] The integral Tate module, in Appendices (\bpslug{proposition-{}\allowbreak{}integral-{}\allowbreak{}tate-{}\allowbreak{}module})
\item[Proposition B.15] No Frobenius-fixed vector, in Appendices (\bpslug{proposition-{}\allowbreak{}no-{}\allowbreak{}frobenius-{}\allowbreak{}fixed-{}\allowbreak{}vector})
\item[Appendix C] The Picard scheme of a curve, in Appendices (\bpslug{chap-{}\allowbreak{}The-{}\allowbreak{}Picard-{}\allowbreak{}scheme-{}\allowbreak{}of-{}\allowbreak{}a-{}\allowbreak{}curve})
\item[Theorem C.3] Riemann–Roch and coherent Serre duality, in Appendices (\bpslug{theorem-{}\allowbreak{}riemann-{}\allowbreak{}roch-{}\allowbreak{}coherent-{}\allowbreak{}serre-{}\allowbreak{}duality})
\item[Corollary C.15] Torsion of the Jacobian and field extension, in Appendices (\bpslug{corollary-{}\allowbreak{}torsion-{}\allowbreak{}jacobian-{}\allowbreak{}field-{}\allowbreak{}extension})
\item[Proposition C.17] $H^1$ of a curve and the Tate module of its Jacobian, in Appendices (\bpslug{proposition-{}\allowbreak{}h-{}\allowbreak{}1-{}\allowbreak{}curve-{}\allowbreak{}tate-{}\allowbreak{}module})
\item[Appendix D] Cohomology of torsion sheaves on curves, in Appendices (\bpslug{chap-{}\allowbreak{}Cohomology-{}\allowbreak{}of-{}\allowbreak{}torsion-{}\allowbreak{}sheaves-{}\allowbreak{}on-{}\allowbreak{}curves})
\item[Proposition D.2] Smooth proper curves, torsion prime to the characteristic, in Appendices (\bpslug{proposition-{}\allowbreak{}smooth-{}\allowbreak{}proper-{}\allowbreak{}curves-{}\allowbreak{}torsion-{}\allowbreak{}prime})
\item[Appendix F] Algebraic groups in characteristic zero, in Appendices (\bpslug{chap-{}\allowbreak{}Algebraic-{}\allowbreak{}groups-{}\allowbreak{}in-{}\allowbreak{}characteristic-{}\allowbreak{}zero})
\item[Lemma F.2] Connected components and images of homomorphisms, in Appendices (\bpslug{lemma-{}\allowbreak{}connected-{}\allowbreak{}components-{}\allowbreak{}images-{}\allowbreak{}homomorphisms})
\item[Lemma F.3] Stable subspaces and the center, in Appendices (\bpslug{lemma-{}\allowbreak{}stable-{}\allowbreak{}subspaces-{}\allowbreak{}center})
\item[Proposition F.8] Lie algebras with a faithful semisimple representation, in Appendices (\bpslug{proposition-{}\allowbreak{}lie-{}\allowbreak{}algebras-{}\allowbreak{}faithful-{}\allowbreak{}semisimple-{}\allowbreak{}representation})
\end{description}
\clearpage\phantomsection%
\section*{References}
\addcontentsline{toc}{section}{References}
\bpbibitem{BBD82}{bp353}{A. Beilinson, J. Bernstein, and P. Deligne. \emph{Faisceaux pervers}, volume 100 of \emph{Astérisque}. 1982.}

\bpbibitem{BS15}{bp354}{B. Bhatt and P. Scholze. The pro-étale topology for schemes. \emph{Astérisque}, 369:99–201, 2015. \href{https://people.mpim-bonn.mpg.de/scholze/proetale.pdf}{Author's manuscript}.}

\bpbibitem{Con07}{bp355}{B. Conrad. Deligne's notes on Nagata compactifications. \emph{Journal of the Ramanujan Mathematical Society}, 22:205–257, 2007. \href{https://math.stanford.edu/\~conrad/papers/nagatafinal.pdf}{Author's manuscript}.}

\bpbibitem{Del74}{bp350}{P. Deligne. La conjecture de Weil. I. \emph{Publ. Math. IHÉS}, 43:273–307, 1974.}

\bpbibitem{Del80}{bp351}{P. Deligne. La conjecture de Weil. II. \emph{Publ. Math. IHÉS}, 52:137–252, 1980.}

\bpbibitem{Den22}{bp356}{C. Deninger. There is no “Weil-”cohomology theory with real coefficients for arithmetic curves. \emph{Ann. Sc. Norm. Super. Pisa Cl. Sci.}, 2022. To appear. arXiv:\href{https://arxiv.org/abs/2204.02714}{2204.02714}.}

\bpbibitem{Deu41}{bp357}{M. Deuring. Die Typen der Multiplikatorenringe elliptischer Funktionenkörper. \emph{Abh. Math. Sem. Univ. Hamburg}, 14:197–272, 1941.}

\bpbibitem{Eke90}{bp358}{T. Ekedahl. On the adic formalism. In \emph{The Grothendieck Festschrift, Vol. II}, volume 87 of \emph{Progress in Mathematics}, pages 197–218. Birkhäuser, 1990.}

\bpbibitem{Has36}{bp359}{H. Hasse. Zur Theorie der abstrakten elliptischen Funktionenkörper. I, II, III. \emph{J. Reine Angew. Math.}, 175:55–62, 69–88, 193–208, 1936.}

\bpbibitem{Hum72}{bp360}{J. E. Humphreys. \emph{Introduction to Lie Algebras and Representation Theory}, volume 9 of \emph{Graduate Texts in Mathematics}. Springer, 1972.}

\bpbibitem{Ill94}{bp361}{L. Illusie. Autour du théorème de monodromie locale. \emph{Astérisque}, 223:9–57, 1994. §4. \href{https://www.numdam.org/item/AST_1994__223__9_0/}{Original article}.}

\bpbibitem{Kat88}{bp362}{N. M. Katz. Travaux de Laumon. \emph{Astérisque}, 161–162:105–132, 1988. Séminaire Bourbaki, exp. 691.}

\bpbibitem{KW01}{bp363}{R. Kiehl and R. Weissauer. \emph{Weil conjectures, perverse sheaves and \(\ell\)-adic Fourier transform}, volume 42 of \emph{Ergebnisse der Mathematik}. Springer, 2001.}

\bpbibitem{Lau87}{bp352}{G. Laumon. Transformation de Fourier, constantes d'équations fonctionnelles et conjecture de Weil. \emph{Publ. Math. IHÉS}, 65:131–210, 1987.}

\bpbibitem{Mil86a}{bp364}{J. S. Milne. \href{https://www.jmilne.org/math/xnotes/AVs.pdf}{Abelian Varieties}. In \emph{Arithmetic Geometry}. 1986. Corrected online version.}

\bpbibitem{Mil86b}{bp365}{J. S. Milne. \href{https://www.jmilne.org/math/xnotes/JVs.pdf}{Jacobian Varieties}. In \emph{Arithmetic Geometry}. 1986. Corrected online version.}

\bpbibitem{Ran39}{bp366}{R. A. Rankin. Contributions to the theory of Ramanujan's function \(\tau(n)\) and similar arithmetical functions. II. \emph{Proc. Cambridge Philos. Soc.}, 35:357–372, 1939.}

\bpbibitem{SGA4}{bp367}{M. Artin, A. Grothendieck, and J.-L. Verdier, editors. \emph{Théorie des topos et cohomologie étale des schémas}, volumes 269, 270, 305 of \emph{Lecture Notes in Mathematics}. Springer, 1972–1973. SGA 4.}

\bpbibitem{SGA4½}{bp368}{P. Deligne. \emph{Cohomologie étale}, volume 569 of \emph{Lecture Notes in Mathematics}. Springer, 1977. SGA 4½; \emph{Théorèmes de finitude}, §3, and \emph{Rapport sur la formule des traces}, §3, and \emph{Sommes trigonométriques}, §1.}

\bpbibitem{SGA5}{bp369}{A. Grothendieck, editor. \emph{Cohomologie \(\ell\)-adique et fonctions \(L\)}, volume 589 of \emph{Lecture Notes in Mathematics}. Springer, 1977. SGA 5.}

\bpbibitem{SGA7II}{bp370}{P. Deligne and N. Katz, editors. \emph{Groupes de monodromie en géométrie algébrique}, volume 340 of \emph{Lecture Notes in Mathematics}. Springer, 1973. SGA 7 II; Exposé XIII, (2.1.7).}

\bpbibitem{Stacks}{bp371}{The Stacks Project Authors. \href{https://stacks.math.columbia.edu}{The Stacks Project}. Individual results are cited by their stable tags in the text.}

\bpbibitem{Var05}{bp372}{Y. Varshavsky. Lefschetz–Verdier trace formula and a generalization of a theorem of Fujiwara. 2005. Version 2. arXiv:\href{https://arxiv.org/abs/math/0505564}{math/0505564}. \href{https://arxiv.org/abs/math/0505564v2}{Author's preprint}.}

\bpbibitem{Wei48a}{bp373}{A. Weil. \emph{Sur les courbes algébriques et les variétés qui s'en déduisent}. Hermann, Paris, 1948.}

\bpbibitem{Wei48b}{bp374}{A. Weil. \emph{Variétés abéliennes et courbes algébriques}. Hermann, Paris, 1948.}

\bpbibitem{Wei49}{bp375}{A. Weil. Numbers of solutions of equations in finite fields. \emph{Bull. Amer. Math. Soc.}, 55:497–508, 1949.}

\end{document}
